% 本文件由 LaTeX 排版 Agent 根据有效 translation.json 自动生成。
% author=Athanasius; work_id=2816; title=驳斥亚略派四论

\chapter{\OriginalTerm{驳斥亚略派四论}{Four Discourses Against the Arians}}

% structure: p0001 source=h1 target=omitted reason=page-title-already-rendered-by-parent

\Emph{第一论}\newline{}\Emph{第二论}\newline{}\Emph{第三论}\newline{}\Emph{第四论}\par

\SourceRecord{Translated by John Henry Newman and Archibald Robertson.；Nicene and Post-Nicene Fathers, Second Series；Vol. 4.；Edited by Philip Schaff and Henry Wace.；Buffalo, NY: Christian Literature Publishing Co.,；1892.；<http://www.newadvent.org/fathers/2816.htm>.}

% structure: p0001 source=h1 target=section reason=page-title

\section{反驳亚略派第一讲}

% structure: p0002 source=h2 target=subsection reason=relative-heading-rank; base=h2

\subsection{第一章 引言。写作缘由；某些人对亚略主义漠不关心；亚略派不是基督徒，因为分裂派总以其创始者之名自称。}

1. 在所有其他背离真理的异端中，人们承认它们不过是设计了疯狂，其不虔已经久为人所共知。因为他们的作者是从我们中间出去的，正如真福\Person{若望}所写的，这明显表明：他们从不曾，也从未与我们同思。因此，正如救主所说：他们既不与我们相聚，就是同魔鬼分散；他们窥视沉睡的人，好借着这第二次撒下自己的致命毒种，在死亡中得同伴。然而，有一个异端，即那最后出现的，现今作为\Person{假基督}的先驱而兴起，被称为\OriginalTerm{亚略派}{Arian}的，见其他异端——她的姐姐们——已被公开禁止，便凭她的诡计和狡猾，如同其父魔鬼一样，装模作样地披上圣经的言辞，正强行夺回教会乐园——以便借着基督徒的伪装，她那油滑的诡辩（因为她没有理性）能欺骗人们，导致对\Person{基督}的错误思想。——更有甚者，她既已诱骗了一些愚昧的人，不仅败坏了他们的耳朵，甚至还让他们与\Person{厄娃}一同吃了，直至他们在随之而来的无知中，以苦为甜，赞赏这可憎的异端。为此，我应你的请求，认为有必要拆开\Quote{它胸甲的褶层}，并暴露其愚蠢的恶臭。这样，那些远离它的人可以继续躲避它，而那些受骗的人可以悔改；并开启心灵的眼睛，明白黑暗不是光明，虚假不是真理，亚略主义也不是善；不仅如此，那些称这些人为基督徒的人，乃在重大而严重的错误中，因为他们既没有研读圣经，也完全不明白基督教，以及其中所包含的信德。\par

2. 他们在这异端中发现了什么与宗教信仰相似之处，使他们徒劳地声称，似乎其追随者并没有说什么邪恶的话呢？这实际上就是把\Person{盖法}也称为基督徒，把叛徒\Person{犹达斯}仍算在宗徒之列，说那些要求释放\Person{巴拉巴}而不要求救主的人没有作恶，并推荐\Person{依默纳约}和\Person{亚历山大}是思想正确的人，仿佛\Person{宗徒}诽谤了他们似的。然而，一个基督徒既不能忍受听这些话，也不能认为敢于如此说话的人是心智健全的。因为对他们来说，\Person{亚略}代替了基督，正如摩尼教徒以\Person{摩尼}代替基督一样；为了取代梅瑟和其他圣者，他们竟然找到\Person{索塔德}这样连外邦人都嘲笑的人，以及\Person{黑落狄雅}的女儿。因为亚略在模仿索塔德放荡且柔弱的语调，仿照他写了\WorkTitle{塔利亚}；同时又在亵渎救主时，模仿了黑落狄雅之女的舞蹈，东倒西歪，嬉戏作乐；直至他的异端的牺牲品丧失理智，变得愚妄，将荣耀的主的名字变成\Quote{可朽坏之人的形象}的模样，并且基督徒竟被称为亚略派，以此作为他们不虔的标记。因为他们不要为自己辩解；也不要将他们的羞辱归咎于那些不像他们的人，按着他们教师的名称来称呼基督徒，以便他们自己似乎同样拥有那名号。当他们为自己的可耻称呼感到羞耻时，也不要以此开玩笑；相反，如果感到羞耻，就让他们掩面，或退避自己的不虔。因为基督徒子民在任何时候都不是从他们中间的主教们取得名号的，而是从我们信德所寄讬的主那里得来的。因此，虽然真福宗徒们成了我们的老师，并传了救主的福音，然而我们不是从他们得名，而是从\Person{基督}，我们是并被称为基督徒。但对于那些从别人领受其所宣认的信德的人，他们当然应当承担那人的名号，因为他们已成为那人的产业。\par

3. 确实如此；我们所有人因基督是、也被称为基督徒。很久以前，\Person{玛尔强}宣扬异端并被驱逐；那些与驱逐他的人继续在一起的人仍是基督徒；但那些跟随玛尔强的人不再称为基督徒，此后称为玛尔强派。同样，\Person{瓦伦廷}、\Person{巴西利得}、\Person{摩尼}和\Person{术士西满}，也都将自己的名号传给追随者；于是有人被称为瓦伦廷派、或巴西利得派、或摩尼派、或西满派；另一些因\Place{夫黎基雅}而被称为\OriginalTerm{卡塔夫黎基雅派}{Cataphrygians}，因\Person{诺瓦提安}而被称为诺瓦提安派。\Person{默勒提乌斯}也是如此，当他被主教兼殉道者\Person{伯多禄}驱逐时，称他的党派不再是基督徒，而是默勒提乌斯派。因此，当值得怀念的\Person{亚历山大}驱逐\Person{亚略}时，那些与亚历山大在一起的人仍是基督徒；而那些与亚略一同出去的，将救主的名留给与亚历山大在一起的我们，至于他们，从此就被称为亚略派。看，即使在亚历山大去世后，那些与他的继承人\Person{亚大纳削}共融的人，以及那些与所说的亚大纳削共融的人，都是同一规则的实例；他们中没有一个以他的名字为名，他也不是以他们得名，而是所有人都同样地，并像通常那样，被称为基督徒。因为尽管我们有连续的教师并成为他们的门徒，然而，由于我们被他们教导基督的事，我们仍然是，也一直被所有人称为基督徒。但那些追随异端者的人，尽管在异端中有无数的继任者，却无论如何都承袭了创始者的名号。因此，尽管亚略已经死去，他的许多党羽继任了他，然而那些与他同思想的人，因着亚略而知名，都被称为亚略派。并且，这一点有一个显著的证据：即使在现今，希腊人中那些进入教会的人，一旦弃绝偶像的迷信，都不以他们的慕道教师的名，而是以救主的名，开始被称为基督徒而非希腊人；而那些投奔异端的人，以及所有从教会转而信奉此异端的人，都放弃了基督的名，从此被称为亚略派，因为他们不再持有基督的信德，而是继承了亚略的疯狂。\par

4. 那么，那些对基督徒而言是\OriginalTerm{阿里乌斯狂信者}{Ario-maniacs}的人，如何能是基督徒呢？或者，那些抛弃了宗徒的信仰、成为新邪恶的始作俑者的人，如何能属于天主教会呢？他们离弃了\OriginalTerm{神圣圣经的神谕}{oracles of divine Scripture}，却称\Person{阿里乌斯}的\WorkTitle{塔利亚}为新智慧。这也不无道理，因为他们正在宣扬一种新的异端。因而令人惊奇的是，尽管许多人撰写了大量关于旧约和新约的论著和丰富的讲道，但其中却找不到一部\WorkTitle{塔利亚}；不仅如此，即便在外邦人中那些较为可敬的人士中间也没有，却只见于那些在酒酣耳热之际，夹杂着欢呼和玩笑，唱出此类调子以引他人发笑的人那里。直到这位不可思议的\Person{阿里乌斯}，不取任何严肃的典范，甚至不知可敬为何物，他从其他异端大肆剽窃，却想在滑稽可笑方面独树一帜，唯有\Person{索塔德}可与之匹敌。当他想要狂舞攻击救主时，还有什么比将他那些不虔不敬的卑劣言辞，编成放荡而松散的韵律更合他身的呢？这样，正如\BibleBook{智慧}{篇}所言：\BibleQuote{人从言语的发表即可被认识}，从那些韵律中也可看出作者柔弱的灵魂和思想的败坏。事实上，那个狡诈之徒并未逃脱被识破；尽管他如蛇一般左右扭动，终究还是陷入了法利塞人的错谬之中。法利塞人为要干犯法律，便佯装关切法律的言词；为要否认那受人期待且已临在的主，他们虚伪地使用天主的名，结果当他们说\BibleQuote{你既是人，却把自己当作天主}，又说\BibleQuote{我与父原是一体} \BibleRef{若 10:30}时，就被断定犯了亵渎之罪。同样，这个伪冒而具\OriginalTerm{索塔德之风}{Sotadean}的\Person{阿里乌斯}，假意谈论天主，引用圣经的语言，但在各方面都被认作不虔的\Person{阿里乌斯}，他否认圣子，并将圣子列为受造物。\par

% structure: p0007 source=h2 target=subsection reason=relative-heading-rank; base=h2

\subsection{第二章 摘自\Person{阿里乌斯}的\WorkTitle{塔利亚}。阿里乌斯主张：天主成为圣父，而圣子并非始终存在；圣子是从无中被造；他曾经不存在；他在受生之前并不存在；他是受造的；他被称为智慧和圣言，乃是按照天主的属性；他被造是为了叫我们得以受造；他是天主众多能力之一；可变的；因天主的预知他将成为怎样的人而被高举；并非真天主；只是如同他人一样，因分享而得名；在本质上与圣父相异；不认识也不看见圣父；不自知。}

5. 阿里乌斯的\WorkTitle{塔利亚}及其轻浮之作的开端，其音韵与本质皆属柔弱，如下所引：——\par

\Quote{按照天主选民的信德，天主的明智之士，\newline{}圣洁的儿女，正确分解，领受天主的圣神，\newline{}我从智慧的分享者那里学得这一切，\newline{}他们成就非凡，受天主训诲，凡事通达。\newline{}我沿着他们的足迹，怀着相同的意见行走。\newline{}我，极为著名者，为天主的光荣多受苦难；\newline{}且受天主指教，我获得了智慧和知识。}\par

而他在其中所发的嘲弄之语，令人厌恶且极不虔敬，诸如：—— \Quote{天主并非永远都是圣父。}但\Quote{天主曾有一段时间是独一的，尚不是圣父，后来才成为圣父。}\Quote{圣子并非始终存在；}因为，既然万有都是从无中造生，一切现存的受造物和工程都是被造的，那么天主圣言本身也是\Quote{从无中造生}，且\Quote{他曾经不存在}，\Quote{他在其起源之前并不存在}，而是如同其他万物一样，\Quote{有一个受造的起源}。他说：\Quote{因为，天主曾独自存在，圣言尚未存在，智慧亦未存在。后来，天主愿意造生我们，于是便创造了一个某一位，并命名其为圣言、智慧和圣子，好藉着他来造生我们。}因此，他说有两种智慧：第一种是与天主同存的属性；其次，圣子在这智慧中受生，并且仅仅是因分享这智慧而被称为智慧和圣言。他说：\Quote{因为，智慧因智慧的天主的旨意，在智慧中有了其存在。}同样，他说在天主内除了圣子之外，还有另一个圣言，而圣子又因分享它，按恩宠被称为圣言和圣子。这亦是他们异端的特有观念，如他们其他著作所显示的，即有很多大能者；其中一个是天主本然和永恒的大能；但基督相反，却不是天主的真大能，而是如同其他所谓的大能者一样，其中之一，即蝗虫和蚱蜢，在圣经中不仅被称为\BibleQuote{大能}，甚至被称为\BibleQuote{大能者}。\Quote{其他的大能者有许多，且如同圣子一样，达味在圣咏中论及他们时说：\InnerQuote{万军的上主}或\InnerQuote{众大能者}。}而且，按本性而言，如同其他所有一样，圣言自身也是可改变的；他因自己的自由意志，在选择时保持善；然而，当他愿意时，便能如我们一样改变，因为他的本性是可改变的。他说：\Quote{因为，由于预知他会是善的，天主便预先赐给他这光荣，这光荣后来他作为人，由德行而获得。正因如此，由于他的行为被预知，天主就使他成为如此，使他得以存在。}\par

6. 更有甚者，他竟敢说，\Quote{圣言并非真正的天主；}\Quote{虽然祂被称为天主，但祂并非真正的天主，}而是\Quote{祂同别人一样，只是由于分享恩宠，而在名义上称为天主。}凡所有受造物在本质上都与天主相异且不同，同样地\Quote{圣言在各方面都与圣父的\OriginalTerm{本质}{essence}和\OriginalTerm{特性}{propriety}疏远且不相似}，反而属于有始之物和受造物，且是其中之一。随后，他仿佛承继了魔鬼的放肆，在其\WorkTitle{塔利亚}中声称，\Quote{即使对圣子，圣父也是不可见的，}以及\Quote{圣言不能完善且确切地看见或认识祂自己的父；}甚至祂所认识和看见的，都是\Quote{按照祂自己的程度，}正如我们按照自己的能力认识一样。因为他还说，圣子不仅不能确切地认识父，因为祂在理解上有缺失，而且\Quote{祂甚至不认识自己的本质；}——并且\Quote{圣父、圣子和\OriginalTerm{圣神}{Holy Ghost}的本质，在本性上是分离的、疏远的、不相连接的、不相干的，且彼此毫无共享；}用他自己的话说，\Quote{在本质和光荣上彼此完全不像，直到无限。}因此，关于\Quote{光荣与本质的相似}，他说圣言与圣父和圣神都全然不同。这位不虔敬者就用这种言辞说话，主张圣子自身是分离的，在任何方面都不分享圣父。这些都是\Person{亚略}那些戏谑作品中出现的荒谬言论的部分。\par

7. 谁听到这一切，甚至听到\WorkTitle{塔利亚}的曲调，能不憎恨，且理应憎恨这个在舞台上戏谑如此重大事情的\Person{亚略}呢？当他佯装称呼天主、谈论天主时，谁不视他为那诱惑女人的蛇呢？谁读了其作品的下文，却不从他的不虔敬教义中看出那谬误，那蛇后来用诡辩引诱女人陷入的谬误呢？听到如此亵渎之言，谁不惊愕呢？\BibleQuote{诸天惊愕，大地战栗}见：\BibleRef{耶 2:12}，正如先知所言，因法律的违犯。但太阳更为惊骇，无法容忍那所有人的共同之主为我们的缘故自愿承受的身体耻辱，就转面而去，收回它的光线，使那日变成无光。面对\Person{亚略}的亵渎，全人类岂不都应无言，捂住耳朵，闭上双眼，以免听见或见到其作者？相反，上主岂不有理由用祂借先知欧瑟亚所发的言语，谴责如此不虔敬、如此忘恩负义的人吗？\BibleQuote{他们有祸了，因为他们远离了我；毁灭临于他们，因为他们背叛了我；我虽拯救他们，他们却向我撒谎}见：\BibleRef{欧 7:13}。随后不久，\BibleQuote{他们对我图谋邪恶；他们转向虚无。}因为背离那真实存在的天主圣言，而去为自己塑造一个不存在的，就是堕入虚无。正是为此，\OriginalTerm{大公会议}{Ecumenical Council}，当\Person{亚略}如此说话时，将他逐出教会，并弃绝他，因为无法容忍这种不虔敬。自此以后，\Person{亚略}的谬误被视为比寻常更甚的\OriginalTerm{异端}{heresy}，被认为是基督的仇敌和\OriginalTerm{敌基督}{Antichrist}的前驱。虽然如此重大的谴责本身就有特别的力量，使人逃离这不虔敬的异端，正如我上面所说，然而，仍有一些被称为基督徒的人，或是出于无知，或是假装，以为它与真理相差无几，并称其信奉者为基督徒；那么，就让我们按己所能，向他们提出一些问题，以此揭露此异端的肆无忌惮。或许，当这样被抓住时，他们会闭口无言，并逃离它，如同逃避蛇的身影。\par

% structure: p0013 source=h2 target=subsection reason=relative-heading-rank; base=h2

\subsection{第三章 主题的重要性。\OriginalTerm{亚略派}{Arians}喜爱圣经的语言，但他们的教义却是新颖且不合圣经的。公教教义的陈述：圣子属于圣父的本体，且是永恒的。与之相反，亚略主义重申：祂是一个有起始的受造物。争论的核心在于：我们所当信为天主的那一位，能否只是名义上的天主，且仅仅是一个受造物？那么，在这场争论中有何借口可以漠不关心呢？亚略派依赖国家的庇护，不敢公开承认他们的信条。}

8. 这样，若按他们看来，借用圣经的某些用语，就能将\WorkTitle{Thalia}的亵渎言辞转变为虔敬语言；那么理所当然，当他们看见现今的犹太人如何研读法律与先知时，他们也该与他们一同否认基督了；或许他们还会像摩尼教徒一样，否认法律与先知，因为后者也读福音书的某些部分。倘若这种困惑和空谈是出于无知，圣经会教导他们，魔鬼——异端的制造者，因着邪恶所固有的恶臭，乃借用圣经的言语，作为外衣，好将自己的毒种播撒其间，并诱惑心地纯朴的人。他如此欺骗了\Person{厄娃}；如此构造了先前的异端；如此在此时说服\Person{亚略}，让他表面上反对那些先前的异端，好将自己的异端不知不觉地引进来。然而，那诡诈之人终究没能逃脱。他因为对天主的圣言不虔敬，便立刻丧失了一切，并向众人暴露了他对其他异端的无知；且在自己的信仰中毫无真理，只是假装有真理罢了。因为，那否认圣子的人，怎能论及圣父讲论真理呢？圣子正是启示圣父的那位；一个人对供给圣神的圣言出言亵渎，又怎能对圣神有正统的认识呢？他既否认基督为从死者中首生者，谁还会在复活一事上信靠他呢？对于圣子由圣父所生的真实而真正的诞生，他全然无知，又如何能在祂降生成人的临在一事上不出错呢？过去，犹太人也是这样，他们否定了圣言，说：\BibleQuote{我们除了凯撒，没有君王}\BibleRef{若 19:15}，于是立刻被剥夺了他们所有的一切，丧失了灯的光亮、膏油的馨香、预言的知识，以及真理本身；直到如今，他们什么都不明白，犹如在黑暗中行走。因为，有谁曾听过这样的道理呢？这异端的同伙和佣工又从何处、从何人获得了它？当他们上学时，谁曾如此向他们讲解过？谁曾告诉他们：\Quote{抛弃对造物的敬拜，然后来崇拜一个受造物和工程}？但若他们自己承认如今是头一次听到此事，又怎能否认这异端是外来的、并非出自我们的教父呢？凡不是出自我们的教父，而在今日才显露出来的，岂不正是\Person{圣保禄}所预言的：\Quote{在最后的时期，有人要背弃信德，听信欺诈的神和魔鬼的训言，这训言是出于那些伪善的说谎者，他们的良心已经烙上了火印，离弃了真理}吗？\par

9. 看，我们拿起圣经，从中自由地宣讲虔敬的信仰，并将它高举在灯台上，如同光明，说道：—— 祂是圣父的真圣子，是本性而真实的，属于圣父的本质，是独生的智慧，是天主的真与独一的圣言；祂不是受造物或工程，而是属于圣父本质的亲生子。因此，祂是真天主，与真圣父是同一性体；至于其它的存在者，祂曾对他们说：\Quote{我曾说过：你们是神}，他们则从圣父那里，借着对圣言的分有，经由圣神，获得这种恩宠。因为祂是圣父位格的表达，是出自光明的光明，是大能，是圣父本质的真肖像。因为主也曾说过：\BibleQuote{谁看见了我，就是看见了父}\BibleRef{若 14:9}。祂一直存在，现今存在，且从未不存在。因为圣父既永恒，祂的圣言和智慧也必然是永恒的。另一方面，这些人从那可耻的\WorkTitle{Thalia}中能给我们显出什么呢？或者，首先让他们自己读一读，模仿作者的语调；至少，他们从别人那里所要遭到的嘲笑，会教导他们已堕落到何等地步；然后再让他们解释自己。因为他们从它里面能说什么呢？无非是说：\Quote{天主并非始终是父，而是后来才成为父；圣子并非始终存在，因为在祂受生之前祂并不存在；祂并非由圣父而来，而是像其他事物一样，从虚无中形成实体；祂不属于圣父的本质，因为祂是受造物和工程。} \newline{}又说：\Quote{基督不是真天主，而是像其他人一样，因着分有才成为天主；圣子没有对圣父的确切知识，圣言也不能完美地看见圣父；祂既不精确地理解，也不认识圣父。祂不是圣父的真与独一的圣言，只是名义上被称为圣言和智慧，并且是因恩宠而被称为子和能力。祂不像圣父那样是不能变更的，而是像受造物一样，在性体上是可变更的，祂无法获臻圣父的完美知识。} \newline{}这异端真是惊人，甚至毫无道理，竟敢妄自臆测那位自有者，说祂不存在，并将亵渎之辞处处冒充为虔敬的话！如果有人查问了双方之后，被问到，在这两方中，他愿意信仰哪一方，或者哪一方所谈论的天主相称呢——毋宁让这些不信者的帮腔人自己说吧，若有人问及天主（因为\Quote{圣言就是天主}），怎样的回答才是合宜的。因为从这一个问题，双方便可定案，何为相称的说法——祂是自有者，或不是；是始终存在，还是在其受生之前存在；或是从彼时彼地才开始存在；或是真实的，还是因义子、分有和在观念上才如此；是称祂为受造物之一，还是将祂与圣父结合；是认为祂在本质上与圣父不同，还是相似并专属于祂；是受造物，还是万物藉祂而成为受造物；祂是圣父的圣言，还是在祂之外另有一个圣言，并且祂是藉这另一位圣言和另一位智慧而形成；祂只是名义上被称为智慧和圣言，并成为这智慧的分有者，而次于它？\par

10. 两种神学中，哪一种是宣认我们的主耶稣基督为天主、为父的圣子呢？是你们所吐出的那种呢，还是我们从圣经中所讲论和坚持的这种呢？如果救主不是天主，不是圣言，不是圣子，那么你们就可以随意说什么，外邦人和现今的犹太人也可以如此。但如果祂是父的圣言、真圣子，出自天主的天主，且\BibleQuote{在万有之上，永远应受赞颂的\BibleRef{罗 9:5}}，那么，难道不应该抹去并删除那些其他的言论以及那部亚略的\WorkTitle{塔利亚}吗？它不过是邪恶的典型，是一切不虔敬的仓库；凡落入其中的，\Quote{不知巨人与她一同丧亡，且下到\OriginalTerm{阴府}{Hades}的深处？}他们自己知道这一点，却狡猾地隐藏起来，没有勇气明说，却另说一套。因为如果他们说出，就会受到定罪；如果他们遭人怀疑，圣经的证据就会从四面八方向他们投来。因此，他们狡猾地，如同今世之子，用野橄榄油点着他们所谓的灯之后，害怕它很快会熄灭（因为经上说：\BibleQuote{恶人的光必要熄灭\BibleRef{约 18:5}}），就把灯藏在自己伪善的斗底下，假装承认别的，自夸有朋友的庇护和\Person{君士坦提乌斯}的权势，好叫他们靠着伪善和所承认的口头信仰，使那些前来的人看不见他们的异端是何等污秽。难道连这一点也不可恶吗？它竟不敢明讲，却被自己的朋友隐藏起来，如同蛇被养大一样。他们是从哪里收集到这些话的呢？或是从谁那里领受了他们竟敢宣讲的呢？无人能指出供应这些言语的任何人。因为在全人类中，无论是希腊人还是蛮人，有谁敢把一位他同时宣认为天主的，列入受造物中，并说祂在受造之前并不存在呢？又有谁，对他所信仰的天主，当天主说：\BibleQuote{这是我的爱子\BibleRef{玛 3:17}}时，竟借口祂不是子，而是受造物，就不肯相信呢？毋宁说，这种疯狂会激起普遍的义愤。圣经也没有给他们任何借口；因为经常已证明，而且现在将要证明，他们的道理是与神圣的圣言相悖的。因此，既然唯一剩下的就是说他们的疯狂来自魔鬼（因为这类意见只有他是撒种者），就让我们着手抵抗他——因为我们的真正对手是他，他们只是工具而已——好叫我们靠着主的助佑，如同平常一样，用论证战胜仇敌，使他们看见在他们内播下这异端的他毫无办法，便自觉羞愧，虽然迟了，也学到：作为亚略派，他们不是基督徒。\par

% structure: p0017 source=h2 target=subsection reason=relative-heading-rank; base=h2

\subsection{第四章 圣子是永恒而非受造的。这些属性是争论的要点，首先以直接的圣经经文加以证明。论及罗 1:20中天主的\Quote{永恒大能}，这被证明是指圣子。评述亚略派的公式：\Quote{曾有一时，圣子不存在}，其支持者不敢明说\Quote{一个圣子不存在的时刻}。}

11. 在他的唆使下，你们就主张并认为，\Quote{曾有一时，圣子不存在}；这是你们教义观点的第一层外衣，必须被剥去。那么，你们这些毁谤而不虔敬的人啊，说说看，在圣子不存在的那\Quote{一时}，有什么存在呢？如果你们说是父，那你们的亵渎就更大了；因为说父是\Quote{一时}存在，或用\Quote{一时}这个词指称祂，是不虔敬的。因为祂常存、现今存在，如同子存在一样，祂也存在，祂自己就是那自有的，是圣子的父。但如果你们说子曾存在于某一时，而那时祂自己并不存在，这回答是愚昧而无意义的。因为祂怎能既存在又不存在呢？在这困难中，你们只能回答说，曾有一段时间，圣言并不存在；因为你们所用的副词\Quote{一时}自然意味着这一点。你们另一句话，\Quote{子在受生之前不存在}，也就等于说，\Quote{曾有一时，祂不存在}，因为这两种说法都意味着在圣言之前有一段时间。那么，你们这一发现从何而来？为什么你们像\Quote{外邦狂怒，构想空言反对上主和祂的基督}那样呢？因为没有任何圣经对救主使用这样的言语，反而使用\Quote{永远}、\Quote{永恒}、\Quote{常与父同在}这些话。因为，\BibleQuote{在起初已有圣言，圣言与天主同在，圣言就是天主\BibleRef{若 1:1}}。在\WorkTitle{默示录}中，他这样说：\Quote{那今在、昔在及将来永在者}。现在，谁能剥夺\Quote{今在}和\Quote{昔在}的永恒性呢？同样，为驳斥犹太人，\Person{保禄}在\WorkTitle{罗马书}中写道：\BibleQuote{论及他的血肉，基督是从他们出来的，他是在万有之上，永远应受赞颂的天主\BibleRef{罗 9:5}}；为使希腊人闭口无言，他说：\BibleQuote{其实，自从天主创世以来，祂那看得见的事物，都可凭祂所造的万物，清楚地看见，即祂永远的大能和祂为神的本性}；至于天主的大能是什么，他在别处亲自教导我们说：\BibleQuote{基督是天主的大能和天主的智慧}。毫无疑问，在这句话里，他并非指父，不像你们彼此低声传说，硬说父是\Quote{祂的永恒大能}。并非如此；因为他不是说\Quote{天主自己是这大能}，而是说\Quote{祂的大能}。这对所有人都十分明显，\Quote{祂的}并非\Quote{祂自己}；但又非某种外来的东西，而是专属于祂的。也要研究上下文，\Quote{转向主}；现在\Quote{主就是那神}；你们就会看出，所指的是子。\par

他既提起了创造，自然也就谈到了在创造中所见的造化者的大能；这大能，我指的是天主的圣言，万物是藉着祂而造成的。如果创造本身，无需圣子，就足以使人认识天主，那么你就该小心，不要因以为创造无需圣子而跌倒。但如果创造是藉着圣子而有的，且\BibleQuote{万有都赖祂而存在}\BibleRef{哥 1:17}，那么必然推论出，谁正确地默观创造，便也是在默观创造它的圣言，并藉着祂开始认识圣父。而且，正如救主也说过：\BibleQuote{除了圣父外，没有人认识子；除了子和子所愿意启示的人外，也没有人认识圣父}\BibleRef{玛 11:27}，并且当\Person{斐理伯}请求：\BibleQuote{把圣父显示给我们罢！}时，祂没有说：\BibleQuote{看哪，受造物}，而是说：\BibleQuote{谁看见了我，就是看见了圣父}\BibleRef{若 14:8}。同样，\Person{保禄}也合理地在指责希腊人默观创造的和谐与秩序，却不省思创造内在的造化圣言（因为受造物为自己的造化者作证），好能藉着创造认识真天主，并放弃对它的崇拜时，合理地说出了\BibleQuote{祂永远的大能和天主性}\BibleRef{罗 1:20}，以此指示圣子。而神圣作者在说\BibleQuote{祂在万世之前就已存在}，以及\BibleQuote{藉着祂造成了万世}\BibleRef{希 1:2}时，他们便以此同样清楚地宣讲了圣子的永恒和永存的存在，甚至就在他们指明是天主自己的时候。因此，既然\Person{依撒意亚}说：\BibleQuote{永生的天主，上主，创造地极的主}\BibleRef{依 40:28}，\Person{苏撒纳}呼求说：\BibleQuote{永生的天主啊！}，\Person{巴路克}写道：\BibleQuote{在我有生之日，我必呼求永生者}，以及稍后：\BibleQuote{我的希望寄托于永生者，祂必拯救你们，喜乐从圣者那里临到我}；然而，正因为宗徒致希伯来人书中说：\BibleQuote{祂是天主光荣的反映，是天主本体的真像}\BibleRef{希 1:3}，并且\Person{达味}在第\BibleBook{圣咏}{八十九篇}中说：\BibleQuote{愿上主的荣光临照我们}，又说：\BibleQuote{在祢的光中，我们看见光}，谁还会这么愚昧，竟怀疑圣子的永恒呢？因为，人什么时候见过光而没有它光辉的照耀，以至于对圣子说\Quote{曾有一时，祂不存在}，或\Quote{在祂受生之前，祂不存在}呢？而在第\BibleBook{圣咏}{一百四十四篇}中对圣子所说的话：\BibleQuote{祢的王国是万代的王国}，禁止任何人想象曾有任何圣言不存在的时段。因为，如果岁月中的每一时段都是可度量的，而圣言却是所有岁月的君王和创造者，那么，由于在祂之前没有任何时段存在，说\Quote{曾有一时永生者不存在}和\Quote{圣子来自虚无}便是疯狂。而既然主亲自说：\BibleQuote{我是真理}，而不是\Quote{我成为真理}，却总是\BibleQuote{我是——我是善牧——我是世界的光}，又说：\BibleQuote{你们不是称呼我为主、为师傅吗？你们说得对，我本是}。谁听到了来自天主、智慧与圣父的圣言的这些话论及自己，还会对真理再犹豫，而不即刻相信\Quote{我是}这一表述是指示圣子是永恒且无始呢？\par

由上可知，圣经宣布了圣子的永恒；同样明显的是，接下来亚略派所用的短语\Quote{祂不存在}、\Quote{之前}和\Quote{当……时}，在同一圣经中却是用来述说受造物的。例如，\Person{梅瑟}在描述我们世界的生成时说：\BibleQuote{地上还没有生出田野的植物，田间还没有长出蔬菜，因为上主天主还没有使雨降在地上，也没有人耕种土地}\BibleRef{创 2:5}。在\BibleBook{}{申命纪}中：\BibleQuote{当至高者为民族分配产业时}\BibleRef{申 32:8}。主亲自说：\BibleQuote{如果你们爱我，就该喜欢我往圣父那里去，因为圣父比我大。如今在事发生前，我就告诉了你们，为叫你们当事发生时能相信}\BibleRef{若 14:28}。论及受造的智慧时，祂藉\Person{撒罗满}的口说：\BibleQuote{大地还没有创造以前，深渊还没有存在以前，我就被生；水泉还没有涌出以前，山岳还没有奠定，丘陵还没有存在以前，我就被生}\BibleRef{箴 8:23}。以及：\BibleQuote{在\Person{亚巴郎}出现以前，我存在}\BibleRef{若 8:58}。论及\Person{耶肋米亚}时，祂说：\BibleQuote{我还没有在母腹内形成你以前，我已认识了你}\BibleRef{耶 1:5}。\Person{达味}在圣咏里说：\BibleQuote{群山尚未形成，大地和世界尚未生出，从永远直到永远，祢就是天主}。在\BibleBook{}{达尼尔}中：\BibleQuote{\Person{苏撒纳}大声呼喊说：永生的天主啊！祢洞察隐秘，事前尽知一切}。由此可见，\Quote{曾不存在}、\Quote{未发生以前}、\Quote{当……时}之类的短语，属于来自虚无的受造物和产物，却与圣言无关。但若圣经对受造物使用这些术语，对圣言却使用\Quote{永远}，那么，天主的仇敌啊，结论便是：圣子并非来自虚无，也根本不属于受造之物的范畴，而是圣父的肖像和永恒的圣言，从没有不存在过，自始至终存在，如同永恒之光的永恒辉耀。那么，为何要设想圣子之前的时日呢？或者，为何要亵渎圣言，说祂在时日之后呢——连万世都是藉着祂造成的？因为，按照你们的说法，当圣言尚未出现时，时光和万世又如何能存在呢？万物是藉着祂而造成的，没有祂，什么也没有造成\BibleRef{若 1:3}。或者，当你们意指时间时，为什么不直说\Quote{曾有一个时间圣言不存在}呢？然而你们却略去\Quote{时间}一词以欺骗单纯的人，但你们完全遮掩不了自己的感受，即使试图遮掩，也终究会被揭露。因为你们说\Quote{曾有一时祂不存在}和\Quote{在祂受生以前祂不存在}时，仍然直接意指时间。\par

% structure: p0021 source=h2 target=subsection reason=relative-heading-rank; base=h2

\subsection{续论主题 反对者以圣子的永恒性使其与圣父同位，由此引入圣子的身份作为其永恒性的第二个证据。\Quote{圣子}一词虽在次要意义上被引入，但应从真实意义上理解。由于万物在\OriginalTerm{分有}{participation}圣子时分有圣父，他就是圣父的完全分有，即就本性而言他是圣子；因为完全被分有便是\OriginalTerm{受生}{generation}。}

14. 当这些论点如此得到证明后，他们的亵渎更进一步。他们说：\Quote{如果从来没有过圣子不存在的时候，但他是永恒的，与圣父同存，你们就不称他为圣父的圣子，而称为弟兄了。}噢，无知又好辩的人啊！因为如果我们只说他是永恒地同圣父在一起，而不是他的圣子，他们假装的疑虑还有些道理；但如果我们在说他是永恒的同时，也承认他是由圣父所生的圣子，那么受生者怎能被视为生他者的弟兄呢？如果我们的信仰在于圣父与圣子，他们之间有什么弟兄关系呢？\OriginalTerm{圣言}{Word}怎能被称为那以他为圣言者的弟兄呢？这不是真正无知之人的反驳，因为他们明白真理何在；这是犹太人的借口，是那些如同撒罗满所说的话，\BibleQuote{因着欲望离群索居}，\BibleRef{箴 18:1}，脱离真理。因为圣父和圣子不是从某个先在的本原受生的，好让我们认为他们是弟兄，而是圣父是圣子的本原，生了圣子；圣父是圣父，不由任何圣子所生；圣子是圣子，而非弟兄。再者，如果他被称作圣父永恒的\OriginalTerm{子嗣}{offspring}，这样称呼是正当的。因为圣父的\OriginalTerm{性体}{essence}从不是不完美的，以致之后需要添加适合它的东西；圣子也不是像人由人那样受生，从而晚于圣父的存在，而是他是天主的\OriginalTerm{子嗣}{offspring}，并且是天主的本有圣子，天主永远存在，他也就永恒存在。因为，人由于本性不完美，因而在时间内生育，而天主的子嗣是永恒的，因为他的本性永远完美。因此，如果他不是圣子，而是从虚无中造成的工程，他们就必须证明这一点；那时他们可以随意，好像想象一个受造物那样，高喊：\Quote{曾有一时，他不存在}；因为受造之物过去不存在，后来才有。但如果他是圣子，正如圣父所说，圣经所宣告的，\Quote{圣子}无非就是从圣父受生者；而从圣父受生的，就是他的圣言、智慧和光辉；那么，坚持\Quote{圣子曾一时不存在}，他们岂不是像掠夺者一样抢夺天主的圣言，并公开声称他曾一度没有自己的圣言和智慧，光曾一度没有光辉，泉源曾一度荒芜干涸吗？因为尽管他们假装对时间的名称感到恐慌，由于那些以此责备他们的人，并说他在万世之前就已存在，然而他们却设定某些时段，在其中想象他不存在，他们仍然是最不虔敬的，因为他们同样暗示时间，并把没有\OriginalTerm{圣言}{Reason}归给天主。\par

15. 另一方面，如果他们虽然与我们一同承认\Quote{圣子}的称号，不愿公开普遍地受到谴责，却否认圣子是圣父的\OriginalTerm{性体}{essence}的本有\OriginalTerm{子嗣}{offspring}，理由是这必然意味着部分与分割；那么这岂不是否认他是真正的圣子，而仅仅在名义上称他为圣子吗？对非物质的东西持有物质的想法，并且由于自己本性的软弱而否认那属于圣父的本性和特征，这难道不是严重的错误吗？这只能导致他们进一步否认他，因为他们不明白天主是怎样的，圣父是什么，这些愚昧的人竟然按自己来衡量圣父的子嗣。对于那种认为不可能有天主之圣子的人，我们应当怜悯；但必须质问并揭露他们，好使他们可能恢复理智。那么，如果如你们所说，\Quote{圣子是从虚无中来的}，并且\Quote{在他\OriginalTerm{受生}{generation}之前不存在}，那么他当然，和其他受造物一样，只能因\OriginalTerm{分有}{participation}而被称作圣子、天主和智慧；因为其他所有受造物也是如此构成，并借着\OriginalTerm{圣化}{sanctification}而得光荣。那么你必须告诉我们，他分享了什么。其他一切受造物都分享了圣神，但他，按照你们的说法，分享了什么呢？分享了圣神？不，倒是圣神自己从圣子领受，正如他亲口说的；若说后者被前者圣化，这并不合理。因此，他所分享的是圣父；因为只剩下这个可说了。但这个被分有的，是什么呢？从哪里来呢？如果它是圣父提供的一个外在的东西，那么他就不是圣父的\OriginalTerm{分享者}{partaker}，而是圣父以外的那个东西的分享者；他也不再是仅次于圣父的了，因为他前面有了那个东西；他也不能被称为圣父的圣子，而要被称为那个东西的圣子，因分享那个东西，他才被称作圣子和天主。如果这是不恰当且不虔敬的，那么当圣父说：\BibleQuote{这是我的爱子}\BibleRef{玛 3:17}，并且圣子说天主是他的圣父时，结论就是，所分有的不是外在的，而是来自圣父的性体。对于这一点，如果它不同于圣子的性体，我们同样会遇到荒谬；因为那样，在这来自圣父的东西和圣子的性体之间，就存在着某种东西，无论它是什么。\par

16. 因此，这些思想既然显然不雅且虚妄，我们就不得不说，凡是出自圣父的性体、且为祂本有的，就完全是圣子；因为说天主完全被分受，与说祂生圣子，是同一回事；而生除了生一个子还能是什么意思呢？这样，万物都靠着从圣子那里来的圣神的恩宠，分享圣子本身；这表明圣子自己并不分享什么，而是凡由圣父所分享的，就是圣子；因为，正如我们分享圣子本身，我们就被称为分享了天主；这就是伯多禄所说的话：\BibleQuote{为使你们能成为有分于天主性体的人}\BibleRef{伯后 1:4}；宗徒也这样说：\BibleQuote{你们不知道，你们是天主的宫殿吗？}\BibleQuote{我们是永生天主的宫殿。}再者，看见了圣子，我们就是看见了圣父；因为对圣子的思想和领悟，就是对圣父的认识，因为祂是出自圣父性体的本有之子。既然没有人会把\Quote{被分受}称为天主性体的受动或分割（因为已经指出并承认，天主是被分受的，而被分受等同于生）；因此，那受生者既不是那真福性体的受动，也不是它的分割。所以，天主应有一个儿子，祂自己性体的子，这并非不可信；我们称\Quote{子}和\Quote{子}时，也不意指天主性体有受动或分割；相反，我们是在承认那真实、真正的、为天主所生的独生子时这样相信的。因此，正如我们已陈述并正在表明的，如果出自圣父性体的子就是圣子，我们就不能犹豫，反而必须确信，这同一位也是圣父的智慧和圣言，圣父在祂内并藉着祂创造并造成万物；也是圣父的光辉，圣父在祂内光照万物，并随意启示给人；也是圣父的肖像与形象，在祂内得被观照和认识，因此，\BibleQuote{祂与父是一体}\BibleRef{若 10:30}，谁看见了祂，就是看见了圣父；而且是基督，在祂内，万物得救赎，新的创造得以完成。另一方面，圣子既是这样的子，那么说祂是从无中造成的受造物，或说祂在受生之前不存在，就不相宜了，甚至充满危险。因为如此谈论圣父性体本有之事的人，已经在亵渎圣父本身了；因为他把关于圣子的虚妄想象，真实地归到了圣父身上。\par

% structure: p0025 source=h2 target=subsection reason=relative-heading-rank; base=h2

\subsection{第6章 继续主题。圣子永恒的第三项证明，即从其他表明祂同性同体的称号；如造物主；真福天主圣三中的一位；如智慧；如圣言；如肖像。如果圣子是圣父的完美肖像，为何祂不也是一位父呢？因为天主是完美的，并非一个族类的起源。唯独圣父是父，因为唯一的父；唯独圣子是子，因为唯一的子。人并非真正的父亲和真正的儿子，而是那真实者的影子。圣子不会成为父，因为祂由圣父领受了永恒不变、永远相同的特质。}

17. 这本身就足以驳斥亚略异端；然而，从以下几点亦可看出其异端道理：——如果天主是创造者与造物主，且藉着圣子创造祂的工程，我们又不能视万物而不经由圣言，那么说天主是创造者时，却说祂那塑造万物的圣言与智慧曾经不存在，岂不是亵渎吗？这等于是说，如果天主没有祂本有的、由祂而来的塑造圣言，而是祂藉以塑造的，是从外界附加于祂，与祂陌异且在性体上不相像，那么天主就不是创造者了。其次，让他们告诉我们——或者不如从其中学习，他们宣称\Quote{祂曾经不存在}以及\Quote{祂在受生之前不存在}是多么不虔敬：因为如果圣言不是自永远就与圣父同在，那么天主圣三就不是永恒的；而是一元先存，后来经由增添才成为天主圣三；这样随着时间推移，似乎我们关于天主的知识在增长并成形。再者，如果圣子不是圣父性体的本有之子，而是从无中产生的，那么天主圣三就是由无构成的，且曾经没有天主圣三，只有一元；且天主圣三先是有所欠缺，然后才得完备；在圣子受生之前有所欠缺，到祂产生后才完备；这样，一个受造之物竟被列入造物主之列，而曾经不存在者竟与永恒存在者同享属神的敬拜与荣耀。不仅如此，更严重的是，天主圣三竟被发现与自己不相像，由不同本质和陌异的性体构成。换言之，这等于是说，天主圣三具有受造的组成。那么，这是一种怎样的宗教呢？它甚至不与自己相似，而是随时间推移而趋于完成，时而如此，时而又如彼。因为它很可能还会接受某种新的增添，如此无限制地继续下去，因为从一开始，它就是由增添而组成的。同样，它无疑也可能减少，因为凡是增添的，显然容许被减去。\par

18. 但这并非如此：绝无此意！\OriginalTerm{圣三}{Triad}并非受造；反之，在\OriginalTerm{圣三}{Triad}内存在着永恒而独一的\OriginalTerm{天主性}{Godhead}，圣三拥有同一的荣耀。你们竟敢将它分割为不同的性体；圣父是永恒的，你们却说坐在祂身旁的圣言\Quote{曾经不存在}；而且，圣子坐在圣父身旁，你们却想将祂置于远离圣父之处。圣三是创造者和塑造者，你们却不怕将它贬低到受造之物；你们不惮将奴仆与圣三的尊贵等同，将君王、万军的上主与臣仆并列。停止这种对不相容事物的混淆，或更确切地说，停止将非存在者与存在者相混淆。这样的言论不会荣耀和尊崇主，反而相反；因为凡不尊敬子的，就是不尊敬父。如果关于天主的教义现今在圣三内是完全的，这是真实而唯一的信仰，这是善和真理，那么它必然始终如此，除非善和真理是后来出现的，而天主的教义是靠添加得以完备的。我要说，它必然永远如此；但如果不是永远如此，那么现今也非如此，而是如同你们所设想的起初那样——我的意思是，现今不是圣三。然而，没有基督徒会容忍这样的异端；引入一个受造的圣三，并将它与受造物等同，这是属于希腊人的；因为他们容许缺陷和增加；但基督徒的信仰承认那有福的圣三是永恒不变的、完全的、始终如它曾所是，既不加添什么，也不归咎于它任何损失（因为这两种想法都是不虔敬的），因此它将圣三与一切受生之物分开，并守护天主性本身的合一，使之不可分割，并加以钦崇；它远避亚略派的亵渎，并且宣认和承认圣子永恒存在；因为祂是永恒的，如同圣父是永恒的，祂是圣父的永恒圣言——现在让我们再次回到这个主题。\par

19. 如果天主是、且被称为智慧的泉源和生命——正如祂藉耶肋米亚所说：\BibleQuote{他们离弃了我这活水的泉源；} \BibleRef{耶2:13}又说：\BibleQuote{自始就光荣崇高的宝座，是我们的圣所；以色列的希望——上主啊！凡离弃祢的，必受羞辱；离开我的，必在地上记录，因为他们离弃了上主，活水的泉源；} 又在巴路克书中记载：\BibleQuote{你们离弃了智慧的泉源，} \BibleRef{巴3:12}——这意味着生命和智慧并非外在于那泉源的本质，而是其固有的，从未曾不存在，而一直存在。现在，圣子就是这一切，祂说：\BibleQuote{我是生命，} \BibleRef{若14:6}又说：\BibleQuote{我智慧以明达为居所。} \BibleRef{箴8:12}那么，说\Quote{圣子曾经不存在}岂不是不信吗？因为这等同于说\Quote{泉源曾经枯干，缺少生命和智慧}。但若如此，它就不再是泉源；因为不从自身生出东西的，就不是泉源。这是何等荒谬！因为天主应许那些承行祂旨意的人将如同永不枯竭的泉源，祂藉依撒意亚先知说：\BibleQuote{上主必在干旱中满足你的心灵，使你的骨头强壮；你必像浇灌的园子，又像水流不绝的泉源。} \BibleRef{依58:11}然而这些人，当天主被称为且是智慧的泉源时，竟敢侮辱祂，说祂是不结果实、缺少自己本有智慧的。但他们的教义是虚假的；真理见证天主是祂本有智慧的永恒泉源；如果泉源是永恒的，智慧也必然永恒。因为在祂内万有受造，如达味在圣咏中说：\BibleQuote{你以智慧造了万有；} 撒罗满说：\BibleQuote{上主以智慧奠定大地，以明达建立诸天。} \BibleRef{箴3:19}而这智慧就是圣言，并且藉着祂，如若望所说：\BibleQuote{万物是藉着祂而造成的}，\BibleQuote{没有一样不是藉着祂而造成的。} 而这圣言就是基督；因为\BibleQuote{我们只有一个天主，就是圣父，万物都出于祂，我们也归于祂；只有一个主，就是耶稣基督，万物藉祂而有，我们也藉祂而有。} \BibleRef{格前8:6}如果万物是藉着祂而有的，祂自己便不应算在那\Quote{万物}之中。因为谁敢将万物藉以有的那一位，称为\Quote{万物}中的一位，那么他对于万物所从出的天主，也必然会有类似的臆想。但如果他因这不合适而退缩，并将天主排除在万物之外，那么他一贯地也应将那独生子排除在万物之外，因为祂属于圣父的性体。如果祂不是万物中的一位，那么关于祂说\Quote{祂曾不存在}，以及\Quote{祂在受生之前不存在}，便是罪过。这样的话可用于受造物；但至于圣子，祂与圣父一样，祂是圣父性体的本有后裔、圣言和智慧。因为这属于子对父的关系，这也显示父属于子；因此我们既不能说天主曾没有圣言，也不能说圣子曾不存在。因为如果子不是从父而来，何以为子？或者，如果圣言和智慧不是始终属于父，何以为圣言和智慧？\par

20. 那么，\Person{天主}何时曾经没有祂所固有的那位呢？或者，一个人怎么能将那固有的那位视为在本质上陌生而异己的呢？因为其他事物，按照受造物的本性，在本质上与创造者没有相似之处；而是外在于祂，是由\OriginalTerm{圣言}{Word}藉着祂的恩宠和旨意造成的，因此如果创造它们的\Person{天主}乐意，它们就可以停止存在；因为这就是受造物的本性。但是，至于\Person{圣父}本体所固有的那位（因为我们已经发现这就是\Person{圣子}），若在不信中竟敢说\Quote{这来自于虚无}，\Quote{在受生之前它不存在}，而是外来的，并且可能再次于某个时候停止存在，这是何等的大胆啊？只要一个人深思这个想法，他就会看出这异端如何损害了\Person{圣父}本体的完美与圆满；然而，如果他考虑到\Person{圣子}是\Person{圣父}的\OriginalTerm{肖像}{Image}、\OriginalTerm{光辉}{Radiance}、\OriginalTerm{言语}{Expression}和\OriginalTerm{真理}{Truth}，他就会更清楚地看到它的不恰当。因为，如果光存在，就同时有它的肖像，即光辉；如果\OriginalTerm{自立体}{Subsistence}存在，就有它完全的言语；如果\Person{圣父}存在，就有祂的真理（即\Person{圣子}）；那么让他们想想，那些以时间度量\OriginalTerm{天主性}{Godhead}的肖像和形式的人，陷入了何等深沉的不信中。因为如果\Person{圣子}在祂受生之前不存在，那么真理就不常在\Person{天主}内，这是有罪的言论；因为自从\Person{圣父}存在，在祂内就常有真理，即\Person{圣子}，祂说：\BibleQuote{我是真理}\BibleRef{若 14:6}。而自立体既然存在，自然立即就有它的言语和肖像；因为\Person{天主}的肖像不是从外部描绘的，而是\Person{天主}自己生出了它；祂在其中看到自己，就喜悦，正如\Person{圣子}自己所说：\BibleQuote{我原是祂的喜悦}\BibleRef{箴 8:30}。那么，\Person{圣父}何时不在自己的肖像中看见自己？或者祂何时没有喜悦，竟有人敢说，\Quote{肖像出于虚无}，以及\Quote{肖像被造之前圣父没有喜悦}？创造者和造物主又怎能在受造且受生的本体中看见自己呢？因为\Person{圣父}是怎样的，肖像也必须怎样。\par

21. 那么，让我们继续思考\Person{圣父}的属性，我们就会知道这肖像是否真是祂的。\Person{圣父}是永恒的、不死的、大能的、光、君王、主宰、\Person{天主}、\Person{主}、创造者和造物主。这些属性必须也在肖像内，才能真实地说，\BibleQuote{那看见了子的，就是看见了父}\BibleRef{若 14:9}。如果\Person{圣子}不是所有这些，而是像亚略派所认为的那样，是受造的且非永恒的，那么这就不是\Person{圣父}真实的肖像，除非他们确实不知羞耻，进而宣称，赋予\Person{圣子}的\OriginalTerm{肖像}{Image}称号，不是一个相似本体的标记，而仅仅是祂的名字。但是，另一方面，基督的仇敌们啊，这不是肖像，也不是言语。因为，那出于虚无的东西，与那使本来虚无者得以存在的祂，有何相似呢？或者，那非存在者，怎能像那存在者呢？它曾在某个时候不存在，并且被列于受造物之中，在这方面不及祂。然而，亚略派愿意祂成为这样的，于是为自己设计出这样的论证：\Quote{如果\Person{圣子}是\Person{圣父}的后裔和肖像，且在凡事上都与\Person{圣父}相像，那么必然的结论就是：正如祂受生，祂也生，祂也成了一个儿子的父亲。再者，从祂所生的，也照样生，如此无穷无尽；因为这样才使受生者与生育祂者相像。}这些天主的仇敌，确实是亵渎的制造者！他们不愿承认\Person{圣子}是\Person{圣父}的肖像，反而对\Person{圣父}本身持有物质的和属世的观念，将分割、流出和影响归于祂。如果\Person{天主}像人一样，那就让祂也像人一样成为父母，好使祂的\Person{圣子}成为另一位的父亲，如此一个接一个地延续，直到他们所想象的系列变成众多的神明。但如果\Person{天主}不像人——祂也确实不像——我们就不应将人的属性归于祂。因为禽兽和人类，在一位创造者开始他们之后，是按着相继的方式出生的；而儿子，既是由一位曾是儿子的父亲所生，也就相应地轮到他自己成为一个儿子的父亲，因为他从父亲那里继承了使他得以存在的方式。因此，在这样的例子中，严格说来，既没有父亲，也没有儿子；父亲和儿子并不保持各自的角色，因为儿子自己成为了父亲，他既是自己父亲的儿子，又是自己儿子的父亲。但在\OriginalTerm{天主性}{Godhead}内并非如此；因为\Person{天主}不像人；因为\Person{圣父}不是出自一位父亲；所以祂所生的不会是成为父亲的一位；\Person{圣子}也不是出自\Person{圣父}的流出，也不是从一位受生的父亲所生；因此祂的受生也不是为了生育。因此，这唯独属于\OriginalTerm{天主性}{Godhead}：\Person{圣父}恰是父亲，\Person{圣子}恰是儿子，而且在祂们内——也唯独在祂们内——\Person{圣父}永远是父，\Person{圣子}永远是子。\par

22. 因此，询问\Person{圣子}为何不生育一个儿子的人，必须探究为何\Person{圣父}没有父亲。但这两种假设都是不恰当的，且充满不虔。因为正如\Person{圣父}永远是父，绝不可能成为子，同样，\Person{圣子}永远是子，绝不可能成为父。因为正是在这一点上，更显出祂是\Person{圣父}的言语与肖像，保持祂之所是，无需改变，却从\Person{圣父}领受了同一的身份。如果\Person{圣父}改变，就让肖像也改变；因为肖像与光辉，相对于生它者的关系就是如此。但如果\Person{圣父}是不改变的，祂永远是祂所是，那么肖像也必然继续保持祂所是，且不会改变。如今祂是来自\Person{圣父}的子；因此祂绝不会变成与\Person{圣父}本体所固有的那位不同的存在。所以那些愚妄人，徒然想出这个异议，企图将肖像与\Person{圣父}分开，好将\Person{圣子}与受造物等同。\par

% structure: p0032 source=h2 target=subsection reason=relative-heading-rank; base=h2

\subsection{第七章 对前述证明的异议。在圣子的受生中，\Person{天主}是创造了已存在者，还是创造了尚未存在者？}

22（\Emph{续}）。按照优西比乌的教导，将祂列在这些受造物之中，视祂如同那些借着祂而有的受造物，亚略及其同伙便背离了真理；他们在这异端开始时，惯于四处游说他们所拼凑的不实之言。不，直到如今，他们中有些人，当他们在集市上碰到孩童，就问他们，不依据圣经，而是像这样，仿佛满心都是\BibleQuote{心里所充满的}\BibleRef{玛 12:34}；—— \Quote{那自有者，是使那非有者从〔非存在〕而成，还是从已有者而成呢？因此，祂是创造子时，子是存在的，还是不存在的呢？} 还有，\Quote{非受生者是一位还是两位？} 以及，\Quote{祂有自由意志，却因本性可变而不按自己的选择改变吗？因为祂不像石头那样自己保持不动。} 然后他们转向愚昧的妇女，以妇人之言对她们说：\Quote{你生子之前有儿子吗？既然你没有，那么生子之前，圣子也没有。} 这些无耻之人以这样的言语戏耍狂欢，将天主比作人，假装是基督徒，却将天主的光荣\Quote{变成为仿佛可朽坏的人像}。\par

23. 如此愚昧迟钝的言论本不配任何答复；然而，为免他们的异端看似有所依据，即便要绕路，在此驳斥他们也是合宜的，尤其是为了那些容易被他们欺骗的愚昧妇女。当他们这样说时，他们本应问问建筑师，不靠材料能否建造；若不能，是否由此推得天主不靠材料就不能创造宇宙。或者他们应问每一个人，人能否没有场所；若不能，是否由此推得天主也在场所之内，好使他们在听众面前也自觉羞愧。再者，当听说天主有子时，他们以自身的类比来否定子；而当听说天主创造和制造时，他们却不再提出属人的观念，这是为何？在创造的事上，他们也当持守同样的想法，为天主提供材料，从而否认祂是创造者，直到他们沦于与摩尼教徒一样卑下。然而，若仅仅天主的观念便超越这些思想，人一听见就相信并知道祂是存在着，不是像我们这样的存在，而是作为天主的存在，祂创造也不是像人创造，而是作为天主创造，那么显然，祂生育也不是像人生育，而是作为天主生育。因为天主并不以人为范本；反倒是我们人，因为天主是真正地、唯一真确地是祂圣子的父，我们才也被称为自己子女的父亲；因为\BibleQuote{上天下地的一切父性都是由祂得名}\BibleRef{弗 3:15}。他们的论点，若不细察，看似有理；但若有人用理性加以细察，便会发现它们招来许多讥笑与嘲弄。\par

24. 首先，论到他们这样的第一个问题，是何等愚钝模糊！他们不说明所问的是谁，使人无法回答，只是抽象地说\Quote{那自有者}，\Quote{那非有者}。亚略派啊，那么谁是\Quote{有}，什么又是\Quote{非有}？或者说，谁\Quote{存在}，谁\Quote{不存在}？什么被称为\Quote{存在}，什么被称为\Quote{不存在}？因为那自有者，能创造那非已有的，那已有的，以及那先前已有的。例如，木匠、金匠、陶匠，各按其手艺，在预先存在的材料上加工，制造他喜欢的器皿；而万有的天主自己，取了已有的、已受造的尘土，塑造了人；然而，那尘土本身，原先并不存在，祂在某一时刻借着自己的圣言创造了它。那么，如果这就是他们问题的含义，一方面受造物在其受造之前显然不存在，另一方面，人则加工现有的材料；因此他们的推理不合逻辑，因为正如这些例子所示，\Quote{已有的}成为，\Quote{非已有的}也成为。但如果他们谈论的是天主和祂的圣言，就让他们把问题说全了再问：那\Quote{自有者}天主，曾没有\OriginalTerm{圣言}{Reason}吗？再者，祂既是光，曾缺少光芒吗？或者，祂向来是圣言的父吗？或者再这样问：那\Quote{自有者}圣父，是创造了那\Quote{非有者}圣言，还是祂一直有自己的圣言，作为祂实体的特生之子？这将向他们表明，他们不过是擅自并斗胆对天主和那出自天主者使用诡辩。真的，谁能忍受听他们说天主曾没有圣言呢？这正是他们再次陷入的谬误，尽管他们徒劳地设法逃避并用诡辩隐藏它。不，考虑到已经引证了大量证据反对他们，人们甚至不愿听到他们争论说天主不是始终为父，而是后来才成为父亲（这对他们的幻想是必要的，即祂的圣言曾不存在）；而若望还说过：\BibleQuote{圣言在起初就有}\BibleRef{若 1:1}，保禄也写道：\BibleQuote{祂是天主光荣的辉耀}\BibleRef{希 1:3}，以及：\BibleQuote{祂是在万有之上，永远可赞颂的天主。阿们}\BibleRef{罗 9:5}。\par

25. 他们最好闭口不言；但既然事实并非如此，那就只能以大胆的反驳来回应他们那无耻的质问。或许，当他们看到同样困扰自己的荒谬之处时，便会停止与真理为敌。在我们多次祈求天主垂怜之后，或许可以转而反问他们：那自有永有的天主，难道曾是如此而成的，而祂原不存在吗？或者，祂在自己存在之先就存在了吗？那自有者，祂是自造的，还是出自虚无，从前本为虚无，忽而自己显现出来呢？这种追问既不庄重，又极为亵圣，却与他们的质问如出一辙；因为他们给出的回答本身就充满了不虔敬。然而，若以这种方式追问天主是亵圣且完全不虔敬的，那么对祂的圣言提出类似的追问也同样是亵圣的。然而，为了揭露这种愚昧迟钝的问题，有必要如此回答：——自有永有的天主，是永恒存在的；因此，圣父既然永在，祂的光辉——也就是祂的圣言——也永在。再者，自有永有的天主，从自己而有了也存在的圣言；圣言并非在先前不存在之后才被加添，圣父也从未没有过圣言。因为，对圣子的这种攻击，使亵渎反落于圣父；仿佛祂从自身之外为自己设计了一位智慧、圣言和圣子；因为无论你使用其中哪一个称号，你所指的都是出自圣父的那位子嗣，如前所述。因此，他们的异议站不住脚；这很自然，因为他们否认\OriginalTerm{圣言}{Logos}，结果提出的问题也不合理性。那么，就像有人看见太阳，然后追问其光芒，说：\Quote{自有者造了那曾有者，还是那曾无者，}他会被视为不明事理，彻底昏聩，因为他竟以为那从光而出的光芒是在光之外，并且追问它是何时、何处、是否被造；同样地，对圣子和圣父如此推测与追问，乃是更大的疯狂，因为这是将圣父的圣言设想为在祂之外，并轻率地称那本体的子嗣为受造物，同时宣称：\Quote{祂在受生之前并不存在。}不，让他们在问题之外再接受这一回答——自有者圣父，造了自有者圣子，因为\BibleQuote{圣言成了血肉}；\BibleRef{若 1:14}；并且，祂本是天主之子，在时间的终末，也使祂成为人子，除非他们真的追随\OriginalTerm{萨摩撒塔派}{Samosatene}，坚称祂在成为人之前根本不存在。\par

26. 我们以上对第一个问题的回答已经足够。现在，\OriginalTerm{亚略派}{Arians}啊，你们也该回想自己的话，告知我们：那位自有者，是需要一位曾无者来塑造宇宙，还是需要一位有者？你们曾说过，祂从虚无中为自己造了圣子，作为用以创造宇宙的工具。那么，哪一方更优越：是需要者，还是供给需求者？或者，双方不是互补对方的不足吗？你们反倒证明了创造者的软弱，因为祂若没有能力凭自己创造宇宙，却从自身之外为自己预备了工具，就像木匠或造船者，没有锛凿和锯子就什么也做不成！还有比这更不虔敬的吗？既然前面的论述已足以表明他们的教义纯属幻想，又何必再细述其可憎之处呢？\par

% structure: p0038 source=h2 target=subsection reason=relative-heading-rank; base=h2

\subsection{第八章 继续反对意见。我们能否以人类子孙——通常晚于其父母出生——的类比来裁定这个问题？不可，因为类比的要点在于\OriginalTerm{同质性}{connaturality}。时间并不包含在\Quote{子}的观念中，而是\OriginalTerm{附加}{adventitious}于它的，并且不适用于天主，因为祂没有部分和情欲。\Quote{圣言}和\Quote{智慧}的称号护卫我们的思想，使关于天主及其圣子的观念不陷入这一误解。天主并非作为创造者而在永恒中\OriginalTerm{潜在}{in posse}为父，因为创造并不像生那样涉及天主的本质。}

26. （\Emph{续}）他们向愚蠢妇女提出的另一个极其简单而愚蠢的问题，也无需回答；或只需重复已有的答案，即用人的本性来衡量天主的生是不合适的。不过，为了让他们像前面那样自我评判，最好在同样的基础上回应他们，如下：——显然，如果他们向父母询问其儿子，就让他们想想那所生的孩子是从何处来的。就算父母在生育之前没有儿子，然而拥有之后，他们拥有他，不是如同外来的或陌生的，而是从自己而来的，属于他（父亲）的本质和他精确的肖像，以致前者在后者中被看见，后者在前者中被瞻仰。那么，如果他们从人的例子认定生育包含时间，为什么不从同样的例子推断出生育也包含\Quote{本然}与\Quote{专属}呢？反而如同毒蛇从地里只吸取那些化为毒液的东西。那些询问父母的人说：\Quote{你在生他之前有儿子吗？}他们应该再加上：\Quote{如果你有儿子，你是否如同购买房屋或其他财产那样，从外面买来的呢？}接着你就会被回答：\Quote{他不是从外面来的，而是从我自身。因为那些从外面来的东西是财产，能从一人转移到另一人；但我的儿子却出自于我，属于我的本质，与我的本质相似，不是从别人那里变成我的，而是由我所生；因此，我也完全在他内，同时我仍然保持我之所是。}确实如此；尽管父母在时间上是区别的，因为他是人，他自己也是时间内形成的，然而他本也能拥有与他永恒共存的孩子，只是他的本性成为限制，使之不可能。因为\Person{肋未}就已经存在在他曾祖父的腰中，早于他自己或他祖父的实际出生。因此，当那人达到本性赋予能力的年龄时，立即，借着不受限制的本性，他就从自己成为儿子的父亲。\par

27. 因此，如果他们询问父母关于子女的事，得到的答复是：本性所生的子女不是从外而来的，而是出自父母，那么让他们同样承认关于天主的圣言，祂完全出自圣父。倘若他们提出时间的问题，就让他们说出什么能限制天主——因为必须在他们发出嘲笑的同一根据上证明他们的不虔敬——让他们告诉我们，有什么能阻止天主永远是圣子的父；因为那受生的必然出自其父，这是无可否认的。再者，他们若将这些事归于天主，就必审断自己；如果像他们就时间问题询问妇女那样，他们向太阳询问其光芒，向源泉询问其流出。他们会发现，这些虽为产物，却永远与产生它们之物同在。如果这样的父母与子女共有本性和延续，那么他们若认为天主比受造之物更低劣，为什么不公开说出他们自己的不虔敬呢？但如果他们不敢公开这样说，并且承认圣子并非来自外边，而是由父所生的本性子嗣，承认没有任何事物能限制天主（因为祂不像人，而是超越太阳，确切地说是太阳的天主），那么结论便是：圣言来自祂，且永远与祂同在，圣父也借着祂使万物从无到有。因此，圣子并非由虚无而来，而是永恒的，且出于圣父，这从事实的本性来看也是确实的；而异端者向父母提出的问题暴露了他们的悖谬；因为他们承认了本性的问题，如今在时间的问题上却被羞辱。\par

28. 如上所述，如今我们重申：天主的生育不应与人的本性相比拟，圣子也不可被视为天主的一部分，生育也不可暗示任何情欲；天主不像人；因为人在生育时是动情的，其本性是可传递的，且因其软弱而等待时间。但天主不能如此；因为祂并非由部分组合而成，而是无情欲且单纯的，祂作为圣父，是以无情欲且不可分的方式生圣子。这一点也有圣经的有力证据和证实。因为天主的圣言是祂的圣子，而圣子是圣父的圣言与智慧；圣言与智慧既非受造物，也非其所出言的那位的一部分，亦非动情受生的子嗣。因此，圣经将两个称号结合起来，谈到\Quote{子}，为的是宣扬祂本性的、出自祂本体的真实子嗣；另一方面，为了让任何人都不会以人的方式理解这子嗣，在表明祂的本体时，圣经又称祂为圣言、智慧和光耀，教导我们这生育是无情欲的、永恒的，且与天主相称。那么，圣言、智慧、光耀，是圣父的哪种情欲呢？或是圣父的哪个部分呢？这点甚至可以给这些愚昧的人留下印象；因为就如他们曾就天主的圣子询问妇女，现在就让他们就圣言询问人，他们就会发现，人所发出的言既非他们的情欲，亦非他们心灵的某个部分。但倘若人就是如此，就是那有情欲且可分割的人，为何他们对于非物质且不可分割的天主，要揣测情欲和部分，以便假借敬畏之名否认圣子的真实而本性的生育呢？上面已经充分说明，由天主所出的子嗣并非情欲；如今又特别证明，圣言并非按情欲受生。对于智慧也可说同样的话；天主不似人；他们也不可在这里以人的方式思考祂。因为，人虽能领受智慧，天主却一无所缺，祂自己是其智慧的父，凡领受这智慧者，就被赋予智慧之名。而这智慧也不是情欲，不是部分，而是圣父所特有的子嗣。因此，祂永远是父，为父的特征对天主来说不是外加的，免得祂显得可变；因为如果祂是父是好的，却并非永远是父，那么善在祂内就不是永恒的了。\par

29. 但是，他们说，请注意，天主一向是造化者，造化之能并非外加于祂的；那么，是否因为祂是万有的创造者，祂的工程也就是永恒的，并且说它们在受造之前并不存在也是邪恶的呢？这些亚略派人士真是愚昧；因为在子和工程之间有什么相似处，竟让他们将父亲的功能与造化者的功能相提并论呢？既然已经指出了\OriginalTerm{子嗣}{offspring}和工程之间的区别，他们怎么还是如此缺乏教导呢？那么，容再说一次，工程是外于性体的，而子却是\OriginalTerm{本质}{essence}所固有的子嗣；所以工程不必常存，因为工人随己意创造它；但子嗣并非随意志而定，而是本质所固有的。一个人可以成为且被称为造化者，尽管工程尚未存在；但除非有子存在，他不能被称为父亲，也不能是父亲。倘若他们好奇地追问，为什么天主虽有创造之能，却不时刻创造（虽然这也是疯子的妄念，因为\BibleQuote{谁知道主的心意？谁作过祂的顾问？}或者\BibleQuote{制造品岂能对制造者说：你为什么这样制作了我？}然而，为了不忽略任何微弱的论证），他们必须被告知，尽管天主一向有创造之能，但受造之物却没有永恒存在的能力。因为它们是出于虚无，所以在它们受造之前并不存在；而在受造之前并不存在的事物，怎能与永恒长存的天主共存呢？因此，天主着眼于它们的好处，当祂看到它们受造后能够存立时，便造了它们。正如，虽然祂能从起初，在亚当、诺厄或梅瑟的时代，就派遣祂的\OriginalTerm{圣言}{Word}，但祂直到时代的末期才派遣了祂（因为祂看这事为整个受造界是好的），同样，受造之物也是当祂愿意，且看为它们好时才造的。但子并非工程，而是圣父的子嗣所固有的，时常存在；因为圣父如何时常存在，祂本质所固有的也必时常存在；这就是祂的圣言和祂的\OriginalTerm{智慧}{Wisdom}。受造物不存在并不贬损造化者；因为祂有权能当祂愿意时创造它们；但子嗣若不与圣父永远同在，便是对祂本质之完美的贬损。因此，祂的工程是当祂愿意时，藉由祂的圣言而造；但子却是圣父本质所时常固有的子嗣。\par

% structure: p0043 source=h2 target=subsection reason=relative-heading-rank; base=h2

\subsection{第九章 诘难续。非受生者是一个还是两个？亚略派使用非圣经词语的不一致；有必要界定其含义。该词的不同含义。若其意为\Quote{无父}，则只有一位非受生者；若意为\Quote{无始或非受造}，则有两位。\Person{阿斯德利乌}的不一致。\Quote{非受生者}是天主的称号，并非与子相对，而是与受造物相对，如同\Quote{全能者}或\Quote{万军之主}。\Quote{父}是更真实的称号，因为它不但是圣经的，还蕴含着一位子，以及我们被收为子嗣。}

30. 这番思考鼓舞了信者，使异端者沮丧，他们看出自己的异端就此被推翻了。此外，他们进一步问道：\Quote{非受生者是一个还是两个？}这显明他们的观点是何等虚妄、奸诈而充满诡诈。他们问这个，并非为了圣父的荣耀，而是为了贬损圣言。因此，若有人不知其诡计，回答说：\Quote{非受生者是一位}，他们立刻吐出自己的毒液说：\Quote{那么子就在受造之物中了}，并且说得好：\Quote{在祂受生之前，祂并不存在。}他们制造各种纷扰和混乱，只要能离间子与圣父，并将万物的造化者归入祂的工程之中。首先，他们可能因此而被定罪，他们指责\Place{尼西亞}的主教们使用了非圣经的用语，虽然这些用语并不有害，反而颠覆了他们的不虔敬，他们自己却陷入了同样的错误，即使用非圣经的词语，而且以之侮辱主，\BibleQuote{他们不知道自己所说的是什么，所论断的是什么}\BibleRef{弟前 1:7}。例如，让他们去问那些曾经教导他们的希腊人（因为这是他们发明的词语，并非圣经的），当他们在其各种含义上受到教导，就会发现自己甚至不能就所从事的主题正确地提问。因为他们使我知道，\OriginalTerm{非受生者}{Unoriginate}的意思可以指尚未形成，但可能形成的事物，例如尚未成为，但能成为器皿的木头；又可以指既非现在有，也永远不能形成的事物，例如四角三角形和偶数奇数。因为三角形既非现在有，也永远不能成为四角形；偶数既非现在有，也永远不能成为奇数。再者，\OriginalTerm{非受生者}{Unoriginate}的意思是指：存在，但不是由任何事物所形成，也根本没有父亲的事物。此外，无原则的诡辩家\Person{阿斯德利乌}，也就是此异端的支持者，在自己的著作中加上说：凡非受造的，而是永恒存在的，就是\OriginalTerm{非受生者}{Unoriginate}。那么，当他们提问时，理应先说明他们是在何种意义上使用\OriginalTerm{非受生者}{Unoriginate}一词，然后被问者才能切中要点地回答。\par

31. 但如果他们仍满足于仅仅提问：\Quote{\OriginalTerm{非受造者}{Unoriginate}是一位还是两位？}那么首先必须告知他们，如同教育不足者，许多事物是如此，而无一是如此：许多事物，即有受生之可能者；而无，即不可能者，如前所述。但如果他们按\Person{阿斯泰略}所规定的那样提问，似乎\Quote{不是受造物，且永恒常在者}就是非受造者，那么必须不断告知他们，在这种意义上，圣子与圣父一样，都必须被称为非受造者。因为祂既不属于受造物之列，也不是受造物，而是永远与父同在，正如已证明的，尽管他们为了反对主而反复多变，提出\Quote{祂出自虚无}与\Quote{祂在被生之前不存在}。当他们处处失败后，转而以问题的另一意义提问，\Quote{存在者，非由任何所生，亦无父亲}，我们将告诉他们，在这种意义上的非受造者只有一位，即圣父；他们的问题一无所得。因为说天主在这种意义上是非受造者，并不表明圣子是受造物，从以上证明已显明，圣言怎样，那生祂者也怎样。因此，如果天主是非受造者，祂的肖像就不是受造的，而是\OriginalTerm{所生者}{Offspring}，即祂的圣言与祂的智慧。因为受造者与非受造者有何相似呢？（我们不可厌倦于重复；）因为如果他们硬要说一者与另一者相像，以致看见这一位就等于看见那一位，他们几乎等于说非受造者是受造物的肖像；其结果是整个论题的混乱，将受造物与非受造者等同，并以受造物衡量非受造者，从而否认非受造者；这一切都是为了将圣子归入他们之列。\par

32. 然而，我想若他们跟随诡辩家\Person{阿斯泰略}，连他们也不愿走到如此极端。因为\Person{阿斯泰略}虽热心拥护\OriginalTerm{亚略异端}{Arian heresy}，并坚持非受造者是一位，但他却与他们相矛盾地说，天主的智慧也是非受造的、无始的。以下是他著作中的一段话：\Quote{真福\Person{保禄}并未说…而是，不用冠词，说\BibleQuote{天主的德能和天主的智慧}\BibleRef{格前 1:24}；因而宣讲说…}稍后他又说：\Quote{然而，祂的永恒德能与智慧，真理论证其为无始且非受造的；这必定是一位。}因为，尽管他误解了宗徒的话，认为有两个智慧；然而，通过仍说有一个智慧与祂共存，他宣告非受造者并非仅是单独一位，而是有另一位非受造者与祂同在。因为共存者，不是与自己共存，而是与另一者共存。因此，如果他们赞同\Person{阿斯泰略}，就永远不要再问：\Quote{非受造者是一位还是两位？}否则他们将不得不在这一点上与他争辩；反之，如果他们不同意他，就不要依赖他的论文，免得\BibleQuote{你们若彼此相咬相吞，你们要小心，免得你们彼此消灭。}\BibleRef{迦 5:15}关于他们的无知，就说到此；但对于他们的狡诈，谁能说得尽呢？谁不憎恨他们，因他们被如此疯狂所附？当他们不再被允许说\Quote{出自虚无}和\Quote{祂在被生之前不存在}时，他们便找到\Quote{非受造者}这个词，以便对纯朴的人说圣子是\Quote{受造的}时，能暗示出同样的短语\Quote{出自虚无}和\Quote{祂曾不存在}；因为在这些短语中，暗示了受造物与受造者。\par

33. 如果他们对自己的立场有信心，就该坚持己见，而非如此多变地改换；但他们不会如此，因为他们认为只要用\Quote{非受造者}这个词作幌子来掩盖其异端，就容易成功。然而，归根结底，这个词并不是与圣子相对照，无论他们如何叫嚣，而是与受造物相对照；同样的用法也可见于\Quote{全能者}与\Quote{万军之主}等词语。因为如果我们说，圣父藉着圣言对万有拥有权能和控制，而圣子掌管圣父的王国，并作为祂的圣言和圣父的肖像，对万有拥有权能，那么很显然，圣子既不被包括在那\Quote{万有}之列，天主也不是因圣子而被称为全能者和主，而是因那些藉圣子而存在、并且祂藉圣言对其行使权能和控制的事物。因此，非受造者这一名称，不是相对于圣子，而是相对于那些藉圣子而存在的事物而言的。而且这样说妙极了：因为天主不像受造物，而是他们的创造者和化工者，藉着圣子。正如\Quote{非受造者}一词是相对于受造物而言，\Quote{圣父}一词则是指示圣子。那称呼天主为化工者、创造者及非受造者的人，是关注并理解受造与形成之物；而那称呼天主为父的人，则由此构想并默观圣子。因此，人们会惊异于他们不信中又加上的顽固：虽然\Quote{非受造者}一词有上述美好的意义，且可被虔诚地使用，他们却在异端中，用它来羞辱圣子，没有读过：\BibleQuote{谁尊敬子，就是尊敬父；谁不尊敬子，就是不尊敬父。}\BibleRef{若 5:23} 如果他们真有一丝关切，要恭敬言谈及给予圣父应得的尊荣，那么他们应当——而且这更好、更高——承认并称呼天主为父，而不是给祂这个名字。因为，在称呼天主为非受造者时，他们正如我之前说过的，是从祂的受造物称呼祂，并且仅以祂为化工者与创造者，臆想着由此可以表示圣言按他们自己的意愿是一个受造物。但是，那称呼天主为父的人，是从圣子指示祂，深知如果有圣子，那么必然藉着那圣子，一切受造物都得以受造。而他们，当他们称呼天主为非受造者时，只是从祂的受造物来命名祂，并不比希腊人更认识圣子；但那称呼天主为父的人，是从圣言命名祂；并且认识圣言，就承认祂是万有的化工者，并明白万有是藉着祂而造成的。\par

34. 因此，由子而指明天主并称祂为父，比仅从祂的工程命名并称祂为\OriginalTerm{非受生}{Unoriginate}，更为虔敬且更为确切。因为后者这一称号，如我已说过的，无非是指一切工程——无论是各别的还是全体的——它们皆因天主的旨意、藉着\OriginalTerm{圣言}{Logos}而得以生成；然而父的称号则唯独源于子才有其意义与关联。而且，圣言既超越\OriginalTerm{受生之物}{originated}，那么称天主为父便远远超越称祂为\OriginalTerm{非受生}{Unoriginate}。因为后者不合圣经，且引人疑虑，因其含有多重意义；以致当人就此被问及时，心思便游移于诸多观念；但父这个词简明、合于圣经，更为确切，且只涵括子。而\Quote{\OriginalTerm{非受生}{Unoriginate}}乃是希腊人的用语，他们不认识子；但\Quote{父}却由我们的主所承认与恩赐。因为祂深知自己是谁的子，便说：\BibleQuote{我在父内，父在我内}；又说：\BibleQuote{看见我的，就看见了父}；以及\BibleQuote{我与父原为一体}。但从未见祂称父为\OriginalTerm{非受生}{Unoriginate}。再者，当祂教导我们祈祷时，祂没有说：\BibleQuote{当你们祈祷时，要说：\InnerQuote{非受生的天主啊}}，而是说：\BibleQuote{当你们祈祷时，要说：\InnerQuote{我们的父，在天上者}}\BibleRef{路 11:2}。而祂也愿意我们信仰的纲领具有同样的指向，吩咐我们受洗时，不是归于\OriginalTerm{非受生}{Unoriginate}者和\OriginalTerm{受生者}{originate}的名，也不是归于创造者和受造物的名，而是归于父、子、圣神的名。因为藉着这样的入门圣事，我们虽被列于工程之中，却得成为儿子；我们既使用父的名，便由此名也在父自身内承认圣言。因此，他们关于\OriginalTerm{非受生}{Unoriginate}这一术语的论证乃是虚妄的，如今已得证明，不过是一场幻梦罢了。\par

% structure: p0049 source=h2 target=subsection reason=relative-heading-rank; base=h2

\subsection{第十章 异议续论。圣言如何拥有自由意志却非\OriginalTerm{可变的}{alterable}；祂因是父的肖像而\OriginalTerm{不可变的}{unalterable}，由经文可证。}

35. 至于他们所提圣言是否\OriginalTerm{可变的}{alterable}这一问题，实属多余，无需细究；只需简单写下他们的话，以揭露其大胆的异端便已足够。他们何等轻率，从以下问题便可见一斑：\Quote{祂是否有自由意志，抑或没有？祂是否凭自由意志选择行善，并且由于本性\OriginalTerm{可变的}{alterable}，祂若愿意，便能改变？或者，祂如同木石，没有自由选择，不能随意移动、倾向这边或那边？}这种说法和想法正与他们的异端相合；因为一旦他们为自己构想出一位从虚无而来的天主和一个受造的圣子，他们当然会采用这类适用于受造物的词汇。然而，当他们在与教会人士辩论时，从后者听到那真实而唯一的父的圣言，却仍胆敢如此议论祂，他们的教义岂非变成了最为可憎的吗？这岂不足以让人仅仅听见就心神迷乱，虽无力答辩，却要掩耳不闻，惊讶于他们所言之新奇——纵然提及便已是亵渎？因为圣言若是\OriginalTerm{可变的}{alterable}且变化不定，祂将止于何处，其进展的终点又是什么？\OriginalTerm{可变的}{alterable}者如何可能相似于\OriginalTerm{不可变的}{Unalterable}者？见到\OriginalTerm{可变的}{alterable}者，怎可视为已见到\OriginalTerm{不可变的}{Unalterable}者？祂必须达到何种状态，我们才能在祂内看见父？显然，我们不会时时都在子内看见父，因为子既不断变化，其本性又属转变。而父却是\OriginalTerm{不可变的}{unalterable}、不能改变的，常处于同一状态和同一本身；但如果像他们所持守的，子是\OriginalTerm{可变的}{alterable}，并非始终如一，其本性又不断变化，那么祂既无父的\OriginalTerm{不可变性}{unalterableness}之肖似，怎能成为父的肖像？祂若意向不定，又如何真的在父内？甚或，正因祂是\OriginalTerm{可变的}{alterable}，且日日进深，祂尚不圆满。但亚略派的这等狂妄，必须摒弃，让真理发光，显明他们的愚拙。与天主同等的，岂不必是圆满的？与父原为一体、又是属于祂本质的亲子的，岂不必是\OriginalTerm{不可变的}{unalterable}？父的本质既是\OriginalTerm{不可变的}{unalterable}，从祂所出的亲子也必是\OriginalTerm{不可变的}{unalterable}。他们若谤毁地将可变性归咎于圣言，就让他们明白，自己的理智何等危殆；因为凭果实即可认识树木。这正是为什么看见子的，就看见了父；认识子，也就是认识父。\par

36. 因此，不可改变的天主的肖像必须是不变的；因为\BibleQuote{耶稣基督昨天、今天、直到永远，常是一样。}\BibleRef{希 13:8} 而圣咏中\Person{达味}论到祂说：\BibleQuote{主，在起初你奠定了下地，上天也是你手的功化。天地必要毁灭，而你永远存在，万物必要如同衣裳一样衰老。\newline{} 你有如更换衣服，你更换它们，它们就被改变；但是你常存在，你的岁月无尽期。} 主自己也藉先知论到自己说：\Quote{现在你们应认清，我就是那一位，}\Quote{我决不改变。} 当然可以说，这话所指的乃是父；但这亦适用于子如此说，尤其因为当祂成为人时，向那些因祂的肉身而认为祂已改变、与从前不同的人，显示了祂的同一性和不变性。圣徒们，或更好说是主，比不虔者的邪僻更可信。因为圣经，在上面所引的 \WorkTitle{圣咏集} 经文章节中，藉着天地的名义，表明一切受造和生成之物的本性是可改变和变化的，却将子排除在这些之外，以此向我们显示，祂绝不是受造之物；反而教导说，祂改变其他一切，而自身并不改变，说：\BibleQuote{你常存在，你的岁月无尽期。}\BibleRef{希 1:12} 这是合理的；因为受造之物是从虚无中来的，在他们受造之前并不存在，因为事实上，他们是从不存在而成为存在的，具有可变的本性；但子，出于父，并属于父的本体，像父一样是不可改变、不可转变的。因为如果说从那不可改变的实体中，生出一个可改变的圣言、一个可变易的智慧，那是亵渎的。如果圣言是可改变的，祂怎能再是圣言呢？或者，可改变的怎能是智慧呢？除非如他们可能主张的，就像在实体中的依附体一样，也就是说，如同在某一特定实体中，偶然存在着一种被称为圣言、子和智慧的德能与习性，可以被取去和增加。因为他们经常表达这种观点，但这并非基督徒的信仰；因为它并未宣认祂真是天主的圣言和圣子，也未宣认祂是那真正的智慧。因为那变化、变迁，不持守同一状态者，怎能是真实的呢？然而主说：\BibleQuote{我就是真理。}\BibleRef{若 14:6} 那么，如果主自己这样论到自己，并宣示祂的不变性，而圣徒们又学习和见证这一点，并且我们对天主的观念也承认这是虔诚的，那么这些不虔的人从何处得出这种新奇说法呢？他们从自己的心中，如同从败坏之穴中呕吐出来。\par

% structure: p0052 source=h2 target=subsection reason=relative-heading-rank; base=h2

\subsection{第十一章 经文讲解；首先是，\WorkTitle{斐理伯书} 2:9,10。是否\Quote{因此天主极其举扬祂}等语证明道德上的考验与进步。首先从\Quote{子}一词的意义予以反驳，此意义与该解释不符。接着，审查该经文。\Quote{极其举扬}、\Quote{赐给}和\Quote{因此}在教会内的意义；即，是就我们的主的人性而言。次要意义；即，暗示圣言透过复活而被\Quote{举扬}，其意义与圣经论及祂在降生成人中的\Quote{下降}相同；此语如何无损于圣言的本性。}

37. 但是，由于他们援引神圣的谕令，并强行按自己的私意曲解，我们必须进行回应，以澄清这些经文，指出它们合乎正统的意义，并揭露对手的错误。他们说，宗徒写道：\BibleQuote{因此，天主极其举扬他，赐给了他一个名字，超越其他所有的名字，好使上天、地上和地下的一切，一听到耶稣的名字，无不屈膝叩拜。}\BibleRef{斐 2:9} 达味也说：\BibleQuote{因此，天主，你的天主，以欢愉的油傅了你，胜过你的同伴。} 然后他们自以为尖锐地主张：\Quote{如果他被举扬并接受了恩宠，是\InnerQuote{因此}，并且是\InnerQuote{因此}被傅油，那么他获得了对他志向的奖赏；但既然他是因志向而行事，他就完全具有可改变的本性。} 这就是\Person{欧瑟比}和\Person{亚略}胆敢说的，甚至写下来的；而他们的党羽竟不避讳在市集中公开谈论，没有看到他们所用的论据何等狂妄。因为如果他拥有的一切是作为他志向的奖赏，假如他不曾需要，并且没有善工可展示，就不会获得；那么从德行和提升中获得了这些，就有理由\Quote{因此}被称为子和天主，而不是本体的子。因为那按本性从另一位而出者，是真正的后裔，如同\Person{依撒格}之于\Person{亚巴郎}，\Person{若瑟}之于\Person{雅各伯}，以及光辉之于太阳；至于那些因德行和恩宠而称为子的人，他们取代本性而拥有的是所获得的恩宠，并且他们本身是赠与之外的另一回事；如同那些因分享而领受圣神的人，经上论到他们说：\BibleQuote{我生育并显扬了子女，他们却背叛了我。} 当然，既然他们不是本性的子女，因此，当他们改变时，圣神就被取去，他们被剥夺了继承权；而当他们悔改时，那位起初赐予他们恩宠的天主，将再次接纳他们，赐予光明，并再称他们为子女。\par

38. 但如果他们也这样论救主，那么祂就既非\OriginalTerm{真天主}{very God}，也非\OriginalTerm{真子}{very Son}，既不肖似圣父，也绝无天主作为其\OriginalTerm{本质}{essence}存在之父，而只是凭着施予祂的纯粹\OriginalTerm{恩宠}{grace}，并有一个按相似于所有受造物的、作为其本质存在之创造者。若果真如他们所主张的，那么就更明显祂并非从一开始就拥有\Quote{子}的名号，倘若这名号是因着祂成为人并取了\OriginalTerm{仆人的形象}{the form of the servant}时所完成的工作及所取得的进步而获得的赏报；而是等到后来，在祂\BibleQuote{听命至死}之后，正如经文所说，祂\BibleQuote{被高举}，并领受那\Quote{名号}作为恩宠，\BibleQuote{致使上天、地上和地下的一切，一听到耶稣的名字，无不屈膝叩拜}\BibleRef{斐 2:8}。那么，若祂是在成为人时才被高举、才开始受崇拜、才被称为子，在此之前祂是什么呢？因为若按他们的错谬，祂是在成为人时才被高举并被称为子，那么祂似乎并未提升血肉，反而自己藉着血肉被提升。那么在此之前祂又是什么呢？我们必须再次追问他们，好使人明白他们的不虔敬教义会得出什么结论。因为如果主是\OriginalTerm{天主}{God}、是\OriginalTerm{子}{Son}、是\OriginalTerm{圣言}{Word}，但在祂成为人之前并非所有这些，那么要么祂是这些之外的某者，后来因着自己的德行而分享了这些，如我们所说；要么他们必须采纳另一种说法（愿这说法落在他们自己头上！），即祂在那一刻之前并不存在，而全然是本性上的人，别无其他。但这并非教会的主张，而是\OriginalTerm{萨莫撒塔派}{the Samosatene}和现今犹太人的主张。那么，如果他们像犹太人一样思想，为什么不也和他们一同受割损，反而装作基督徒，实际上却是基督徒的仇敌呢？因为如果祂并不存在，或虽存在但后来才被提升，那么万物怎会藉着祂而受造？如果祂并非完美的，圣父又怎会在祂内喜悦呢\BibleRef{箴 8:30}？而另一方面，如果祂现在才被提升，那么祂先前如何在圣父面前欢跃呢？又如果祂是在死后才接受崇拜，那么亚巴郎如何在帐幕中向祂敬拜，梅瑟又如何向在荆棘中的祂敬拜呢？而且，如达尼尔所见，有千万的臣仆服事祂？如果如他们所说，祂现在才被提升，那么圣子自己怎会提到祂在世界之前和超越世界的荣耀，说：\BibleQuote{父啊！现在，在你面前光荣我罢！赐给我在世界未有以前，我在你前所有的光荣罢}\BibleRef{若 17:5} 呢？如果如他们所说，祂那时才被高举，那么祂先前如何\Quote{使天低垂而下降}，又如何\Quote{至高者发出自己的雷霆}呢？因此，如果甚至在创世之前，圣子就已拥有那光荣，且是光荣之主和至高者，并自天降下，且永受崇拜，那么由此可知，祂并非因降临而获得提升，反倒是祂自己提升了需要提升的事物；而如果祂降临是为了完成它们的提升，那么祂并非作为赏报而领受了子和天主的名号，反倒是祂自己使我们成为圣父的子女，并藉由自己成为人而将人神化\OriginalTerm{神化}{deified}。\par

39. 因此，祂并非先是人，然后才成为天主，而是祂原是\OriginalTerm{天主}{God}，然后才成为人，为的是将我们神化。因为如果祂成为人时，才被称为子和天主，但在祂成为人之前，天主已称古代的子民为子女，并使梅瑟作法郎的天主（且圣经论多人说：\BibleQuote{\OriginalTerm{天主站在天主的会中}{God stands in the congregation of Gods}}），那么很显然，祂被称为子和天主要比他们晚得多。那么万物如何藉着祂而有，而祂又如何在万有之先呢？或者，如果祂之前有其他被称为子和神的人，祂又怎能是\Quote{一切受造物的首生者}呢？那些最初分享者又怎能不分享圣言呢？这种看法并非真理，而是我们当代\OriginalTerm{犹太化派}{Judaizers}的诡计。因为在这种情形下，任何人怎能认识天主为父呢？因为如果没有真子，就不可能有\OriginalTerm{收养}{adoption}，而真子说：\BibleQuote{除了子和子所愿意启示的人以外，没有人认识父}\BibleRef{玛 11:27}。而如果没有圣言、且在圣言之先，又怎能有神化呢？然而，祂却对他们的同族犹太人说：\BibleQuote{如果那些承受天主圣言的人，天主都称他们为神}\BibleRef{若 10:35}。如果所有被称为子和神的人，无论在地上还是天上，都是藉着圣言而被收养和神化的，而圣子自己就是圣言，那么很明显，万物都是藉着祂而有，祂自己在万有之先，或者说，唯有祂自己是\OriginalTerm{真子}{very Son}，唯有祂是从真天主而来的\OriginalTerm{真天主}{very God}，并非因自己的德行而获得这些殊荣作为赏报，也不是与这些不同的另一位，而是按\OriginalTerm{本性}{nature}和\OriginalTerm{本质}{essence}就是这一切。因为祂是圣父本质的\OriginalTerm{所生者}{Offspring}，因此无人能怀疑，仿效那不可变更的圣父，圣言也是不可变更的。\par

40. 到目前为止，我们已用\Quote{子}这个词语所包含的正确观念回应了他们非理性的想法，这是主亲自给我们的。但接下来最好是引用天主的谕旨，好能更充分地证明\OriginalTerm{子}{Son}永不改变的特性，以及祂那属于父的、不能变更的本性，同时也证明他们的悖谬。于是，宗徒致斐理伯人书说：\BibleQuote{你们该怀有基督耶稣所怀有的心情：他虽具有天主的形体，并没有以自己与天主同等，为应当把持不舍的，却使自己空虚，取了奴仆的形体，与人相似，形状也一见如人；他贬抑自己，听命至死，且死在十字架上。为此，天主极其举扬他，赐给了他一个名字，超越其他所有的名字，致使上天、地上和地下的一切，一听到耶稣的名字，无不屈膝叩拜；一切唇舌无不明认耶稣基督是主，以光荣天主\OriginalTerm{圣父}{Father}。}\BibleRef{斐 2:5}还有什么比这更清楚明白的呢？祂并非由低微的境况擢升，相反，祂本为天主，却取了奴仆的形体，在取这形体时，祂不但没有被擢升，反而贬抑了自己。那么，这里哪有什么德行的赏报，或是在屈辱中的擢升与高升呢？因为，倘若祂本是天主，而成为人，自高天降下后，还说是被举扬，祂既是天主，又怎能被举扬呢？这同时也很清楚：既然天主是至高无上的，祂的\OriginalTerm{圣言}{Word}也必然同样是至高无上的。那么，那位在父内、且在各样事上都与父相似的，又怎能在更高处被举扬呢？所以，祂不需要任何加添；也并非如亚略派所设想的那样。因为尽管圣言降下是为被举扬，且经如此记载，但祂何须贬抑自己，好像要去追求那本已拥有的呢？祂本是\OriginalTerm{恩宠}{grace}的施予者，又领受了什么恩宠呢？祂本因自己的名而常受崇拜，又怎样领受那受敬拜的名呢？的确，在祂成为人之前，圣经作者就呼求祂：\BibleQuote{天主，求你因你的名拯救我}；又\BibleQuote{有人依靠战车，有人依靠战马，我们却只依靠主我们天主的名。} \OriginalTerm{圣祖}{Patriarchs}们敬拜祂，论及天使时经上记着说：\BibleQuote{天主的众天使都要崇拜祂。}\BibleRef{希 1:6}\par

41. 倘若如达味在第七十一篇圣咏中所说：\BibleQuote{他的名必要永世受颂扬，并要在太阳前，万古长存}，他怎会领受那本在领受之前就已永远拥有的呢？或者在受举扬之前他已是至高者，又如何被举扬呢？或者那在领受之前就常受崇拜的，又如何领受受崇拜的权利呢？这不是隐晦的话，而是一个神圣的奥迹。\BibleQuote{在起初已有圣言，圣言与天主同在，圣言就是天主}；但为了我们，圣言成了血肉，正如经上所说。至于此处所说的\Quote{极其举扬}，并不表示圣言的本体被举扬，因为祂从来就是，且仍是\BibleQuote{与天主同等}\BibleRef{斐 2:6}，而是指人性的举扬。因此，这话是在圣言成了血肉之后才说的；这样就能清楚看出，\Quote{贬抑}和\Quote{举扬}都是指祂的人性而言；因为哪里有卑微的境地，那里也就能有举扬；如果因为祂取了血肉而记载了\Quote{贬抑}，那么显然，\Quote{极其举扬}也正因为此。因为人性正需要这个，由于肉身的卑微和死亡的缘故。既然圣言是\OriginalTerm{圣父}{Father}的肖像，是不死不灭的，却取了奴仆的形体，并作为人，在肉身内为我们经历了死亡，好藉着死亡把自己为我们奉献给父；因此，作为人，祂也是为我们的缘故，并代表我们，而被极其举扬，好使我们因基督的死亡，都死于基督，同样也在基督内，从死者中复活，上升天庭，极其举扬，正如经上说：\BibleQuote{作前驱的耶稣已为我们进入了那帐幔的内部，并非进入了一座象征性的圣所，而是进入了天堂本境，如今出现在天主面前，为我们转求}，然而，如果基督如今为我们进入了天堂本境，尽管祂在起初并一直是主及诸天的造主，这如今的举扬也是为我们的缘故而记载的。正如那圣化万物的，也为我们的缘故说自己圣化自己奉献给父，并非为使圣言变为圣的，而是为在自身内圣化我们众人。同样，我们也必须如此理解此处\Quote{祂极其举扬了祂}这句话：并非使祂自己被举扬，因为祂本是至高者，而是为使我们成为祂的正义，我们能在祂内被举扬，并能进入天门，这天门祂也已为我们打开，正如前驱们所说：\BibleQuote{城门，请提高你们的门楣，古老的门户，请加大门扉，因为要欢迎光荣的君王}，因为在这里，并非祂作为万有的主和造主，被拒于门外，而是为我们这些被关闭了乐园之门的人，才记录了这话。所以，从人的关系来说，因着祂所承担的肉身，经上论祂说：\Quote{请提高你们的门楣}和\Quote{将进来}，仿佛是一个人进入；但从神性关系来说，既然\BibleQuote{圣言就是天主}，祂便是\BibleQuote{主}和\BibleQuote{光荣的君王}。关于我们这样的举扬，\OriginalTerm{圣神}{Holy Spirit}在第八十九篇圣咏中预言说：\BibleQuote{他们因你的名而欢欣踊跃，因你的正义而挺胸昂首，因为你是他们所夸耀的力量}。如果\OriginalTerm{圣子}{Son}就是正义，那么祂的举扬并非祂自己需要，而是我们在这正义中得以举扬，这正义就是祂\BibleRef{格前 1:30}\par

42. 同样，\Quote{赐给了他}这些话，也并非因为圣言自身而写的；因为甚至在他\OriginalTerm{降生成人}{became man}之前，他已受到崇敬，正如我们所说，由天使及整个受造界，因着他本是属于父的而受崇敬；而是为了我们并且因了我们的缘故，才在圣言身上如此记述。就如\Person{耶稣基督}作为人而死并被高举，同样，作为人，圣经说他取用了那些作为\OriginalTerm{天主}{God}他向来所拥有的，好叫这\OriginalTerm{恩宠}{grace}的赐予也能临于我们。因为\OriginalTerm{圣言}{Word}在取得了身体时并未受损，因而需要寻求恩宠的赐予，相反，他\OriginalTerm{神化}{deified}了那他所穿上的，并且不仅如此，还恩赐地\Quote{赐给了}人类。因为他既然是圣言，具有\OriginalTerm{天主的形体}{form of God}，向来受到崇敬，那么当他仍然是他素来所是的，虽然降生成人并被称为\Person{耶稣}，他依然将整个受造界置于脚下，使万有因这名向他屈膝，并且宣认：圣言成为血肉，在肉身内经历死亡，并非有损他\OriginalTerm{天主性体}{Godhead}的光荣，而是\Quote{为光荣\OriginalTerm{圣父}{God the Father}}。因为这是圣父的光荣，就是：受造而又失丧的人能再次被寻回；当他死了，能再生活，并成为\OriginalTerm{天主的圣殿}{God's temple}。因为天上的大能者，无论是天使和总领天使，一向崇敬\OriginalTerm{主}{Lord}，就如他们现今以\Person{耶稣}之名崇敬他，这正是我们的恩宠和崇高地位，即当\OriginalTerm{圣子}{Son of God}降生成人的时候，他仍受崇敬，并且天上的大能者见我们这些与他\OriginalTerm{同属一个身体}{one body with Him}的人，均被引入他们的境域，也必不惊讶。而这事若非那具有天主形体的一位，取了\OriginalTerm{奴仆的形体}{servant's form}，并贬抑自己，交出他的身体以至于死，才得以实现，否则就不可能。\par

43. 那么请看，人因着\OriginalTerm{十字架}{Cross}而视为\Quote{天主的愚妄}，竟成为万有中最受尊崇的。因为我们的\OriginalTerm{复活}{resurrection}积蓄在其中；不再只是\Place{以色列}，而是从此万民，诚如\OriginalTerm{先知}{Prophet}所预言的，抛弃了\OriginalTerm{偶像}{idols}，承认了\OriginalTerm{真天主}{true God}，\OriginalTerm{基督的父}{Father of the Christ}。\OriginalTerm{恶魔}{demons}的幻象已归于乌有，唯有那真实是天主的一位，在我们\OriginalTerm{主}{Lord}\Person{耶稣基督}的名内受崇敬。因为主即使在降生成人，取了一个人性的身体，并被称为\Person{耶稣}时，仍受崇敬，并被相信是天主之子，而且借着他，父也被认识，这事实表明，正如已说过的，并非圣言本身，就其为圣言而言，领受了这如此大的恩宠，而是我们领受了。因为由于我们与他\OriginalTerm{身体}{Body}的联系，我们也成了天主的圣殿，因而成了\OriginalTerm{天主的子女}{God's sons}，以致即使在我们的身上，主也受到崇敬，并且看见的人传报说，正如\OriginalTerm{宗徒}{Apostle}所说的，天主确实在他们中间。正如\Person{若望}在\WorkTitle{福音}中也说：\BibleQuote{凡接受他的，他给他们权能，好成为天主的子女。} \BibleRef{若 1:12}他在\WorkTitle{书信}中写道：\BibleQuote{我们所以知道他住在我们内，是借他赐给我们的\OriginalTerm{圣神}{Spirit}。} \BibleRef{若一 3:24}这也是他向我们施恩的一个明证：我们因\OriginalTerm{至高的上主}{Highest Lord}在我们内而被高举，并且为了我们的缘故，恩宠已经赐给了他，只因那供应恩宠的主成了像我们一样的人；而另一方面，他，\OriginalTerm{救主}{Saviour}，贬抑自己，取了我们这\BibleQuote{我们卑贱的身体} \BibleRef{斐 3:21}，取了奴仆的形体，穿上那被\OriginalTerm{罪}{sin}奴役的血肉。而他确实没有从我们获得任何对自己晋升的好处：因为\OriginalTerm{天主的圣言}{Word of God}无乏无缺，圆满无缺；反倒是我们从他获得了晋升；因为他是\BibleQuote{那普照每人的\OriginalTerm{真光}{Light}，正进入这世界} \BibleRef{若 1:9}。\OriginalTerm{亚略派人士}{Arians}徒劳地强调那个连接词\Quote{为此}，因为\Person{保禄}说过：\Quote{为此，天主极其举扬他}。因为他这样说，并非意指任何德行的奖赏，也非由于进境的提升，而是指出将举扬赐予我们的原因。这原因无非就是：那具有天主形体的，\OriginalTerm{尊贵父的圣子}{Son of a noble Father}，贬抑了自己，代替我们并为我们成了奴仆。因为如果主没有降生成人，我们便不能被救赎脱离\OriginalTerm{罪}{sin}，不能从死者中复活，而仍倒卧地下；不能升入\OriginalTerm{天堂}{heaven}，而仍躺卧在\OriginalTerm{阴府}{Hades}。那么，\Quote{极其举扬}和\Quote{赐给}这些话，正是为了我们并因着我们的缘故而说的。\par

44. 因此，我认为这段经文的意义就在于此，并且这是一种极为合乎教会信仰的意义。不过，还可以用另一种方式加以评述，以平行的方式给出相同的意义；也就是说，虽然经文并未提及圣言自身的举扬，就其是圣言而言（因为，正如刚才所说，祂至高无上且与祂的圣父相似），但借着祂降生成人，它指示了祂从死者中的复活。因为在说了\BibleQuote{祂自谦自卑，直至死亡}之后，他立即加上\BibleQuote{为此，天主极其举扬祂}；旨在表明，虽然就人而论祂据说死了，但祂作为生命，却在复活中被举扬；因为\BibleQuote{那下降的，就是那上升的}。祂以身体下降，也因祂是身体中的天主本身而复活。这也是按照这层意义，他引入连接词\Emph{为此}的另一个原因；这并不是作为德行的赏报或地位的提升，而是为表明复活发生的原因；并且说明，为何从亚当至今的所有其他人，都死去并仍然处于死亡中，只有祂完完整整地从死者中复活了。原因正是祂自己已经教导我们的：祂是天主，却成了人。因为其他所有人，仅是从亚当而生，都死了，死亡统治了他们；但祂，\Emph{第二人}，却来自天上，因为\BibleQuote{圣言成了血肉}\BibleRef{若 1:14}，而这个人就被称为来自天上、属于天上的，因为圣言从天下降；因此，祂并没有被死亡所拘禁。因为虽然祂自谦自卑，舍弃自己的身体去接受死亡，因这身体能死，但祂因着是在肉体中的天主子，却从地上被高举。因此，这里所说的\BibleQuote{为此，天主也极其举扬祂}，正与\WorkTitle{宗徒大事录}中伯多禄的话相符：\BibleQuote{天主解除了祂死亡的苦痛，举扬了祂，因为祂不可能被死亡所拘禁}\BibleRef{宗 2:24}。因为正如保禄所写的：\BibleQuote{祂虽具有天主的形体，却成了人，自谦自卑，直至死亡，因此天主也极其举扬祂}，所以伯多禄也说：\BibleQuote{因为祂是天主，却成了人，并以奇迹异事向观看者证明祂是天主，所以祂不能被死亡拘禁}。为人，这不可能成功；因为死亡属于人；因此，圣言，既是天主，就成了血肉，好使祂在血肉内被处死，能以祂自己的大能使所有人获得生命。\par

45. 但是，既然据说祂自己是被\Emph{举扬}的，而天主\Emph{赐予}了祂，而异端者们认为这乃是圣言本质中的一个缺陷或一种感情，那么就有必要解释这些话语是如何运用的。祂被说成是从地的低下之处被举扬，因为甚至死亡也被归于祂。这两件事都被算作祂的，因为正是祂的身体，而非其他任何人的身体，从死者中被举扬，并被接到天上。再者，由于身体是祂的，而圣言并不在身体之外，当身体被举扬时，祂，作为人，因为身体的缘故，被言说是被举扬的，这是很自然的。那么，如果祂没有成为人，就不要对祂说这话：但是，既然圣言成了血肉，复活和举扬，就必然像在人身上那样归给祂，好使那归给祂的死亡，能成为人类罪恶的救赎，并废除死亡，而复活和举扬，为了祂的缘故，能为我们确保稳固。在这两方面，经上论到祂都说：\BibleQuote{天主极其举扬了祂}，以及\BibleQuote{天主赐予了祂}；从而他还要表明，不是圣父成了血肉，而是祂的圣言，成为人，并以人的方式从圣父领受，且如已说过的，由圣父举扬。并且很清楚，也没有人会争辩，圣父所赐予的，是透过圣子而赐予的。这的确是奇妙且令人敬畏的；因为圣子从圣父那里赐予的恩宠，圣经说圣子自己领受了；并且圣子从圣父那里赐予的举扬，圣子自己也因此被举扬。因为祂既是天主子，自己却成了人子；并且，作为圣言，祂从圣父那里赐予，因为凡圣父所做和赐予的一切，祂都是透过祂而做并供给的；而作为人子，祂自己以人的方式被言说成从祂那里领受了这些，因为祂的身体不是别的，正是祂的身体，并且如已说过的，是恩宠的自然承受者。因为祂领受恩宠，就祂的人性被举扬而言；这举扬就是使人性\OriginalTerm{被神化}{deified}。但这样的举扬，圣言自己照圣父的天主性和完美始终拥有，这是祂本有的。\par

% structure: p0062 source=h2 target=subsection reason=relative-heading-rank; base=h2

\subsection{第十二章 经文释义；其次，\BibleBook{圣咏}{}四十五篇七至八节。——\Quote{因此}、\Quote{受傅}等字眼是否暗示圣言获得了酬报。首先从\Quote{同伴}或\Quote{有分者}一词加以反驳。祂在人性中受圣神傅油，以圣化人性。因此，圣神在祂于\Place{约旦河}受洗时降临在祂身上。祂也被说成是为我们圣化自己，并将祂所领受的光荣赐给我们。\Quote{因此}一词暗示祂的神性。\Quote{你爱慕正义}等，并不暗示试探或选择。}

46. 对宗徒话语的这种解释足以驳斥那些不虔诚的人；而圣咏作者的话也可作同样正统的理解，他们妄加曲解，但圣咏作者的意思显然是虔敬的。他这样说：\BibleQuote{天主，你的宝座是永远和永远的；正义的权杖是你王国的权杖。你喜爱正义，憎恨邪恶，因此天主，你的天主，用喜乐之油给你傅油，胜过你的同伴。} 看，你们亚略派的人，由此承认真理吧。圣咏作者称我们众人都是主的\Quote{同伴}或\Quote{分享者}：但如果他是由虚无中而来的事物和受造物之一，他自己也是分享者之一。然而，既然他赞颂他为永生的天主，说：\BibleQuote{天主，你的宝座是永远和永远的}，并且宣告一切其他事物都分享他，我们还能得出什么结论，除非他不同于受造物，并且唯独他是父的真圣言、光辉和智慧，一切受造物都分享他，并借着他而在圣神内受圣化呢？因此，他在这里是\Quote{受傅者}，并不是要他成为天主，因为他本来就是；也不是要他成为君王，因为他拥有永恒的王权，作为天主的肖像而存在，正如神圣的神谕所显示的；而是为了我们的缘故，如前面所写。因为以色列的君王，在受傅之后才成为君王，此前并不是，如\Person{达味}、\Person{希则克雅}、\Person{约史雅}和其余的人；但救主相反，他本是天主，永远在父的王国中为王，并且是那位赐予圣神者，然而这里却说他是受傅者，好使他像先前那样，作为人而受圣神的傅油，好能为我们人类预备，不仅是高举和复活，更是圣神的内住和亲密。主自己也在\BibleBook{若望}{福音}中亲口说明这事：\BibleQuote{我已将他们派遣到世界上，我又为他们的缘故圣化我自己，好叫他们也因真理被圣化。} 他说这话，显示他不是被圣化者，而是圣化者；因为他不是由其他受圣化，而是他圣化自己，好叫我们能在真理中被圣化。那圣化自己的，就是圣化的主。这是如何发生的呢？他的意思是：\Quote{我，作为父的圣言，当我成为人时，我将圣神赐给我自己；而我自己成为人，就在祂内圣化我自己，好使在我——真理（因为\InnerQuote{你的圣言就是真理}）内，众人都得圣化。}\par

47. 那么，如果祂是为了我们的缘故圣化自己，且是在成为人时这样做，那么很明显，圣神在约旦河降在祂身上，乃是降在我们身上，因为祂负荷了我们的肉体。这并非为了圣言的提升，而是为了我们的圣化，好使我们能分享祂的傅油，而关于我们能说：\BibleQuote{你们不知道你们是天主的圣殿，天主圣神住在你们内吗？}\BibleRef{格前 3:16} 因为主作为人，在约旦河受洗时，乃是我们藉着祂并在祂内受洗。当祂领受圣神时，正是我们藉着祂而成为祂的接受者。而且，为此缘故，祂不像\Person{亚郎}或\Person{达味}或其他人那样受油傅，而是以超越一切同伴的方式，\Quote{用喜乐之油}受傅，祂自己将此解释为圣神，藉先知说：\BibleQuote{上主的神临到我身上，因为上主给我傅了油}\EditorialNote{依撒意亚 61:1}；正如宗徒也说过：\BibleQuote{天主怎样以圣神傅了纳匝肋的耶稣}\BibleRef{宗 10:38}。那么，这些事何时论及祂呢？正是在祂以肉躯降来，在约旦河受洗，圣神降在祂身上之时。主自己也说：\BibleQuote{圣神要由我领受；} 且说：\BibleQuote{我要派遣祂来；} 又对门徒说：\BibleQuote{你们领受圣神吧。} 然而，祂作为父的圣言和光辉，施予他人者，如今却说是受圣化，因为此刻祂已成为人，那受圣化的身体是祂的。那么，我们从祂开始领受了傅油和印记，\Person{若望}说：\BibleQuote{你们从那圣者领受了傅油；} 宗徒也说：\BibleQuote{你们受了恩许圣神的印记。} 所以这些话语是因我们且为了我们而说的。那么，从我们主的这个实例中，能证明什么提升的进步、德行的赏报或一般品行的回报呢？因为若祂不是天主，后来成为天主，若非王而升为王，你们的论证或许会有一点表面的合理。但若祂是天主，且祂王国的宝座是永恒的，天主如何能进步呢？坐在父宝座上的祂，还有什么缺欠的呢？而且如果正如主自己所说，圣神是祂的，也从祂领受，且祂派遣祂，那么受傅以自己施予的圣神者，并不是作为圣言和智慧的圣言，而是祂所取的肉躯，这肉躯在祂内并藉着祂受傅；好使那临于作为人的主的圣化，能从祂临于众人。因为祂说，圣神不是凭自己说话，而是圣言将祂赐予配得的人。这类似于前面讨论的经文；正如宗徒写道：\BibleQuote{他虽具有天主的形体，并没有以自己与天主同等为应把持不舍的，却空虚自己，取了奴仆的形体，} 同样\Person{达味}赞颂主，为永生的天主和君王，却被派遣到我们这里，并取了我们有死的肉躯。这就是他在圣咏中的意思：\BibleQuote{你的衣冠散布没药、芦荟和肉桂的芬芳；} 而\Person{尼苛德摩}和\Person{玛利亚}的同伴正代表了这一点，因为前者带着\BibleQuote{大约一百斤没药和芦荟的混合物}而来，后者则带了\BibleQuote{为安葬主身体所预备的香料}\BibleRef{若 19:39}。\par

那么，对不死者来说，取了有死的人性，有什么进步呢？对永恒者来说，穿上了暂世之物，又有什么提升呢？在圣父怀中的永恒天主和君王，能得到什么伟大的赏报呢？你们难道看不出，这也是为了我们、由于我们而做并记下来的，为叫我们这些有死且暂世的人，主成为人，能使我们成为不朽的，并领入永恒的天国吗？你们散布谎言反对神圣的圣言，难道不脸红吗？因为当我们的主耶稣基督在我们中间时，我们确实得到了提升，如同从罪中得到解救；但祂始终如一；祂成为人时也并没有改变（重复说我已说过的），而是如经上所写：\BibleQuote{天主的圣言永远常存。} 确切地说，正如祂在成为人之前，圣言分赐圣神给圣徒，当作祂自己的，同样，当祂成为人时，祂也借着圣神圣化所有人，并对祂的门徒说：\BibleQuote{领受圣神吧。} 祂也曾赐给梅瑟和其他七十人；达味也曾借着祂向圣父祈祷说：\BibleQuote{不要从我身上收回祢的圣神。} 另一方面，当祂成为人时，祂说：\BibleQuote{我要给你们派遣护慰者，就是真理之神；}\BibleRef{若 15:26} 而祂确实派遣了祂，身为天主圣言的祂，是忠信的。因此，\BibleQuote{耶稣基督昨天、今天、直到永远，常是一样，}\BibleRef{希 13:8} 保持不变，同时给予又领受，以天主圣言的身份给予，以人的身份领受。那么，受提升的并不是圣言，即从圣言角度看的圣言；因为祂本来就拥有一切，并且永远拥有一切；而是人，人是在祂内并借着祂才有领受这些的根源。因为，当如今说祂以人的方式受傅时，那实在是我们在祂内受傅；同样，当祂受洗时，也是我们在祂内受洗。然而，救主在这一切事上大大地光照我们，祂对圣父说：\BibleQuote{祢赐给我的光荣，我已赐给了他们，为使他们合而为一，如同我们合而为一一样。}\BibleRef{若 17:22} 那么，祂是为了我们而祈求光荣，经上出现的\Quote{领受}、\Quote{赐予}和\Quote{高举}这些词，是为使我们能领受，为能赐给我们，并使我们在祂内被高举；同样，祂也是为了我们而祝圣自己，为使我们在祂内被祝圣。\par

但是，如果他们利用\Quote{因此}这个字，就是圣咏中与\BibleQuote{因此天主，就是祢的天主，给祢傅了油}这句相连的\Quote{因此}，来为自己的目的服务，那么让这些对圣经一知半解却精通不虔敬的人知道，如前所说，\Quote{因此}这个字并不表示对圣言的德行或行为的赏报，而是说明祂降临到我们这里的原因，以及为了我们而在祂内发生的圣神傅油的原因。因为祂并没有说：\Quote{因此祂给祢傅油，为使祢成为天主、君王、圣子或圣言}——因为祂以前就如此，且永远如此，正如已经证明的；而是说：\Quote{正因为祢是天主和君王，所以祢才受傅，因为除祢以外，无人能将人与圣神结合；祢是圣父的肖像，我们在起初是按照它受造的；因为圣神也属于祢。} 因为受造物的本性无法为此提供保证，天使背命了，人违命了。因此，需要一位天主，而圣言就是天主；为使那些处于诅咒之下的人，祂自己能释放他们。那么，如果祂是从虚无中来的，祂就不会是基督或受傅者了，因为祂只不过是众人中的一位，与其他人一样有分。然而，正因为祂是天主，因为祂是天主子，是永恒的君王，且是圣父的\OriginalTerm{光荣的反映和本体的真像}{Radiance and Expression}\BibleRef{希 1:3}，所以祂正是人们所期待的基督，圣父通过启示给祂的圣先知们，而向人类宣布了祂；为叫我们怎样借着祂而有存在，同样也能在祂内使众人从自己的罪中得蒙救赎，并由祂统御万有。这就是在祂内发生的傅油的原因，也是圣言降生成人的临在的原因；圣咏作者预见此事而加以歌颂，他首先以这些话称扬祂的天主性和国，那是属于圣父的：\BibleQuote{天主，祢的宝座永远常存；祢治国的权杖是公正的权杖}；然后如此宣布祂降临到我们这里：\BibleQuote{因此天主，就是祢的天主，以喜乐的油给祢傅了油，胜过祢的同伴。}\par

50. 如果那位赐予圣神的主，在此被说成是受圣神傅油，这有什么可惊讶的？有什么不可信的？当时，由于需要，祂也未曾因自己的人性而拒绝称自己低于圣神。因为当犹太人说祂仗赖\Person{贝耳则步}驱魔时，祂为揭露他们的亵渎，回答他们说：\BibleQuote{但我若仗赖天主的神驱魔}\BibleRef{玛 12:28}。看，圣神的赐予者在此说祂仗赖圣神驱魔；但这话之所以说出，无非是因祂的肉身。因为人的本性自身不足以驱魔，唯有赖圣神的能力，因此祂作为人这样说：\BibleQuote{但我若仗赖天主的神驱魔}。当然，祂也表明，冒犯圣神的亵渎，比针对祂人性的亵渎更为重大，当祂说：\BibleQuote{凡出言干犯人子的，可得赦免}；就像那些说\BibleQuote{这人不是那木匠的儿子吗？}的人一样；但是那些亵渎圣神，将圣言的作为归于魔鬼的人，必受不可免的惩罚。这就是主作为人对犹太人所说的话；但对门徒，为显示祂的天主性和尊威，并表明祂并不低于圣神、而是与圣神同等，祂赐下圣神，说：\BibleQuote{你们领受圣神吧！}，又说\BibleQuote{我派遣祂}，\BibleQuote{祂要光荣我}，\BibleQuote{凡祂所听见的，都要说出来}。因此，正如在这一处，主自己——圣神的赐予者——并不拒绝说祂作为人，仗赖圣神驱魔；同样，同一位主，圣神的赐予者，也不拒绝说：\BibleQuote{上主的神临于我身上，因为祂给我傅了油}\BibleRef{依 61:1}，这是基于祂的降生成人，正如\Person{若望}所说的；为的是在这两点上都显明：我们是在圣化中需要圣神恩宠的人，又是没有圣神的能力便不能驱魔的人。那么，圣神应藉着谁、从谁而来呢？岂不是藉着圣子，圣神也是祂的吗？我们何时得以领受圣神呢？岂不是在圣言降生成人之时吗？而且，正如宗徒的经文所指明的，若非那位具有天主形体的取了奴仆的形体，我们就不会得蒙救赎并被高举；同样，\Person{达味}也指明，我们若非如此，便不能分享圣神并蒙圣化，而只有圣神的赐予者——圣言自己——曾亲自说祂为我们受了圣神的傅油。因此，我们稳妥地领受了圣神，因为祂被说成是在肉身内受傅油；因为肉身先在祂内受圣化，而祂作为人，被说成为肉身而领受圣神，于是我们便蒙受随之而来的圣神恩宠，\BibleQuote{从祂的满盈中}\BibleRef{若 1:16} 我们都领受了。\par

51. 至于\BibleBook{}{圣咏}中所加上的以下话语——\BibleQuote{你爱慕正义，憎恨邪恶}——也并非如你们再次所猜想的那样，显示圣言的本性是可变的，反倒正因其力量而表明祂的不变性。因为受造物的本性是可变的，其中一部分曾犯罪过，另一部分背命不从，正如已经说过的；而且它们将如何行动是不确定的，往往现今是善的，后来却改变而变成不同的，以致刚才公义的人，不久就被发现为不义；因此，在此也需要一位不变的，好使人们能以圣言正义的不变性作为德行的肖像和楷模。这一思想强烈地吸引着正直的人。因为第一个\Person{亚当}改变之后，死亡便经由罪进入了世界；因此，第二亚当必须是不变的；这样，若蛇再次攻击，蛇的诡计便可被挫败，而主既是不变不易的，蛇在攻击众人时便无能为力了。就如\Person{亚当}犯罪之后，他的罪殃及了众人，同样，当主成为人、并推翻了蛇之后，祂那如此伟大的能力也要伸展到众人，使我们每人都能说：\BibleQuote{因为我们不是不知道他的诡计。}\BibleRef{格后 2:11}所以，主——祂在本性上永远是不变的，爱慕正义，憎恨邪恶——理所当然应受傅油并被派遣，为的是，祂虽保持原有的身份，却取了这可变的肉身，好\BibleQuote{在肉身上定了罪恶的案}，并确保肉身的自由，且使它从此能在自身内\BibleQuote{满全法律所要求的正义}，以致能说：\BibleQuote{但我们已不属于肉性，而是属于圣神，只要天主的圣神住在你们内。}\BibleRef{罗 8:9}\par

52. 所以，亚略派啊，你们在这里又一次徒然作了这种臆测，也徒然地引用了圣经的话；因为天主的圣言是不变更的，且常处于同一状态，并非可以随意变动，而是如同圣父一样；因为如果他不如此，又怎能肖似圣父呢？再说，倘若他没有圣父的不变更性和不易性，圣父所有的一切，又怎能也全是圣子的呢？他爱此一，恨彼一，并非如同受法律约束、有所偏倚一样，以免我们因为他是由于怕堕落而选择此一，从而承认他在其他方面也是可变更的；不，他既是天主，又是圣父的圣言，所以他是一位公义的审判者，也是德行的爱慕者，或者更好说，是德行的施与者。因此，他因本性即是公义和圣善的，故而经上说，他喜爱正义，憎恨邪恶；这等于说，他喜爱并选择有德之人，弃绝并憎恨不义之人。天主的圣经对圣父也说了同样的话：\BibleQuote{正义的上主喜爱正义；你憎恨所有作恶的人}，以及\BibleQuote{上主喜爱熙雍的城门，胜过雅各伯所有的帐棚}；还有，\BibleQuote{我爱了雅各伯，而恨了厄撒乌}\BibleRef{拉 1:2}；在\BibleBook{}{依撒意亚}中，又有天主的声音说：\BibleQuote{我上主喜爱正义，憎恨不义的抢劫}\BibleRef{依 61:8}。那么，就让他们像解释后面这些经文一样，去解释前面那些经文吧；因为前面那些也是论及天主的肖像而写的：否则，若是对这些经文也像解释那些经文一样妄加曲解，他们就会设想圣父也是可变更的了。然而，连听到别人这样说也并非没有危险，所以我们最好这样想：经上说天主喜爱正义、憎恨不义的抢劫，并非好像他有所偏倚，能够接受相反的事，以致选择后者而不选择前者，因为那属于受造之物；而是说，他如同一位审判者，喜爱义人，收纳义人，远离恶人。因此，我们关于天主的肖像也当这样想，即他的爱与恨不会以别的方式发生。因为天主的肖像必须与父有同样的本性，尽管亚略派在盲目之中既看不见那肖像，也看不见神圣言语的任何其他真理。因为，他们受到自己内心观念，或更好说是误解的驱使，便求助于圣经的段落，而由于缺乏理解，一如他们惯常的那样，辨别不出其真义；反倒把自己的不虔敬立为某种解经的规矩，强行将所有神圣的话语都曲解为迎合它。因此，他们一提出这类学说，就只配得到这样的回答：\BibleQuote{你们错了，不明白圣经，也不了解天主的能力}\BibleRef{玛 22:29}；如果他们执迷不悟，就该以下面的话使他们缄默：\BibleQuote{将属于人的归于人，将属于天主的归于天主}。\par

% structure: p0070 source=h2 target=subsection reason=relative-heading-rank; base=h2

\subsection{第十三章 经文解释；第三点，\BibleBook{希伯来}{书}一章四节。引用的附加以作反对的经文；例如 \BibleBook{希伯来}{书} 1:4; 7:22。\OriginalTerm{更优越}{better}一词是否暗示与天使相似；\OriginalTerm{受造}{made}或\OriginalTerm{成为}{become}是否意味受造。需要考量圣经发言时的情境。\OriginalTerm{更优越}{better}与\OriginalTerm{更大}{greater}之间的区别；用作证据的经文。\OriginalTerm{受造}{made}或\OriginalTerm{成为}{become}是一般性的词语。在 \BibleBook{希伯来}{书} 1:4 中，在本性上将圣子与受造物作对比。两个盟约下惩罚的差异显示了圣子与天使本性的差异。\OriginalTerm{成为}{become}不是相关于圣言的本性，而是相关于祂的人性、职分以及对我们的关系。平行段落中，这一用语被应用于永恒的圣父。}

53. 但是他们说，在\BibleBook{}{箴言}中写着：\BibleQuote{上主自始即拿我作祂行动的起始，作祂作为的开端}；在\BibleBook{希伯来}{书}中，宗徒说：\BibleQuote{祂所承受的名分既超越天使，祂就成了比天使更优越的}；随后又说：\BibleQuote{因此，有分于天上召叫的圣徒弟兄们！你们应细心想想我们宣认的宗徒和大司祭——基督耶稣，祂对那委派祂的是忠信的}；在\BibleBook{宗徒}{大事录}中：\BibleQuote{所以，以色列全家应确切知道：天主已把你们所钉死的这位耶稣，立为主，立为默西亚了}。他们反复提出这些经文，却误解其意义，以为它们证明了天主的圣言是一个受造物、一件工程，且是受造之物中的一员；他们就这样欺骗了那些不加思考的人，以圣经的语言为借口，却在那里播下他们自己异端的毒种，取代了真实的意义。假使他们知道，就不会对\BibleQuote{光荣的主}\BibleRef{格前 2:8}不虔敬，也不会曲解圣经的善言。如果他们从此以后公开采取\Person{盖法}的做法，决意犹太化，并且不懂得那经文，即天主真要住在地上，那么，就让他们不要探究宗徒们的言论罢；因为这不是犹太人的作风。但是，如果他们与不虔敬的摩尼教徒勾结，否认\BibleQuote{圣言成了血肉}以及祂降生成人的临在，那么，他们就别提出\BibleBook{}{箴言}，因为这与摩尼教徒格格不入。不过，如果他们为了利禄，以及随之而来的贪婪之财和渴望获得美名，不敢否认\BibleQuote{圣言成了血肉}这句经文，因为经上是这样写的，那么，要么让他们正确地解释圣经上关于救主人性临在的那些话，要么，如果他们否认真义，就让他们彻底否认主曾成为人。因为，一方面承认\BibleQuote{圣言成了血肉}，另一方面又以论及祂的记载为耻，并因此扭曲其意义，这实属不妥。\par

54. 因为经上记载说：\BibleQuote{他远超众天使之上}。我们就先来查考这话。如今，如同对待一切圣经一样，在此也有必要忠实地阐明宗徒所写的时代、人物和要点；以免读者因无知而忽略了这些或类似的细节，偏离了真义。那位好问的太监明白了这一点，便如此恳求\Person{斐理伯}说：\BibleQuote{请问，先知说这番话是指谁呢？是指自己，还是别人呢？}\BibleRef{宗 8:34} 因为他担心，如果解经时引用的人物不当，便会远离真义。而门徒们为了要得知所预言之事发生的时间，便恳求主说：\BibleQuote{请告诉我们：这些事什么时候发生？你来临的征兆是什么？}\BibleRef{玛 24:3} 再者，当他们从救主口中听到末世的事件后，便渴望得知其时间，好使自己不陷于谬误，也能教导别人；例如，当他们明白之后，便纠正了走上错路的\Person{得撒洛尼人}。所以，谁若正确了解这些要点，他对信仰的理解便是正确而健全的；但若误解了类似之点，便立即陷入异端。因此，\Person{依默纳约}、\Person{亚历山大}及其同伙，就是在时代上错了，他们说复活已经发生；\Person{迦拉达人}则是在时间之后，他们现今仍执着于割损。至于误认对象，正是犹太人的通病，至今如此；他们认为\BibleQuote{看，一位贞女将怀孕生子，人将称他的名字为\Person{厄玛奴耳}，意思是：天主与我们同在}\BibleRef{依 7:14}；而认为\BibleQuote{上主你们的天主，要从你们当中给你们兴起一位先知}\BibleRef{申 18:15}，是指某一位先知；至于\BibleQuote{他如同被牵去宰杀的羊}\BibleRef{依 53:7}，他们非但不向\Person{斐理伯}求教，反而臆测是指\Person{依撒意亚}或先前某一位先知说的。\par

55. (3.) 这正是基督的仇敌陷入其可憎异端时的心态。因为，如果他们认识宗徒这话所指的人物、主题和时机，就不会把属于基督人性之事曲解为基督的天主性，也不会在愚昧中犯下如此大不敬的行为。如今，只要正确解释这段选经的开头，这一点便显而易见了。因为宗徒说：\BibleQuote{天主在古时，曾多次并以多种方式，藉着先知对我们的祖先说过话；但在这末期内，祂藉着自己的儿子对我们说了话}\BibleRef{希 1:1}，稍后他又说：\BibleQuote{当他独自洁除了我们的罪恶之后，便在高天之上至大者的右边坐下，他所承受的名字既然比天使的名字更尊贵，他便远超众天使之上。}由此看来，宗徒的话提到那个时代——天主藉着自己的儿子向我们说话，以及洁除罪恶发生之时。那么，祂何时藉着儿子向我们说话，洁除罪恶又在何时发生呢？祂又在何时成了人？不就是在先知之后的末世吗？接着，当他讲述那关乎我们的\OriginalTerm{经世}{economy}，并论及末世之时，他自然提及，即使在古时，天主也未曾对人缄默，而是藉着先知向他们说话。既然先知执行职务，法律是由天使颁布的，而圣子也来到地上，且是为了执行职务，他便不得不加上\BibleQuote{远超众天使之上}这话，为要表明：圣子既远超仆人，同样，圣子的职务也远胜于仆人的职务。宗徒将旧的职务与新的职务加以对照，便坦然对犹太人写道：\BibleQuote{远超众天使之上}。正因如此，他通篇没有使用\Quote{成为更伟大的}或\Quote{更尊荣的}这类比较词，免得我们以为基督与天使属于同类；他特意使用了\Quote{更好}一词，为标明圣子的性体与受造之物全然不同。对此我们有圣经为证；例如，\Person{达味}在圣咏中说：\BibleQuote{在你宫庭内逗留一日，远胜过在别处逗留千日}；\Person{撒罗满}也高呼：\BibleQuote{你们应听取我的教训，而不要银子；应吸收知识，而不要纯金；因为智慧胜过任何珍珠，任何可贪恋的事都不能与她伦比。}\BibleRef{箴 8:10}智慧与地上的宝石在本质和性体上岂非截然不同？天上的庭院与地上的房屋岂有任何相似？永生和属灵之事与暂时和可朽之事岂有可比之处？\Person{依撒意亚}也这样说：\BibleQuote{因为上主这样说：那些遵守我的安息日、拣选我所喜悦的事、持守我盟约的太监，我要在我的屋内、在我的墙垣内，赐给他们比子女更美的一份和名号；我要赐给他们一个永久不能取消的名号。}\BibleRef{依 56:4}同样，圣子与天使之间毫无相类之处；因此，\Quote{更好}一词并非用于比较，而是用于对照，因为祂的性体与天使的性体截然不同。因此，宗徒本人解释\Quote{更好}一词时，其含义无非是指圣子远超一切受造之物：他称前者为\Quote{子}，后者为\Quote{仆}；前者以子的身分与父同坐于右边；后者则以仆役的身分侍立于前，奉差遣，尽职务。\par

56. 圣经这样说话，阿里乌派啊，暗示圣子并非受造，反而异于受造之物，且为圣父所专有，住在祂的怀中。（四）就连此处的\Quote{\OriginalTerm{成为}{become}}这个说法，也不像你们所设想的，表明圣子是受造的。如果单单是\Quote{成为}而无其他，那么阿里乌派或许还有理由；然而经上通篇已先有\Quote{圣子}一词，表明祂异于受造之物，同样地，\Quote{成为}一词也绝非单独出现，而是紧接着加上了\Quote{\OriginalTerm{更好}{better}}。因为作者认为这说法无关紧要，知道对于一位公认的真正儿子，说\Quote{成为}就等于说祂\Quote{受造}，并且是\Quote{更好}的。因为即使我们说受生者是\Quote{成为}或\Quote{受造}也无妨；但是反过来，\OriginalTerm{受造之物}{things originate}却不能被称作受生者，因为它们是天主的化工，除非在受造之后，它们分享了受生的圣子，因而也可说它们已经受生，这完全不是按它们的本性，而是由于它们在圣神内分享了圣子。这点圣经也再次确认；因为论到受造之物，它说：\BibleQuote{万物是藉着祂而造成的；凡受造的，没有一样不是藉着祂而造成的}\BibleRef{若 1:3}，又说：\BibleQuote{你以智慧造成了万有}；但论到受生的儿子，它说：\BibleQuote{约伯生了七个儿子和三个女儿}\BibleRef{约 1:2}，\BibleQuote{亚巴郎一百岁时，他生了儿子依撒格}\BibleRef{创 21:5}；梅瑟说：\BibleQuote{若有人生了儿子}\BibleRef{申 21:15}。所以，既然圣子异于受造之物，单独是圣父\OriginalTerm{本质}{essence}的亲生子，阿里乌派关于\Quote{成为}一词的申辩就毫无价值。\par

（五）此外，如果他们至此仍被驳倒，却还强行坚持该语言是用于比较的，并因而说比较意味着种类相同，以至圣子具有天使的本性，那么他们首先要蒙受与\Person{瓦伦廷}、\Person{卡尔波克拉特}和其他异端者同流合污的耻辱，瓦伦廷曾说天使与基督是同类的，而卡尔波克拉特则说天使是世界的创造者。或许他们正是在这些师傅的教导下，将\OriginalTerm{天主圣言}{Word of God}与天使相比较。\par

57. 不过，就在这样的思辨中，他们必定会被圣咏作者的话所感动，他说：\BibleQuote{众神之中谁能相似上主？}，又说：\BibleQuote{上主啊，众神之中没有谁可比祢。}然而，必须要回答他们，或许能让他们从中获益，那就是：比较固然应用于同一种类的事物，而非不同种类的事物。例如，没有人会将天主与人相比，或将人与兽相比，或将木头与石头相比，因为它们的本性不同；而是天主超乎一切比较，人与人相比，木头与木头相比，石头与石头相比。在这些情况中，我们不该用\Quote{更好}，而该用\Quote{更}和\Quote{更加}；比如\Person{若瑟}比他的兄弟们更俊美，\Person{辣黑耳}比\Person{肋阿}更美；星星并不比星星更好，而是在光荣上更加卓越；然而当把种类不同的事物放在一起时，就用\Quote{更好}来标记差异，正如之前论智慧和珠宝时所说的。因此，倘若宗徒说\Quote{圣子如此领先于天使}，或说\Quote{如此更大}，你们就有借口，仿佛圣子与天使相比；但事实上，宗徒说祂\Quote{更好}，且与仆人们的差别如同儿子与仆人，这就表明圣子在\OriginalTerm{本性}{nature}上异于天使。\par

（六）此外，他说是祂\BibleQuote{奠定了大地的根基}\BibleRef{希 1:10}，这也就表明祂异于一切受造之物。但如果祂在本质上与它们的本性相异，祂的本质又能与它们有什么比较，或什么相似之处呢？然而，即便他们这样想，\Person{保禄}也会加以驳斥，他切中了要点，说：\BibleQuote{论到天使又说\InnerQuote{天主使自己的天使为神，使自己的仆役为火焰}}\BibleRef{希 1:7}。\par

58. 在此要注意，\Quote{制造}一词适用于受造物，他称它们为受造物；但论到子时，他不说制造，也不说成为，而是论及永恒、王权与创造者的职能，他宣告：\BibleQuote{你的宝座，天主，是永远永远的；} 以及 \BibleQuote{上主，你在起初奠定了大地的基础，诸天是你手中的工程；它们必将毁灭，你却常存。}从这些话中，他们若愿意，就能看出创造者不同于受造之物，前者是天主，后者是受造物，从无中被造。因为所说\Quote{它们必将毁灭}，并非说受造界注定要毁灭，而是借此表明受造物的本性及其趋向。因为那些能够毁灭的事物，虽然因创造者的\OriginalTerm{恩宠}{grace}而不致毁灭，但却是从无中而来的，本身见证它们曾经不存在。正是基于此，由于它们的本性如此，论到子时说\Quote{你却常存}，以显示主的永恒；因为他不像受造物那样有毁灭的可能性，却具有永恒的延续，因此说他\Quote{在受生之前不存在}对他来讲是陌生的，而常存、与父同在才是他所固有的。即使宗徒在 \BibleBook{希伯来}{书} 中没有这样写，他的其他书信以及整部圣经也必定会禁止他们对圣言抱有这种观念。但既然他在这里明确地写出来了，并且，如上所示，子是父\OriginalTerm{本质}{essence}的子，他是创造者，其他万物是借着他受造的，他是父的\OriginalTerm{辉耀}{Radiance}、\OriginalTerm{圣言}{Word}、\OriginalTerm{肖像}{Image}与\OriginalTerm{智慧}{Wisdom}，而受造物则在\OriginalTerm{天主圣三}{Triad}之下各按其位站立侍奉，因此，子在种类和本质上都与受造物不同，相反，他是父的本质所固有，与父本质同一。因此，子也不是说：\BibleQuote{我的父比我更好} \BibleRef{若 14:28}，免得我们以为他与父的本性相异，而是说\Quote{更大}，并非在伟大或时间上，而是因为他从父自身的\OriginalTerm{受生}{generation}；而且，在说\Quote{更大}时，他再次表明自己是父的本质所固有的。\par

59. (7). 而且，宗徒之所以说\Quote{远超过众天使}，其首要意图并非比较圣言与受造物的本质（因为他无法比较，毋宁说两者不可度量），而是针对圣言在肉身内的临在，以及他所实施\OriginalTerm{经世计划}{Economy}的，他要表明他并不像先前那班人；因此，正如他的本性远胜过他先前所派遣的使者，同样，由他而来并借着他而来的\OriginalTerm{恩宠}{grace}，也远胜过透过天使的服役。因为仆人的职责是要求果实，仅此而已；而儿子兼主人的职责则是免除债务，并移交葡萄园。\par

(8.) 当然，宗徒继续说下去的话，显示了子远超受造物的卓越性：\BibleQuote{因此，我们必须更紧密切记所听到的事，免得随流失。因为借着天使所传的话既是确实，凡犯法干纪的，都受了应得的报复；那么，我们若忽略了这么大的救恩，怎能逃罚？这救恩原是主亲自开始宣讲，由听信的人给我们证实的。} \BibleRef{希 2:1}。但如果子属于受造物之列，他就不比他们更好，悖逆他也不会导致加重的惩罚；正如在天使的服役中，违背不同的天使并不会导致犯罪者罪的轻重不同，法律只有一个，对犯法者的报复也只有一个。但正因为圣言不属于受造物之列，而是父的子，所以，正如他自身更善、他的作为也更善且超绝，同样，惩罚也更严厉。因此，让他们仔细思量那借着子而来的\OriginalTerm{恩宠}{grace}，并承认他甚至借着他的作为所见证的：他与受造物不同，唯独是父内的真实子，父也在他内。法律是由天使颁布的，却未成全任何人 \BibleRef{希 7:19}，它需要圣言的临在，如\Person{保禄}所说的；但那临在成全了父的工程。并且，从\Person{亚当}到\Person{梅瑟}，死亡为王 \BibleRef{罗 5:14}；但圣言的来临废除了死亡 \BibleRef{弟后 1:10}。我们不再在\Person{亚当}内全死去 \BibleRef{格前 15:22}；而是在基督内全复活了。那时，从\Place{丹}到\Place{贝尔舍巴}均宣报了法律，只有在\Place{犹太}才认识天主；但如今，他们的声音传遍普世，大地充满了对天主的认识，门徒们使万民成为门徒 \BibleRef{玛 28:19}，现在应验了经上所载：\BibleQuote{众人都要蒙天主的教导} \BibleRef{若 6:45}。那时所启示的不过是\OriginalTerm{预像}{type}；但如今真理已彰显出来。这事宗徒本人在后文又更清晰地描述说：\BibleQuote{耶稣因此成了更美\OriginalTerm{盟约}{covenant}的保证。} 又说：\BibleQuote{如今他领受了更卓越的职分，因为他作了那基于更美恩许所制定的更美\OriginalTerm{盟约}{covenant}的中保。} 以及：\BibleQuote{法律本来就不能完善什么，接着所引进的更美的希望则能。} 他又说：\BibleQuote{按理，天上事物的模型必须这样用这些\OriginalTerm{祭献}{sacrifices}来洁净，但那天上事物的本身，却需用比这些更美好的\OriginalTerm{祭献}{sacrifices}。} 因此，无论是在我们面前的这节经文，还是贯穿始终，他都将\Quote{更美}一词归于主，主比受造物更美且与之不同。因为借着他的祭献更美，在他内的希望更美；借着他的恩许也更美，这不仅是大事与小事的比较，而是本性上的差异，因为实施这\OriginalTerm{经世计划}{Economy}的那位，远比受造物\Quote{更美}。\par

60.（9.）此外，\Quote{祂\OriginalTerm{成了}{become}中保}这句话，是指祂为我们所提供的保证。因为，作为\Quote{\OriginalTerm{圣言}{Word}}，祂\BibleQuote{成了血肉}\BibleRef{若 1:14}，而我们将\Quote{\OriginalTerm{成了}{become}}归于血肉，因为血肉是有起初的和受造的；同样，我们在这里对\Quote{祂成了}这一表述，按照第二种意义加以解释，即：因为祂成了人。要让这些好争辩的人知道，他们这悖谬的图谋已告失败；要让他们知道，保禄并非意指祂的性体\Quote{成了}，因为他知道，祂是\OriginalTerm{圣子}{Son}、\OriginalTerm{智慧}{Wisdom}、\OriginalTerm{光荣}{Radiance}和\OriginalTerm{圣父}{Father}的\OriginalTerm{肖像}{Image}；但在这里，他同样把\Quote{成了}这个词，应用于那盟约的\OriginalTerm{职务}{ministry}，在那盟约中，一度作王的死亡已被废除。因为在这里，通过祂而来的职务也变得更好，因为\BibleQuote{法律因了肉性的软弱所不能行的，天主却行了：祂派遣了自己的儿子，带着罪恶肉身的形状，当作赎罪祭，在这肉身上定了罪恶的罪案}\BibleRef{罗 8:3}，从而除去了过犯；这肉身在过犯中不断被囚禁，以致不接纳天主的心思。祂使血肉能够承载圣言，使我们不再随从肉性，而随从圣神行事，并不断说：\Quote{但我们不在肉性内，而在圣神内}，以及\BibleQuote{因为天主派遣子到世界上来，不是为审判世界，而是为叫世界藉着祂而获救}\BibleRef{若 3:17}。从前，世界因为有罪，处在法律的审判之下；但现在，圣言亲自承担了审判，并在肉身中为众人受苦，将救恩赐给所有人。有鉴于此，若望大声说：\BibleQuote{法律是藉梅瑟传授的，恩宠和真理却是由耶稣基督而来的}\BibleRef{若 1:17}。恩宠比法律更美善，真理也比影像更美善。\par

61.（10.）所以，正如已经说过的，\Quote{更美善}的事，除了坐在\OriginalTerm{圣父}{Father}右边的\OriginalTerm{圣子}{Son}之外，没有任何其他能够成就。这除了表明圣子的真实无伪，以及圣父的\OriginalTerm{天主性}{Godhead}与圣子的相同之外，还能表明什么呢？因为圣子在圣父的国度里为王，与圣父同坐一个宝座，并在圣父的天主性内被瞻仰，所以\OriginalTerm{圣言}{Word}就是天主，凡看见圣子的，就是看见了圣父；因此只有一个天主。这样，祂坐在右边，却并未将圣父放在左边；凡是圣父内一切宝贵和美好的，圣子也同样拥有，正如祂所说：\BibleQuote{凡父所有的一切，都是我的}\BibleRef{若 16:15}。因此，圣子虽然坐在右边，却也看见圣父在右边，尽管祂是以成了人的身份说：\Quote{我常将主置于我眼前，因为祂在我右边，我决不动摇。}这更表明圣子在圣父内，圣父在圣子内；因为圣父在右边，圣子也在右边；当圣子坐在圣父的右边时，圣父也在圣子内。天使上去下来服役；但论到圣子，经上说：\BibleQuote{天主所有的天使都要朝拜祂}\BibleRef{希 1:6}。天使服役时，会说：\Quote{我被派遣到你们这里}，以及\Quote{主命令了}；但圣子，尽管以人性方式说\Quote{我被派遣}，并且来完成那工程并服役，却仍以圣言和肖像的身份说：\Quote{我在父内，父在我内}；又说：\Quote{谁看见了我，就是看见了父}；以及\Quote{那住在我内的父，作祂自己的工程}；因为我们在那肖像内所看见的，乃是父的工程。\par

（11.）前面所说的，应当使那些与真理本身争辩的人感到羞赧；然而，如果因为经上记着\Quote{成了更美善的}，他们就拒绝把用于圣子的\Quote{成了}，理解为\Quote{过去如何，现今亦然}；也不愿像我们所说的，把它理解为那更美善的盟约已经成立；反而从这句话断定，圣言被称为受造的，那么，就让他们再听一遍简明的陈述，因为他们已经忘记了前面所说的话。\par

62. 如果圣子属于天使之列，那就让\Quote{成了}这个词用在祂身上，如同用在天使身上一样，让祂在本性上与天使毫无分别；要么让天使与祂同为子，要么让祂与天使一样成为天使；让天使与祂一同坐在圣父的右边，要么让圣子与众天使一同站立，如同服役的神体，被派遣去像他们一样服役。但是，如果保禄却将圣子与受造之物区分开来，说：\Quote{天主曾向哪一位天使说过：你是我的儿子？}，并且圣子造了天地，而天使是藉祂而受造的；圣子与圣父同坐，天使却侍立服役；谁还看不出，保禄没有把\Quote{成了}这个词用于圣言的性体，而是用于那通过祂而来的职务呢？因为，作为\Quote{圣言}，祂\Quote{成了血肉}；同样，当祂成了人时，祂在职务上变得比由天使而来的职务更美善得多，就如圣子远胜仆人，创造者远胜受造物一样。所以，让他们不要再把\Quote{成了}这个词用于圣子的本体，因为祂不属于受造之物；要让他们承认，这个词是指祂的职务和那已经实现的\OriginalTerm{救世计划}{Economy}。\par

(12.) 然而，祂如何在职务上成为更美，祂的本性本优于\OriginalTerm{受生}{originate}之物，这从前文所述已可看出，我认为这本身足以令他们羞愧。但如果他们继续争论，面对他们鲁莽的胆大妄为，与他们交锋并引用论及父自身的类似话语来反驳他们，这是恰当的。这或可使他们羞愧而勒住舌头不说邪恶，或可教导他们认识自己愚昧的深重。因为经上记着说：\BibleQuote{求你作我坚固的磐石，救我的堡垒，为拯救我。}又说：\BibleQuote{主成了受压迫者的保障}，以及圣经中类似的记载。如果他们将这几段话用于子，这或许最接近真理，那么让他们承认，圣作者祈求祂——并非受生的那一位——成为他们的\BibleQuote{坚固的磐石和堡垒}；而且今后，让他们将\Emph{成为}、\Emph{祂造}、\Emph{祂创造}等词语，理解为指祂降生成人的临在。因为当祂在木架上，以自己的身体承担我们的罪，并说：\BibleQuote{凡劳苦和负重担的，你们都到我这里来吧，我要使你们安息}\BibleRef{玛 11:28}时，祂就成为了\BibleQuote{坚固的磐石和堡垒}。\par

63. 但如果他们将这几段话用于父，当此处同样记着\Emph{成为}和\Emph{祂成了}时，他们胆敢竟然断言天主是受生的吗？是的，他们胆敢如此，正如他们论及祂的圣言时所争论的；因为他们论证的进程驱使他们，将论及圣言时所设想的，同样猜想到父身上。但让这样的观念远离所有信徒的思想吧！因为圣子既不在受生之物的行列，所讨论的圣经经文中的\Emph{成为}和\Emph{祂成了}，也不表示存在的开始，而是表示赐予穷乏者的援助。因为天主常存不改，始终如一；而人却是后来借着圣言而出现的，当父自己愿意时；并且天主对受生之物是不可见、不可接近的，尤其对地上的人。因此，当人在软弱中呼求祂，在迫害中求帮助，在受冤屈时祈祷，那不可见者，作为热爱人类者，便以祂的\OriginalTerm{恩惠}{beneficence}照耀他们，这恩惠是祂借着并在祂固有的圣言内施行的。于是，神圣的显现立即按每个人的需要而赐下，给软弱者健康，给受迫害者成为\Quote{避难所}和\Quote{堡垒}；对受冤屈者祂说：\BibleQuote{你呼喊时，我必说：我在这里}\BibleRef{依 58:9}。因此，凡借着圣子临到各人的护佑，各人就说天主已成为他自己的，因为这援助是来自天主自己，借着圣言。而且人类的习惯也认可这点，人人都会承认其恰当性。通常，人与人之间互相帮助；有人为受冤屈者劳苦，如\Person{亚巴郎}为\Person{罗特}；有人为受迫害者敞开家门，如\Person{敖巴狄雅}接待先知的门徒；有人款待旅客，如\Person{罗特}接待天使；有人周济穷乏者，如\Person{约伯}帮助那些向他乞求的人。那么，如果这些受益者中，有人这样说：\Quote{某人成了我的援助}，另一人说\Quote{也成了我的避难所}，第三个人说\Quote{成了某人的供应}，他们在这样说时，并不是在谈论施恩者的原始生成或本质，而是在谈论从他们那里临到自己的恩惠；同样，当圣人论及天主说\Emph{祂成了}和\Emph{求你成为}时，他们并不表示任何原始的生成，因为天主是无始的、\OriginalTerm{非受生}{unoriginate}的，而是表示从祂那里施于人的救恩。\par

64. 既如此理解，关于圣子的情况也类似，即无论多少次说诸如\Emph{祂成了}和\Emph{成为}之类的话，都应始终作同一理解：以致当我们听到所讨论的\Quote{成了比天使更美}和\Quote{祂成了}时，我们不应对圣言设想任何原始的生成，也绝不应从这些词语臆想祂是受生的；而应将\Person{保禄}的话理解为指祂成为人时的职务和\OriginalTerm{救世计划}{economy}。因为当\BibleQuote{圣言成了血肉，寄居在我们中间}\BibleRef{若 1:14}，并来服务及赐予众人救恩时，那时祂便成了我们的救恩，成了生命，成了赎罪祭；那时，祂为我们的救世计划便远远超越天使，祂成了道路，成了复活。正如\Quote{求你作我坚固的磐石}这话并不表示天主本身的本质有所生成，而是表示祂的\OriginalTerm{慈爱}{lovingkindness}，如已说过的；同样，此处\Quote{成了比天使更美}、\Quote{祂成了}、\Quote{耶稣成了更美\OriginalTerm{盟约}{covenant}的中保}这些词语，也并不表示圣言的本质是受生的（绝无此意！），而是表示借着祂成为人而临到我们的恩惠；尽管异端者忘恩负义，顽固坚持他们的不虔。\par

\SourceRecord{Translated by John Henry Newman and Archibald Robertson.；Nicene and Post-Nicene Fathers, Second Series；Vol. 4.；Edited by Philip Schaff and Henry Wace.；Buffalo, NY: Christian Literature Publishing Co.,；1892.；<http://www.newadvent.org/fathers/28161.htm>.}

% structure: p0001 source=h1 target=section reason=page-title

\section{驳斥亚略派第二讲}

1. 我确实认为，针对那些空洞宣扬亚略狂妄的人，已经说了足够多的话，无论是为了驳斥他们，抑或为了维护真理，以确保他们停止并忏悔自己对救主的邪恶念头与言论。然而，不管出于什么原因，他们依然不屈服；反倒像猪狗那样在自己所吐的秽物和自己所陷的泥沼中打滚，反而为自己不信虔诚的邪念想出新的花招。他们就这样误解了\BibleBook{}{箴言}中的话：\BibleQuote{上主自始即造了我，以作祂工程的开端，为祂事业之始}，以及宗徒的话：\BibleQuote{祂曾忠于创造祂者}\BibleRef{希 3:2}，并随即推论说，天主之子是一个工程和一个受造物。但是，虽然他们本可以从上述所论中学到——如果他们尚未完全丧失领会能力的话——圣子并非出自虚无，也绝非诸受造物之数，真理为此作证（因为，祂既是天主，就不能是一个工程，称祂为受造物是亵渎不敬的；而我们说\Quote{出自虚无}以及\Quote{它在受造之前不存在}这样的话，是针对受造物和工程的），然而，既然他们仿佛害怕放弃自己的虚构，惯于援引前述的圣经段落——这些段落本有美好的意义，却被他们滥用了——那么，让我们重新着手探讨这些段落的意义问题，以提醒忠信者，并从每段经文指明，他们对基督信仰根本一无所知。若非如此，他们就不会把自己封闭在现今犹太人的无信之中，而会去探询、学习到：既然\BibleQuote{在起初已有圣言，圣言与天主同在，圣言就是天主}，因此，当圣父以祂的美意使圣言成了人时，关于祂才说了这样的话，譬如由\Person{若望}说的：\BibleQuote{圣言成了血肉}\BibleRef{若 1:14}；由\Person{伯多禄}说的：\BibleQuote{天主已立祂为主，为基督了}\BibleRef{宗 2:36}——如同借\Person{撒罗满}以主自己的位格说的：\BibleQuote{上主自始即造了我，以作祂工程的开端，为祂事业之始}\BibleRef{箴 8:22}；由\Person{保禄}说的：\BibleQuote{祂所承受的名字既然超越众天使}\BibleRef{希 1:4}；又说：\BibleQuote{祂空虚了自己，取了奴仆的形体}\BibleRef{斐 2:7}；又说：\BibleQuote{因此，蒙受天上召选的圣善弟兄们，你们该留心思想那位我们宣认的宗徒和大司祭耶稣，祂曾忠于创造祂者}\BibleRef{希 3:2}。因为所有这些经文都有同样的力量和意义，一种虔诚的意义，宣示圣言的天主性，即使是那些按人而言论及祂的经文，也是就祂成了人子而言。然而，尽管这一分别足以驳斥他们，但既然他们出于对宗徒话语的误解（首先提及此事），因着经上写着\BibleQuote{祂曾忠于创造祂者}，而将天主圣言视为诸工程之一，我便认为有必要着手处理他们进一步的论点，像先前那样，来驳斥他们的陈述。\par

2. 如果祂不是\OriginalTerm{圣子}{Son}，就称祂为一个工程，凡关于工程所说的，都说到祂身上，而且不可单称祂为圣子，也不可称祂为\OriginalTerm{圣言}{Word}或\OriginalTerm{智慧}{Wisdom}；也不可称天主为圣父，只可称祂为万物的塑造者和创造者，万物都是借着祂而有的；让受造物成为祂创造意志的\OriginalTerm{肖像}{Image}和表现，并且按照他们的说法，让祂没有生育的本性，以致在祂的固有实体中，既无圣言，也无智慧，更没有肖像。因为如果祂不是圣子，那么祂也不是肖像。可是如果没有圣子，那么你怎么能说天主是一位创造者呢？既然万有都是借着圣言并在智慧内而有，没有这一位便一无所有，而你却说祂没有那位在祂内并借着祂创造万物的那位。若是天主性本质自身不结果实，而是贫瘠的，如同他们所主张的那样，好似不发光的明灯，枯干的泉源，难道他们不羞于宣称祂具有创造的能吗？他们既否认出于本性的，难道不脸红把出于意志的置于本性的之前吗？然而，祂若是意愿那些原不存在、自身之外的事物存在，因而成为它们的造作者，那么祂便更会首先成为出自其固有本质的后裔之父了。因为假如他们归给天主的是关于非存在事物的意志，为何不承认天主内那超越意志的呢？那超越意志的，就是祂凭本性而是，并是祂固有圣言的父。倘若那在先的、依本性而有的，如他们愚昧所想的，并不存在，那么那在后的、依意志而有的，又怎能出现呢？因为圣言在先，然后才是受造界。相反，无论那些不虔敬的人怎样断言，圣言是存在的；因为万有都是借着祂而受造，而天主，既是造作者，显然拥有祂的创造圣言，这圣言不在祂之外，而是祂自己固有的——这一点必须一再重申。如果祂拥有意志的能力，而祂的意志是有效的，足以使万物得以存在，且祂的圣言是有效的，是造作者，那么，那圣言必是圣父活的意志，是本质性的能，是真实的圣言，万物都在祂内得以维系，并被卓越地统辖。谁也无法怀疑，安排者先于安排本身及被安排的事物。因此，如我所说，天主的创造次于祂的生发：因为\Quote{子}意味着某种祂所固有、真实地出自那真福而永恒的性体者；而出于祂意志的，则从外部进入存在，并借着那出自这性体的固有后裔而塑造形成。\par

3. 我们已经指出，那些声称主不是天主子，而是\OriginalTerm{受造物}{work}的人，实在是犯了大错；因此，我们都必须宣认祂是圣子。而如果祂确是圣子——事实上祂正是——且圣子被宣认不是外在于父，而是由父而来，那么就让他们不要像我们先前所指出的，质疑圣经作者们论及圣言本身时所用的措辞，即不是\Quote{对生祂者}，而是\Quote{对造祂者}；既然祂的\OriginalTerm{本性}{nature}已获宣认，在这种情况下用词如何便无需质疑。因为措辞不会贬低祂的本性；相反，本性将其纳入自身并改变它们。因为措辞并不先于\OriginalTerm{本质}{essence}，而是本质在先，措辞在后。因此，当本质是受造物时，那么\Quote{祂造}、\Quote{祂成为}和\Quote{祂创造}这些词恰当地用于它，并指称受造物。但当本质是\OriginalTerm{后裔}{offspring}和圣子时，那么\Quote{祂造}、\Quote{祂成为}和\Quote{祂创造}便不再恰当地属于它，也不指称受造物；而是我们用\Quote{祂造}一词无争议地表示\Quote{祂生}。因此，父亲们常称由他们所生的儿子为仆人，却并不否认他们本性的真实；他们也常亲切地称自己的仆人为儿女，却并未忘记他们原本是买来的；因为他们以父亲的权威使用一种称呼，以慈爱之情使用另一种称呼。\Person{撒辣}称\Person{亚巴郎}为主，尽管她不是仆人而是妻子；宗徒将身为仆人的\Person{敖乃息摩}，作为兄弟与主人\Person{费肋孟}联合在一起；\Person{巴特舍巴}虽是母亲，却称自己的儿子为仆人，对他父亲说：\BibleQuote{你的仆人\Person{撒罗满}}\BibleRef{列上 1:19}；——随后\Person{纳堂}先知也进来，向\Person{达味}重复了她的话：\BibleQuote{\Person{撒罗满}你的仆人}。他们并不介意称儿子为仆人，因为\Person{达味}听见时，认识到\Quote{本性}，而他们说话时，并未忘记\Quote{真实性}，祈求让他成为他父亲的继承者，他们却称他仆人之名；因为对\Person{达味}来说，他按本性是儿子。\par

4. 因此，当我们读到这些时，我们公正地解释，并不因听到\Person{撒罗满}被称为仆人，就认为他是仆人，而是按本性、真实的儿子。同样，论及救主——祂被宣认确实是圣子，且按本性是圣言——若圣徒们说：\Quote{祂忠于造祂的主}，或祂论及自己说：\Quote{主创造了我}，以及\Quote{我是你的仆人，你婢女的儿子}等等，任何人都不可因此否认祂属于圣父并由父而来；相反，应像在\Person{撒罗满}和\Person{达味}的例子中那样，正确理解圣父与圣子。因为，如果他们听到\Person{撒罗满}被称为仆人，却承认他是儿子，那么那些人，不在主的例子中保持同样的解释，每当听到\Quote{后裔}、\Quote{圣言}和\Quote{智慧}等称谓，就强行曲解并否认圣子由圣父所生的、本性而真实的生，反倒在听到适用于\OriginalTerm{受造物}{work}的词语时，立即堕入认为祂本性是受造物的观念，并否认圣言；然而，既然祂已成了人，完全可以将所有这些词语归因于祂的人性——这样的人岂不该死？而且，他们岂不是被证明在上主眼中是\Quote{可憎恶的}，因为他们随身带着\BibleQuote{不同的砝码}\BibleRef{箴 20:23}，用一套衡量别的事例，用另一套亵渎主？但或许他们承认\Quote{仆人}一词是在某种理解下使用的，却强调\Quote{造祂的}作为他们异端的重大支柱。然而，他们这根支柱也只是折断的芦苇；因为他们若明白圣经的风格，就必立刻定自己的罪。因为正如\Person{撒罗满}虽是儿子，却被称为仆人，同样，重复上面所说，尽管父母称由自己所出的儿子为\Quote{被造的}、\Quote{被创造的}和\Quote{成为的}，他们并不因此否认儿子的本性。因此，\Person{希则克雅}——正如\WorkTitle{依撒意亚}所载——在祈祷中说：\Quote{从今日起，我必生养儿女，他们要宣扬你的正义，我救恩的天主。}他说了\Quote{我必生养}；但那位先知在同一卷书和\WorkTitle{列王纪下}中这样说：\BibleQuote{从你而出的儿子}\BibleRef{列下 20:18}。这里他用\Quote{生养}代替\Quote{生}，并称由他而出的后裔为\Quote{被造的}，没有人会质疑这词语是否指自然的后裔。又如，\Person{厄娃}生了\Person{加音}后说：\Quote{我赖上主获得了一个男子}；她也是用\Quote{获得}代替\Quote{生出}。因为，她先生下了孩子，然后却说\Quote{我获得了}。没有人会因她说\Quote{我获得了}，就认为\Person{加音}是从外面买来的，而非由她所生。此外，圣祖\Person{雅各伯}对\Person{若瑟}说：\Quote{如今，你在\Place{埃及}所生的两个儿子，\Person{厄弗辣因}和\Person{默纳舍}，在我来到\Place{埃及}你这里以前，他们是我的。}圣经论及\Person{约伯}说：\Quote{他有七个儿子和三个女儿。}同样，\Person{梅瑟}在\WorkTitle{法律书}中说：\Quote{若人有了儿子}，以及\Quote{若人生了儿子}。这里，他们再次将所生的说成\Quote{成了}和\Quote{造了}，因为他们知道，只要他们被承认是儿子，我们就不必质疑\Quote{他们成了}、\Quote{我获得了}或\Quote{我造了}这些说法。因为本性和真理将意义引向自身。\par

5. 既然如此，当有人询问主是否受造物或工程时，首先宜问他们：他是否是\OriginalTerm{圣子}{Son}和\OriginalTerm{圣言}{Word}和\OriginalTerm{智慧}{Wisdom}。因为若证明了这一点，关于工程与受造的猜测便立刻瓦解消亡。工程绝不可能是\OriginalTerm{圣子}{Son}与\OriginalTerm{圣言}{Word}；\OriginalTerm{圣子}{Son}也不可能是工程。再者，既然事实如此，为众人明显的是\Quote{对于造成他的那一位}这一句并不支持他们的异端，反成为其定罪。因为经上业已证明\Quote{他造了}这一说法即使对真正由亲生而来的子女也适用；由此，既然主被证明是\OriginalTerm{圣父}{Father}亲生的真实\OriginalTerm{圣子}{Son}、\OriginalTerm{圣言}{Word}和\OriginalTerm{智慧}{Wisdom}，纵使论及他时用了\Quote{他造了}或\Quote{他成了}，这并非视他为工程，而是圣人们毫不迟疑地使用这些说法——例如在\Person{撒罗满}和\Person{希则克雅}的子女身上。因为尽管父亲由自己生育了他们，经上仍写着：\Quote{我造了}、\Quote{我获得了}并\Quote{他成了}。所以，天主的仇敌虽然一再搬弄那些说法，但如今，纵使为时已晚，在所述之言后，应摒弃他们不敬的思想，而视主为\OriginalTerm{圣父}{Father}的真\OriginalTerm{圣子}{Son}、\OriginalTerm{圣言}{Word}和\OriginalTerm{智慧}{Wisdom}，而非工程、非受造物。因为倘若\OriginalTerm{圣子}{Son}是受造物，他自己又是凭何\OriginalTerm{圣言}{Word}、何\OriginalTerm{智慧}{Wisdom}造成的呢？原来万有都是藉著\OriginalTerm{圣言}{Word}和\OriginalTerm{智慧}{Wisdom}而造，如经上所载：\BibleQuote{你以智慧造了他们一切}，并\BibleQuote{万有是藉著他而造成的；没有他，就没有什么能够造成}。但他若就是那万有赖以生成的\OriginalTerm{圣言}{Word}和\OriginalTerm{智慧}{Wisdom}，便推论他不在工程的数目之内，简言之，不属于受造之物，而是\OriginalTerm{圣父}{Father}之\OriginalTerm{圣子}{Son}。\par

6. 试想，将天主的\OriginalTerm{圣言}{Word}称作工程，这是何等严重的错误。\Person{撒罗满}在\WorkTitle{训道篇}某处说：\BibleQuote{天主将把一切工程，连同一切隐密的事，无论善恶，都带入审判}\BibleRef{训 12:14}。因此，如果\OriginalTerm{圣言}{Word}是一项工程，难道你的意思是他也如其他事物一样将被带入审判吗？但若审判者受审，审判又有什么余地呢？那时主（按你们的假定）也与众人一同受审，谁来给予义人祝福，谁又给予不配之人惩罚呢？赐予律法者自身又要按照哪一法律受审呢？这些事本属于工程：受审、由\OriginalTerm{圣子}{Son}赐予祝福与惩罚。所以，现在当敬畏审判者，让\Person{撒罗满}的话说服你们。因为倘若天主将一切工程都带入审判，而\OriginalTerm{圣子}{Son}却不在受审之列，他反而是一切工程的审判者，那么，\OriginalTerm{圣子}{Son}并非工程，而是\OriginalTerm{圣父}{Father}的\OriginalTerm{圣言}{Word}，一切工程在他内得以生成并受审，这证明岂不比太阳更明亮吗？再者，如果\Quote{谁是忠信的}这句话令他们困惑，因他们认为\Quote{忠信}一词用于他如同用于别人，仿佛他凭信心而获得信德的赏报，那么他们势必也要指责\Person{梅瑟}说\Quote{天主是忠信且真实的}，并指责\Person{圣保禄}写道：\BibleQuote{天主是忠信的，他决不许你们受那超过你们能力的试探}\BibleRef{格前 10:13}了。但圣人们这样说时，并非以人的方式思考天主，而是承认圣经中\Quote{忠信}一词有两种含义：一是\Quote{信}，二是\Quote{可信赖}；前者属于人，后者属于天主。所以\Person{亚巴郎}有信德，是因为他信了天主的话；而天主是忠信的，因为如同\Person{达味}在\WorkTitle{圣咏}中所说：\BibleQuote{主的一切言语都是忠信的}（或译为可信赖的），他绝不说谎。再说，\BibleQuote{若信主的妇人中有寡妇}\BibleRef{弟前 5:16}，她被称为信主的，是因她正当的信德；而\Quote{这话是可信的}，则是因为他所言配得上我们的信德，真实无伪，绝无别样。因此，\Quote{他对那造成他的是忠信的}这句话，并非暗示他与别人并列，也不意味着他因有信德而获得悦纳；而是说，身为真天主的\OriginalTerm{圣子}{Son}，他本身亦是忠信的，凡他所说所行都当被相信，而在他为人的\OriginalTerm{经世计划}{Economy}与血肉的\OriginalTerm{临在}{presence}中，他自己始终保持不变易、不更改。\par

因此，这样我们就可以应对这些厚颜无耻的人了；从\Quote{造了}这单单一个词上，就可以指出他们错以为天主的\OriginalTerm{圣言}{Logos}是一件受造之物。进一步，既然上下文的趋向也是正统的，指明了这个说法所指向的时间和关系，我也应当由此指出异端者如何缺乏理智；也就是说，要像我们上面所做的那样，考虑它是何时使用的，为了什么目的。如今，宗徒如此说话时，不是在讨论创造之前的事，而是讨论当\BibleQuote{圣言成了血肉}的时候；因为经上这样记载：\BibleQuote{因此，同蒙天上召选的圣善弟兄们，应想想我们宣认的宗徒和大司祭耶稣，祂对那造了祂的忠心可信。}那么，祂何时成了\Quote{宗徒}呢？不正是当祂穿上我们的血肉之时吗？祂何时成了\Quote{我们宣认的大司祭}呢？不正是当祂为我们祭献了自己，使祂的身体从死者中复活，并如今亲自使那些怀着信德接近祂的人靠近并奉献于圣父，救赎众人，并为众人平息天主的义怒之时吗？宗徒说\Quote{祂对那造了祂的忠心可信}，并非愿意指明圣言的本体，也不是指明祂由圣父的受生——（这种想法绝不应有！因为圣言不是被造的，而是创造者）——而是为了指明祂降来到人间、以及祂\Quote{成为}的大司祭职，正如人可容易地从关于法律和\Person{亚郎}的记载中看到的。我是说，\Person{亚郎}并非生来就是大司祭，而是一个人；随着时间进程，当天主愿意时，他才成了大司祭；然而他并非简单地成为大司祭，也不是穿著平常的衣裳来指示这点，而是在这些衣裳之上，穿戴上了厄弗得、胸牌\BibleRef{出 29:5}，长袍——这长袍是妇女们按天主的命令织成的——身穿这些走进圣所，为人民奉献祭献；他就这样借着这些，如同在天主的显现与人的祭献之间作中保。因此，主也是如此：\BibleQuote{在起初已有圣言，圣言与天主同在，圣言就是天主}；但是，当圣父愿意为众人并为了一切支付赎价，赐下恩宠时，圣言确实，如同\Person{亚郎}穿上他的长袍那样，也穿上了尘世的血肉，以\Person{玛利亚}为祂身体的母亲，犹如未开垦的童贞之土一样，好使祂作为大司祭，如同其他人那样有所奉献，能将祂自己祭献给圣父，并以自己的血洗净我们众人的罪，且从死者中复活。\par

因为古时发生的事是这事的影子；而救主在降来时所行的，就是\Person{亚郎}按照法律所预表的。那么，\Person{亚郎}并未因穿上大司祭的礼服而改变自己，仍是原先的\Person{亚郎}，只是穿上了礼服；因此，若有人看见他献祭，并说：\Quote{看，\Person{亚郎}今天成了大司祭}，这人并非暗示\Person{亚郎}那时才生而为人，因为他在成为大司祭之前已经是一个人，而是说他在职务上成了大司祭，穿上了为司祭职所制作和预备的礼服。同样，对于主的情形，也可能有正确的理解：祂取了血肉，并未变成与祂自己不同的一位，而是和先前一样，只是穿上了血肉；而\Quote{祂成为}和\Quote{祂被造}这些说法，不能理解为，圣言就其作为圣言而言是被造的，而是说，圣言本是万有的创造者，后来才成了大司祭，借着穿上了那受造且被制成的身体，这身体是祂能为我们奉献的；因此祂被称为\Quote{被造}。果真如此，如果主没有成为人，那是亚流派争辩的事；但若\BibleQuote{圣言成了血肉}，那么论到祂成为人时，除了说\Quote{祂对那造了祂的忠心可信}之外，还能说什么呢？因为对于圣言，论到祂宜说\BibleQuote{在起初已有圣言}；对于人，就宜说\Quote{成为}和\Quote{被造}。有谁看见主像人一样行走，却又从祂的作为显出是天主，会不问：谁使祂成了人？面对这样的问题，又有谁不会回答说，是圣父使祂成了人，并派遣祂到我们这里作为大司祭？而这意义、时机和特性，乃是那写下\Quote{祂对那造了祂的忠心可信}这话的宗徒，只要我们留意这些话的上下文，就能最好地向我们说明的。因为有一个连贯的思路，整段读经都是关于同一位。他在\BibleBook{希伯来}{书}中这样写道：\BibleQuote{那么，既然孩子们有血有肉，祂也同样取了血肉，为能借着死亡，毁灭那握有死亡权势的——就是魔鬼，并解救那些因死亡的恐怖，一生当奴隶的人。其实并没有援助天使，而是援助了\Person{亚巴郎}的后裔。因此，祂应当在各方面相似弟兄们，好在关于天主的事上，成为仁慈和忠信的大司祭，以补赎人民的罪。祂既然亲自经过试探受了苦，也必能扶助受试探的人。所以，同蒙天上召选的圣善弟兄们！你们应细心想想，我们公认的宗徒和大司祭耶稣；祂对那造了祂的忠心可信。}\par

9. 谁能读完这整段经文，而不谴责\OriginalTerm{亚略派}{Arians}，且不钦佩那位说话得体的真福\OriginalTerm{宗徒}{Apostle}呢？因为，基督在什么时候被\Quote{立}，又在什么时候成为\Quote{宗徒}，若不祂像我们一样\BibleQuote{分享了血肉}呢？而祂又是什么时候成为\BibleQuote{慈悲而忠信的大司祭}，若不\BibleQuote{在各方面都与祂的弟兄相似}呢？而那时祂\Quote{被造为相似}，就是当祂成为人，穿上我们的肉躯时。因此，\Person{保禄}论及\OriginalTerm{圣言}{Word}在人间的\OriginalTerm{经世}{Economy}时，说：\Quote{祂对委任祂的那位是忠信的}，而不是论及祂的\OriginalTerm{本体}{Essence}。所以，不要再胡言乱语，说天主\OriginalTerm{圣言}{Word}是工程；祂因性体是\OriginalTerm{独生子}{Only-begotten}\OriginalTerm{圣子}{Son}，而后来祂有了\Quote{弟兄}，是在祂穿上和我们一样的肉躯之后；此外，因祂自献己身，祂被命名且成为\Quote{慈悲而忠信的}——慈悲，是因祂对我们大发慈悲，为我们奉献了自己；忠信，并非祂与我们分享信德，也非祂像我们一样对某人有信心，而是祂所言所行配得众人信赖，且祂奉献了忠信的\OriginalTerm{祭献}{sacrifice}，这祭献常存不废。因为按照法律所献的祭品，并无这种忠信，它们随日子过去而消逝，且需要再次洁净；但\OriginalTerm{救主}{Saviour}的\OriginalTerm{祭献}{sacrifice}，一次完成，便成全了万有，且因永远常存而成为忠信的。\Person{亚郎}有继承者，总之，法律下的\OriginalTerm{司祭职}{priesthood}随时间和死亡而更替最初的任职者；但\OriginalTerm{主}{Lord}拥有不经转换、无继承的\OriginalTerm{大司祭}{High Priest}职，祂成了\Quote{忠信的大司祭}，因为祂永远长存；且祂因许诺也是忠信的，好能俯听而不误导那些前来投靠祂的人。这也可以从大\Person{伯多禄}的书信中学得，他说：\BibleQuote{愿照天主圣意受苦的人，将他们的灵魂托付给忠信的创造者}\OriginalTerm{创造者}{Creator}\BibleRef{伯前 4:19}。因为祂是忠信的，祂不改变，永远常存，并履行祂的许诺。\par

10. 至于希腊人所谓的诸神，名不副实，无论就其本质或其所许诺，都不忠信；因为它们并非处处皆在，相反，地方性的神明随时间过去而归于虚无，且经历自然的分解；为此，\OriginalTerm{圣言}{Word}大声斥责他们，说\Quote{他们的信靠不坚固}，他们是\Quote{干涸的河床}，\Quote{在他们内毫无信实}。但万有的天主，是独一、实在且真实的，是忠信的，祂永远一样，且说：\BibleQuote{你们要认清，我就是那位，我就是我，}\BibleQuote{我决不改变}；因此，祂的\OriginalTerm{圣子}{Son}也是\Quote{忠信的}，永远一样，永不改变，无论在本质或应许上都绝不欺蒙——正如\OriginalTerm{宗徒}{Apostle}致书\Place{得撒洛尼}人时又说：\BibleQuote{召叫你们的那位是忠信的，祂必实行此事}\BibleRef{得前 5:24}；因为祂履行所许，\Quote{祂对自己的话语是忠信的}。而他也曾论及这词的意义为\Quote{不可改变的}，致书\Place{希伯来人}如此说：\BibleQuote{我们即使不信，祂仍是忠信的，祂不能否认自己}\BibleRef{弟后 2:13}。因此，\OriginalTerm{宗徒}{Apostle}论及\OriginalTerm{圣言}{Word}的有形体临在时，理所当然地说祂是\Quote{\OriginalTerm{宗徒}{Apostle}，并对委任祂的那位忠信的}，向我们表明，即使在祂成为人时，\BibleBook{耶稣}{基督}仍是\BibleQuote{昨日、今日、直到永远，是一样的}\BibleRef{希 13:8}，是永不改变的。正如\OriginalTerm{宗徒}{Apostle}在信中论及祂的\OriginalTerm{大司祭职}{High Priesthood}时提到祂成为人，他同样也没有长久沉默关于祂的\OriginalTerm{天主性}{Godhead}，反而立即提到，从各方面给我们提供保障，尤其在他论及祂的谦卑时，好使我们能立刻认出祂的崇高以及属于圣父的尊威。例如，他说：\BibleQuote{\Person{梅瑟}为仆，而基督为子}\BibleRef{希 3:5}；前一人\Quote{在祂的家中尽忠}，后一人\Quote{却是在家之上}，因为祂亲自建造了这家，且是它的主和\OriginalTerm{塑造者}{Framer}，并作为天主圣化着它。因为\Person{梅瑟}，按本性是人，因着信赖藉自己的圣言向祂说话的天主而成为忠信的；但\OriginalTerm{圣言}{Word}并非像由身体所出的受造物那样，也非受造物中的受造物，而是作为\OriginalTerm{天主}{God}在肉躯内，是万有的\OriginalTerm{塑造者}{Framer}，且是在祂所建造的家中建造的那一位。而人穿上肉躯是为了存在并存活；但天主\OriginalTerm{圣言}{Word}成为人，是为圣化肉躯，且祂虽是\OriginalTerm{主}{Lord}，却取了\OriginalTerm{仆人的形象}{form of a servant}；因为整个受造界都是\OriginalTerm{圣言}{Word}的仆人，是藉着祂而得以存在并被造成的。\par

11. 由此可知，\OriginalTerm{宗徒}{Apostle}所用的\Quote{祂立了}一词，并不证明\OriginalTerm{圣言}{Word}被造，而是证明那与我们相同的肉躯是祂所取的；因此祂既成了人，就称为我们的弟兄。但即使已证明，纵然\Quote{造成}一词可应用于\OriginalTerm{圣言本身}{Very Word}，它也是作\Quote{生}解，那么，现在既然讨论已从各方面将此词意义清楚阐明，并指出\OriginalTerm{圣子}{Son}并非\OriginalTerm{工程}{work}，在\OriginalTerm{本体}{Essence}上确系\OriginalTerm{圣父}{Father}的苗裔，而在\OriginalTerm{经世}{Economy}中，按\OriginalTerm{圣父}{Father}的\OriginalTerm{旨意}{good pleasure}，祂为我们而被造，且成了人，他们还能想出什么邪恶的伎俩呢？为此之故，\OriginalTerm{宗徒}{Apostle}才说：\Quote{祂对委任祂的那位是忠信的}；而在\BibleBook{箴言}{}中，甚至提及\OriginalTerm{创造}{creation}。因为只要我们承认祂成了人，那么如先前所述，无论说\Quote{祂成为}、\Quote{祂被造}、\Quote{受造}、\Quote{被形成}、\Quote{仆人}、\Quote{婢女之子}、\Quote{人子}、\Quote{被设立}、\Quote{启程}、\Quote{新郎}、\Quote{侄子}或\Quote{弟兄}，都毫无问题。所有这些词汇恰当地属于人的构成；这些词汇所指的，并非\OriginalTerm{圣言}{Word}的\OriginalTerm{本体}{Essence}，而是祂成了人。\par

% structure: p0013 source=h2 target=subsection reason=relative-heading-rank; base=h2

\subsection{第十五章 经文解释；第五点，\BibleRef{宗 2:36}。必须遵守\OriginalTerm{信德准则}{Regula Fidei}；\Quote{造成}一词适用于主的\OriginalTerm{人性}{manhood}；也适用于祂的显现；以及祂对我们所有的职分；且与犹太人相关。\BibleRef{创 27:29, 37}有平行的例子。上下文反驳了\OriginalTerm{亚略派}{Arian}的解释。}

11（\Emph{续}）。他们引用的\BibleBook{宗徒}{大事录}篇章，其意义也是如此；在那篇里，\Person{伯多禄}说：\BibleQuote{天主已把你们所钉死的这位耶稣，立为主，立为基督了。}因为在这里也没有写着\BibleQuote{祂为自己造了一个儿子}，或\BibleQuote{祂使自己成了圣言}，好让他们有这种想法。那么，如果他们还没有忘记，他们所谈论的是关于天主子，就让他们查考，在什么地方写着\BibleQuote{天主为自己造了一个儿子}，或\BibleQuote{祂为自己创造了一个圣言}；或者，再查考，什么地方明确地写着\BibleQuote{圣言是工程或受造物}；然后，让他们再去论证他们的主张吧，这些无知的人，好让他们在这里也得到回答。但如果他们提不出任何这类证据，而只是抓住像\BibleQuote{祂造了}和\BibleQuote{祂被造成了}这类零散的表述，那么我担心，他们听了\BibleQuote{在起初天主造了天地}、\BibleQuote{祂造了日月}、\BibleQuote{祂造了海}这些话，迟早会把圣言称为天、把第一天出现的光、地以及各个具体受造物都称为圣言，最终变得像所谓\OriginalTerm{斯多葛派}{Stoics}那样：一派把他们的神延伸进一切事物，另一派把天主的圣言与每一样具体工程等同起来；而他们几乎已经这样做了，因为他们说圣言是祂的工程之一。\par

12. 但在这里他们必须得到与前面相同的答案，首先要告诉他们，圣言是子，如上所述，而不是工程，并且这类用语不应理解为指祂的天主性，而应探究其理由和方式。对于那些如此探究的人，祂为了我们而承担的人性\OriginalTerm{安排}{Economy}就会清楚呈现。因为\Person{伯多禄}在说了\BibleQuote{祂立了主和基督}之后，立刻又说\BibleQuote{就是你们所钉死的耶稣}；这清楚地向所有人，甚至向他们，表明——只要他们注意上下文——不是圣言的本质，而是按祂的人性而言，祂被造成了。因为被钉死的是什么，不就是身体吗？除了用\BibleQuote{祂造了}的说法，又怎能表达圣言身上属肉体的部分呢？尤其是，\BibleQuote{祂造了}这个说法，具有与正统信仰一致的意义；因为正如我在前面所观察到的，他并没有说\BibleQuote{祂使祂成为圣言}，而是说\BibleQuote{祂使祂成为主}，并且不是泛指，而是\Quote{向着}我们，并\Quote{在}我们中间，意思相当于\Quote{祂彰显了祂}。而且，\Person{伯多禄}本人在开始他这第一要理讲授时，就小心地表达了这一点，当时他对他们说：\BibleQuote{你们以色列人，请听这些话：纳匝肋人耶稣，是天主用德能、奇迹和征兆——即天主藉祂在你们中间所行的，一如你们自己所知道的——向你们所证明的}\BibleRef{宗 2:22}。因此，他在结尾所用的\BibleQuote{造了}一词，他在开头就用了\Quote{彰显了}来解释，因为主所行的征兆和奇迹，彰显了祂不仅仅是人，而是以身体显出的天主，也是主，是基督。\BibleBook{若望}{福音}中的那一段也是如此：\BibleQuote{为此犹太人越发想要杀害祂，因为祂不但犯了安息日，而且又称天主是自己的父，使自己与天主平等。}因为主当时并没有把自己塑造成天主，事实上，\Quote{被造的天主}是无法想象的，而是祂藉着工程彰显了这一点，说：\BibleQuote{你们纵然不信我，也该信我的工程，好使你们知道并认识：父在我内，我在父内。}因此，父就这样\Quote{使}祂在我们中间，并为我们这些曾经悖逆的人，成为主和君王；很明显，那现在显为主和君王的那一位，并不是在那时才开始了为王为主，而是开始彰显祂的\OriginalTerm{主位}{Lordship}，并将其扩展到悖逆的人身上。\par

13. 因此，如果他们以为救主甚至在成为人和忍受十字架之前，也不是主和君王，而是那时才开始成为主，那么让他们知道，他们是公然在重提\Person{撒摩撒塔的保禄}的言论。但如果我们如上所引证和宣告的，祂是永恒的主和君王，因为\Person{亚巴郎}以祂为主而敬拜祂，并且\Person{梅瑟}说：\BibleQuote{当时，上主使硫磺和火，从天上上主那里，降于\Place{索多玛}和\Place{哈摩辣}}\BibleRef{创 19:24}；\Person{达味}在\BibleBook{}{圣咏}中则说：\BibleQuote{上主对我主说：坐在我右边}；又说：\BibleQuote{天主，你的宝座是永永远远的；你王国的权杖是正义的权杖}；又说：\BibleQuote{你的王国是万代的王国}；很明显，在祂成为人之前，祂就是永恒的君王和上主，是父的\OriginalTerm{肖像}{Image}和\OriginalTerm{圣言}{Word}。圣言既然是永恒的上主和君王，那么同样非常清楚的是，\Person{伯多禄}并没有说子的本质被造了，而是说祂对我们的\OriginalTerm{主位}{Lordship}，当祂成为人，并藉着十字架救赎一切时，成为了万有的上主和君王。但如果他们继续以经上写着\BibleQuote{祂造了}为依据，而不愿将\BibleQuote{祂造了}理解为\Quote{祂彰显了}的意思，无论是由于缺乏理解，还是出于他们反对基督的意图，那么让他们聆听对\Person{伯多禄}话语的另一种健全解释。因为，如果一个人成为别人的主，他是在取得对已经存在的存在者的拥有；然而，如果上主是万物的\OriginalTerm{造物主}{Framer}和永恒的君王，并且在祂成为人时，才取得了对我们拥有，那么，这里也看出，\Person{伯多禄}的话语显然不是表示圣言的本质是一项工程，而是指万有的后来顺服，以及救主对一切所拥有的主位。这与我们前面所说的一致；因为当时我们引用了这些话：\BibleQuote{求你做我的天主和保障}，\BibleQuote{上主成了受压迫者的避难所}，并且理所当然的是，这些表达并不表示天主是受造物，而是表示祂对每一个人的\OriginalTerm{恩惠}{beneficence}\Quote{成为}了；\Person{伯多禄}的表达也具有同样的意义。\par

14. 因为天主子自己就是\OriginalTerm{圣言}{Logos}，是万有的主；而我们曾一度从起初就屈服于败坏的奴役和法律的诅咒，然后逐渐为自己塑造了不存在之物，正如真福宗徒所说，我们侍奉了\BibleQuote{他们本性不是神}\BibleRef{迦 4:8}，因不认识真天主，我们选择了虚无之物而不是真理。但后来，如同古代人民在埃及受压迫时呻吟一样，当我们也有了那\BibleQuote{嫁接}\BibleRef{雅 1:21}的律法在我们内，并按照圣神那无法言喻的叹息\BibleRef{罗 8:26}献上我们的转求：\BibleQuote{主我们的天主，求你拥有我们}，那时，正如\BibleQuote{他成了避难所}和\BibleQuote{天主与保障}一样，他也成了我们的主。他并非那时才开始存在，而是我们开始以他为主。因此，天主既是良善的，又是主之父，出于怜悯，并切愿被众人认识，便使自己的圣子穿戴上人的身体，成为人，并被称为\Person{耶稣}，好使他在这身体内将自己献为众人牺牲，拯救众人脱离虚假的崇拜和败坏，并使自己成为万有的主和君王。所以他以这种方式成为主和君王，这正是\Person{伯多禄}所说\BibleQuote{他立他为主}与\BibleQuote{派遣基督}的意思。这等于说，父在使他成为人时（因为\Quote{成为}属于人），不是简单地使他成为人，而是使他成为人，好让他成为所有人的主，并藉着傅油圣化一切。因为圣言虽具有天主的形体，却取了奴仆的形体，然而，采用人性并未使那本性是主的圣言成为奴仆；相反，不仅整个人类由圣言获得解放，而且那本性是主、当时成为人的圣言，藉着奴仆的形体被立为万有的主和基督，也就是为藉着圣神圣化万有。就如天主，当他\BibleQuote{成为天主与保障}并说\BibleQuote{我要作他们的天主}时，他并非那时比以往更成为天主，也并非那时才开始成为天主，而是他自永远所是者，在蒙他喜悦时，为了那些需要他的人而成为那样的；同样，本性是主和永恒君王的基督，在他被派遣时，并没有比以前更成为主，也没有那时才开始成为主和君王，而是他永远所是者，那时根据肉体被立为（主和君王）；并且，在救赎了众人之后，他由此再次成为生者和死者的主。从此，万物都服事他，这就是\Person{达味}在圣咏中所说的意思：\BibleQuote{主对我主说：你坐在我右边，等我使你的仇敌作你的脚凳}。因为通过那本性是主的那一位完成救赎是合宜的，以免我们虽由圣子所造，却称呼另一位为主，而陷入\OriginalTerm{亚略派}{Arian}和\OriginalTerm{希腊人}{Greek}的愚妄，侍奉受造物而不侍奉那创造万有的天主。\par

15. 这至少根据我的卑微理解，是这段经文的含义；而且，\Person{伯多禄}这番话对犹太人而言具有真实而美好的意义。因为犹太人偏离了真理，固然期待基督来临，却不承认他承受苦难，他们说着自己不明白的话：\BibleQuote{我们知道当基督来临时，他永远存在，你怎么说人子必须被举起来呢？}其次，他们设想他不是那降生成人的圣言，而仅仅是一个凡人，如同所有的君王一样。于是，主训诫\Person{克罗帕}和另一个人，教导他们基督必须先受苦；并对其他犹太人表明天主已经来到他们中间，说：\BibleQuote{如果那些承受天主话的人，圣经尚且称他们为神，而圣经是不能废弃的，那么，对于父所圣化并派遣到世界上来的那一位，你们却说：你亵渎了，因为我说过：我是天主子？}\BibleRef{若 10:36}\par

16. \Person{伯多禄}已从救主那里学到了这一点，便在两个方面纠正了犹太人，说：\Quote{诸位犹太人！神圣的圣经宣告基督要来，你们却视祂为出自\Person{达味}后裔的纯粹凡人，然而关于祂所记载的，表明祂并非如你们所说的那样，反而宣告祂是主、天主、不死者及生命的施予者。因为\Person{梅瑟}曾说过：\BibleQuote{你的生命悬于你眼前。}\Person{达味}在第一百零九篇圣咏中说：\BibleQuote{上主对我主说：你坐在我右边，等我使你的仇敌作你的脚凳。}又在第十五篇中说：\BibleQuote{你绝不会将我遗弃在阴府，也不让你的圣者见到腐朽。}这些章节并非以\Person{达味}为指归，他亲自作证，承认那位要来的正是他自己的主。况且你们自己也晓得他已死了，他的遗骸尚在你们那里。因此，基督必定如圣经所言的那一位，你们自己也会公开承认。因为那些宣告出自天主，其中不可能有虚假。倘若你们能说，先前曾有如此一位来临，并能凭他所行的奇迹异事证明他是天主，你们便有理据坚持争辩；但若你们无法证明祂已来临，且仍在期盼这样的一位，就当由\Person{达尼尔}先知认识真正的时期，因为他的话正关涉现今的时刻。可是，若这现今的时期正是古时所预言的，而你们也已看见发生在我们中间的事，就当确信这位\Person{耶稣}，就是你们所钉死的，正是所期待的基督。因为\Person{达味}与众先知都已死去，他们所有人的坟墓都在你们那里，但现今发生的复活，已显明这些章节的指归正是\Person{耶稣}。因为钉十字架由\BibleQuote{你的生命悬着}所表示，肋旁被枪刺伤，应验了\BibleQuote{他如同被牵去宰杀的羊}\BibleRef{依 53:7}，而复活，更且是远古死者从坟墓中起来（这些你们大半都亲眼见过），这正是\BibleQuote{你绝不将我遗弃在阴府}，以及\BibleQuote{他永远吞噬了死亡}\BibleRef{依 25:8}，又说：\BibleQuote{天主要擦去……}因为实际发生的奇迹，表明那在肉身内的是天主，也是生命和死亡的主。因为基督在赐予别人生命时，理当自己不被死亡束缚；但若祂如你们所想的仅是一个凡人，这事便不能发生。其实，祂是天主圣子，因为众人皆隶属于死亡。所以，谁也不可疑惑，但愿以色列全家确实知道：这位\Person{耶稣}，你们曾见祂具有人的形体，行过从未有人行过的奇迹与工程，祂就是基督，是万有的主。因为祂虽成为人，被称为\Person{耶稣}，如我们从前所说，却并未因那人的受苦而受损，反而在成为人时，彰显出祂是生死的主。既然如宗徒所言：\BibleQuote{因为世人没有凭自己的智慧认识天主，天主遂以自己的智慧，决意以愚妄的宣讲来拯救那些相信的人。}\BibleRef{格前 1:21}因此，既然我们人类不愿藉祂的圣言认识天主，也不愿事奉我们的天然主——天主圣言，天主就乐意在人身上显示自己的主权，如此吸引众人归向祂。但藉一个纯粹的人行这事并不相宜，免得我们以人为我们的主，而成为人的敬拜者。所以，圣言亲自成了血肉，圣父给祂起名叫\Person{耶稣}，这样\OriginalTerm{立}{made}祂为主、为基督，无异于说：祂立祂为王，治理一切；好使众人因\Person{耶稣}的名字——就是你们所钉死的——无不屈膝，我们得以承认子、并藉着子承认父，为主与君王。}\par

17. 犹太人听了这些话，大半醒悟过来，当即承认了基督，正如\BibleBook{宗徒}{大事录}所记载的。然而，那些\Person{阿里乌}派狂热者却宁愿继续作犹太人，与\Person{伯多禄}争辩；因此，让我们把一些类似的语句摆在他们面前；也许会对他们产生些影响，使他们认清圣经的通用语式。基督是永远的主与君王，已由前文的论述显明，无人可对此怀疑；因为祂是天主圣子，必定与父相似，既相似，就定然是主、是君王，因为祂亲口说过：\BibleQuote{谁看见了我，就是看见了父。}另一方面，\Person{伯多禄}的这句话——\BibleQuote{天主已立祂为主，为基督了}——并不意味子是受造物，这可从\Person{依撒格}对\Person{雅各伯}的祝福中看出，尽管这个例子对我们的主题而言只是一个微弱的比喻。他对\Person{雅各伯}说：\BibleQuote{愿你作你兄弟的主}；又对\Person{厄撒乌}说：\BibleQuote{看，我已立他作你的主}。然而，就算\OriginalTerm{立}{made}这个字意味着\Person{雅各伯}的本质和产生，即便如此，他们也不可将之设想为同样适用于天主圣言，因为天主圣子并非如同\Person{雅各伯}一样的受造物；况且，他们本可查问而摆脱那种狂妄。若他们不把这理解成他的本质，也不理解成他的产生——虽然\Person{雅各伯}按本性是受造物与工程——那么，难道他们的疯狂比魔鬼更甚吗？因为他们在类似用语上连对于本性受造之物都不敢做的事，竟加于天主圣子，宣称祂是受造物！因为\Person{依撒格}说\Quote{愿你成为}和\Quote{我已立了}，既不指\Person{雅各伯}的产生，也不指他的本质（因为这话是他出生三十多年后说的）；而是指他后来实行的对兄弟的权柄。\par

18. 因此，\Person{伯多禄}这样说，更不意味着圣言的\OriginalTerm{性体}{Essence}是一项工程；因为他认识他是天主的圣子，宣认道：\BibleQuote{你是基督，永生天主之子}\BibleRef{玛 16:16}；但他所说的乃是他的国度和主权，那是按照恩宠而形成并出现的，且是相对于我们而言的。因为他在说这话时，并没有缄默不言天主子那永远的神性，亦即圣父的神性；但他先前已经说过，他已经把圣神倾注在我们身上了。如今，以权柄赐予圣神，并非受造物或工程的能力所能做到的，而圣神是天主的恩赐。因为受造物是由圣神所圣化的；而圣子，既然不是由圣神所圣化，相反，他本身就是将圣神赐予众人的施予者，因此不是受造物，而是圣父真正的圣子。然而，那赐予圣神者，同样也被称为被建立者；也就是说，因着他的人性，他在我们中间被立为主，而同时因他是天主的圣言，他又赐予圣神。因为他永远存在、且现在仍是，作为圣子，同样也是万有的主和主宰，在万有上都相似于圣父，并拥有圣父所有的一切，正如他自己所说过的。\par

% structure: p0022 source=h2 target=subsection reason=relative-heading-rank; base=h2

\subsection{第十六章 引论：\BibleBook{箴言}{}8:22，圣子不是受造物。亚略派公式：一个受造物却不像受造物中的一个；但每个受造物都与其他一切受造物不同；且没有受造物能创造。因此圣言不同于所有受造物之处，正是在于它们虽在其他方面彼此相异，但作为受造物却共同一致的那一方面；就是说，在于他是\OriginalTerm{动力因}{efficient cause}；在于他是创造中的唯一中介或\OriginalTerm{工具性动因}{instrumental agent}；此外，在于他是圣父的启示者；并在于他是受敬拜的对象。}

18. \Emph{续前}。接下来，让我们思考\BibleBook{箴言}{}中的经文：\BibleQuote{上主创造了我，作他行动的起始，作他作为的开端}；尽管在证明圣言不是工程时，也就证明了圣言不是受造物。因为说工程或受造物是同一个意思，所以证明他不是工程也就是证明他不是受造物。然而，人们可能会对这些人大感惊奇，他们如此设法为不虔敬找借口，并且在每一个论点上都遭到反驳时，却毫不畏缩。因为起初他们用这些问题欺骗头脑简单的人：\Quote{他是从无中造出那非有的，还是那有的？}以及\Quote{你在生子之前就有子了吗？}当这些问题被证明毫无价值之后，他们又提出了新的问题：\Quote{\OriginalTerm{非受生者}{Unoriginate}是一位还是两位？}然后，当他们在这方面被驳倒后，紧接着又另造一个问题：\Quote{他拥有\OriginalTerm{自由意志}{free-will}和可改变的本性吗？}但在被迫放弃这个问题之后，他们接下来又开始说：\Quote{他被造得远胜于天使}；而当真理揭露了这个伪装之后，如今他们又把这些全都收罗在一起，想用\Quote{工程}和\Quote{受造物}来推荐他们的异端。因为他们无非是在重复同样的意思，并且忠实地坚持他们自己的邪恶，将同样的错误改换成各种形式，翻来覆去，以便借这种多变欺骗一些人。尽管上文已经充分证明他们这鲁莽的伎俩，可是，既然他们到处宣扬\BibleBook{箴言}{}中的这段经文，而且对许多不懂得基督徒信仰的人来说，他们似乎能说点什么，所以，有必要分别来审查\Quote{他创造了}以及\Quote{他是忠信的，对待那位造了他的}这两处；以便，如同在其他所有经文上一样，在这一段经文上也可以证明他们只不过是幻想而已。\par

19. 首先，让我们看看他们在异端开始形成之初，给有福者\Person{亚历山大}主教所作的答复。他们写道：\Quote{他是受造物，但并非如同受造物之一；是工程，但并非如同工程之一；是后裔，但并非如同后裔之一。}请每个人思考一下这异端的无耻与狡猾；因为它知道自己本身的邪恶之苦涩，所以努力用漂亮话来装点自己，它确实说了它的本意——他是受造物，却自以为能够借加上\Quote{但并非如同受造物之一}来掩盖自己。然而，他们这样写，反而定了他们不虔敬的罪；因为，如果按你们的意见，他干脆就是个受造物，为什么还要加上\Quote{但并非如同受造物之一}这个借口呢？如果他干脆是一项工程，又怎么说是\Quote{并非如同工程之一}呢？从中我们可以看出异端的毒素。因为当他们说\Quote{后裔，但并非如同后裔之一}时，他们便计入了许多儿子，并宣称主是其中之一；这样，按他们的说法，他便不再是\OriginalTerm{独生子}{Only-begotten}，而是众弟兄中的一位，被称为后裔和儿子。那么，说他是受造物而又不是受造物的这种借口，又有什么用呢？因为尽管你们说，\Quote{并非如同受造物之一}，我却要证明你们的这种诡辩是愚蠢的。因为你们仍然宣称他是受造物之一；而且，一个人对其它受造物所能说的任何话，你们也照样用来论及圣子，你们真是\BibleQuote{瞎子和盲者！}\BibleRef{玛 23:19}。难道有哪一个受造物恰好与另一个受造物一样，以至于你们竟把这一点当作某种特权而归于圣子呢？而且，整个可见的创造是在六天之内造成的：——第一天，造了那被他称为\Quote{昼}的光；第二天，造了穹苍；第三天，聚水，显露旱地，并生出地上的各种果实；第四天，造了太阳、月亮和一切星辰；第五天，创造了水中的生物和空中的飞鸟；第六天，造了地上的走兽，最后造了人。并且\BibleQuote{其实，自从天主创世以来，他那看不见的美善，即他永远的大能和他为神的本性，都可凭他所造的万物，辨认洞察出来}\BibleRef{罗 1:20}；并且，光不同于夜，太阳不同于月亮；无灵之物不同于有灵之人；天使不同于座天使，座天使不同于权天使，然而它们全都是受造物，每一个受造之物都各按其类而存在，并保持其自身的本质，就如当初被造时那样。\par

20. 因此，应将圣言从工程中分别出来，作为造物主归于圣父，并承认祂是本性的圣子；或者，倘若祂仅是一个受造物，便应赋予祂与其馀受造物相同的境况，并使它们每一个都如祂一样被称为 \Quote{一个受造物，却非如同受造物之一；子嗣或工程，却非如同工程或子嗣之一。} 因为你们声称子嗣等同于工程，写道 \Quote{受生或受造}。因为尽管在比较上，圣子超出其馀者，但祂仍然如同它们一样是个受造物；因为在那些本性是受造物的事物中，可以发现有些比其它更卓越。例如，星与星在光荣上有别，其馀的当互相比较时也各有差异；然而，这并不意味其中某些是主，而其它的则服事较优者，也不意味某些是 \OriginalTerm{动力因}{efficient cause}，其它则藉由它们而产生，而是所有都拥有一个生发和受造的本性，在它们自身内承认它们的 \OriginalTerm{造主}{Framer}：如 \Person{达味} 在 \BibleBook{圣咏}{集} 中说：\BibleQuote{高天陈述天主的光荣，穹苍宣扬祂手的化工；} 又如贤哲 \Person{则鲁巴贝耳} 说：\BibleQuote{全地呼求真理，高天赞美它：一切工程都因它而震动战栗} \BibleRef{厄上 4:36}。但如果整个大地赞美 \OriginalTerm{造主}{Framer} 和真理，赞美并敬畏它，而它的造主是圣言，并且祂亲口说：\BibleQuote{我是真理} \BibleRef{若 14:6}，那么，圣言就不是受造物，而是惟独属于圣父，在祂内万物得以安排，并且祂被万有赞颂为 \OriginalTerm{造主}{Framer}；因为 \BibleQuote{我曾在祂身旁作安排者；} 以及 \BibleQuote{我圣父到现在一直在工作，我也工作} \BibleRef{若 5:17}。而 \Quote{到现在} 一词表明祂作为圣言永远存在于圣父内；因为圣言的本分就是作圣父的工程，而非在祂之外。\par

21. 但如果圣父所作的，圣子也作，并且圣子所创造的，就是圣父的创造，而圣子却是圣父的工程或受造物，那么，要么祂将作出自己，成为自己的造物主（因为圣父所作的，也是圣子的工程），这是荒谬且不可能的；要么，在祂创造并作出圣父的事物时，祂自己便不是工程也不是受造物；因为否则祂自己既是 \OriginalTerm{动力因}{efficient cause}，便可能使被造之物变得与祂自己所成为的一样，或者更确切地说，祂可能根本没有能力作任何因。\par

因为，若按你们所持的，祂是从无中出来的，祂如何能将无中事物造成实有呢？或者，如果祂自身是受造物，却仍能制造受造物，那么同样的能力在每一个受造物中也应能够设想，即制造其他事物的能力。如果这令你们满意，那么圣言还有什么必要呢？因为低等的事物可以由高等的事物产生。或者无论如何，每一被产生之物都可以在起初聆听天主的话——\Quote{有} 和 \Quote{造成}——如此便被制造。但经上并非如此记载，也不可能如此。因为一切被产生之物中没有一个 \OriginalTerm{动力因}{efficient cause}，万物都是藉着圣言而造成的；若圣言自己也在受造物之列，祂就不会造成万物。因为天使也无法制造，他们同样是受造物，尽管 \Person{瓦伦廷}、\Person{马西翁} 和 \Person{巴西利得} 这样想，而你们是他们的仿效者；太阳既是受造物，也绝不能将虚无变成实有；人不能塑成人，石头不能设计石头，木头不能使木头生长。但天主是在母胎中形成人、安置山岳、使树木生长的那位；而人，由于具有知识，只是组合和安排材料，并按所学的方式加工现存之物；只要它们得以产生，他便满足，并且意识到自己的本性之后，若有任何需要，便知道向天主祈求。\par

22. 那么，如果天主也用材料制作和组合，这实在是一种\OriginalTerm{外邦}{Gentile}思想，按照这种思想，天主只是一个工匠，而非创造者；但即使在这种情形下，也让\OriginalTerm{圣言}{Word}按天主的吩咐并在为天主服务中，去加工那些材料吧。如果祂藉着祂自己的圣言，召叫那原本不存在的物进入存在，那么圣言就不属于那些本不存在而被召叫的物之列；否则我们就必须寻找另一个圣言，藉着祂那一位也被召叫；因为万有原是藉着圣言从无中产生了。如果天主藉着祂创造并制作，祂自己就不属于受造和制作之物；祂反而是创造者天主的圣言，并由父的工程（这些工程也是祂自己所行的）而被人认识为\BibleQuote{祂在父内，父在祂内}，和\BibleQuote{谁看见了祂，就是看见了父}，因为子的\OriginalTerm{性体}{Essence}是父所专有的，祂在各方面都与父相似。那么，祂怎能藉着祂创造呢，除非祂是祂的圣言和智慧？而祂又怎能是圣言和智慧，除非祂是祂性体的真正子嗣，且不像其他万有那样，出自虚无？既然万有都出自虚无，都是受造物，而圣子，如他们所说的，也是受造物之一，且属于那曾经不存在的物之列，那么祂怎能单独启示父，且除了祂以外无人认识父呢？因为，如果作为工程的祂可能认识父，那么父也必按各人程度的比例被所有受造物认识：因为所有受造物和祂一样都是工程。然而，如果说那出于无而有的受造物，既不能看见也不能认识，因为对祂的看见和认识超越万有（因为天主自己说过：\BibleQuote{人见我的面，不能再存活}），但圣子却宣告：\BibleQuote{除了子，没有人认识父}\BibleRef{玛 11:27}，因此圣言与一切出于无而有的受造物不同，因为唯独祂认识父，唯独祂看见父，正如祂所说的：\BibleQuote{这不是说有人看见过父，只有那从父来的，祂看见过父}，以及\BibleQuote{除了子，没有人认识父}，尽管亚略想法不同。那么，祂怎能独自认识呢，除非唯独祂是父所专有的？而如果祂是受造物，而非出自父的真子，又如何是专有的呢？（为了信仰的缘故，我们不应介意经常重复同样的话。）因此，认为圣子属于万有之一是反宗教的；称祂为\Quote{受造物，但不是如同受造物之一；工程，但不是如同工程之一；子嗣，但不是如同子嗣之一}，这是亵渎且毫无意义的。因为如果像他们所说的，祂在受生之前并不存在，又怎能不是如同其中之一呢？因为，受造物和工程的特性，就是在其产生之前并不存在，且从虚无中得以成立，即使它们在光荣方面超越其他受造物；因为这种个别差异存在于一切受造物中，在可见之物中也可看到。\par

23. 而且，如果像异端者所主张的，圣子是受造物或工程，但因祂在光荣上超越其他受造物，故并非如同受造物之一，那么圣经自然地应透过与其他工程比较而描述并展示祂的优势；例如，圣经应说祂大于总领天使、比座天使更尊贵、比日月更明亮、比诸天更高。但事实上圣经并没有这样论及祂；相反，父表明祂是祂自己专有的独生子，说：\BibleQuote{你是我的儿子}，以及\BibleQuote{这是我的爱子，我所喜悦的}。因此，众天使侍奉祂，视祂为超越他们自身的一位；他们朝拜祂，并非因为祂在光荣上更伟大，而是因为祂是超越一切受造物、超越他们自身、且按性体说，唯独是父专有的子。因为如果祂只因在光荣上超越而受朝拜，那么每一种臣服之物都该朝拜那超越自己的。但事实并非如此；因为受造物不朝拜受造物，而是仆人朝拜主人，受造物朝拜天主。因此，宗徒伯多禄阻止那想朝拜他的科尔乃略，说：\BibleQuote{我自己也是人}\BibleRef{宗 10:26}。而在默示录中，一位天使，当若望想朝拜他时，也阻止他，说：\BibleQuote{切不可这样！我只是你和你众弟兄先知，以及那些遵守本书预言者的同仆；你要朝拜天主！}\BibleRef{默 22:9}。因此，唯有天主堪受朝拜，这一点连众天使都知道，即：他们虽然超越其他存在的荣耀，但都是受造物，不可受朝拜，而只应朝拜主。因此，三松的父亲玛诺亚，想向那位天使献祭，立即被天使阻止，说：\BibleQuote{不要向我献祭，只应向天主}。另一方面，主甚至也受众天使朝拜；因为经上记载：\BibleQuote{天主的众天使都要朝拜祂}\BibleRef{希 1:6}；也受一切外邦人朝拜，如依撒意亚所说：\BibleQuote{埃及的劳碌、雇士的买卖，以及身材高大的色巴人，都要过来归于你，他们要作你的奴隶；他们要向你跪拜，向你哀求说：天主确在你们中间，除祂以外，再没有别的神，没有别的天主}\BibleRef{依 45:14}。祂也接受门徒的朝拜，并向他们证实祂是谁，说：\BibleQuote{你们称我师傅、主，你们说得对，我实在是}。当多默对祂说：\BibleQuote{我主！我天主！}时，祂容许这样的话，甚至接纳了他，而非阻止。因为祂就是，正如其他先知所宣告，达味在圣咏中所说：\BibleQuote{万军的主，撒巴耳的主}，意即\Quote{军旅的主}，且是真实而全能的天主，即使亚略派闻之暴怒。\par

24. 然而，如果祂仅仅是受造物，就绝不会如此受朝拜，也绝不会如此被论及。但如今，正因为祂不是受造物，而是那位受朝拜天主的\OriginalTerm{性体}{Essence}所生的真正子嗣，且按本性就是祂的子，所以祂受朝拜，并被相信是天主，且是万军的主，具有权威，全能如同父；因为祂亲口说过：\BibleQuote{凡父所有的一切，都是我的}\BibleRef{若 16:15}。因为，拥有父的一切，并且让人在祂内看见父，而且万有都是藉着祂而造成的，万民的救恩在祂内实现并得以完成：这都是子所专有的。\par

% structure: p0031 source=h2 target=subsection reason=relative-heading-rank; base=h2

\subsection{第十七章 对\BibleBook{}{箴言} 8:22的引介（续）。假设为了创造其他受造物而创造一个圣子或圣言是荒谬的；至于受造物无法承受天主的直接之手，天主纡尊降贵至最低处。此外，如果圣子是受造物，祂也不能承受天主之手，就需要无限系列的媒介。反对说，如同带领以色列人出离的梅瑟是人一样，我们的主也是人；但梅瑟不是创造中的行动者：——又反对说，统一性出现在受造的服事中，但所有这些服事都是有缺陷且依赖的：——又反对说，祂学习创造，然而天主的智慧需要教导吗？而且祂为什么要学习，如果圣父工作直到如今？如果圣子被创造是为了创造我们，那么祂是为了我们，而不是我们为了祂。}

24.（续）在此，为更清楚地驳斥他们的异端，我们最好也向他们提出这个问题——既然万物都是受造物，所有一切皆从虚无中得以造存，而根据你们的说法，圣子自己也是受造物与工程，曾一度不存在，那么为何唯独\BibleQuote{万物是藉著祂而造成的}，\BibleQuote{凡受造的，沒有一樣不是藉著祂而造成的}\BibleRef{若 1:3}？或者为何当提及\Quote{万物}时，无人认为圣子被包括在内，而只指受造之物；但当圣经论及圣言时，它并不视祂为\Quote{万物}之一，却将祂与圣父并列，作为万物藉由圣父管理并实现眷顾与救恩的那一位？然而万物本当也可藉着那道命令——就是圣子由天主单独造成的那命令——而得以存在吗？因为天主不会因命令而疲倦，祂的能力也足以创造万物，祂无需单独创造唯一圣子，再需要祂的辅助与合作来构造其余万物。祂不让任何祂愿意做的事迟延；祂一愿意，万物便存在了，无人\BibleQuote{抗拒了祂的意志}\BibleRef{罗 9:19}。那么，为什么万物不是由天主藉着那同一命令——即圣子由以受造的命令——单独造成的呢？或者让他们告诉我们，为什么万物要藉着祂——祂自己不过是受造物——而得以存在呢？何等毫无理性！然而，关于祂，他们这样说：\Quote{天主愿意创造受造的本性，但见它无法承受圣父未经调剂的手，也无法由祂直接创造，于是首先单独创造了一个，称祂为子与圣言，好叫万物借着祂作为中介，随后得以造成。}他们不仅这样说了，而且胆敢将其笔之于书，这些人就是：欧瑟比乌斯、亚略和曾献祭的阿斯忒里乌斯。\par

25. 这岂不正是他们那种不虔敬的充分证据吗？他们以极度的疯狂麻醉自己，以致厚颜无耻地醉心于反对真理。因为如果他们以创造万物的劳苦作为天主只创造圣子的理由，那么整个受造界将要呼喊反对他们，因为他们在说有关天主的不相称的话；依撒意亚也在圣经中说：\BibleQuote{上主是永生的天主，是创造地极的主，他并不疲倦，也不困乏，他的智慧高深莫测。}\BibleRef{依 40:28}而如果天主只创造了圣子，因为不屑于创造其余万物，而把它们交给圣子作为助手，这另一方面也不相称于天主，因为在祂内毫无骄傲。实际上，主谴责这种思想，当祂说：\BibleQuote{两只麻雀不是卖一个铜钱吗？}和\BibleQuote{没有你们的天父许可，它们中连一只也不会掉在地上。}又説：\BibleQuote{不要为你们的生命忧虑吃什么，或喝什么；也不要为你们的身体忧虑穿什么。难道生命不是贵于食物，身体不是贵于衣服吗？你们仰观天空的飞鸟，它们不播种，也不收获，也不在粮仓里屯积，你们的天父还是养活它们；你们不比它们更贵重吗？你们中谁能运用思虑，使自己的寿数增加一肘呢？关于衣服，你们又忧虑什么？你们观察一下田间的百合花怎样生长：它们既不劳作，也不纺织；可是我告诉你们：连撒罗满在他极盛的荣华时代所披戴的，也不如这些花中的一朵。田里的野草今天还在，明天就投在炉中，天主尚且这样装饰，信德薄弱的人哪，何况你们呢？}因此，既然天主对如此微小的事物——头发、麻雀、田野的草——施行眷顾，并不有失身份，那么创造它们也不有失身份。因为凡祂眷顾的事物，祂也藉自己的圣言造成了。然而，对于如此说话的人，还有一个更糟糕的荒谬之处；因为他们把受造物与构造工程区分开来；认为后者是圣父的工程，而受造物是圣子的工程；其实，要么万物都须由圣父与圣子共同创造，要么如果所有受造物都藉圣子而存在，我们就不该称祂为受造物之一。\par

26. 其次，他们的愚昧可以如此揭露：如果\OriginalTerm{圣言}{Word}也是\OriginalTerm{受造的本性}{originated nature}，那么这种本性既过于软弱，不足成为天主的亲手化工，祂又怎能独自从一切中忍受由你们所说的\OriginalTerm{非受生且无掺杂的本体}{unoriginate and unmitigated Essence}所造呢？因为如果祂能忍受，则所有人都能忍受；或者，既然无人能忍受，那么圣言也不能忍受，因为你们说祂是受造物之一。再者，如果因为受造的本性不能忍受成为天主的亲手化工，而产生了需要\OriginalTerm{中保}{mediator}，那么必然得出，圣言既是受造物，在祂的形成中也需要\OriginalTerm{中保}{medium}，因为祂也属于那种不能忍受由天主所造而需要中保的受造本性。但如果为祂找到某个存在作中保，那么为此第二者，又需要一个新中保，这样追溯下去，我们将发明出大量累积的中保；于是创造便不可能存在，因为永远缺少中保，而那个中保若不藉另一位中保也不会出现；因为按照你们的说法，它们全都属于那种不能忍受唯独由天主所造的受造本性。这愚昧何其多！它迫使他们认为已经存在的事物不容许生成！或者他们可能以为它们甚至还未生成，因为仍在寻找它们的中保；根据他们如此不虔而又虚妄的观念，存在之物由于缺少中保便不会存在。\par

27. 但是，他们又提出这一点：\Quote{看，祂也曾借着\Person{梅瑟}带领百姓出离\Place{埃及}，并借着他颁赐法律，然而他不过是人；所以，相似者可以借相似者而生。}当他们说这话时，该蒙上脸，以遮其耻。因为梅瑟并非受遣去创造世界，不是去呼唤那不存在者进入存在，也不是去塑造同他一样的人，而只是去向百姓和\Person{法郎}王传达言语的执行者。这两者完全不同：执行是属于受造物、如同仆役的事，而创造与形成则唯独属于天主，并属于祂固有的\OriginalTerm{圣言}{Word}和\OriginalTerm{智慧}{Wisdom}。所以，在创造的事上，我们只能找到天主的圣言；因为\BibleQuote{万物皆在智慧中造成}，并且\BibleQuote{没有圣言，就没有造成一物}。但至于执行任务者，则不是只有一位，而是主愿派遣谁，就从全体人中派遣谁。因为有许多\OriginalTerm{总领天使}{Archangels}、许多\OriginalTerm{上座者}{Thrones}、\OriginalTerm{宰制者}{Authorities}和\OriginalTerm{率领者}{Dominions}，成千上万，无量无数，侍立在天主面前，奉命服务并随时被派遣。还有许多先知、十二宗徒和\Person{保禄}。梅瑟本人也不是独一的，\Person{亚郎}与他同在，随后又有七十人充满圣神。梅瑟之后是\Person{农的儿子若苏厄}，他之后是民长，而他们也不是由一人接替，乃是由众多君王接替。那么，如果圣子是受造物，是受造者之一，就必然会有许多这样的儿子，好让天主有许多这样的执行者，正如前面那些数目众多一样。但如果这并非我们所见，反而受造物虽多，圣言却只有一位，那么任何人由这一点便可推断：圣子与众受造物不同，不与受造物同级，而是属于圣父的。因此，圣言不是许多，而是只有一父的唯一圣言，以及唯一天主的唯一肖像。但他们又说：\Quote{看，太阳只有一个，地球也只有一个。}他们既然这样愚蠢，不如也主张只有一水、一火，然后就可以被告知：凡被造之物，在其各自的本体上确是一个；但为赋予它的执行与服务，单靠自己既不合宜也不充足。因为天主说过：\BibleQuote{在天空中要有光体，以照耀大地，分别明与暗；并作记号，定节令、日、年。}然后又说：\BibleQuote{于是天主造了两个大光体：较大的统治白昼，较小的统治夜晚，又造了星宿，并将它们安置在天空，照耀大地，统治白昼与黑夜。}\par

28. 看，光体很多，不只有太阳，也不只有月亮，但各自在本质上是一，然而所有光体的服务却是一致共通的；它们各自的欠缺，由其它光体补充，照明的功能由全体共同完成。这样，太阳有权在白日普照，仅此而已；月亮则在夜间；星辰与它们一起完成季节与年岁，并成为征兆，各自应所需而显。同样，大地并非为了一切，而是仅为产果，并为居住其上的生物提供立足之地。穹苍则是为了分开水与水，并为安置星辰的场所。火与水，以及其他事物，也都是为了成为物体的构成部分而受造；简言之，没有一物是孤立的，一切受造之物，好像彼此互为肢体，共同组成一个身体，即宇宙。那么，如果他们如此设想圣子，就让众人向他们扔石头吧，因为他们将圣言视为这个宇宙的一部分，而且是没有其余部分便不足以完成所受委托之职分的部分。但若这显然是不虔敬的，就让他们承认，圣言不属于受造物之列，而是圣父唯一本有的圣言，是他们的形成者。他们说：\Quote{但是，虽然祂是受造物，属于受造之类，但祂如同从师傅和工匠那里学会了形成，如此服事了那位教导祂的天主。}正因为如此，智者\Person{阿斯特里乌斯}，由于学会了否认主，竟敢如此写，没有注意到由此产生的荒谬。因为如果形成是可教授的事，他们便当心，以免说天主本身形成万物不是由于本性，而是由于学识，这样就有失去能力之虞。此外，如果天主的智慧是通过教导而达到形成的，当祂需要学习时，祂又如何仍是智慧呢？那么，在学习之前祂是什么？如果祂需要教导，就不是智慧；肯定只是一个空洞之物，而不是\OriginalTerm{本质的智慧}{essential Wisdom}，只是由于进步才得到智慧之名，并且只有在能保守所学时才持续是智慧。因为凡不是由本性，而是由学习得来的，都可能有一日失学。但如此谈论天主的圣言，不是基督徒的事，而是希腊人的事。\par

29. 因为如果形成的能力是通过教导而赋予某人的，这些不明事理的人便将嫉妒与软弱归咎于天主——嫉妒，因为祂没有教导许多人如何形成，以致在祂周围，如同有许多总领天使和天使，也有许多形成者；软弱，因为祂不能独自创造，而需要一位同工或下手；而且，已经证明受造的本性只可由天主单独创造，他们却认为圣子具有这样的本性并被造成这样。但天主一无所缺：愿此念消逝！因为祂亲自说过：\BibleQuote{我已饱享。}\BibleRef{依 1:11}圣言也不是因教导而成为万有的形成者；而是作为圣父的肖像和智慧，祂行圣父的事。祂也不是为了创造万物而造了圣子；因为看，虽然圣子存在，圣父仍被看见在工作，正如主亲自说的：\BibleQuote{我父到现在一直工作，我也工作。}\BibleRef{若 5:17}然而，如果如你所说，圣子受生是为了造祂之后的事物，而圣父甚至在圣子之后还被看见在工作，那么即便如此，你也必须认为这样的圣子受生是多余的。此外，当祂要创造我们时，为何还要寻求中保，仿佛祂的旨意不足以成就祂所喜悦的一切呢？然而圣经说：\BibleQuote{祂凡所喜悦的，无不实行}，又说：\BibleQuote{谁能抗拒祂的旨意？}\BibleRef{罗 9:19}如果仅仅祂的旨意就足以形成万物，你便使中保的职分成为多余了；因为你所举的梅瑟、太阳和月亮的例子，已证明是不恰当的。这里又有使你无言以对的论证。你说，天主愿意创造受生的本性，并为此筹谋，便设计和创造了圣子，为能藉着祂形成我们；若果如此，想想你竟敢说出了多么不虔敬的话。\par

30. 首先，圣子似乎更是为我们而得以存在，而不是我们为祂；因为我们并非为祂受造，而是祂为我们所受造；因此祂应感谢我们，而非我们感谢祂，正如女人之于男人。因为圣经上说：\BibleQuote{男人不是为女人受造，女人乃是为男人受造。}所以，正如\BibleQuote{男人是天主的肖像和光荣，女人是男人的光荣}\BibleRef{格前 11:7}，同样，我们成为天主的肖像，为祂的光荣；但圣子是我们的肖像，为我们的光荣而存在。我们受造是为了存在；但天主的\OriginalTerm{圣言}{Word}，如你们所必须主张的，并非为了祂自己存在而受造，而是作为满足我们需要的一个工具，因此不是我们从祂而来，乃是祂从我们的需要而被构成。难道能设想这种思想的人，不是比麻木不仁更甚吗？因为如果圣言是为我们而受造，那么祂在天主面前就不居于我们之先；因为祂并非在自己内拥有祂而与我们商议，而是在自己内拥有我们，如他们所说，便商议关于祂自己的圣言。但若如此，或许圣父甚至根本对圣子没有意愿；因为祂并非因有意于祂而创造祂，而是有意于我们，为了我们的缘故而造了祂；因为祂在计划我们之后，才计划了祂；以致於按照这些不虔敬者的说法，如今那些被造的人既已存在，作为工具而被造的圣子就成了多余的了，因为祂是为他们而受造的。但如果圣子由天主独自所造，是因为祂能承受这创造，而我们却不能，因此我们由圣言所造，那么祂为什么不首先商议关于那能承受祂创造的圣言，反而商议我们呢？或者祂为什么更重视那强壮的祂，而不重视软弱的我们呢？或者为什么先造了祂，却不先商议关于祂呢？或者为什么先商议我们，却不先造我们呢？祂的旨意原本足以构成万有。但祂先造了祂，却先商议我们；且祂愿意我们，是在中保之前；并且当祂愿意创造我们，并商议我们时，祂称我们为受造物；但祂为我们所塑造的那位，祂却称为圣子和合法继承人。但为了我们而造祂的我们，反更应该被称为儿子；或者至少，祂的圣子，更是祂预先思想和旨意的对象，为了祂，祂才造了我们众人。这就是异端者的病态，这就是他们的呕吐物。\par

% structure: p0039 source=h2 target=subsection reason=relative-heading-rank; base=h2

\subsection{第十八章 续论\BibleBook{}{箴言}8:22。圣父在圣子内直接而自然地运作，借受造物作工具；圣经用语说明此点。这些说明应依教会教义解释；亚略派所赋予的曲解，予以驳斥。神圣生育的奥迹。详述天主圣言与人之言的区别。\Person{阿斯特留斯}被陷入主张两个\OriginalTerm{无始者}{Unoriginate}；其不一致。洗礼如何由圣子和圣父共同施行。论异端者的洗礼。为何亚略派比其他异端更坏。}

31. 但关于此事的真理情操，决不可隐藏，而必须高声宣扬。因为天主的圣言并非为我们而受造，反而是我们为祂而受造，并且\BibleQuote{万物都是在祂内受造的}\BibleRef{哥 1:16}。也不是因为我们软弱，祂才强壮，并由圣父独自所造，以便祂借祂为工具塑造我们；绝无此念！并非如此。因为即使天主定意不创造任何受造之物，圣言依然与天主同在，圣父在祂内。同时，受造之物若无圣言便不能存在；因此它们是通过祂而受造——这是合理的。因为圣言既是按本性专属于祂本体的天主圣子，且是从祂而来，并在祂内，如祂自己所说，那么受造物若非通过祂，便无法存在。正如光借其光芒照亮万物，没有光芒，任何事物都不能被照亮；同样，圣父借着圣言，如同借一只手，创造万有，离了祂，什么也不造。例如，正如\Person{梅瑟}所记载，天主说：\BibleQuote{要有光}，\BibleQuote{让水聚集在一起}，\BibleQuote{让旱地出现}，以及\BibleQuote{让我们造人}；至圣的\Person{达味}在圣咏中亦说：\BibleQuote{祂一发言，万有造成，祂一出命，万物化生}。祂说话，并不像人间那样，某个副手听见，得知发言者的旨意后，离去执行；因为那是属于受造物的方式，但以这种方式去思想或谈论圣言则是不恰当的。因为天主的圣言是塑造者和创造者，祂就是圣父的旨意。因此，圣经并没有说，有一位听见并回答，如同论及祂所愿造之物的方式或本性；而是天主只说了：\BibleQuote{要有}，便接着说：\BibleQuote{事就这样成了}；因为凡祂所乐意和计划的，圣言便立刻开始实行并完成。因为当天主命令其他受造物时，无论是天使，或是与\Person{梅瑟}对话，或是命令\Person{亚巴郎}，那时，听者会回答；一个说：\BibleQuote{我凭什么能知道}\BibleRef{创 15:8}？另一个说：\BibleQuote{请派别人去}\BibleRef{出 4:13}；又：\BibleQuote{他们若问我，祂叫什么名字，我该对他们说什么}？天使对\Person{匝加利亚}说：\BibleQuote{主这样说}；他又问主：\BibleQuote{万军的主，你无情对待耶路撒冷要到何时呢}？并等候聆听慰藉之言。因为所有这些，都有中保圣言，以及彰显圣父旨意的天主智慧。但当圣言亲自工作与创造时，便没有询问与回答，因为圣父在祂内，圣言也在圣父内；只要愿意，便完成了；所以\Quote{祂说}一词，是为我们的缘故，作为旨意的标记，而\Quote{事就这样成了}，则表明通过圣言和智慧所完成的事工，在这智慧中，也蕴涵着圣父的旨意。而\Quote{天主说}即在\Quote{圣言}中得到阐释，因为经上说：\BibleQuote{你以智慧造成万物}；和\BibleQuote{因主的一句话，诸天造成}；以及\BibleQuote{只有一个主，就是耶稣基督，万物借祂而有，我们也借祂而有}。\par

32. 由此清楚可见，亚略派并非为他们的异端与我们相争；当他们假意针对我们时，他们真正的争战乃是反对\OriginalTerm{天主性}{Godhead}本身。因为，倘若说\BibleQuote{这是我的爱子}的声音是出自我们，我们对他们的指责便无足轻重；但这既是圣父的声音，门徒们也听到了，且圣子自己也论到祂说：\BibleQuote{在群山被造之前，祂已生了我}，他们岂不正是与天主为敌，如同神话中的巨灵，照圣咏作者的话，他们的舌头是攻击信仰的利剑吗？因为他们既不畏惧圣父的声音，也不尊敬救主的话，又不信从圣人，其中一位写道：\BibleQuote{祂是光荣的光辉和祂本体的真像}，又说\BibleQuote{基督是天主的德能和天主的智慧}；另一位在圣咏中说：\BibleQuote{在祢那里有生命的泉源，在祢的光中我们必得见光}，又说\BibleQuote{祢以智慧造成万物}；先知们说：\BibleQuote{上主的圣言传给我说}\BibleRef{耶 2:1}；若望说：\BibleQuote{在起初已有圣言}；路加说：\BibleQuote{照那些自起初亲眼见过并为圣言服役的人，传递给我们的人}\BibleRef{若 1:1}；而达味又说：\BibleQuote{祂发出自己的圣言，治好了他们}。所有这些经文从各方面都明令禁止亚略异端，并表明\OriginalTerm{圣言}{Word}的永恒性，且祂并非外来的，而是圣父\OriginalTerm{性体}{Essence}本有的。因为谁曾见过无\OriginalTerm{光辉}{Radiance}的\OriginalTerm{光}{Light}呢？或者谁敢说\OriginalTerm{真像}{Expression}能与\OriginalTerm{本体}{Subsistence}不同呢？一个人若竟设想天主曾一度没有\OriginalTerm{理智}{Reason}和\OriginalTerm{智慧}{Wisdom}，岂不是自己丧失理智了吗？因为\OriginalTerm{圣经}{Scripture}提出这样的比喻和这样的影像，正是体谅人性不能领悟天主，我们才能藉着这些，无论多么贫弱模糊，尽可能形成一些概念。正如受造界含有丰富的材料，使人认识天主的存在和\OriginalTerm{上智安排}{Providence}（\BibleQuote{因为从受造物的伟大和美丽，可以比拟地看到它们的创造者}\BibleRef{智 13:5}），我们无需寻问声音，通过它们学习，而是聆听圣经，我们相信；观察万物的秩序与和谐，我们承认祂是万有的创造者、主宰和天主，并领悟祂对万物的奇妙\OriginalTerm{上智安排}{Providence}和治理；同样，关于圣子的\OriginalTerm{天主性}{Godhead}，上文所述已足够，再去争论，或者说以异端的方式询问：圣子怎能从永恒而生？或者祂怎能出自圣父的\OriginalTerm{性体}{Essence}，却不是一部分？这便成了多余，甚至极为疯狂。因为凡说出于某者的，便是他的一部分；而分开的，就不是完整的。\par

33. 这些就是异端分子邪恶的诡辩；然而，虽然我们已指出其浅薄，这些经文本身的准确含义和这些比喻的力量，将足以表明他们所持可憎信条之毫无根据。因为我们看到，\OriginalTerm{理智}{Reason}是永存的，并且发自那理智所归属者，且是其性体所本有的，它不容许有先后。同样，我们看见出自太阳的\OriginalTerm{光辉}{Radiance}是太阳本有的，太阳的性体并不因此被分割或减损；其性体是完整的，而其光辉也是完美而完整的，却无损于光的性体，且是它的真正产物。我们同样理解，圣子并非从外而生，而是由圣父所生，圣父保持完整，祂\OriginalTerm{本体}{Subsistence}的\OriginalTerm{真像}{Expression}便永存，并保存了圣父的肖像和不变的\OriginalTerm{肖像}{Image}，以致凡看见祂的，也在祂内看见那\OriginalTerm{本体}{Subsistence}，而祂正是那\OriginalTerm{本体}{Subsistence}的\OriginalTerm{真像}{Expression}。并且，由\OriginalTerm{真像}{Expression}的运作，我们明了\OriginalTerm{本体}{Subsistence}的真实\OriginalTerm{天主性}{Godhead}，正如救主亲自教导说：\BibleQuote{住在我内的父，行祂自己的事业}\BibleRef{若 14:10}，即我所作的；\BibleQuote{我与父原是一体}，又说\BibleQuote{我在父内，父在我内}\BibleRef{若 10:30}。因此，让这反对基督的异端先尝试分割那些在受造物中找到的比喻，说：\Quote{太阳曾一度没有光辉}，或\Quote{光辉并非光之性体本有的}，或\Quote{它确实是本有的，但分割后成为光的一部分}；然后，让它分割\OriginalTerm{理智}{Reason}，宣称理智与心灵相异，或曾经不存在，或并非其性体本有的，或分割后成为心灵的一部分。至于祂的\OriginalTerm{真像}{Expression}和\OriginalTerm{光}{Light}和\OriginalTerm{德能}{Power}，也让它如对待\OriginalTerm{理智}{Reason}和\OriginalTerm{光辉}{Radiance}那样加以粗暴处理；反而凭己意任意想象。但如果他们无法如此狂妄，难道他们不是大大失常，擅自闯入超越受造物和他们自身本性的事物，企图行那不可能之事吗？\par

34. 因为如果在这些受造和非理性的事物中，发现有子嗣，它们既不是其所出自之本体的部分，也不是带着受动而存在，更不会损害其来源的本体，那么他们竟然又在非物质而真实的天主身上寻求并揣测部分和受动，将分裂归于那超越一切受动与变化者，企图扰乱纯朴人的耳朵，使他们偏离真理，这岂不是疯狂之举吗？因为谁听见\Quote{子}，不想到那属于父亲本体的事物呢？谁在初受教理讲授时，听见天主有一位圣子，并藉着自己的圣言造了万物，却按照我们现在所意指的那种理解去领会的呢？在这可恨的亚略异端兴起时，谁听到他所听到的，不立刻感到惊异，视之为怪异，且是在那从起初就撒下的圣言之外，又撒的另一种种子呢？因为从起初在每一个灵魂中撒下的是：天主有一位圣子，祂是圣言、智慧、大能，也就是祂的肖像和辉耀；由此立刻得出：祂永远存在；祂出自圣父；祂与圣父相似；祂是圣父本体的永恒子嗣；在这些观念中，完全没有受造物或工程的含义。但当人们睡觉的时候，仇敌来，又撒下了稗子，如\Quote{祂是受造物}，\Quote{曾有一个时期祂不存在}，以及\Quote{怎么可能？}之类的话；从此，基督仇敌的邪恶异端就如稗子般生长出来，他们随即丧失了一切正当的思想，像强盗般干预，竟敢说：\Quote{圣子怎可能常与圣父同在？}因为人是由人生出来的，是儿子，是在一段时间之后才有的；父亲三十岁时，儿子才开始出生，由父所生；总而言之，对于每一个人的儿子，他出生之前并不存在，这是确实的。他们又低声议论说：\Quote{圣子怎能是圣言，或者圣言怎能是天主的肖像？因为人的言语是由音节组成的，仅仅表达说话者的意志，然后就结束并消逝了。}\par

35. 然而他们再次，仿佛忘记了已向他们提出的证据一般，用这些不虔敬的绳索\Quote{将自己刺透}，并如此争论。但真理的言语却如此驳斥他们：如果辩论是针对某个人，那么关于祂的圣言和圣子，就让他们以人的方式推理；但如果是关于那位创造了人的天主，那么他们就不应再怀有人性的想法，而应怀有那些超越人性本性的想法。因为怎样的人，所生的也必然怎样；圣言的父如何，圣言也必然如何。人，是在时间内受生的，也在时间内生子；而且人是从虚无中形成的，所以他的言语也就消逝而不延续。但天主并不像人，正如圣经所说；祂是自有者，且永远存在；因此祂的圣言也是自有者，且永远与圣父同在，如同光明中的辉耀。人的言语由音节组成，既不能生活，也不能运作什么，仅只是表示说话者的意向，发出后便消逝，不复再现，因为它在被说出以前根本就不存在；所以人的言语既不能生活，也不能运作什么，简言之，它并不是人。而人的言语之所以如此，如我前面所说，是因为生出它的人，其本性出自虚无。但天主的圣言并非如人们所说的，仅是被说出的言语，也不是音节的声音，所称的圣子也不是指祂的命令；而是如同光明中的辉耀，是由完美者所生的完美子嗣。因此祂也是天主，因为祂是天主的肖像；因为经上说：\BibleQuote{圣言就是天主}\BibleRef{若 1:1}。人的言语不能产生效果；因此人并非靠言语行事，而是靠双手，因为双手是实有的，而人的言语却非实存。但正如宗徒所说，\BibleQuote{天主的圣言却是生活的，有效的，比各种双刃的剑更锐利，直穿入灵魂和神魂，关节与骨髓的分离点，且可辨别心中的感觉和思念。没有一个受造物，在天主面前不是明显的；万物在祂眼前都是袒露敞开的，我们必须向祂交账。}\BibleRef{希 4:12}祂就是万物的造成者，\BibleQuote{凡受造的，没有一样不是藉着祂而造成的}\BibleRef{若 1:3}，而没有祂，也没有任何事物能被造成。\par

36. 我们也不应探究\OriginalTerm{天主圣言}{the Word of God}为何不像我们的言语，因为正如先前所说，天主并不像我们一样。再者，我们也不宜寻求圣言如何出自天主，祂如何是天主的光辉，天主如何生，以及祂的生是怎样的方式。因为凡敢涉足这类问题的人，必定是心智失常；这等于是要求用言语说明那不可言说、专属于天主本性、只有祂和圣子知道的事。这就好像探问天主在哪里、天主怎样、圣父有什么本性一样。探问这类问题既是不虔敬的，也显出对天主的无知；同样，贸然探问天主圣子的生发，或用我们自己的本性及软弱去度量天主和祂的智慧，也是不神圣的。任何人都不能因此就有偏离真理的自由，也不应因为对这类追问感到困惑，就怀疑经上所写的。因为在困惑中保持沉默并相信，胜过因困惑而不信：困惑的人也许还能获得某种怜悯，因为他虽提出了问题，但仍保持静默；但若有人因困惑而为自己捏造不堪的道理，说出有损于天主的言论，这样的胆大妄为只会招致无怜悯的判决。因为遇到这类困惑时，圣经能够提供某种解救，使人正确地领受所记载的，并可以拿我们自己的言语作为比方：正如我们的言语是属于我们的、由我们发出，而不是外在于我们的工程；同样，天主的圣言也属于祂、由祂发出，而不是一件工程；然而祂的圣言并不像人的言语，否则我们就必须设想天主也是人了。请观察：人的言语多种多样，逐日消逝；后来的言语并不延续前面的，而是消失不见。这种情况之所以发生，是因为发言者是会过去的人，他们有时间，意念也前后相继；他们先想到什么、后想到什么，就说出什么；因此他们有许多言语，过后却完全不留下什么；因为说话的人一停，他的言语立刻就消尽了。但是天主的圣言却是同一不变的，如经上所载：\BibleQuote{天主的圣言永远常存}\BibleRef{依 40:8}，祂不是变化或被替代的，而是始终一样地存在。因为既然天主是唯一的，祂的肖像也该是唯一的，祂的圣言是唯一的，祂的智慧也是唯一的。\par

37. 因此我感到惊奇：天主既是唯一的，这些人却凭自己私下的想法，引入了许多肖像、智慧和圣言，并说圣父本有的和本性的圣言并非圣子，圣父甚至借着这个圣言造了圣子；又说那真正是圣子的，只是名义上被称为圣言，如同葡萄树、道路、门和生命树一样；还说他被称为智慧也只是名号，而圣父本有的和真实的智慧，以\OriginalTerm{非受生}{ingenerate}方式与祂共存，却是另一位，并非圣子，圣父甚至借着它造了圣子，又因圣子分享了它而称圣子为智慧。他们不仅限于口头说说，\Person{亚略}在\WorkTitle{塔利亚}中，以及诡辩家\Person{阿斯泰里乌斯}，都曾如上所述写道：\Quote{真福保禄并没有说他宣讲的是基督，即天主的德能或天主的智慧，而是不加冠词地说\BibleQuote{天主的德能}和\BibleQuote{天主的智慧}\BibleRef{格前 1:24}，以此宣讲天主本有的、属于祂本性的、以\OriginalTerm{非受生}{ingenerate}方式与祂同在的德能，乃是另一位，它生成了基督，并创造了整个世界；关于这德能，他在致罗马人书中如此教训说——\BibleQuote{其实，自从天主创世以来，祂那看不见的美善，即祂永远的大能和祂为神的本性，都可凭祂所造的万物，辨认洞察出来}\BibleRef{罗 1:20}。正如没有人会说那里所提的天主性是指基督，而是指圣父自己；同样，我想，\InnerQuote{祂永远的大能和天主性也不是独生子，而是生祂的圣父}。并且他教导说，还有另一个天主的德能和智慧，借着基督彰显出来。}稍后，同一个\Person{阿斯泰里乌斯}又说：\Quote{然而祂那永远的德能和智慧，真理论证是无始无终、非受生的，必定是同一个。因为有许多智慧，是祂一个一个造的，而基督是其中的首生者和独生子；不过所有这些智慧都同样依靠拥有它们的那位。而一切德能，正确地被称为创造并运用它们的那位的德能——正如先知所说，那降来作为人的罪恶之天罚的蝗虫，被天主自己不只称为一种德能，而是称为大德能；真福达味在大多数圣咏中，不仅邀请天使，也邀请德能赞美天主。}\par

38. 现在，他们仅仅说出这样的话，难道不应该招致所有人的憎恨吗？因为，如果他们坚持认为，祂是圣子，不是因为祂由圣父所生，且是圣父本体的专属者，而是因为祂被称为圣言，仅仅是由于那些有理性的存在物；被称为智慧，是由于那些被赋予智慧者；被称为大能，是由于那些被赋予能力者，那么，祂必定也是由于那些被立为子女的人，才被称为圣子：而且，或许因为有了存在物，祂才有了自己的存在，而这存在仅存在于我们的观念中。那么，祂究竟是什么呢？如果这些只不过是祂的名号，祂自身就什么都不是：祂只有一个存在的假象，并且从我们这里获得了这些名号的粉饰。如果他们并不是不愿意承认自己真实存在，却认为天主的圣言仅仅是名号而已，这简直是魔鬼的狂妄，或更糟。说智慧与圣父同在，却不承认那就是基督，反而说有许多受造的能力和智慧，而主只是其中之一，他们甚至将其比作毛虫和蝗虫，这岂不是荒谬绝伦吗？当他们听到我们说圣言与圣父同在时，立即嘀咕说：\Quote{你们不是在说两个\OriginalTerm{非受生者}{Unoriginate}吗？}然而，当他们自己谈论\Quote{祂非受生的智慧}时，却没有看到自己已经招致了他们如此轻率地用来指责我们的罪名，这难道不是放荡无耻吗？此外，他们那种认为那非受生的智慧与天主同在就是天主本身的想法，是多么愚蠢啊！因为凡是同在者，不是与自己同在，而是与另一位同在，正如圣史们论及主时所说，祂与祂的门徒同在；因为祂不是与祂自己同在，而是与门徒同在---除非他们要说天主具有复合的本性，智慧是祂本体的一个组成部分或补充，并且像祂自己一样也是非受生的，而他们更装作这智慧就是世界的创造者，以便剥夺圣子对世界的创造之功。因为，他们为了不持守关于主的真理，什么话不愿坚持呢？\par

39. 他们究竟在圣经的何处找到了，或者从谁那里听说过，除了这位圣子之外，还有另一位圣言和另一位智慧，以至于他们竟为自己构造出这样的教义呢？确实，经上记载说：\BibleQuote{我的话岂不是像火，又像打碎岩石的大锤吗？}\BibleRef{耶 23:29}又如在\BibleBook{}{箴言}中：\BibleQuote{我要将我的话传给你}\BibleRef{箴 1:23}；但这些乃是规诫和命令，是天主藉着祂自己专属的、唯一真实的圣言，向圣者们所说的；关于这圣言，圣咏作者说：\BibleQuote{我禁止我脚走一切邪路，为要遵守你的话。}因此，救主指明这样的话是与祂自己不同的，当祂亲自说：\BibleQuote{我对你们所说的话}\BibleRef{若 6:63}。因为毫无疑问，这样的话不是什么后裔或子，也不会塑造世界，也不是唯一真天主的多重肖像，也不是许多为我们成了人者，也不是仿佛从许多这样的圣言中有一位成了血肉，如同若望所说的那样；相反，若望宣讲祂，因祂是天主的唯一圣言：\BibleQuote{圣言成了血肉}，并且\BibleQuote{万物都是藉着祂而造成的}。因此，只有关于祂---我们的主耶稣基督，以及祂与圣父的合一，才有那些证言被写下并公布出来，既有圣父的证言，指明圣子是一位的；也有圣者的证言，他们意识到这一点，并说圣言是一位的，且祂是独生子。祂的工程也公布出来；因为一切可见与不可见的事物都是藉着祂而有的，\BibleQuote{没有一样不是藉着祂而造成的}。然而，关于另一位或任何其他者，他们根本不曾思想，也不曾为自己构造出什么圣言或智慧，这些在圣经中既无名号也无行动的证据，只是被这些人如此称呼罢了。因为这是他们的虚构，是反对基督的妄想，他们极度利用圣言与智慧的名号；为自己虚构了另一些，却否认了天主的真实圣言，以及圣父的真正且唯一的智慧，从而这些可怜虫与摩尼教徒一争高下。因为摩尼教徒也是如此，当他们看到天主的化工时，却否认唯一真实的天主，并为自己虚构另一位天主，而这位天主他们既无法用工程来展示，也无法用任何取自神圣记载的证言来证明。\par

40. 因此，如果在神圣的记载中，除了这位圣子之外找不到别的智慧，我们也没有从教父那里听说过任何这样的智慧，然而他们却承认并写道，有一种智慧，非受生地...与圣父同在，是祂的专属者，且是世界的创造者，那么这必定就是那位即便按他们的说法也与圣父永远同在的圣子。因为祂是万物的创造者，正如经上所载：\BibleQuote{你以智慧造成万物。}不仅如此，\Person{阿斯特里乌斯}本人，仿佛忘记了先前所写的，后来却像\Person{盖法}那样，当劝勉希腊人时，不由自主地不说许多智慧和毛虫，反而只承认一位智慧，用这些话---\Quote{天主圣言是一位，但理性的事物众多；智慧的本体和本性是一位，但智慧和美善的事物众多。}不久之后他又说：---\Quote{他们用天主儿女的称号尊崇的那些人是谁？因为他们不会说那些也是圣言，也不会主张有许多智慧。因为，既然圣言是一位，智慧也被宣布为一位，就不可能将圣言的本体分配给众多的儿女，并将智慧的名号赐予他们。}因此，亚略派人士与真理争战，一点也不奇怪，因为他们与自己的原则相冲突，彼此矛盾：有时说有许多智慧，有时又主张只有一位；有时把智慧与毛虫同列，有时又说它与圣父同在且是祂的专属；现在说只有圣父是非受生的，转眼又说祂的智慧和祂的大能也是非受生的。他们因我们说天主的圣言永远存在而攻击我们，却忘记了自己的教义，自己却说智慧非受生地...与天主同在。他们在这一切事上如此晕头转向，否认真实的智慧，却发明一个不存在的智慧，就像摩尼教徒一样，他们在否认那真实存在者之后，为自己制造另一位神。\par

41. 然而，让其他异端和摩尼教徒\OriginalTerm{摩尼教徒}{Manichees}也知道，基督的父是唯一的，且是藉着祂本有的\OriginalTerm{圣言}{Word}而为受造界的主和创造者。尤其让\OriginalTerm{亚略狂徒}{Ario-maniacs}知道，\OriginalTerm{天主}{God}的圣言是唯一的，是独生子，出自祂的本体\OriginalTerm{本体}{Essence}，本真无伪，并与祂的父拥有不可分的、同一的\OriginalTerm{天主性}{Godhead}，正如我们多次说过，且由救主自己教导我们的。因为，若非如此，何以父要藉着祂创造，并在祂内将自己启示给祂所愿意的人，并光照他们呢？又何以在圣洗的祝圣\OriginalTerm{祝圣}{consecration}中，圣子与父一同被呼名呢？因为若他们说父并非全备自足，回答便是不虔敬的\OriginalTerm{不虔敬}{irreligious}；若父是全备自足的——这当然是对的——那么，为了创造世界，或为了\OriginalTerm{圣洗}{holy laver}，又何需圣子呢？受造物与创造者之间，有何相通呢？在我们众人的祝圣中，何以被造之物竟与创造者并列？又何以如你们所持守的，交付给我们的信仰，是信仰一位创造者和一位受造物呢？因为若为使我们结合于天主性，何需受造物？若为使我们结合于圣子——一位受造物，那么依你们说，圣洗中呼求圣子之名便是多余的，因为那位使祂成为子的天主，同样能使我们成为子。此外，若圣子是受造物，\OriginalTerm{理性受造物}{rational creatures}的本性既同为一体，受造物无从得受造物的扶助，因为万有都需要\OriginalTerm{恩宠}{grace}。我们刚才略说了万物应藉祂而造的合宜性；但因讨论进程也使我们提及圣洗，故必须陈明——如我所想所信的——圣子与父一同被呼名，并非因为父不全备自足，也非无意且偶然；而是因为祂是天主的圣言和本有的智慧\OriginalTerm{智慧}{Wisdom}，又是父的\OriginalTerm{光耀}{Radiance}，常与父同在，因此，若父施予恩宠，祂不可能不在圣子内施予，因为圣子在父内，如同光耀在光明中。因为，天主并非出于需要，而是作为父，在祂本有的智慧中奠定了大地，并在出自祂的圣言内创造万有，且在圣子内坚定了圣洗。因为父在哪里，子就在那里；光在哪里，光耀就在那里。父所行的，都藉着子而行，正如主亲口说：\BibleQuote{我看见父所作的，我也照样作。}同样，施洗时，父所洗的，子也洗他；子所洗的，便是在\OriginalTerm{圣神}{Holy Ghost}内受祝圣。再者，犹如太阳照耀时，可以说光耀在照明，因为光是唯一而不可分的，也不能分离；同样，哪里有父，或父在哪里被呼名，那里就有子；圣洗中既呼父的名，也就必须与祂一同呼子的名。\par

42. 因此，当祂向\OriginalTerm{圣人}{saints}许下恩许时，祂这样说：\BibleQuote{我与我父要来到他那里，并在他那里作我们的住所}；又说：\BibleQuote{好使他们合而为一，如同我们合而为一}。所赐的恩宠是一，由父在子内赐下，正如\OriginalTerm{保禄}{Paul}在每封\OriginalTerm{书信}{Epistle}中所写：\BibleQuote{愿恩宠与平安，由我们父天主和主耶稣基督赐与你们。}因为光必须与光线同在，光耀也必须与其本光一同被瞻仰。因此，\OriginalTerm{犹太人}{Jews}怎样否认子，也就怎样没有父；因为，正如\OriginalTerm{巴路克}{Baruch}所责备的，他们离弃了\BibleQuote{智慧之泉源}\BibleRef{巴 3:12}，也就拒绝了由泉源涌出的智慧——我们的主耶稣基督（因为\OriginalTerm{宗徒}{Apostle}说，基督是\BibleQuote{天主的德能和天主的智慧}\BibleRef{格前 1:24}），那时他们说：\BibleQuote{我们除了凯撒，没有君王。}\BibleRef{若 19:15}这样，犹太人因他们的否认而受惩罚；他们的城邑与他们的理性一同归于乌有。这些人也危及了\OriginalTerm{奥迹}{mystery}的圆满，我指的即是圣洗；因为若祝圣是奉父及子之名授予我们的，而他们不宣认一个真父，因为他们否认那出自父、且与父本体相似的，也否认一个真子，却另立一位自己所虚构、从虚无中造出的，那么，他们所施行的礼仪岂非全然空洞无益，徒有虚表，却无补于\OriginalTerm{虔敬}{religion}吗？因为\OriginalTerm{亚略派}{Arians}并非因父及子施洗，而是因创造者和受造物、因制造者和作品施洗。正如受造物不同于圣子，同样，他们自以为施行的圣洗，也不同于真理，尽管他们假借\OriginalTerm{圣经}{Scripture}的字句，装作呼求父与子的名。因为，并非单单说\Quote{主啊}的人就施洗；而是那有名却也有\OriginalTerm{正信}{right faith}者。为此，我们的救主也没有单单命令施洗，而是先教导说：\BibleQuote{教导}，然后说：\BibleQuote{因父及子及圣神之名给他们施洗}；好使正信随学习而来，并与\OriginalTerm{信德}{faith}一同到来的是圣洗的\OriginalTerm{祝圣}{consecration}。\par

43. 还有许多其它异端，正如我所言，只使用言辞，却非以正确的意思，亦非以健全的信德，结果他们施行的水毫无益处，因缺乏虔诚，以致被他们洒水者，与其说是得救，不如说是被不虔所污染。外邦人也是这样，虽然口称天主的名，却遭\OriginalTerm{无神论}{Atheism}的指控，因为他们不认识真实而真正的天主，即我们的主耶稣基督的父。同样，摩尼教徒、弗里吉亚人、\Place{撒摩撒塔}的弟子，虽然使用这些名号，仍然是异端者，而亚略派也循着同样的路径，虽然他们阅读圣经的文字，使用这些名号，却嘲笑那些从他们领受圣洗的人，比其它异端更不虔，且更甚于它们，因自己言语的放肆而使那些异端显得无辜。因为这些其它的异端在某个方面违背真理，或关于主的身体犯错，好像祂没有从\Person{玛利亚}取肉身，或祂根本没有死，也没有成为人，只是显现，并非真实，似乎有身体其实没有，似乎有人形，如梦境中的幻影；但亚略派毫不掩饰地对父本身不虔。因为他们从圣经听到，祂的\OriginalTerm{天主性}{Godhead}在圣子内表现如在\OriginalTerm{肖像}{Image}中，他们就亵渎说，那是一个受造物，并且关于那肖像，他们到处带着\Quote{祂不存在}这句话，如同钱包里的泥土，像毒蛇吐毒一样吐出来。然后，当他们的教义令所有人生厌时，立刻，为了支持它免于倒塌，他们用人的庇护支撑异端，以便头脑简单的人，因眼见或甚至因惧怕，而忽略他们悖谬的祸害。的确应怜悯受他们欺骗的人；为他们哭泣是应当的，因为他们牺牲自己的益处换取那由快乐提供的眼前幻象，而丧失了未来的希望。他们以为领受圣洗归于一个不存在者的名，就将一无所获；将自己与受造物同列，就得不到受造界的帮助；相信一个在本质上与父不同且陌生者，他们就无法与父结合，因为按本性他们没有祂自己的圣子，这圣子是出于祂、在父内、父也在祂内的，正如祂自己所说；但因被他们引入歧途，这些可怜的人从此失去了天主性，赤贫如洗。因为这种属世财物的幻象，在他们死后不会跟随他们；当他们看见他们所否认的主，坐在祂父的宝座上，审判生者死者时，他们也不能呼求现在欺骗他们的任何人帮助他们；因为他们也将看见这些人在审判座前，为自己的罪恶和不虔的行径而懊悔。\par

% structure: p0053 source=h2 target=subsection reason=relative-heading-rank; base=h2

\subsection{第19章。经文解释；第六，\BibleBook{}{箴言} 8:22。箴言具有比喻性质，必须照此解释。我们必须根据\OriginalTerm{信仰规则}{Regula Fidei}解释它们，特别是这段经文。\Quote{祂造了我}不等同于\Quote{我是一个受造物}。智慧就其人性身体而言是受造物。再者，如果祂是一个受造物，那是作为\Quote{道路的开端}，这是一个职分，虽非属性，却是更高神性本质的结果。并且是\Quote{为了工程}，这暗示工程存在，因此远在祂被造之前，祂就已存在。还有，\Quote{主}并非父\Quote{造了}祂，这意味着那次创造是仆人的创造。}

44. 我们已经在\BibleBook{}{箴言}的那段经文前讲了这么多，抵制他们内心发明的昏愚寓言，好让他们知道天主的圣子不应被称为受造物，并学会轻松地阅读那些实有正确解释的经文。因为经上写着：\BibleQuote{上主造我作祂道路的开端，为了祂的工程}；然而，既然这些是箴言，并以箴言的方式表达，我们不可按字面意思直接解释，而必须探究位格，从而虔诚地赋予其意义。因为在箴言中所说的，不是直白道出，而是隐晦地提出，正如主自己在\BibleBook{若望}{福音}中教导我们的，说：\BibleQuote{这些事我对你们说了，是用比喻；时候将到，我不再用比喻对你们说话，而要明明地告诉你们} \BibleRef{若 16:25}。因此，必须揭示所说之事的含义，如同寻求隐藏之物，不可径直解释，好像那意义是\Quote{明明}说出的，免得因错误的解释而偏离真理。那么，如果所写的既指天使，或任何别的受造物，如同关于我们这些工程中的一员，就可以说\Quote{造了我}；但如果那是天主的智慧，万物都是藉着祂而造成的，而祂是在论说自身，我们应该理解什么？就是\Quote{祂造了}并不与\Quote{祂生了}相悖。智慧并不忘记自己是造物主和形成者，也不忽略造物主与受造物的区别，而将自己列于受造物中；而是暗示某种意义，如同在箴言中，不是\Quote{明明}的，而是隐晦的；这智慧曾启示圣人们在预言中使用，随后祂自己又以其他类似的表述，亲自解释了\Quote{祂造了}的意义，说：\BibleQuote{智慧建造了自己的房舍} \BibleRef{箴 9:1}。显然，我们的身体就是智慧的房舍，这身体祂取了来成为人；因此，\Person{若望}也一致地说：\BibleQuote{圣言成了血肉} \BibleRef{若 1:14}；而\Person{撒罗满}借智慧之口，以谨慎的精确性论说自身，不是\Quote{我是一个受造物}，而只是\Quote{上主造我作祂道路的开端，为了祂的工程}，却不是\Quote{造我使我存在}，也不是\Quote{因为我有受造物的开端和根源}。\par

45. 因为在这一段经文中，\OriginalTerm{圣言}{Word}借着\Person{撒罗满}说话，并不是指明祂天主性本体的实质，也不是指明祂由圣父而来的永恒而本真的生成；恰恰相反，是指明祂的人性，以及祂为我们所行的\OriginalTerm{经世}{Economy}。而且，如我之前所说，祂并没有说\Quote{我是受造物}或\Quote{我成了受造物}，而只是说\Quote{祂造了}。因为受造物具有受造的实质，是有起源的，且被说成是受造的，当然受造物就是受造的；然而，\Quote{祂造了}这个单纯的词语并不必然表示实质或生成，而是指明在它所谈论的那一位身上所发生的其他事，并不单纯是说那位被说成受造的，在其本性和实质上立刻就是受造物。这一区别，圣经也识别出来，论及受造物说：\BibleQuote{大地充满了祢的化工}，又说：\BibleQuote{一切受造之物都一同叹息，同受产痛}；在\BibleBook{若望}{默示录}中则说：\BibleQuote{海中一切有生命的受造物死了三分之一}；同样\Person{保禄}也说：\BibleQuote{天主所造的样样都好，如以感恩的心领受，没有一样是可摈弃的}；而\BibleBook{智慧}{篇}上亦记载：\BibleQuote{祢以祢的智慧造了人，叫他统治祢所造的万物}\BibleRef{智 9:2}。而这些既是受造物，也就被称为受造的，正如我们还可从主那里听到的话，祂说：\BibleQuote{那造了他们的，造了男造了女}；也可从\Person{梅瑟}在歌中听到的话，他写道：\BibleQuote{你且试问过去的日子，就是在你以前的日子，自从天主在地上造了人那天起，从天的这一边到天的那一边}\BibleRef{申 4:32}。\Person{保禄}在\BibleBook{哥罗森}{书}中说：\BibleQuote{祂是不可见的天主的肖像，是一切受造物的首生者，因为在天上和在地上，可见的与不可见的，无论是上座者，或是宰制者，或是率领者，或是掌权者，都是在祂内受造的：一切都是借着祂，并且是为了祂而受造的。祂在万有之先就有}。\par

46. 那么，被称为受造物以及受造，这属于按本性具有受造实质之物，这些章节足以提醒我们，尽管圣经中类似的章节比比皆是；另一方面，\Quote{祂造了}这个单纯的词语并不简单表示实质及其生成方式，\Person{达味}在圣咏中指明说：\BibleQuote{愿此事为未来的世代记下，愿受造的子民赞美主}；又说：\BibleQuote{天主，求祢给我造一颗纯洁的心}；\Person{保禄}在\BibleBook{厄弗所}{书}中说：\BibleQuote{废除了那由规条所组成的诫命律法，为将二者在自己身上造成一个新人}\BibleRef{弗 2:15}；又说：\BibleQuote{穿上新人，这新人是按照天主的肖像，在正义和真实的圣德中受造的}。因为\Person{达味}既不是指某个在实质上受造的民族，也不是祈求得到一颗与他已有的不同的心，而是指按照天主的更新和重生；\Person{保禄}并非指主内有两个在实质上受造的人，也不是劝告我们穿上另一个人；而是他把合乎德行的生活称为\Quote{按照天主的新人}，并用在基督内\Quote{受造}指在祂内获得更新的两个民族。\BibleBook{耶肋米亚}{书}中的话也是如此：\BibleQuote{主创造了一个新的救恩，为了一项种植，使人们在这救恩中往来行走。}；他说这话，并不是指某个受造物的实质，而是预言在人间救恩的更新，这救恩已经在基督内为我们完成了。那么，\Quote{受造物}和单纯的\Quote{祂造了}之间既有这样的区别，如果你在圣经的任何地方能找到主被称为\Quote{受造物}，就拿出来证明并争辩吧；但如果没有任何地方写着祂是受造物，只是祂自己在\BibleBook{}{箴言}中论到自己说：\Quote{主造了我}，你们就该羞愧，既由于上述的区别，也因为那是箴言式的措辞；因此，\Quote{祂造了}应当被理解为，不是指祂是受造物，而是指那成了祂的人性，因为这人性才属于受造。的确，当\Person{达味}和\Person{保禄}说\Quote{造了}时，你们就不理解为实质和生成，而理解为更新；但当主说\Quote{造了}时，你们却把祂的实质列于受造物之中，这难道不是明显的不公正吗？再者，当圣经说：\BibleQuote{智慧建造了房舍，立起了七根支柱}\BibleRef{箴 9:1}，你们就把\Quote{房舍}作寓意的解释，却按字面意思理解\Quote{祂造了}，并把受造物的观念强加给它。不但祂为万有创造者的身份对你们毫无分量，而且你们也不畏惧祂是父独有且本真的后裔，却鲁莽行事，好像你们已征兵对抗祂一样，竟如此争战，视祂不如人。\par

47. 因为这同一段经文证明了，称主为受造物只是你们自己的捏造。因为深知自己的\OriginalTerm{性体}{Essence}是\OriginalTerm{独生智慧}{Only-begotten Wisdom}和圣父的\OriginalTerm{后裔}{Offspring}，与一切受造的和自然的事物不同的主，出于对人的爱，说：\Quote{\BibleQuote{上主造了我，作祂道路的起始}}，仿佛是说：\Quote{\InnerQuote{我父给我预备了一个身体，并为人的得救而创造了我}}。因为，就如\Person{若望}说：\Quote{\BibleQuote{\OriginalTerm{圣言}{Word}成了血肉}\BibleRef{若 1:14}}，我们并不认为整个圣言本身是血肉，而是取了血肉而成为人；在听到\Quote{\BibleQuote{基督为我们成了诅咒}}和\Quote{\BibleQuote{那不认识罪的，祂替我们成了罪}}时，我们并不简单地认为，整个基督变成了诅咒和罪，而是祂背负了那加在我们身上的诅咒（正如\OriginalTerm{宗徒}{Apostle}所说：\Quote{\BibleQuote{由诅咒中赎出了我们}}，又如\Person{依撒意亚}所说：\Quote{\BibleQuote{背负了我们的罪}}，以及\Person{伯多禄}所写：\Quote{\BibleQuote{在木架上，亲自承担了我们的罪}}）；同样，如果在\BibleBook{箴}{言}中说\BibleQuote{祂造了}，我们不可认为，整个圣言在性体上是受造物，而是祂穿上了受造的身体，天主为了我们而创造了祂，给祂预备了那受造的身体，如同经上所载，为我们，好使我们在祂内得以更新和被\OriginalTerm{被神化}{deified}。那么，愚妄的人啊，是什么欺骗了你们，竟称\OriginalTerm{创造主}{Creator}为受造物呢？你们是从哪里买来了这个新思想，到处炫耀呢？因为《\BibleBook{箴}{言}》说\BibleQuote{祂造了}，却未称圣子为受造物，而是后裔；而且，根据前面所说圣经中\BibleQuote{祂造了}与\Quote{受造物}的区别，他们承认圣子本性所固有的：祂是独生智慧和受造物的\OriginalTerm{创造者}{Framer}；而当他们说\BibleQuote{祂造了}时，并不是针对祂的性体说的，而是表明祂成为\OriginalTerm{众多道路的起始}{beginning of many ways}；因此，\BibleQuote{祂造了}是与\Quote{后裔}相对，而称祂为\BibleQuote{道路的起始}，则与祂是独生圣言相对。\par

48. 因为如果祂是\OriginalTerm{后裔}{Offspring}，你们怎能称祂为\OriginalTerm{受造物}{creature}呢？因为没有人会说，祂所创造的，是祂生出的，也不会称自己的亲生后裔为受造物；再者，如果祂是\OriginalTerm{独生者}{Only-begotten}，又怎能成为\BibleQuote{\OriginalTerm{道路的起始}{beginning of the ways}}呢？因为如果祂被创造为万物的起始，那祂就不再是独一的了，因为有那些在祂之后而来的存在。比如\Person{勒乌本}，当他成为子孙的起始时，并不是独生子，虽然在时间上为第一，但在性体和关系上，却是后来者之一。因此，如果\OriginalTerm{圣言}{Word}也是\BibleQuote{道路的起始}，那么祂必定像那些道路一样，那些道路也必定像圣言一样，尽管在时间上祂是其中首先被造的。因为一座城的起始或开端，就像城的其他部分一样，各部分连接起来，使城市完整合一，如同一身的许多肢体；并不是一部分创造，另一部分产生并从属于前者，而是全城同样都由其建造者获得治理和构造。那么，如果主就是以这样的意义而被造为万物的\Quote{起始}，结论便是：祂与所有其他事物共同组成\OriginalTerm{受造界}{creation}的统一体；祂虽成为万物的\Quote{起始}，却与其他事物无异；祂虽在时间上较古，却不是它们的主；而是与其余事物有着同样的受造方式，同一位主。不止如此，如果如你们所坚持，祂是受造物，祂又如何能独独地、首先地被造，以至成为万物的起始呢？从以上所述可见，在受造物中，没有一个具有永恒的性体及先于其他的受造，而是各自与所有其他事物一同起源，尽管在荣耀上可能超过其他。就如那分离的星辰或大光，并不是这个先出现，那个后出现，而是在同一天，由同一命令，它们一起蒙召而存在。四足兽、飞鸟、鱼类、牲畜、植物的最初形成，也是如此；按照\OriginalTerm{天主的肖像}{God's Image}而造成的人类，也是如此，即人类；因为虽然只有亚当是由泥土形成的，但在他内，却包括了整个人类的延续。\par

49. 从可见的受造界，我们清楚地分辨出，祂那不可见的事物，\BibleQuote{借着所造的万物，可被辨见}\BibleRef{罗 1:20}，也不是各自独立的；因为并非先有一个，然后另一个，而是按类同时被组成。因为宗徒并没有一一列举，如说\Quote{或\OriginalTerm{天使}{Angel}、或\OriginalTerm{上座者}{Throne}、或\OriginalTerm{宰制者}{Dominion}、或\OriginalTerm{掌权者}{Authority}}，而是按其类别一起提及：\Quote{或\OriginalTerm{众天使}{Angels}、或\OriginalTerm{总领天使}{Archangels}、或\OriginalTerm{率领者}{Principalities}}：因为受造物的起源便是如此。那么，如我所言，如果圣言是受造物，祂就必不会首先于它们受造，而是与所有其他\OriginalTerm{异能者}{Powers}一起受造，尽管在荣耀上祂远超其他。因为我们发现它们的情况正是如此：它们是一同出现，既无第一亦无第二，且彼此在荣耀上有所区别，有的在宝座右边，有的环绕四周，有的在左边，但全体一致赞颂侍立于主前。因此，如果圣言是受造物，祂就不会是其他一切的第一或起始；然而，如果祂在万有之先，正如实际所是，且唯独祂是\OriginalTerm{首生者}{First}和子，这并不表示，就其\OriginalTerm{性体}{Essence}而言祂是万有的起始，因为万有的起始本身就在万有之数内。而如果祂不是这样的起始，那么祂也不是受造物；相反，这是非常明显的：祂在性体和本性上与受造物不同，且与它们相异，是唯一真天主的\OriginalTerm{肖似}{Likeness}和\OriginalTerm{肖像}{Image}，祂自己也是唯一的。因此，在圣经中祂不被列于受造物中，而是\Person{达味}责备那些竟敢如此想祂的人，说：\Quote{\BibleQuote{众神之中，谁能像主？}}和\Quote{\BibleQuote{在天主的众子中，谁能像主？}}；\Person{巴路克}也说：\Quote{\BibleQuote{这是我们的天主，再没有别的可同祂相比}\BibleRef{巴 3:35}}。因为那位创造，其余的都是受造；那位是圣父性体的本有\OriginalTerm{圣言}{Word}和\OriginalTerm{智慧}{Wisdom}，藉着这圣言，先前不存在的万物得以造成。\par

50. 那么，你们那著名的断言——即圣子是受造物——并不真实，只是你们的幻想而已；不仅如此，\Person{撒罗满}证明你们多次诽谤了他。因为他并没有称他为受造物，而是称他为天主的后裔和智慧，说：\BibleQuote{天主以智慧奠定了大地}，\BibleQuote{智慧建造了自己的房舍。}而所讨论的这段经文本身就证明了你们不虔敬的精神；因为经上写着：\BibleQuote{主自始即造了我，作为他道路的起始，为他的工程。}因此，如果他在万有之前，却说\BibleQuote{他造了我}（不是\Quote{为使我造作工事}，而是）\BibleQuote{为工程}，除非\BibleQuote{他造了}是指后来于他自己的事，否则他会显得晚于那些工程，在他的受造时，那些工程已先于他存在，而他就是为此被造出来的。如果这样，他又如何在万有之前呢？又如何万有是借着他并在内而存在的呢？看哪，你们说那些工程先于他存在，而他就是为了这些工程而被造和受差遣的。但是事实并非如此；断不可存此想念！异端者的假设是虚假的。因为天主圣言不是受造物，而是创造者；他以\BibleBook{}{箴言}的方式说\BibleQuote{他造了我}，是指他穿上了受造的肉身。从这段经文本身也可以理解出别的东西；因为，作为圣子，有天主为他的圣父——因为他是父本有的后裔——然而在这里他却称圣父为主；并不是因为他是仆人，而是因为他取了仆人的形象。因为作为出自父的圣言，他称天主为父是合宜的：这是儿子对父亲当有的事；但当他来完成那工程，并取了仆人的形象时，称父为主也是合宜的。他自己通过一个恰当的区别教导了这一分别，在福音书中说：\BibleQuote{父啊！我感谢你}，然后\BibleQuote{天地的主宰！}\BibleRef{玛 11:25}。因为他称天主为他的父，但对受造物，他称他为主；由此清楚表明，当他穿上了受造物时，便称父为主。因为在\Person{达味}的祈祷中，圣神也标明了同样的分别，在圣咏中说：\BibleQuote{求你赐给你的孩子力量，拯救你婢女的儿子。}因为那本性且真正的天主子只有一位，而婢女的子女——即出自受造物本性的——则是另一些。因此，那作为圣子的那一位，拥有父的权能；而其余的则需要救恩。\par

51. （但如果他们因为他被称为儿子而妄加议论，要让他们知道，\Person{依撒格}也被称为\Person{亚巴郎}的儿子，\Place{叔能}妇人的儿子也被称为小孩子。）因此，我们是仆人，当他成为我们一样时，他也像我们一样称父为主；他这样做是出于对人的爱，为使我们这些本为仆人的人，因领受圣子的神，能有信心藉着恩宠称那本性是我们主的天主为父。但是，正如我们称主为父，并不否认我们本性的奴仆身份（因为我们是他的工程，是\BibleQuote{他造了我们，不是我们自己}），同样，当圣子取了仆人的形象，说\BibleQuote{主自始即造了我，作为他道路的起始}时，他们就不应否认他天主性的永恒，也不应否认\BibleQuote{在起初已有圣言}、\BibleQuote{万有是借着他而造成的}、\BibleQuote{在他内万有都受造了。}这些真理。\par

% structure: p0062 source=h2 target=subsection reason=relative-heading-rank; base=h2

\subsection{第二十章 经文解释；第六，\BibleBook{}{箴言} 8:22 续。我们的主被称为受造\BibleQuote{为工程}，即有一个特定的目的，这是任何单纯的受造物从未被说过的。与 \BibleBook{依撒意亚}{} 49:5 等相平行。当提到他的人性时，会加上原因；提到他的天主性时则不然。证明的经文。}

51（\Emph{续}）。因为\BibleBook{}{箴言}中的那段经文，正如我之前说过的，所指的并非圣言的本质，而是他的人性；因为如果他说他受造是\BibleQuote{为工程}，他就表明他的用意不是指他的本质，而是指那\BibleQuote{为他的工程}而发生的、次于存在的\OriginalTerm{经世}{Economy}。因为那些处于形成和受造中的事物，特别被造是为了使它们存在和生存，然后它们才去做圣言所吩咐的事，正如在万物身上可以看到的。因为\Person{亚当}被造，不是为了工作，而是为了首先成为人；在这之后，他才领受了工作的命令。\Person{诺厄}被造，不是为了方舟，而是为了首先存在并成为人；在这之后，他才领受了准备方舟的命令。同样的情况，在调查之后可以发现于每一事例中——因此伟大的\Person{梅瑟}首先被造为人，然后才被托付治理百姓的职责。所以，这里我们也必须作相似的假定；因为你们看，圣言并非被造成存在，而是\BibleQuote{在起初已有圣言}，他随后为了工程和朝向它们的经世而被差遣。因为在工程被造之前，圣子已经存在，也不需他被造；但当工程被造，之后又有了为修复它们的经世之需要时，圣言才取了这种\OriginalTerm{屈尊就卑}{condescension}和与工程的同化；他通过\BibleQuote{他造了}这句话向我们显示了这一点。而且通过先知\Person{依撒意亚}，为表达相似的意思，他又说：\BibleQuote{如今主这样说：他自母胎形成我作他的仆人，为把\Person{雅各伯}和\Person{以色列}聚集到他那里——我必在主前被聚集并受荣耀。}\par

52. 在此也要注意，祂被形成，并非为了存在，而是为了聚集那些在祂被形成之前就已存在的支派。因为如前一段经文说\BibleQuote{祂创造}，此处则说\BibleQuote{祂形成}；并且如前段说\BibleQuote{为了工程}，此处则说\BibleQuote{为聚集}；以致从各个角度看，\BibleQuote{祂创造}与\BibleQuote{祂形成}都是在\BibleQuote{圣言已存在}之后才说的。正如在祂被形成之前，支派已经存在，而祂正是为了它们而被形成的；同样，那些工程也已存在，祂正是为了它们而被创造的。当\BibleQuote{在起初已有圣言}时，工程尚不存在，正如我先前说过的；但当工程已被造成，且需要迫切时，那时才说\BibleQuote{祂创造}；这就像某个儿子，当仆人们因自身的疏忽而落入敌手，且情势危急时，被父亲派遣去救助和收复他们；他出发时，穿上与他们相同的衣服，使自己打扮得与他们一样，以免俘获者认出他是主人而逃跑，阻止他下到那些被他们隐藏在地下的人那里；若有人问他为什么这样做，他回答说：\Quote{我父亲这样形成并预备我，为他的工程}，他这样说，既不表明自己是仆人，或工程之一，也不提及他原本受生的起始，而是指此后交付给他治理工程的责任——同样，主也披上了我们的血肉，\BibleQuote{具有人的形状}，如果那些看见祂如此并感到惊奇的人问祂，祂会说：\BibleQuote{主创造了我，作祂工程的起始，为祂的工程}，以及\BibleQuote{祂形成我，为聚集以色列}。圣神在\WorkTitle{圣咏}中再次预言说：\BibleQuote{祢立祂于祢手所造的工程之上} \BibleRef{希 2:7}；而主在别处指自己说：\BibleQuote{我被祂立为君王，在祂的圣山\Place{熙雍}之上}。正如祂在肉身中显现于\Place{熙雍}时，并非那时才开始存在或为王，而是身为天主的圣言和永恒的君王，祂甘愿让祂的王权以人的方式彰显于\Place{熙雍}，好将那些受罪恶辖制的人及我们从罪恶中救赎出来，带他们进入祂父的国度；同样，当祂被立\BibleQuote{为了工程}时，不是为尚未存在的，而是为已经存在且需要恢复的工程。\par

53. 因此，\BibleQuote{祂创造}、\BibleQuote{祂形成}和\BibleQuote{祂立}，具有相同意义，并不表示祂存在的开始，或祂受造的本质，而是指祂为我们成就的仁慈的更新。据此，虽然祂这样说，但祂也教导说，祂自己在此前就已存在；当祂说：\BibleQuote{在\Person{亚巴郎}出现以前，我就有} \BibleRef{若 8:58}；又说：\BibleQuote{当祂建立高天时，我已在场}，又说：\BibleQuote{我在祂身旁，作技师}。正如祂自己在\Person{亚巴郎}出现之前就已存在，而以色列是在\Person{亚巴郎}之后才存在的，显然祂先存在，后被形成；祂的被形成不表示祂存在的开始，而是指祂摄取人性，在其中祂也聚集以色列的支派；同样，由于\BibleQuote{常与圣父同在}，祂自己是受造界的塑造者，祂的工程显然晚于祂自己，\BibleQuote{祂创造}所指的，不是祂存在的开始，而是为工程而施行的\OriginalTerm{经世}{Economy}，即祂在肉身中完成的。因为，祂既与工程有别，更是它们的塑造者，由祂来承担工程的更新，原是合宜的；这样，由于祂为我们而被创造，万有便可在祂内受造了。因为当祂说\BibleQuote{祂创造}时，立即加上理由，指明\BibleQuote{工程}，以便祂为工程而被创造，正表明祂为更新它们而成为人。这是圣经的惯例；因为当它表明圣子肉身的起源时，也加上祂成为人的原因；但每当祂自己或祂的仆人论及祂的天主性时，一切都以简明的言辞和绝对的意义说出，不加任何理由。因为祂是圣父的光辉；正如圣父存在，却不需要任何理由，我们也无需探求那光辉的理由。因此经上记载：\BibleQuote{在起初已有圣言，圣言与天主同在，圣言就是天主} \BibleRef{若 1:1}；未指明原因；但当\BibleQuote{圣言成了血肉} \BibleRef{若 1:14} 时，才加上原因说：\BibleQuote{寄居在我们中间}。宗徒也说：\BibleQuote{祂虽具有天主的形体}，并未引入理由，直到\BibleQuote{祂取了奴仆的形体}；然后才继续：\BibleQuote{祂贬抑自己，听命至死，且死在十字架上} \BibleRef{斐 2:6}；因为正是为此，祂才成了血肉，并取了奴仆的形体。\par

54. 主自己说了许多比喻；但当他论及自己时，却说得很绝对：\BibleQuote{我在父内，父在我内}，\BibleQuote{我与父是一体}，\BibleQuote{看见我的，就是看见了父}，\BibleQuote{我是世界的光}，\BibleQuote{我是真理}，却没有每次都说明理由和原因，免得他显得比那些他为之而被造的事物次要。因为理由本来必须先于他，没有这理由，连他自己也不会存在。例如，\Person{保禄} \BibleQuote{被分别为宗徒，为了福音——这福音是天主先前借众先知所应许的}\BibleRef{罗 1:1}，他就因此隶属于福音，成了福音的仆役；\Person{若望}被选为先驱预备主的道路，就隶属于主。但主却并不隶属于任何他之所以成为圣言的理由，只除了他是圣父的后裔和独生智慧这一点；当他成为人时，就说明了他即将取肉身的原因。因为人类的急需先于他降生成人，若非如此，他就不会穿上肉身。他自己这样指出他成为人的原因：\BibleQuote{我自天降下，不是为行我的旨意，而是为行派遣我来者的旨意。派遣我来者的旨意就是：凡他交给我的，叫我连一个也不失掉，且在末日要使他复活。因为这是我父的旨意：凡看见子并信从他的，必得永生，在末日我要使他复活。}又说：\BibleQuote{我来到世界上，作为光，使凡信我的，不住在黑暗中。}又说：\BibleQuote{我为此而生，也为此而来到世界上，为给真理作证。}\Person{若望}写道：\BibleQuote{天主子所以显现出来，是为消灭魔鬼的作为。}\BibleRef{若一 3:8}\par

55. 所以，救主来，是为作证，并为我们忍受死亡，使人类复活，并消灭魔鬼的作为；这就是他降生亲临的理由。因为若不先有死亡，就不会有复活；而他若不取一个可死的身体，又怎能死亡呢？宗徒从他那里学到这一点，这样阐明：\BibleQuote{那么，儿女既然都有同样的血肉，他也照样取了同样的血肉，为能藉著死亡，消灭那握有死亡权势者——就是魔鬼，并解救那些因死亡的恐惧，一生当奴隶的人。}\BibleRef{希 2:14} 又说：\BibleQuote{死亡既因一人而来，死者的复活也因一人而来。}\BibleRef{格前 15:21} 又说：\BibleQuote{法律因了肉性所不能行的，天主却行了：他派遣了自己的儿子，取了罪恶的肉身的形状，当作赎罪祭，在这肉身上定了罪的罪案，为使法律所要求的正义，成全在我们今后不随从肉性，而随从圣神生活的人身上。}\BibleRef{罗 8:3} \Person{若望}说：\BibleQuote{因为天主没有派遣子到世界上来审判世界，而是为叫世界藉着他而获救。}\BibleRef{若 3:17} 救主也亲口说过：\BibleQuote{我是为了审判，才到这世界上来，叫看不见的，看得见；叫看得见的，反而成了瞎子。}因此，他来，不是为了自己，而是为了我们的救恩，为消灭死亡，定罪罪过，使瞎子看见，使众人从死者中复活。如果他不为自己而来，而是为我们，那么，他也并非为自己，而是为我们而被造。然而，如果他并非为自己而被造，而是为我们，那么他自己并不是受造物，而是因为穿上了我们的肉躯，才用这样的言语。这正是圣经的意思，我们可以从宗徒学到，他在\BibleBook{厄弗所}{书}中说：\BibleQuote{他是我们的和平，他拆毁了我们中间阻隔的墙壁，就是仇恨，并在他自己的肉身上，废除了那法律上诫命所规定的规条，为把双方在自己身上造成一个新人，而成就和平。}\BibleRef{弗 2:14} 如果双方在他内被造成，而这双方是在他的身体内，那么，他既然在自己内肩负了双方，也就可以说好像自己是被造的；因为他将在他内受造的双方造成一个，他也在他们内，如同他们在自己内一样。这样，双方既在他内受造，他就可以恰当地说：\BibleQuote{上主造了我}。因为正如他承担了我们的软弱，自己也被称为软弱的，他自己并非软弱，因为他是天主的大能；他为我们成了罪，成了诅咒，自己却没有犯罪，只因他亲自背负了我们的罪和诅咒；同样，因着在我们内创造我们，他可以说：\BibleQuote{他为工程造了我}，尽管他自己不是受造物。\par

56. 因为，若如他们所持守，\OriginalTerm{圣言}{Word}的本质属于受造的本性，因此祂说：\BibleQuote{主造了我}，那么祂既是\OriginalTerm{受造物}{creature}，就不是为我们而受造；但若祂不是为我们而受造，我们也就不是在祂内受造的；而若我们不是在祂内受造的，我们便不是内在地拥有祂，而是以外在的方式，例如，像从一位老师那里接受教导一样。这在我们身上既是如此，罪便没有失去对肉身的统治，因为它仍内在其中，未被逐出。但\Person{宗徒}稍前一点便反对这种主张，他说：\BibleQuote{因为我们是祂的化工，是在基督耶稣内受造的}\BibleRef{弗 2:10}；而若我们是在基督内受造的，那么就并非祂受造，而是我们在祂内受造；因此\BibleQuote{祂创造}这话是为了我们的缘故。因为出于我们的需要，圣言虽是创造者，却忍受了用于受造物的言语；这些言语并不与身为圣言的祂相称，而是相称于我们这些在祂内受造的人。而且正如圣父如何永在，祂的圣言也如何永在，永远存在着，便永远说：\BibleQuote{我天天是祂的喜乐，不断在祂前欢跃}\BibleRef{箴 8:30}，以及\BibleQuote{我在父内，父在我内}\BibleRef{若 14:10}；同样，当祂为了我们的需要而成为人时，祂便一贯地使用我们人的语言，如\BibleQuote{主造了我}，好使因着祂居住于肉身，罪得以完全从肉身逐出，我们便能拥有自由的心灵。因为当祂成为人后，祂该说什么呢？\Quote{在起初我是人}吗？这既与祂不相称，也不真实；而既然不宜这样说，那么在人的情形中说\BibleQuote{祂创造}和\BibleQuote{祂造了}祂，便是自然且适当的。正因如此，才加上了\BibleQuote{祂创造}的理由，即工作的需要；而既然加上了理由，这理由便是正确地解释了经文。因此在这里，当祂说\BibleQuote{祂创造}时，祂给出了原因——\BibleQuote{工作}；另一方面，当祂绝对地表明从圣父的生育时，便立即加上：\BibleQuote{在山岳未有之前，祂生了我}\BibleRef{箴 8:25}；但祂并没有如在\BibleQuote{祂创造}的情形中那样加上\Quote{为何}，说\BibleQuote{为了工作}，而是绝对地说\BibleQuote{祂生了我}，正如经上的话：\BibleQuote{在起初已有圣言}\BibleRef{若 1:1}。因为，即使没有任何工作受造，天主的\BibleQuote{圣言}仍然\BibleQuote{存在}，且\BibleQuote{圣言就是天主}。而祂成为人这件事，若非人的需要成为原因，便不会发生。所以，子不是受造物。\par

% structure: p0069 source=h2 target=subsection reason=relative-heading-rank; base=h2

\subsection{第二十一章 经文解释；第六，续\BibleBook{}{箴言}8:22。我们的主在圣经中并没有被称为\Quote{受造的}，或工作被称为\Quote{受生的}。\Quote{在起初}在工作的事上是指\Quote{从起初}。解释圣经章节。我们先是被天主所造，然后受生；本性上是受造物，借恩宠成为子。基督先受生，后被造或受造。\Quote{死者中的首生者}、\Quote{众多弟兄中的首生者}、\Quote{一切受造物的首生者}的意义，与\Quote{独生子}对比。进一步解释\Quote{道路的开端}和\Quote{为了工作}。为何受造物不能救赎；为何救赎根本是必要的。对比圣言与工作的经文。}

57. 因为，如果祂是\OriginalTerm{受造物}{creature}，祂就不会说：\Quote{祂生了我}，因为受造物来自外界，是制造者的作品；而子并非来自外界，也不是作品，而是出自父，并与父的本质相合。因此，它们是受造物；而祂是天主的圣言和\OriginalTerm{独生子}{Only-begotten}。例如，\Person{梅瑟}论及创造时，并没有说\Quote{在起初祂生了}，也没有说\Quote{在起初已有}，而是说\BibleQuote{在起初天主创造了天地}\BibleRef{创 1:1}。\Person{达味}在圣咏中也没有说：\Quote{你的双手 \InnerQuote{生了我}}，而是说\Quote{创造了我，形成了我}，将\Quote{造}这个词遍用于受造物。但对圣子则相反；因为他并没有说\Quote{我造了}，而是说\Quote{我生了}，以及\Quote{祂生了我}，还有\Quote{我的心灵涌溢出一句好圣言}。而在创造的事例中，是\Quote{在起初祂造了}；但在圣子的事例中，则是\BibleQuote{在起初已有圣言}\BibleRef{若 1:1}。而且有这样的差别：受造物是在起初受造，有一个与时间段相连的存在之始；因此，论及它们时所说的\Quote{在起初祂造了}，就等于说\Quote{从起初祂造了}——正如主了解祂所造的，在使法利塞人缄默时教导说：\BibleQuote{那从起初制造他们的，造了他们男和女}\BibleRef{玛 19:4}；因为从某个起初，当它们还不存在时，起源之物便被带入存在并受造。\OriginalTerm{圣神}{Holy Spirit}在圣咏中也表达了这一点，说：\Quote{主啊，你起初奠定了大地的基础}；又说：\Quote{求你纪念你自起初所获得的会众}；明显可知，在起初发生之事，有一个创造的起始，并且天主从某个起初获得了祂的会众。而关于\Quote{在起初祂造了}，从他说\Quote{造了}便意谓\Quote{开始造}，\Person{梅瑟}自己在完成一切之后说：\BibleQuote{天主降福了第七日，定为圣日，因为在这一天内，天主停止了祂所开始造的一切工作}\BibleRef{创 2:3}，这便表明了这点。因此，受造物开始被造；但天主的圣言，既没有存在之始，肯定不会开始存在，也不会开始形成，而是永恒存在的。它们的工作在受造中有其开始，且开始先于其形成；但圣言，不属于那些形成之物，反而自身成了那些有开始之物的形成者。起源之物的存在，是由其生成来度量的，天主通过圣言从某个起始开始造它们，为使人知道它们在被造之前并不存在；但圣言的存在，除了在无始的父内之外，别无其他起始，而他们承认父是无始的，因此祂同样在父内无始存在，是父的子，而不是受造物。\par

58. 因此，圣经认识到了\OriginalTerm{受生者}{Offspring}与受造之物之间的区别，并指明受生者是\Emph{子}，并非从任何开端开始，而是永恒的；但受造之物，作为造物主的外在工作，则开始存在。因此，\Person{若望}传授关于\Emph{圣子}的神圣道理，并知道这些用语的区别，他没有说\Quote{在起初成了}或\Quote{被造}，而是说\BibleQuote{在起初\Emph{有}圣言}，好使我们借着\Quote{有}理解\Quote{\OriginalTerm{受生者}{Offspring}}，不以时间间隔来度量祂，而是相信\Emph{圣子}始终且永恒地存在。有了这些证据，亚略派的人士啊，你们误解了\BibleBook{申命}{纪}中的经文，怎敢对主做出新的不敬行为，说\Quote{祂是一件工程}或\Quote{受造物}，甚或\Quote{受生者}呢？因为你们把受生者和工程视为同一回事；但在这里，你们也将被证明既无学识又不虔敬。你们的第一段经文是：\BibleQuote{祂岂不是你的父，赎买了你？岂不是祂造了你，创造了你？}很快在同一首颂歌中他又说：\BibleQuote{生你的天主，你竟离弃，忘了那养育你的天主。}这些经文所传达的意义非常值得注意；因为他并没有先说\Quote{祂生了}，唯恐那用语与\Quote{祂造了}相混淆，这些人就可以借口说：\Quote{梅瑟确实告诉我们，天主起初说了，\BibleQuote{让我们造人}\BibleRef{创 1:26}}但他在后文自己立即说：\Quote{生你的天主，你竟离弃}，仿佛这些用语毫无区别；因为受生者和工程是同一回事。但在说了\Quote{赎买}和\Quote{造了}之后，他最后加上了\Quote{生了}，好使这句话带有自己的解释；因为在\Quote{造了}这个词中，他精确地指出了属于人性的事，即他们是工程和受造之物；但在\Quote{生了}这个词中，他显示了天主在创造了他们之后向人施行的慈爱。既然他们对此表现出忘恩负义，于是\Person{梅瑟}就责备他们，先说：\BibleQuote{你们竟这样报答主吗？}接着加上：\BibleQuote{祂岂不是你的父，赎买了你？岂不是祂造了你，创造了你？}\BibleRef{申 32:6}接着他又说：\BibleQuote{他们祭祀邪魔，不祭祀天主，祭祀他们素不认识的神，就是近来新兴，你们祖先所不敬畏的神；生你的天主，你竟离弃。}\par

59. 因为天主不仅创造他们成为人，还召叫他们成为儿女，就如祂生了他们。因为\Quote{生了}一词在这里，如同在别处一样，是表达\Quote{子}的，就如祂借着先知说：\BibleQuote{我生了儿子，抚养了他们；}一般地，当圣经要指明一个儿子时，它不是用\Quote{受造}一词，而是无疑地用\Quote{生了}。而\Person{若望}似乎这样说：\BibleQuote{祂给了他们权能，好成为天主的子女，就是那些信祂名字的人；他们不是由血气，也不是由肉欲，也不是由男欲，而是由天主生的。}\BibleRef{若 1:12}在这里，谨慎的区别也被小心保存，因为他先说\Quote{成为}，因为他们不是本性上称为儿女，而是由于收养；然后他说\Quote{是生的}，因为他们无论如何的确已领受了儿子的名分。但是，正如先知所说，那百姓\Quote{藐视}了他们的恩主。但这是天主对人的慈爱：祂是他们的造主，但依照恩宠，祂后来也成为父亲；就是说，当人——祂的受造物——在心中接受了，正如宗徒所说，\Quote{祂儿子的圣神，呼喊说：\InnerQuote{阿爸，父啊！}}。这些人就是，接受了圣言，从祂获得了权能，成为天主的子女；因为他们本性是受造物，除非接受那本性和真正的圣子的圣神，否则就不能成为子女。因此，为成就这事，\BibleQuote{圣言成了血肉}，好使祂能使人有分于天主性。同样的意义也可从\Person{玛拉基亚}先知获得，他说：\BibleQuote{不是一个天主造了我们吗？我们岂不都共有一个父吗？}\BibleRef{拉 2:10}因为他先说\Quote{造}，后说\Quote{父}，为指明，如同其他作者所行的，从起初我们本性是受造物，天主借着圣言是我们的造主；但后来我们被造成为子女，从此天主造主也成为我们的父。因此，\Quote{父}是\Emph{圣子}所专有的；不是\Quote{受造物}，而是\Quote{圣子}是\Emph{父}所专有的。因此，这段经文也证明，我们本性上不是子女，而是那在我们内的圣子；并且再次证明，天主本性上不是我们的父，而是我们内圣言的父，在祂内并因着祂，我们\Quote{呼喊：\InnerQuote{阿爸，父啊！}}\BibleRef{迦 4:6}同样地，父在所有人中，只要祂看见自己的圣子，就称他们为儿子，并说：\Quote{我生了。}因为生是表示子，而造是表示工程。因此，我们不是首先生，而是受造；因为经上写着：\BibleQuote{让我们造人}\BibleRef{创 1:26}；但后来，在接受圣神的恩宠时，从那时起我们也说被生；正如伟大的\Person{梅瑟}在他的颂歌中，以贴切的意义先说\Quote{祂赎买}，然后说\Quote{祂生了}，以免人听见\Quote{祂生了}，可能忘记自己的原本本性；而是要他们知道，从起初他们是受造物，但当他们依照恩宠被说成生，作为子女时，他们仍然不亚于从前，是按本性而言的工程。\par

60. 至于\OriginalTerm{受造物}{creatura}与\OriginalTerm{子}{gennema}并非同一，而是本性及词义上彼此有别，主自己在《箴言》中也已表明。因为他说了：\BibleQuote{上主造了我，作他行动的开端}；之后又加上：\BibleQuote{但在一切丘陵未有之先，他已\OriginalTerm{生}{gennao}了我。}因此，倘若圣言按其本性及本质是受造物，且子与受造物没有区别，他就不会加上\BibleQuote{他生了我}，而只会满足于说\BibleQuote{他造了我}，好像那个词已蕴含\BibleQuote{他生了我}之意；但事实上，他说了\BibleQuote{上主造了我，作他行动的开端，为他的工程}之后，又加上\BibleQuote{他生了我}，并非简单地加，而是以连词\Quote{但}来连接，以护卫那\OriginalTerm{造}{ktizo}字，他说：\BibleQuote{但在一切丘陵未有之先，他已生了我。}因为\BibleQuote{他生了我}紧接在\BibleQuote{他造了我}之后，使意思合一，并显示\Quote{造}是为了某个目的而说的，但\BibleQuote{他生了我}却先于\BibleQuote{他造了我}。因为，假若他反过来先说\BibleQuote{上主生了我}，然后接上\BibleQuote{但在丘陵未有之先，他造了我}，那么\Quote{造}必然先于\Quote{生}；同样，他既然先说了\Quote{造}，然后加上\BibleQuote{但在一切丘陵未有之先，他已生了我}，这就必然表明\Quote{生}先于\Quote{造}。因为他说：\BibleQuote{在万有之先他生了我}，暗示他异于万有；本书前面已经证明，没有一个受造物先于另一个受造物而被造，万有一同凭着同一命令在起源上立即存在。因此，紧接在\Quote{造}之后的那些话，并不也接在\BibleQuote{生了我}之后；而是对于\Quote{造}，加上了\OriginalTerm{行动的开端}{arche hodon}，但对于\BibleQuote{生了我}，他没有说\BibleQuote{他生了我作为开端}，而是说\BibleQuote{在万有之先他生了我}。然而那在万有之先的，并非万有的开端，而是异于万有；而如果异于万有（这\Quote{万有}包括了一切的开端），则他必然是异于受造物的；这就清楚地表明，圣言既是异于万有，且在万有之先，随后才被\Quote{造}为\BibleQuote{工程之行动的开端}，因为他成了人，正如宗徒所说，\BibleQuote{他是元始，是死者中的\OriginalTerm{首生者}{prototokos}，为使他在万有之上独占首位}\BibleRef{哥 1:18}。\par

61. \Quote{造}与\BibleQuote{生了我}，以及\OriginalTerm{行动的开端}{arche hodon}与\BibleQuote{万有之先}之间既然有这样的区别，那么，天主作为最初的创造者，其次，正如已说过的，因祂的圣言居住在人们内，而成为人类之父。但在圣言身上，情况却相反；因为天主是祂由本性所生的父，随后，当圣言穿上那受造且被造成的血肉，成为人时，祂也成了祂的创造者和制造者。因为，如同人接受了圣子的\OriginalTerm{神}{pneuma}，借着他而成为子女，同样，天主圣言，当祂亲自穿上人的血肉时，那时才说祂\BibleQuote{受造}且\BibleQuote{被造成}。因此，如果我们是由本性而为儿子，那么祂就由本性而为受造物和工程；但如果我们是由\OriginalTerm{义子名分}{huiothesia}和\OriginalTerm{恩宠}{charis}而成为子女，那么，圣言也因着对我们的恩宠而成为人时，说了：\BibleQuote{上主造了我。}其次，当祂穿上了受造的本性，在身体上与我们相似时，祂被合理地称为我们的\BibleQuote{兄长}和\BibleQuote{\OriginalTerm{首生者}{prototokos}}。因为，虽然祂是在我们之后，为了我们而成为人，因身体相似而成为我们的兄长，但祂仍被称为且是\BibleQuote{我们中的首生者}，因为，在众人因亚当的过犯而丧亡之后，祂的肉身先于所有其他人而得救和释放，因为那是圣言的身体；从此以后，我们与这身体合为一体，便按照它的模式而得救。因为在祂内，主成为我们通往天国和自己父的向导，祂说：\BibleQuote{我是道路}和\BibleQuote{门}，以及\BibleQuote{凡经过我进来的，必得安全}。因此，祂也被称为\BibleQuote{死者中的首生者}\BibleRef{默 1:5}，这并非说祂先于我们而死，因为我们先死了；而是因为，祂为我们经历了死亡并废除了死亡，作为人，为了我们，首先使祂自己的身体复活。从此，祂既然复活了，我们也从祂并由祂，按次序从死者中复活。\par

62. 但祂若也被称为\Quote{一切受造物的首生者}，这并非说祂与受造物等同，只是在时间上先于它们而已——（因为祂既是\Quote{独生子}，如何可能如此？）——而是因为圣言对受造物的屈尊就卑，祂依此成为\Quote{众多}的\Quote{弟兄}。因为\Quote{独生子}一词用于没有弟兄之处，而\Quote{首生者}则因有弟兄而使用。因此，圣经中从未记载\Quote{天主的首生者}，或\Quote{天主的受造物}；而是用\Quote{独生子}、\Quote{子}、\Quote{圣言}、\Quote{智慧}指称祂，作为专属于父的。因此，\BibleQuote{我们见了祂的光荣，正如父独生子的光荣}\BibleRef{若 1:14}；\BibleQuote{天主派遣了祂的独生子}\BibleRef{若一 4:9}；\BibleQuote{上主，祢的圣言永远常存}；\BibleQuote{在起初已有圣言，圣言与天主同在}；\BibleQuote{基督是天主的德能，是天主的智慧}\BibleRef{格前 1:24}；\BibleQuote{这是我的爱子}；\BibleQuote{你是基督，永生天主之子}。但\Quote{首生者}暗示降临于受造界；因为正是为受造界，祂才被称为首生者；而\Quote{祂受造}则暗示祂对工程的恩宠，因为祂是为它们而受造的。因此，如果祂是独生子，正如事实那样，\Quote{首生者}就需要一些解释；但如果祂确实是首生者，那么祂就不是独生子。因为同一位不可能既是独生子又是首生者，除非从不同关系而言——即，独生子是因祂由父所生，如前所述；而首生者则因祂屈尊就卑，降临于受造界，并使众人成为祂的弟兄。当然，这两个词既彼此不相容，应当说，就圣言而言，独生子的属性确实优先，因为没有别的圣言，也没有别的智慧，唯独祂是父的真子。此外，如前所述，论及祂说这话，并非基于任何理由，而是绝对的：\BibleQuote{在父怀里的独生子}\BibleRef{若 1:18}；而\Quote{首生者}一词则又与受造界有关，作为其理由，\Person{保禄}接着说：\BibleQuote{因为在天上、在地上，一切都是在祂内受造的}\BibleRef{哥 1:16}。但如果一切受造物都是在祂内受造的，祂就与受造物有别，而非受造物，而是万物的创造者。\par

63. 因此，祂之被称为\Quote{首生者}，并非因为祂由父而来，而是因为在祂内受造界得以形成；正如在受造之前，祂已是儿子，万物藉祂而受造，同样，在祂被称为一切受造物的首生者之前，圣言本就与天主同在，祂就是天主。但那些不虔敬的人既不明白这点，便到处说：\Quote{如果祂是一切受造物的首生者，显然祂也是受造物之一。}愚昧的人！如果祂仅是\Quote{一切受造物的首生者}，那么祂便与一切受造物有别；因为并未说\Quote{祂是其他受造物之上的首生者}，免得祂被算作受造物之一，而是记着\Quote{一切受造物的}，为显示祂与受造物有别。例如，\Person{勒乌本}并非被称为\Person{雅各伯}所有子孙的首生者，而是\Person{雅各伯}本人及其弟兄的首生者；免得他被人以为与\Person{雅各伯}的子孙不同。不仅如此，即使论及主自己，宗徒也没有说\Quote{为使祂成为一切的首生者}，免得祂被认为取了与我们不同的身体，却说\Quote{在众多弟兄中作长子}\BibleRef{罗 8:29}，这是因为肉身的相似。如果圣言也是受造物之一，圣经论及祂也会说祂是其他受造物的首生者；但事实上，圣人们说祂是\Quote{一切受造物的首生者}\BibleRef{哥 1:15}，便清楚表明天主子与一切受造物有别，且并非受造物。因为如果祂是受造物，祂将是自己的首生者。那么，亚略派啊，祂怎么可能在自己之前又在自己之后呢？其次，如果祂是受造物，而整个受造界又藉着祂形成，且在祂内得以维系，祂怎能既创造受造界，又是那在祂内得以维系的事物之一呢？既然这种想法本身就不适当，真理便驳倒了他们：祂被称为\Quote{在众多弟兄中作长子}，是因肉身的关联；被称为\Quote{死者的首生者}，因为死者的复活是从祂而来且在祂之后；被称为\Quote{一切受造物的首生者}，则是因为父对人的慈爱，这慈爱使得在祂的圣言内，不仅\BibleQuote{万有都得以维系}，而且受造界本身，正如宗徒所说，\BibleQuote{期待天主子女的显扬}，将有一天\BibleQuote{脱离败坏的控制，得享天主子女的光荣自由}。对于如此得到解救的受造界，主将是它的首生者，也是所有成为子女者的首生者，为使那在后的，因被称为首的祂而得以存立，犹如依靠圣言这元始。\par

64. 我认为，不虔敬的人自己必因这样的想法而感到羞愧；因为若情况并非如我们所说，他们反而断定祂在本质上\Quote{是整个受造物的首生者}——不过是受造物中的一份子，就让他们想想，他们这样想，就是视祂与没有理性和生命之物为兄弟和同伴了。因为那些没有理性和生命之物也是整个受造物的一部分；而\Quote{首生者}若只是时间上的首先，且在种类和相似性上与其他所有受造物相同，他们这样说，怎能不逾越一切不虔敬的限度呢？若他们说出这样的话，谁能容忍他们呢？谁又不会厌恶他们竟构想这般事呢？因为众人都清楚，祂之被称为受造物的\Quote{首生者}，既不是为祂自己（好像祂是受造物），也不是因为祂与整个受造物在本质上有什么关联；而是因为圣言在起初创造万物时，\OriginalTerm{屈尊就卑}{condescended}，降与受造之物，好使它们得以存在。因为受造物无法承受祂的本性，即那未加调和的、圣父的光辉，除非祂因圣父对人的爱而屈尊就卑，托住它们，扶持它们，并将它们引入存在；其次，因为这圣言的屈尊就卑，受造物也藉由祂而成为子嗣，好使祂如所说的，在各方面都是受造物的\Quote{首生者}，既在于创造，也在于为万民而被引入这世界之中。正如经上所载：\BibleQuote{当祂引领首生者进入世界时，祂说，\InnerQuote{天主的所有天使都要朝拜祂}}\BibleRef{希 1:6}。让基督的仇敌听见而自相撕裂吧，因为正是祂来到世界这件事，使祂被称为万有的\Quote{首生者}；因此，圣子是圣父的\Quote{唯一圣子}，因为唯独祂由圣父所出；而祂是\Quote{受造物的首生者}，则是由于万有都被\OriginalTerm{收养}{adoption}为子。正如祂是弟兄中的首生者，并从死者中复活，成为\BibleQuote{睡了之人初熟的果子}\BibleRef{格前 15:20}；所以，既然祂理当\BibleQuote{在万有中居首位}\BibleRef{哥 1:18}，因此祂被造为\Quote{道路的开端}，好使我们行走其上，并通过那位说\Quote{我就是道路}和\Quote{我是门}的祂而入内，分享对圣父的认知，也能听到经上的话：\Quote{行为成全，遵行上主法律的人，真是有福}，以及\Quote{心里洁净的人是有福的，因为他们要看见天主。}\par

65. 因此，既然真理宣称圣言并非本性上是受造物，现在便适合说明，祂在何种意义上被称为\Quote{道路的开端}。因为当那藉\Person{亚当}而来的第一条路失落，我们偏离乐园而走向死亡，并听见经上的话：\BibleQuote{你既是灰土，你还要归于灰土}\BibleRef{创 3:19}，于是那爱人的天主圣言，便按圣父的旨意，给自己穿上了受造的血肉，好使那因原祖的过犯而死亡的人性，能在祂自己肉身的宝血中获得生命，并给我们开了一条\Quote{又新又活的路}，正如宗徒所说：\BibleQuote{即经过帐幔，就是祂的肉身}\BibleRef{希 10:20}；他在别处也表明了这一点：\BibleQuote{所以谁若在基督内，他就是一个新的受造物，旧的已成过去，看，都成了新的}\BibleRef{格后 5:17}。但如果一个新的受造物已经出现，那么这受造物中必有一位是首先的；然而，那仅仅由土造成的、如同我们因过犯而变成的人，却不能作这首先的。因为在最初的受造物中，人已变得不忠，而最初的受造物也因他们而失落了；所以需要另一位来更新最初的受造物，并保存那已经出现的新受造物。因此，出于对人的爱，那新受造物的\Quote{开端}——主，被造为\Quote{道路}，并且一如祂所说的：\Quote{上主造我作祂道路的开端，为祂的工程}；这样，人就不再依照那最初的受造物而行，而是有一个新受造物的开端，并同基督一起作为\Quote{祂道路的开端}，从此我们便跟随那位对我们说\Quote{我就是道路}的祂——正如真福宗徒在\WorkTitle{致哥罗森人书}中教导的：\Quote{祂又是身体——教会的头，祂是元始，是死者中的首生者，为使祂在万有之上独占首位。}\par

66. 因为，如所言，由于从死者中复活，祂被称为开端，而复活发生是在祂带着我们的肉身为我们舍身受死之时，所以显然，祂的话\BibleQuote{祂造了我，作为道路的开端}并非指示祂的\OriginalTerm{性体}{essence}，而是指示祂\OriginalTerm{肉身的临在}{bodily presence}。因为死亡是专属于肉身的；同样，\BibleQuote{主造了我，作为祂道路的开端}这话也是专属于肉身的临在。因为既然救主按肉身如此受造，并成了新受造之物的开端，又拥有我们的\OriginalTerm{初果}{first fruits}，即祂所取的人性肉身，因此，理所当然地，在祂之后，将来的子民也受造了，正如\Person{达味}所说：\BibleQuote{愿这事为后代写下，将来受造的子民要赞美主。} 又说，在第二十一篇 \BibleBook{}{圣咏}中：\BibleQuote{将来的世代要向主传述，他们要向主表明祂的正义，向那将要出生的子民，即主所造的子民传扬。} 因为我们不再听见：\BibleQuote{你吃的那日，你必定死} 见：\BibleRef{创 2:17}，而是\BibleQuote{我在那里，你们也在那里}；这样我们就可以说：\BibleQuote{我们是祂的化工，在基督耶稣内受造，为行善工} 见：\BibleRef{若 14:3}。再者，既然天主的化工，即人，虽受造完美，却因过犯而变得有缺欠，因罪而死，而天主的化工不应该停留在不完美的状态（因此所有圣人都为此祈祷，例如在第一百三十七篇 \BibleBook{}{圣咏}中说：\BibleQuote{主啊，祢必为我偿报；求祢不要鄙视祢双手的化工}），于是，\OriginalTerm{完美的圣言}{perfect Word of God}穿上了不完美的肉身，并被说成是\Quote{受造，为完成工程}；好使祂代替我们偿付罪债，亲自成全人所欠缺的。他所欠缺的，正是不朽和通往天堂的道路。这就是救主所说的话：\BibleQuote{我在地上光荣了祢，完成了祢交给我做的工程} 见：\BibleRef{若 17:4}；又说：\BibleQuote{父交给我要我完成的工程，就是我所行的这些工程，为我作证}；但这里祂说父交给祂完成的\Quote{工程}，就是祂为此而受造的工程，正如在\BibleBook{}{箴言}中所说：\BibleQuote{主造了我，作祂道路的开端，为完成祂的工程}；因为说\Quote{父交给我工程}和\BibleQuote{主造了我为完成工程}完全是一回事。\par

67. 那么，天主的敌人啊，祂何时接受了要完成的工程？因为从这一点也可以理解\Quote{祂受造}的含义。如果你们说：\Quote{在起初，当祂从无中使它们存在时}，这是不真实的；因为工程那时尚未造成；然而祂似乎是说，祂接受的是已经存在的事物。把这话归于\OriginalTerm{圣言降生成人}{Word's becoming flesh}之前的时期，也是不虔敬的，免得祂的来临因此显得多余，因为正是为了这些工程，祂才来临。因此，我们只能说，当祂成为人时，祂接受了工程。因为那时祂通过治愈我们的创伤，并恩赐我们从死者中复活，而完成了这些工程。但是，如果圣言成为血肉之时，工程才交给祂，那么很显然，祂成为人之时，也是祂为工程而受造之时。所以，正如已多次说过的，\Quote{祂受造}所指的不是祂的\OriginalTerm{性体}{essence}，而是祂\OriginalTerm{肉身的生发}{bodily generation}。因为那时，工程因过犯变得不完美且有缺陷，所以祂就肉身而言被说成是受造；好借着成全并完善它们，使祂能将教会呈献给圣父，正如宗徒所说：\BibleQuote{没有瑕疵，没有皱纹，或其他类似的缺陷，而是圣洁无玷的} 见：\BibleRef{弗 5:27}。于是，人类在祂内得以成全和恢复，如同起初受造时一样，且获得了更大的\OriginalTerm{恩宠}{grace}。因为从死者中复活后，我们将不再畏惧死亡，而要在天上永远偕同基督为王。这一切得以成就，是由于那出自圣父的天主自己的圣言，穿上了血肉，成了人。因为，如果祂本是一个受造物，而成了人，人就仅保持原状，而不能与天主结合；因为一个受造物怎能借着另一个受造物与造物主结合呢？或者，同类怎能帮助同类，既然两者都需要帮助呢？再说，如果圣言是受造物，祂怎有能力废除天主的判决，并赦免罪恶呢？因为经上天主的先知记载说，这是天主的作为：\BibleQuote{有谁像祢一样，赦免罪过，宽恕过犯的天主？} 见：\BibleRef{米 7:18}。因为，天主既说过：\BibleQuote{你本是尘土，仍要归于尘土} 见：\BibleRef{创 3:19}，人就成了会死的；那么，出自受造之物怎能解除罪呢？然而，主正是解除罪的那一位，正如祂亲口所说：\BibleQuote{除非圣子使你们自由}；而使人自由的圣子，在真理中显示了祂自己并非受造物，也不是出自受造之物的一员，而是\OriginalTerm{本有的圣言和圣父性体的肖像}{proper Word and Image of the Father's Essence}，祂在起初颁布判决，并且唯独祂赦免罪过。因为，既然在圣言中说：\BibleQuote{你本是尘土，仍要归于尘土}，那么，自由和定罪之取消，便恰当地借着圣言本身并在祂内实现了。\par

68. \Quote{然而，}他们说，\Quote{即使救主是受造物，天主也能只说一句话，就解除诅咒。}另一个人也要同样对他们说：\Quote{即使完全不降到我们中间，天主也能只说一句话，就解除诅咒。}但我们当考虑的是，什么对人类有益，而不只是天主可能做到什么。在诺厄的方舟之前，他本可消灭那时的犯罪者；但他却在方舟之后才做。他也能不经梅瑟，只说一句话，就将百姓领出埃及；但借梅瑟来做，更有裨益。天主也能不经民长而拯救他的子民；但为子民有益的是，在一段时期内为他们兴起民长。\OriginalTerm{救主}{the Saviour}也能从起初就降到我们中间，或者在他降来时不被交付给比拉多；但他却在\BibleQuote{时期一满}\BibleRef{迦 4:4}时来了，并在人们寻找他时说道：\BibleQuote{我就是}\BibleRef{若 18:5}。因为他所做的，就是对人有裨益的，而任何其他方式都不适宜；凡有裨益且适宜的事，他就如此供应。因此他来，不是\Quote{为受人服事，而是为服事人}，为成就我们的\OriginalTerm{救恩}{salvation}。他当然也能从天上颁布\OriginalTerm{法律}{the Law}，但他看出，为人类有益的是他从西乃山发言；他便这样做了，好让梅瑟能上去，也让人们就近听到他的话，更容易相信。此外，他如此行事的良善动机，可以这样看出来：如果天主只凭其能力说一句话，就这样解除了诅咒，那么说话者的能力是彰显了，但人却变成了如同亚当在犯罪以前的样子，从外面承受了\OriginalTerm{恩宠}{grace}，却没有使它与自己身体合而为一；（当亚当被置于乐园中时，就是如此）不但如此，他或许还变得更坏，因为他学会了犯罪。那么，人既处于这种状况下，如果再被蛇诱惑，就又需要天主出命解除诅咒；这样需求就永无止境，人陷入罪债的程度将不减从前，因为成了罪恶的奴隶；既不断犯罪，就要不断需要有人赦免他们，永远得不到自由，因为自己只是\OriginalTerm{肉躯}{flesh}，且由于肉躯的软弱，总被法律所克胜。\par

69. 再说了，如果\OriginalTerm{圣子}{the Son}是受造物，人就会如先前一样仍然是会死的，没有与天主结合；因为一个受造物并不会使受造物与天主结合，因为它自己还需要寻求那能使它结合的；受造界的一部分，也不能成为受造界的救恩，因为它自己也需要救恩。为预防此事，他派遣了自己的儿子，圣子也成了\OriginalTerm{人子}{Son of Man}，取了受造的肉身；这样，既然众人都处在死亡的判决之下，他与他们完全不同，便能亲自为众人将自己的身体献于死亡；而且从今以后，仿佛众人都借着他死了，那判决的话得以应验（因为众人都\BibleQuote{在基督内死了}\BibleRef{格后 5:14}），众人借着他便得以脱离罪恶和随之而来的诅咒，并真正永远存留，从死者中复活，穿上不朽和不能败坏。因为\OriginalTerm{圣言}{the Word}以肉身为衣，正如多次解明过的，蛇的一切啮咬便开始从它那里被彻底止住；凡由肉身的情欲所生出的邪恶，都被割除，同时死亡——罪恶的伴侣——也被废除，正如主亲口所说：\Quote{这世界的首领要来，他在我内一无所能}；\BibleQuote{天主子所以显现出来，是为消灭魔鬼的作为}\BibleRef{若一 3:8}。这些既从肉身中被消灭，我们众人便因肉身的亲缘关系而得自由，从此以后，我们也都与圣言结合。既与天主结合，我们就不再居留于世；而是，如他亲口所说，他在哪里，我们也在哪里；从此我们不再畏惧蛇，因为当\OriginalTerm{救主}{the Saviour}在肉身中攻击他，并向他喊道：\BibleQuote{\OriginalTerm{撒殚}{Satan}，退到我后面去！}\BibleRef{玛 16:23}时，他便被彻底击败，如此从乐园被投入永火。我们也不必提防女人诱惑我们，因为\BibleQuote{在复活的时候，他们也不娶，也不嫁，好像天使一样}\BibleRef{谷 12:25}；在基督耶稣内，将是\Quote{\OriginalTerm{新的受造物}{new creation}}，\Quote{既无男女之分，唯有基督是一切，并在一切内}；既有基督在，还有什么恐惧、什么危险能发生呢？\par

70. 然而，如果圣言是\OriginalTerm{受造物}{creature}，这一切就不会发生；因为魔鬼本身也是受造物，与受造物相争，战争就将永远继续下去，而人处于两者之间，将永远处于死亡的威胁中，没有任何人能在其中并通过他结合于天主、脱离一切恐惧。因此真理向我们表明，圣言不属于\OriginalTerm{受造之物}{originate}，反倒是它们的\OriginalTerm{创造者}{Framer}。正因如此，祂才取了一个\OriginalTerm{受造}{originate}的人的身体，为的是作为创造者更新了它以后，能在自己内使它\OriginalTerm{神化}{deify}，并这样带领我们众人依祂的肖像进入天国。因为人若结合于受造物，或者除非圣子是\Emph{真天主}，就不能神化；人也不能进入圣父面前，除非取了身体的祂是天主本然的、真实的圣言。同样，如同我们如果不是由于圣言所取的本然的人性肉躯，就不能脱离罪和诅咒（因为我们对异类之物毫无共通之处），同样，人如果不是由于成为肉身的圣言从圣父而生、本性真实且为圣父本有，就不能神化。因此，这样的结合正是为了把本性为人者与具有\OriginalTerm{天主性体}{Godhead}的那位结合起来，使他的救恩和\OriginalTerm{神化}{deification}得以稳固。所以，凡否认圣子是由圣父所生、本性具有圣父的本体者，也当否认祂由\OriginalTerm{卒世童贞玛利亚}{Mary Ever-Virgin}取得了真实的人肉；因为在这两种情形下，对我们人都是无益的：如果圣言不是真实而本然地是天主之子，或者祂所取的肉躯不是真实的。然而，祂确实取得了真实的肉躯，尽管\Person{瓦伦廷}胡言乱语；同样，圣言确实是本性\Emph{真天主}，尽管\OriginalTerm{亚略狂徒}{Ario-maniacs}胡言乱语；并且，在这肉躯内开始了我们的\OriginalTerm{新造}{new creation}，祂为了我们而被造为人，正如所说过的，为我们开辟了那条新道路。\par

71. 因此圣言既不是受造物也不是\OriginalTerm{工程}{work}；因为受造物、被造的事物、工程，都是同一个意思；如果祂是受造物和被造的事物，祂也就是工程。所以祂没有说\Quote{祂创造我为工程}，也没有说\Quote{祂造我与众工程}，免得祂显得在本性和本质上是一个受造物；也没有说\Quote{祂创造我去做工程}，免得另一方面，按照不信者的谬见，祂似乎是为了我们而被做的工具。祂也没有说\Quote{祂在众工程之前创造了我}，免得祂虽有如\OriginalTerm{圣子}{Offspring}在万有之前，但如果也在众工程之前受造，就会给\Quote{圣子}和\Quote{祂创造}赋予同一含义。而是极其精确地说\Quote{为了工程}；这等于说：\Quote{圣父使我成为肉身，好使我成为人}，这再次表明祂不是工程而是圣子。因为正如进到一所房子里的人，不是房子的一部分，而是与房子不同，同样，为了工程而被创造的那位，必然在本质上与工程不同。但是，如果不然，就如你们\OriginalTerm{亚略派}{Arians}所主张的，天主的圣言是工程，那么祂自己是凭哪只手和哪个智慧而产生的呢？因为凡产生的事物，都是藉天主的\OriginalTerm{手}{Hand}和\OriginalTerm{智慧}{Wisdom}而产生的，天主自己说：\BibleQuote{这一切都是我手所造的}~\BibleRef{依 66:2}；\Person{达味}也在圣咏中说：\BibleQuote{主，你在起初奠定了大地，诸天是你手的工程}；又在第一百四十三篇圣咏中说：\BibleQuote{我回忆往事，默念你的一切工程，思虑你手所作的一切}。因此，如果众工程是藉天主的手而成就的，而经上记载\BibleQuote{万物是藉着圣言而造成的，凡受造的，没有一样不是藉着祂造成的}~\BibleRef{若 1:3}，又说\BibleQuote{只有一个主耶稣，万物藉祂而有}~\BibleRef{格前 8:9}，并且\BibleQuote{万物在祂内而存在}~\BibleRef{哥 1:17}，那么很清楚，圣子不可能是工程，相反，祂是天主的手和智慧。认识到这一点的\Place{巴比伦}的\OriginalTerm{殉道者}{martyrs}，\Person{阿纳尼雅}、\Person{阿匝黎雅}和\Person{米沙耳}，谴责了亚略派的不虔。因为当他们说\BibleQuote{上主的一切工程，请赞美上主}时，他们列举了天上的事物、地上的事物和整个受造界，作为工程；却未提圣子。因为他们没有说\Quote{圣言，请赞美；智慧，请颂扬}，为表明其他一切事物都在赞美，也都是工程；但圣言不是工程，也不属于赞美者，而是与圣父同受赞美、同受钦崇、同被宣认为天主，因为祂是天主的圣言和智慧，是众工程的\OriginalTerm{创造者}{Framer}。\OriginalTerm{圣神}{Spirit}也在圣咏中以极贴切的区别宣告：\BibleQuote{上主的言语是正直的，祂的一切工程都是忠实可信的}；在另一篇圣咏中又说：\BibleQuote{上主，你的工程何其繁多，皆系你以智慧所造成}。\par

72. 但如果\OriginalTerm{圣言}{Word}是一件工程，那么祂也如同其他万物一样，在\OriginalTerm{智慧}{Wisdom}内受造；圣经也不会将祂与工程区分开，也不会在称它们为工程的同时，宣讲祂是圣言，并宣认祂是天主的智慧。然而事实上，圣经将祂与工程区分开，显示智慧是工程的创作者，而非工程。\Person{保禄}在\BibleBook{致希伯来人}{书}中也注意到这区分，写道：\BibleQuote{天主的话是生活的，是有效力的，比各种双刃的剑还锐利，直穿入灵魂和神魂，关节与骨髓的分离点，且可辨别心中的感觉和思念。没有一个受造物，在天主面前不是明显的，万物在他眼前都是袒露敞开的，我们必须向他交账。}\BibleRef{希 4:12} 看，他将受造之物称为\Quote{\OriginalTerm{受造物}{creature}}；但\OriginalTerm{圣子}{Son}他认作天主的圣言，似乎祂是与受造物不同的另一位。他又说，\BibleQuote{万物在他眼前都是袒露敞开的，我们必须向他交账。} 这表明祂与万物不同。因此，祂施行审判，而一切受造物都须向祂交账。同样，当整个受造界与我们一同叹息，渴望脱离败坏奴役时，圣子便因此显明与受造物不同。因为如果祂是受造物，祂也会是那些叹息者之一，也需要一位如对其他人的一样，将义子身份和救恩带给祂。但若整个受造界一同叹息，为获得脱离败坏奴役的自由，而圣子并非那些叹息者，亦非需要自由者，祂却是将儿子的名分和自由赐予万有的那一位，对当时的犹太人说：\BibleQuote{奴隶不能永远住在家里，儿子却永远居住；那么，如果天主子使你们自由了，你们的确是自由的。}\BibleRef{若 8:35} 从这些思考中亦比光更清楚地表明，天主的圣言并非受造物，而是\OriginalTerm{真圣子}{true Son}，按本性是圣父的亲生子。那么，关于\BibleQuote{上主造我于他道路之始}，我认为这寥寥数语已足以为博学者提供材料，以更完备地驳斥\OriginalTerm{亚略异端}{Arian heresy}。\par

% structure: p0086 source=h2 target=subsection reason=relative-heading-rank; base=h2

\subsection{第二十二章 经文阐释；第六，\BibleBook{}{箴言}8:22的上下文，及其应用于受造的智慧——如在工程中所见。圣子通过工程启示圣父，然后通过降生成人。}

但是，由于异端者读下一节经文时，对那节也持有歪曲的看法，因为经上记载：\BibleQuote{在创世之前，祂就奠定了我}\BibleRef{箴 8:23}，他们认为这是针对圣言的\OriginalTerm{天主性}{Godhead}所说，而非针对祂\OriginalTerm{降生成人的临在}{incarnate Presence}，所以有必要也阐释这节经文，以指出他们的错误。\par

73. 经上记载：\BibleQuote{上主以智慧奠定了大地}\BibleRef{箴 3:19}；如果大地是由智慧所奠定，那么奠定者又如何被奠定呢？其实，这话也是按\BibleBook{}{箴言}的体裁说的，我们也必须以同样方式探讨其意义，好叫我们知道，当圣父藉\OriginalTerm{智慧}{Wisdom}架构并奠定大地，使其稳固坚立时，\OriginalTerm{智慧}{Wisdom}本身也为我们而被奠定，好能成为我们新创造和更新的起始与根基。因此这里如同先前一样，祂并没有说\Quote{在世界之前，祂造了我为圣言或圣子}，免得好像祂的受造有个开端。因为我们必须首先寻求这一点，即祂是否是\OriginalTerm{圣子}{Son}，\Quote{并在这一点上特别查考圣经}；因为正是这个问题，当初宗徒们被质问时，\Person{伯多禄}回答说：\BibleQuote{你是基督，永生天主之子}\BibleRef{玛 16:16}。这也是\OriginalTerm{亚略异端}{Arian heresy}之父作为他的首要问题之一提出的；\BibleQuote{你若是天主子，就……}\BibleRef{玛 4:3}；因为他知道这是真理，是我们信德的至高原则；也知道如果祂本身就是圣子，魔鬼的暴政就必终结；但如果祂是\OriginalTerm{受造物}{creature}，祂也不过是那个被魔鬼欺骗的\Person{亚当}的后裔之一，魔鬼就无需忧虑。出于同样的理由，当时的犹太人也恼怒，因为主说祂是天主子，且天主是祂自己的父。因为如果祂称自己是受造物之一，或说\Quote{我是一件工程}，他们不会因这信息而震惊，也不会认为这些话是亵渎，因为他们知道连天使也曾临在于他们的祖先中；但既然祂称自己为子，他们就察觉到这并非受造物的特征，而是\OriginalTerm{天主性}{Godhead}与\OriginalTerm{圣父的本性}{the Father's nature}的特征。因此，亚略派的人即使效法他们自己的父魔鬼，也当在这一点上特别下功夫；如果祂曾经说过\Quote{祂奠定我为圣言或圣子}，那么就如他们那样思想；但祂若没有这样说过，就不要为自己捏造无中生有的事。\par

74. 祂并没有说：\BibleQuote{在世界以先，祂奠定我做\OriginalTerm{圣言}{Word}或圣子}，而单单说\BibleQuote{祂奠定我}，正如我说过的，为要再次表明，祂在此也是以箴言的方式，并非为祂自己，而是为那些被建造在祂身上的人说话。为此知道，宗徒也写道：\BibleQuote{除已奠立的根基，就是耶稣基督外，任何人不能再奠立别的根基；但各人要注意怎样在上面建造。}\BibleRef{格前 3:11} 而且，根基必须与建在其上的物相匹配，好使它们能紧密结合。祂作为圣言，并没有像自己一样能与之结合的；因为祂是\OriginalTerm{独生子}{Only-begotten}；但成为人之后，祂就有了与祂相似的，就是那些祂取了其血肉相似性的人。因此，按照祂的人性，祂被奠立，好使我们如同珍贵的宝石，得以被建造在祂之上，并能成为那居住在我们内的圣神的圣殿。正如祂是根基，我们是被建造在祂上面的石头；同样，祂又是葡萄树，我们如同枝条连结于祂——并非按照\OriginalTerm{天主性体}{Essence of the Godhead}，因为这确实是不可能的；而是按照祂的人性，因为枝条必须与葡萄树相似，我们既按血肉与祂相似。此外，既然异端者持有这些属人的观念，我们就可恰当地用他们所坚持之事中包含的属人类比来驳斥他们。因此，祂并没有说\Quote{祂造我为根基}，免得祂好像是被造的、具有存在的开端，他们由此找到亵渎的卑劣借口；而是说\BibleQuote{祂奠定我}。凡被奠定的，是为那在上建造的石头之缘故；这不是一个随意的过程，而是一块石头首先从山上被运下来，放置在地的深处。当石头在山里时，它尚未被奠定；但当需要时，它被运送、安置在地的深处，那时候，如果石头能说话，它立刻会说：\Quote{那从山上把我带到这里来的，如今奠定我了。}因此，主在被奠定时，并没有获得存在的开端；因为祂在此之前就是圣言；但当祂穿上我们的身体，就是祂从玛利亚所分离并取来的身体时，祂便说\BibleQuote{祂奠定我}；这如同说：\Quote{我是圣言，祂用尘土的身体将我包裹。}因为如此，祂为我们的缘故被奠立，取了我们所有的，好使我们藉着血肉的相似性，被结合、连结、团聚在祂内，得以成为完全的人，并常存不朽不坏。\par

75. 不要让\BibleQuote{在世界以先}、\BibleQuote{祂造大地以前}、\BibleQuote{山岳未曾奠定}这些话困扰任何人；因为它们与\BibleQuote{奠定}和\BibleQuote{造成}十分相符；因为这里再次暗示了那属血肉的\OriginalTerm{救世计划}{Economy}。因为，虽然那由救主临到我们的恩宠，正如宗徒所说，是在现今才显现，并在祂居住于我们中间时来到；然而这恩宠远在我们出现之前，甚至在世界奠基之前就已预备好了，这原因是仁慈而奇妙的。天主不宜在事后才为我们的命运定主意，免得祂显得好像不知道我们的终向。万有的天主——藉着祂自己的圣言创造了我们，比我们更清楚我们的终向，并且预见到，我们被造原是\BibleQuote{好的} \BibleRef{创 1:31}，却会违背诫命，因悖逆而被赶出乐园——祂既仁慈又良善，便预先在祂自己的圣言内，就是那造了我们的圣言内，预备了我们救恩的救世计划；使虽然我们因蛇的欺骗而背离了祂，却不至于永久死亡，而是在圣言内拥有那早已为我们预备的救赎和救恩，得以复活并常存不朽；那时，祂要为我们被造为\BibleQuote{道路的起点}，那\BibleQuote{一切受造物的首生者}要成为\BibleQuote{弟兄中的首生者}，又要复活成为\BibleQuote{死者中的初果}。真福宗徒\Person{保禄}在他的书信中教导了这道理；因为，在解释\WorkTitle{箴言}中\BibleQuote{在世界以先}和\BibleQuote{大地未有以前}这些话时，他对\Person{弟茂德}这样说：\BibleQuote{你要为福音同受苦难，赖天主的能力，祂拯救了我们，以圣召召叫了我们，不是按我们的行为，而是按祂自己的决意和恩宠：这恩宠是万世以前在基督耶稣内赐予我们的，如今却藉着我们救主耶稣基督的显现彰显出来；祂毁灭了死亡，并彰显了生命与不朽。}\BibleRef{弟后 1:8} 又对厄弗所人说：\BibleQuote{愿天主，我们的主耶稣基督的父受赞美！祂在天上，在基督内，以各种属神的祝福，祝福了我们，因为祂于创世以前，在基督内已拣选了我们，为使我们在他面前，成为圣洁无瑕疵的；又由于爱，按照自己旨意的决定，预定了我们藉着耶稣基督获得义子的名分，而归于祂。}\BibleRef{弗 1:3}\par

76. 然而，祂如何在我们尚未存在时就拣选了我们？岂不是正如祂自己所说，在祂内我们预先蒙受\OriginalTerm{预像}{represented beforehand}吗？并且，在人类受造以前，祂如何就预定我们得\OriginalTerm{义子}{adoption}的名分，岂不是因为圣子自己\BibleQuote{在创世之前已被奠定}，承当了那为我们而设的\OriginalTerm{救世工程}{economy}吗？再者，正如宗徒接下去说的，我们如何能够\BibleQuote{预定承受基业}，岂不是因为主自己在\BibleQuote{创世之前}已被奠定，祂怀着为我们而来的目的，通过血肉承当那本来攻击我们的全部审判的罪债，使我们从此在祂内成为儿子呢？我们又如何能够在\BibleQuote{创世之前}领受它呢？那时我们尚不存在，而是在时间中后来才存在的；岂不是因为在基督内早已储存了那临到我们的恩宠吗？因此，在审判之时，当人人都要按照自己的行为受报时，祂要说：\BibleQuote{我父所祝福的，你们来吧！承受自创世以来为你们预备的国度吧！}\BibleRef{玛 25:34}。那么，这国度在我们尚未存在之前在谁内、如何预备好了呢？岂不是在那为此目的\BibleQuote{在创世以前}已被奠定的主内吗？好使我们作为建基于祂的、紧密结合的石头，得以分享从祂而来的生命与恩宠。这事的成就，正如虔诚的心灵自然领悟到的，即如我所说的，我们这由土所造的人，在经过我们短暂死亡后的复活中，能够承受那永生；这本是我们所不能承受的；而这乃是因为在\BibleQuote{创世之前}，就已在基督内为我们预备了生命与救恩的盼望。因此，圣言来到我们的血肉中，并在其中受造，作为\BibleQuote{祂一切工道的起始}，按照圣父在创世之前、如所言，在未有大地、山岳未定、水泉未涌之前，在祂内所定的旨意，被立为根基；这样，虽然大地、山岳和可见自然的万象将在这时代的终结时逝去，我们却不会随它们一样衰败，反而在它们逝去之后仍能存活，因为我们拥有按照\OriginalTerm{拣选}{election}，在这一切之先，就已在圣言自身内为我们预备的属灵生命与祝福。因为这样，我们才能承受非暂时的生命，而是从此以后永远居于基督内生活；因为甚至在此之前，我们的生命就已在基督耶稣内被奠定和预备好了。\par

77. 我们的生命以任何其他方式被奠定都不适宜，而只应奠基在那万世之前已存在、且万世由祂而造的主内；这样，由于在祂内，我们也就能够承受那永恒的生命。因为天主是美善的；祂既常是美善的，就定意如此，因为祂知道我们脆弱的性体需要从祂而来的救助与救恩。又如一位明智的建筑师，计划建造一座房屋，同时也会筹划一旦房屋建成后破损时的修葺事宜；他为此咨询，预备并给予工匠修复的材料；这样，修葺的材料便先于房屋备好；同样，在我们之先，我们救恩的修复工程便在基督内奠定了，好使我们在祂内甚至被新创造。这旨意和定策在\BibleQuote{创世之前}就已预备好，但当需要时便生效，而救主来到了我们中间。因为主自己将在祂接我们进入永生时，为我们代作在天上的一切。这便足以证明天主圣言并非受造物，而经文的正确意义当是如上述。然而，既然这经文从各方面仔细考察都有正确的意义，也许最好说明它是什么；或许详尽的解释能让那些无知的人羞愧。现在，我必须重提前面所说过的话，因为我要说的与同一篇箴言和同一位智慧相关。圣言并未自称本性为受造物，而是在\BibleBook{}{箴言}中说：\BibleQuote{主创造了我}；祂清楚指示的，不是一种\OriginalTerm{显白}{plainly}的意义，而是隐晦的意义，这意义当我们揭开箴言的帷幕时方能寻见。因为，有谁听见那\OriginalTerm{创造者智慧}{Framing Wisdom}说：\BibleQuote{主创造了我作祂一切工道的起始}，会不立刻追问其意义，沉思那创造性的智慧如何能受造呢？有谁听见天主的独生子说，祂受造为\BibleQuote{工道的起始}，会不探究其意义，惊奇那独生子如何能成为众多其他者的起始呢？因为这原是一句\OriginalTerm{隐语}{dark saying}；然而正如他所说的，\BibleQuote{明哲人必懂得箴言与解释，智慧人的言辞与他们的隐语。}\BibleRef{箴 1:5}\par

78. 然而，天主的独生圣子和真智慧乃是万物的造物主与塑造者；因为他说：\BibleQuote{你以智慧创造了万物}，又说：\BibleQuote{大地充满了你的造化}。然而，为使受造之物不仅存在，而且成为美好的，天主乐意使自己的智慧\OriginalTerm{屈尊就卑}{condescend}于受造物，以便将祂自己肖像的\OriginalTerm{印记}{impress}和模样普遍地印在万有和每一物上，使所造之物成为明显充满智慧的工程，堪当天主。因为正如圣子被视为圣言时，我们的话语是肖像；同样，同一位圣子被视为智慧时，那植根于我们内的智慧也是肖像；在这智慧中，我们因拥有认知与思想的能力，成为\OriginalTerm{创造万物的智慧}{All-framing Wisdom}之接受者；并借着这智慧，我们得以认识祂的圣父。祂说：\BibleQuote{凡有子的，也有父}；又说：\BibleQuote{接纳我的，便是接纳派遣我来的那一位}。既然智慧的印记如此被造在我们内，并存在于一切工程中，那\OriginalTerm{真实而创造的智慧}{true and framing Wisdom}便理所当然地将属于自己印记的话语归于自身，说：\BibleQuote{上主造了我，以作祂的工程}；因为我们内智慧所说的话，主自己就仿佛当作自己的话来说；并且，虽然祂本身并非受造，而是造物主，但因为在受造工程中形成了祂的肖像，祂便仿佛就自身而言说了这话。正如主自己曾说：\BibleQuote{谁接纳你们，就是接纳我}\BibleRef{玛 10:40}，因为祂的印记在我们内；同样，虽然祂不在受造物之列，但因为祂的肖像和印记被造在工程中，祂便仿佛以自身的位格说：\BibleQuote{上主造了我作祂道路的开端，为祂的工程}。因此，智慧在工程中的印记被造成，正如我先前所说，为叫世界在其中认出自己的造物主圣言，并借着祂认出圣父。这正是\Person{保禄}所说的：\BibleQuote{因为那能由天主所认识的，已为他们显明了，因为天主已给他们显示出来了：因为祂那看不见的美善，即祂永远的大能和祂为神的本性，从受造世界以来，都可凭着受造之物，辨认洞察出来}\BibleRef{罗 1:19}。但若如此，圣言在性体上并非受造物；而在\BibleBook{}{箴言}这一段中所说的，乃是我们内所谓的智慧。\par

79. 但如果这也未能说服他们，就让他们自己告诉我们，受造物中究竟有没有智慧？如果没有，宗徒为何抱怨说：\BibleQuote{既然世人凭自己的智慧，在依天主的智慧之下，未能认识天主}\BibleRef{格前 1:21}？或者，如果没有智慧，圣经中怎能找到\Quote{众多智者}呢？因为\BibleQuote{智慧人敬畏并远离邪恶}\BibleRef{箴 14:16}；\BibleQuote{因智慧，房屋得以建成}；训道者说：\BibleQuote{人的智慧使他容光焕发}；他也如此责备刚愎之人：\BibleQuote{不要说：先前的日子胜过如今的日子，是什么缘故？因为你这样问，不是出于智慧。}但是，正如\Person{息辣之子}所说：\BibleQuote{祂将智慧倾注在祂的一切化工上；照祂的恩赐，智慧临在于一切血肉之中；祂将智慧赐予爱慕祂的人}\BibleRef{德 1:9}，而这倾注并非真智慧及独生圣子的性体之特征，而是那在世界中映射之智慧的特征，那么，\OriginalTerm{创造万物的真智慧}{All-framing and true Wisdom}本身——其印记正是倾注于世间的智慧与知识——如我已经解释过的那样，仿佛就自身而言说：\BibleQuote{上主造了我以作祂的工程}，这有什么不可信的呢？因为世界中的智慧并非创造的智慧，而是受造于工程中的智慧，依照这智慧，\BibleQuote{高天陈述天主的光荣，穹苍宣扬他手造的化工}。人若内里有这智慧，就会承认天主的真智慧，并且知道自己真是按照天主的肖像受造的。比如，某位君王的儿子，当父王想要建造一座城时，可能会在每项兴建的工程上印上自己的名字，一方面借着他的名字出现在每件事物上，确保工程稳固长存，另一方面也使人从这名字记起他和他的父王；在城建成之后，有人问起如何建成，他便回答说：\Quote{这城建造得稳固，因为依照我父的旨意，我在每项工程上都留下了肖像，我的名字已被刻在工程上}；但他这样说，并非表示自己的性体是受造的，而是指借着自己名字所留下的印记；同样地，若将这比喻应用起来，对于那些惊叹受造物中智慧的人，真智慧回答说：\BibleQuote{上主造了我以作工程}，因为我的印记在其中；我如此屈尊就卑，是为了塑造万有。\par

80. 此外，\Person{圣子}论及我们内之印记，仿佛这印记就是祂自己，此事不应令任何人惊讶，因为（我们不应回避重复）当\Person{扫禄}迫害\Person{教会}时——\Person{教会}有祂的印记和肖像——祂说，如同自己遭受迫害一般：\BibleQuote{主，你为什么迫害我？}\BibleRef{宗 9:4} 因此（如前所述），倘若假设在工程中的智慧之印记说：\BibleQuote{主为他的工程创造了我，}没有人会惊讶；同样，如果那真实的、形成万物的智慧，即天主的独生子\Person{圣言}，使用属于祂肖像的话，如同论及自己，即\BibleQuote{主为他的工程创造了我，}任何人都不要忽视那受造于世界和工程中的智慧，而认为\BibleQuote{祂创造}是指智慧本身的本体而说的，免得他以水掺酒，被判为真理的诈骗者。因为祂是创造者与形成者；而祂的印记是受造于工程中，如同肖像的复制品。祂又说\BibleQuote{道路的开端}，因为这样的智慧成为某种开端，仿佛认识天主的初步；因为一个人就这样先踏上了这条路，且存留在敬畏天主内（正如\Person{撒罗满}所说：\BibleQuote{敬畏主是智慧的开端}），然后在思想中向上探索，察觉那受造界中的形成之智慧，也将在其中察觉它的\Person{圣父}，正如\Person{主}亲口说的：\BibleQuote{看见我的，就是看见了\Person{圣父}，}以及如\Person{若望}所写的：\BibleQuote{凡承认子的，也有\Person{圣父}。}祂又说：\BibleQuote{在世界之前祂就立定了我，}因为在祂的印记中，工程得以稳固并永恒存立。这样，免得任何人听到关于在工程中如此受造的智慧，便以为真智慧，即天主子，本性就是受造物，祂发现有必要加上：\BibleQuote{山岳未成，大地未有，洋海未造，一切山陵未生以先，祂就生了我，}这样，当祂说\BibleQuote{在一切受造物之先}（因为祂用这些名目包括了整个受造界），祂可以表明，按本质祂并非与工程一同受造。因为如果祂是\BibleQuote{为这些工程而受造，}却是在它们之前，那么结论是，祂在受造之前就已经存在了。因此，祂按本性及本质并非受造物，而是如祂自己所说的，是\OriginalTerm{所生者}{Offspring}。然而受造物与所生者有何区别，以及在本性上如何不同，已在前文说明。\par

81. 但是，由于祂继续说：\BibleQuote{当祂建立诸天时，我已在祂身旁，}我们应当知道，祂这样说，并非仿佛\Person{圣父}是在没有智慧的情况下建立了天和云彩（因为毫无疑问，万物都是在智慧内受造的，离了它，没有一物受造\BibleRef{若 1:3}）；祂的意思乃是：\Quote{万物都在我内并借着我而成就，当需要在工程中创造智慧时，我确实在本质与\Person{圣父}同在，但借着俯就受造物，我安排我的印记于工程之上，使整个世界就像一个身体，不致分歧，而能协调一致。}因此，所有那些以正直的理解，按所赐给他们的智慧，来静观受造物的人，都能为自己说：\BibleQuote{万物因你的安排而常存；}但那些轻视这事的人，必被告知：\BibleQuote{他们自负为智者，反而成为愚妄；}因为\BibleQuote{关于天主，人所本来能知的一切，都在他们内原是明显的，因为天主已给他们显示了；因为那看不见的天主性——他永恒的大能和天主性——自创世以来，凭着可理解的受造物，都可辨认出来，以致人无可推诿。因为他们虽然认识了天主，却不以祂为天主而光荣祂，也不感谢祂，反而去敬拜事奉受造物，而非敬拜事奉那万物的造主——祂是可赞美的，直到永远。阿们。}他们听到下列话时，必然羞愧：\BibleQuote{因为，在智慧的安排中（如上文所解释的），世人凭自己的智慧不认识天主，天主遂决意以愚妄的道理来拯救那些相信的人}\BibleRef{格前 1:21}。因为，不再像以往那样，天主愿意借着智慧的肖像和影子，即受造物中的智慧，叫人认识祂，而是使真智慧本身降生成人，并受十字架的死刑；为叫从此以后，凡信祂的人都能获得救恩。然而，仍是同一天主智慧，借着祂在受造物中的自己的肖像（也因此祂被称为受造的），首先显出祂自己，并借着自己显出祂的\Person{圣父}；然后，祂本身就是\Person{圣言}，正如\Person{若望}所说，\BibleQuote{成了血肉}\BibleRef{若 1:14}，并在毁灭死亡、拯救人类之后，更彰显了祂自己和藉着祂的\Person{圣父}，说：\BibleQuote{求祢赐给他们认识祢，唯一的真天主，和祢所派遣来的耶稣基督。}\par

82. 因此，整个大地充满了对他的认识；因为通过子认识父和从父认识子是一样的，父喜悦子，子也以同样的喜乐在父内欢跃，说：\BibleQuote{我天天在他面前喜悦，不断欢跃。}\BibleRef{箴 8:30}。这再次证明子并非外来的，而是属于父的本质。因为看，并非如那些不虔敬的人所说，他是为了我们而存在，也不是从虚无中来的（因为天主并非从自身之外为自己找了一个喜悦的缘由），这些话所指的是他本有的、相似的。那么，何时父曾不喜悦呢？但如果他曾在某时喜悦，那他所喜悦的也必恒在。而父喜悦谁呢？岂不是看见自己在自己的肖像中，即他的圣言中吗？虽然他在造成世界时也喜悦世人，正如这同一部\WorkTitle{箴言}所记载的\BibleRef{箴 8:31}，然而这也有一致的意义。因为即使如此，他的喜悦并非因为喜乐加于他，而是再次看见按自己肖像所造的工作；因此天主的这种喜悦也是因他的肖像。而子又怎能喜悦，除非看见自己在父内？因为这就等于说：\BibleQuote{看见了我，就是看见了父，我在父内，父在我内。}\BibleRef{若 14:9}。因此，基督的仇敌啊，你们的自夸在各方面都已显明是空虚的，你们到处传扬并散布你们的经文：\Quote{上主创造了我作为他道路的开端}，歪曲其意义，并传播的并非撒罗满的原意，而是你们自己的注解。因为看，你们的理解已被证明只是幻想；但\WorkTitle{箴言}中的那段话，以及以上所说的一切，都证明子按本性和本质并非受造物，而是父的亲生子，真正的智慧和圣言，\BibleQuote{万物是借着他而造成的，凡受造的，没有一样不是借他而造成的。}\BibleRef{若 1:3}。\par

\SourceRecord{Translated by John Henry Newman and Archibald Robertson.；Nicene and Post-Nicene Fathers, Second Series；Vol. 4.；Edited by Philip Schaff and Henry Wace.；Buffalo, NY: Christian Literature Publishing Co.,；1892.；<http://www.newadvent.org/fathers/28162.htm>.}

% structure: p0001 source=h1 target=section reason=page-title

\section{反驳亚略异端第三讲}

% structure: p0002 source=h2 target=subsection reason=relative-heading-rank; base=h2

\subsection{第二十三章 解释经文；第七点，\BibleBook{若望}{福音}十四章十节导论。\OriginalTerm{互寓}{coinherence}的教义。圣父与圣子各自都是完整而完美的天主。他们彼此内居于对方，因为他们的\OriginalTerm{性体}{Essence}是同一的。他们各自是完美的，并拥有同一个性体，因为第二位是第一位之子。\Person{Asterius}对所论经文的逃避解释；被驳斥。因为圣子拥有圣父所有的一切，他是圣父的\OriginalTerm{肖像}{Image}；圣父是唯一的天主，因为圣子住在圣父内。}

1. 那些\OriginalTerm{亚略异端的狂热分子}{Ario-maniacs}，显然，一旦决心违犯并背叛\OriginalTerm{真理}{Truth}，便竭力引用圣经的话：\BibleQuote{恶人陷入灾祸的深渊时，便肆无忌惮；}因为反驳不能阻止他们，困惑也不能使他们羞愧；相反，如同拥有\BibleQuote{淫妇的额头}，他们\BibleQuote{不肯羞惭}\BibleRef{耶 3:3}，在众人面前毫不掩饰自己的不敬虔。因为尽管他们所引用的那些经文——\BibleQuote{上主造了我}、\BibleQuote{被立高于众天使}、\BibleQuote{首生者}，以及\BibleQuote{对造他的那位忠信}——都有正确的意义，并且教导人应向\OriginalTerm{基督}{Christ}怀有虔敬，然而这些人，仿佛被蛇的毒液浸润，不看自己当看的，也不明白自己所读的，犹如从他们不虔敬之心深处呕吐出来一般，竟转而去诋毁我们主的言论：\BibleQuote{我在父内，父在我内}\BibleRef{若 14:10}；他们说：\Quote{一者怎能被包含在另一者内，另一者又怎能被包含在这一者内？}或者说：\Quote{父是较大的，怎能被包含在较小的子内？}或者说：\Quote{子若在父内，有何奇妙呢？}，考虑到经上也论及我们：\BibleQuote{我们生活、行动、存在，都在他内}？而这种心态正与他们邪恶的本性相符，他们以为天主是物质的，不明白何为\OriginalTerm{真父}{True Father}、\OriginalTerm{真子}{True Son}，也不明白那\OriginalTerm{不可见的光明}{Light Invisible}与\OriginalTerm{永恒的光明}{Eternal}，以及它的\OriginalTerm{不可见的光辉}{Radiance Invisible}，更不明白\OriginalTerm{不可见的自立存在}{Invisible Subsistence}，以及\OriginalTerm{非物质的表达}{Immaterial Expression}和\OriginalTerm{非物质的肖像}{Immaterial Image}。因为他们若知道，就不会羞辱并嘲笑荣耀的主，也不以物质的方式解释非物质的事物，从而扭曲美善的言语。的确，只要听到是主的话语，就立刻相信，这便已足够，因为单纯的信仰胜于精密的推理说服；然而，既然他们甚至企图将此段经文也亵渎为自己的异端所用，那么就有必要揭露他们的邪恶，并展示真理的心意，至少为保障信友之安全。因为当说到\Quote{我在父内，父在我内}时，并不像这些人所想象的那样，他们彼此倾注于对方，如同空容器一般互相充满，以至于圣子充满圣父的虚空，圣父又充满圣子的虚空，以致他们各自本身并不完整和完美（因为这属于物体，所以仅仅如此主张便已充满不虔），因为圣父是圆满而完美的，圣子则是\OriginalTerm{天主性的圆满}{Fulness of Godhead}。再者，也并非如同天主进入圣徒中坚固他们，他也同样在圣子内。因为他本身就是圣父的\OriginalTerm{能力}{Power}与\OriginalTerm{智慧}{Wisdom}，受造之物借着分受他，而在\OriginalTerm{圣神}{Holy Spirit}内被圣化；但圣子本身并非由于分受而成为圣子，乃是圣父的亲生者。也不是在经文\BibleQuote{我们在他内生活、行动、存在}的意义上，说圣子在圣父内；因为，他作为来自\OriginalTerm{圣父的泉源}{Fount of the Father}的，就是生命，万物在他内获得活力和存立；因为生命不在生命中生活，否则他就不是生命了；他反而是将生命赐给万物。\par

但我们且来看看那被异端收买的辩士，\OriginalTerm{智者}{Sophist}\Person{Asterius}怎么说。他效法犹太人到此地步，这样写道：\Quote{他说他在父内，父又在他内，这话显然是因为：他所论述的话语，并非如他所说的，是他自己的，而是父的；那些工作也不属于他，而是属于那赐给他权能的父。}这样话，若是由一个小儿随口说出，还可因其年幼而恕之；但当一位自称\OriginalTerm{智者}{Sophist}、宣称通晓一切知识的人这样写时，该当受到何等严厉的谴责啊！他岂不是显明自己不认识宗徒的话\BibleRef{格前 2:4}，因为自恃智慧诱人的言词而自高自大，并想凭此成功进行欺骗，\BibleQuote{不明白自己所讲论、所断定的}\BibleRef{弟前 1:7}吗？因为，圣子所说的、作为唯独属于圣子——他是\OriginalTerm{圣言}{Word}、\OriginalTerm{智慧}{Wisdom}和圣父\OriginalTerm{性体}{Essence}的\OriginalTerm{肖像}{Image}——的适当话语，他却将其拉平到一切受造物，使之成为圣子和受造物共有的；而且，这无法无天的人竟说，圣父的能力领受了能力，由此他的不虔便可进一步说，圣子被立为子是在一个儿子内，圣言领受了话语的权威；他远非承认基督是以圣子的身分说这话，反而将他与一切受造之物并列，仿佛他也是受教学习，如同它们一般。因为，如果圣子说\BibleQuote{我在父内，父在我内}\BibleRef{若 14:10}，是因为他的言论不是他自己的话，而是圣父的，他的工作也是如此——那么，既然\Person{达味}说：\Quote{我要听上主天主在我内所说的话}，\Person{撒罗满}也说：\Quote{我的话是出自天主}，并且\Person{梅瑟}是传述来自天主之话语的仆人，每位先知所说的也不是自己的，而是来自天主的：\Quote{上主这样说}，而圣徒们的工作，正如他们自己表白的，并不是他们自己的，而是赐予能力的天主的，例如\Person{厄里亚}和\Person{厄里叟}呼求天主亲自复活死人，\Person{厄里叟}在治愈\Person{纳阿曼}的癞病时对他说：\BibleQuote{好叫你明白在以色列有天主}\BibleRef{列下 5:8}，\Person{撒慕尔}在收割时节祈求天主赐雨，宗徒们也声明他们行奇迹不是凭自己的能力，而是凭主的恩宠——那么显然，根据\Person{Asterius}，这样的话必然为一切所共有，以至于他们每人都能说：\Quote{我在父内，父在我内}；其结果便是：他不再是唯一的\OriginalTerm{圣子}{Son of God}、\OriginalTerm{圣言}{Word}与\OriginalTerm{智慧}{Wisdom}，而是如同其他许多人一样，不过其中之一而已。\par

3. 但是，如果主说了这话，祂的话就不应该是 \BibleQuote{我在父内，父在我内}，而应该是 \Quote{我也在父内，父也在我内}，这样祂就可以与圣父相对，没有任何属于祂自己的、作为圣子的特权，而是与众人有同一的恩宠。但事情并非如他们所想的一样；因为他们不明白祂是从圣父所生的真实圣子，就诬蔑了这位具有这一身份、唯独合宜说 \BibleQuote{我在父内，父在我内} 的圣子。因为圣子在圣父内——这是我们可以知道的——是由于圣子的整个 \OriginalTerm{存有}{Being} 都是圣父 \OriginalTerm{性体}{Essence} 本有的，如同来自光的光辉，来自泉源的溪流；因此，谁看见了圣子，就看见了圣父本有的一切，并且知道圣子的存有，因为来自圣父，所以就在圣父内。因为圣父在圣子内，既然圣子是从圣父而来、并且是圣父本有的，如同在光辉中见到太阳，在言语中见到思想，在溪流中见到泉源：因为谁这样瞻仰圣子，就瞻仰到圣父性体本有的一切，并且知道圣父在圣子内。因为圣父的 \OriginalTerm{形像}{Form} 和 \OriginalTerm{天主性}{Godhead} 就是圣子的存有，所以圣子在圣父内，圣父在圣子内。\par

4. 因此，这是合理且恰当的：祂先前说了 \BibleQuote{我与父是一体的}，接着又说 \BibleQuote{我在父内，父在我内}\BibleRef{若 10:30}，为要显示天主性的同一与 \OriginalTerm{性体}{Essence} 的唯一。因为他们是一，不是如同一物分为两部分，这些部分无非是一，也不是如同同一事物有两个名称，以至于同一位有时是圣父，有时又是祂自己的圣子，因为 \Person{撒伯里乌} 持此说，被判定为异端者。然而他们是两位，因为圣父是父，而不是子；圣子是子，而不是父；但 \OriginalTerm{本性}{nature} 却是一；（因为子嗣并非不像其父，因为祂是父的肖像），并且凡圣父所有的一切，都是圣子的。因此，圣子并不是另一位天主，因为祂并非从外获得，否则若一天主性是外在于圣父而获得的，那就会有许多位天主；因为若圣子作为子嗣是另一位，祂仍然是与天主同一位；并且祂与圣父在 \OriginalTerm{本性}{nature} 的专有与特性上，以及同一的天主性上，是一体的，正如前面所说。因为光辉也是光，并不次于太阳，也不是别的光，也不是由于分沾光，而是太阳完整而本有的子息。而如此生出的必然是同一的光；没有人会说它们是两个光，而是太阳与光辉是两位，而那从太阳发出的在光辉中照亮万物的光却是一个。同样，圣子的 \OriginalTerm{天主性}{Godhead} 也是圣父的；因此也是不可分的；如此便只有一位天主，除他以外再无别位。所以，既然他们是一，且天主性本身也是一，那么论及圣子的事，同样论及圣父，除了说祂是圣父这一点之外：——例如，祂是天主：\BibleQuote{圣言就是天主}\BibleRef{若 1:1}；全能者：\BibleQuote{那今在、昔在及将来永在的全能者这样说}\BibleRef{默 1:8}；主：\BibleQuote{唯一的主耶稣基督}\BibleRef{格前 8:6}；祂是光：\BibleQuote{我是光}\BibleRef{若 8:12}；祂赦免罪过：祂说，\BibleQuote{为叫你们知道人子在地上有权赦罪}\BibleRef{路 5:24}；以及其他属性也是如此。因为圣子亲自说：\BibleQuote{凡父所有的一切，都是我的}；又说：\BibleQuote{我的，都是你的}。\par

5. 当我们听到圣父的属性被论及于圣子时，我们便因此在圣子内看见圣父；当我们发现论及圣子的事也被论及于圣父时，我们也就在圣父内瞻仰圣子。圣父的属性为何归于圣子，岂不是因为圣子是出自祂的子嗣吗？而圣子的属性为何属乎圣父，岂不是也因为圣子是祂性体本有的子嗣吗？而圣子既是圣父性体本有的子嗣，祂便合理地说，圣父的属性也是祂自己的；因此，祂恰当地与所说 \BibleQuote{我与父是一体的} 一致，接着说道 \BibleQuote{好叫你们知道我在父内，父在我内}。此外，祂还又加上了 \BibleQuote{谁看见了我，就是看见了父}；这三段经文具有同一意义。因为谁若这样理解圣子与圣父是一体的，就明白祂在父内，父在子内；因为圣子的天主性是圣父的，且在圣子内；谁领会了这一点，就确信 \BibleQuote{谁看见了子，就是看见了父}，因为在圣子内所瞻仰到的，乃是圣父的天主性。我们立刻可从皇帝的肖像这一比喻中明白这一点。因为在肖像中有着皇帝的形状和形像，而皇帝内也有着那在肖像中的形状。因为肖像中皇帝与皇帝的肖似丝毫不差；因此，注视肖像的人，就在肖像中看见了皇帝；而看见皇帝的人，也就认出他正是在肖像中的那一位。由于这肖似并无不同，肖像或许会对那愿意通过肖像看见皇帝的人说：\Quote{我与皇帝是一体的；因为我在他内，他在我内；你在我内所见的，你在他内也看见；你在他内所见过的，你在我内也拥有。}因此，朝拜肖像的人，就在肖像中朝拜皇帝；因为肖像就是他的形状和显现。既然如此，圣子也是圣父的 \OriginalTerm{肖像}{Image}，必须理解，圣父的天主性和专有属性就是圣子的 \OriginalTerm{存有}{Being}。\par

6. 这就是所说的：\BibleQuote{祂虽具有天主的\OriginalTerm{形象}{Form}} \BibleRef{斐 2:6}，以及\BibleQuote{父在我内}。也不是这\OriginalTerm{天主性}{Godhead}形象的\OriginalTerm{天主性}{Godhead}仅仅是部分的，而是圣父天主性的圆满是圣子的\OriginalTerm{存有}{Being}，而圣子是完全的天主。因此，祂既然与天主同等，也\BibleQuote{不以自己与天主同等为僭越}；再者，既然圣子的天主性和形象无非是父的，这就是祂所说的：\BibleQuote{我在父内}。因此，\BibleQuote{天主在基督内使世界与自己和好；} \BibleRef{格后 5:19}因为父的\OriginalTerm{本体}{Essence}的\OriginalTerm{专有属性}{propriety}就是那个圣子，在祂内受造界当时就与天主和好了。因此，圣子当时所做的一切，就是父的工作，因为圣子就是那行事的天主性父的形象。所以，谁看见了圣子，就看见了父；因为在父的天主性内，圣子存在并被默观；而父在祂内的形象，在祂内彰显父；这样，父就在圣子内。而且那从父在圣子内的专有属性和天主性，在父内显示圣子，以及祂与父的不可分离性；谁若听到并看到，凡论及父的，也同样论及圣子，并非作为借着\OriginalTerm{恩宠}{grace}或分享而附加于祂的本体，而是因为圣子的存有本身就是父的本体特有后裔，就会恰当地理解这些话，正如我先前说过的：\BibleQuote{我在父内，父在我内}和\BibleQuote{我与父原是一体}。因为圣子如何，父也如何，因为祂拥有父的一切。因此，祂也与父一同被蕴含。因为，若没有子，就不能说有父；然而，当我们称天主是创造者时，我们不必同时提及那被造的事物；因为创造者先于他的工程。但当我们称天主为父时，我们随即与父一同意指圣子的存在。因此，谁信圣子，也就信父：因为他信的是父的本体所专有的；如此，\OriginalTerm{信德}{faith}在唯一的天主内是同一的。而谁朝拜并尊敬圣子，在圣子内也朝拜并尊敬父；因为天主性是一个；所以向父在并藉由圣子所献上的尊敬和朝拜，也是同一的。而谁这样朝拜，就是朝拜唯一的天主；因为天主是唯一的，除祂之外，没有别的。因此，当父被称为唯一的天主，我们读到\BibleQuote{天主是唯一的} \BibleRef{谷 12:29}，以及\BibleQuote{我是}，\BibleQuote{除我以外，没有别的神}，和\BibleQuote{我是元始，我是终末}，这有恰当的意义。因为天主是唯一、独一且是首先的；但这不是为否认圣子才说的，断乎不可！因为圣子就在那唯一、独一且首先者之内，作为出自那唯一、独一且首先者的唯一圣言、智慧和\OriginalTerm{光辉}{Radiance}。而且祂也同样是首先的，因为祂是那首先且独一者的天主性的圆满，是完全且完整的天主。那么，这话不是针对祂而说的，而是要否认有别的像父和祂的圣言那样的存在。\par

% structure: p0009 source=h2 target=subsection reason=relative-heading-rank; base=h2

\subsection{经文解释；第八点，\BibleBook{若望}{福音} 17:3及同类经文。我们的主的天主性不能侵犯圣父作为唯一的天主的特权，因为这特权是由圣子热切维护的。\Quote{唯一}是在与假神和偶像的对比中，而不是在反对藉以发言的圣子时使用的。我们的主将自己的名字加于父的名字中，因为祂包含在父之内。圣父是第一位的，并非好像圣子不是第一位，而是圣父是本源。}

7. 现在，先知这话的意思对所有的人都是清楚明显的；但是，由于那些不虔敬的人，同样引用这些章节，侮辱主并谴责我们，说：\Quote{看哪，天主被称为唯一、独一且是首先的；你们怎么说圣子也是天主呢？因为如果祂是天主，祂就不会说，\BibleQuote{惟独我}，也不会说\BibleQuote{天主是唯一的}；}因此必须进一步说明这些话的意思，尽我们所能，好让所有人由此也知道亚略派实际上是在与天主争辩。如果那时圣子与圣父之间有争竞，那么这样攻击圣子的话才可说出；而且如果如同论到达味关于阿多尼雅和阿贝沙隆的话那样，圣父也如此看待圣子，那就让圣父对自己说出并坚持这样的话，免得身为圣子的祂，自称为天主，导致任何人反叛圣父。但是相反地，如果认识圣子的，就认识圣父，是圣子亲自将父启示给他，并且在圣言内他更要看见圣父，正如已经说过的，并且如果圣子降来，不是光荣自己，而是光荣父，祂对前来见祂的人说：\BibleQuote{你为什么称我为善？没有谁是善的，只有一位是善的，那就是天主}；对询问法律中最大的诫命是什么的人，祂回答说：\BibleQuote{以色列，你要听！上主我们的天主是唯一的天主} \BibleRef{谷 12:29}；对群众说：\BibleQuote{我从天降下，不是为执行我的旨意，而是为执行派遣我来者的旨意}；并教导门徒说：\BibleQuote{父比我大}和\BibleQuote{谁尊敬我，就是尊敬派遣我来的那位}；如果圣子对自己父亲尚且如此，那么有什么困难，使得人们非得对这些经文有那样的看法呢？另一方面，如果圣子是父的圣言，除了这些反对基督的人之外，谁如此狂妄，竟以为天主是这样说话，如同在诽谤和否认自己的圣言呢？这不是基督徒的想法；断乎不可！因为这些经文并非针对圣子而写的，而是要否认那些由人所捏造、虚假称为神的东西。\par

8. 这种解释这些经文的意义是令人满意的；因为既然那些敬拜伪称之神的人背叛了\OriginalTerm{真天主}{True God}，因此，那位良善并眷顾人类的天主，为召回迷途者，说道：\Quote{我是\OriginalTerm{唯一天主}{Only God}}，\Quote{我是}，以及\Quote{除我以外，没有别的天主}等等；为能谴责那些虚无之物，并使一切人归向自己。正如假想在白昼太阳照耀之时，有人粗率地在一块木头上作画，那木头连光的样子都没有，却称那肖像为光的起因，而太阳针对它说道：\Quote{唯独我是白昼的光，除我以外，白昼再无别的光}，他说这话并非针对自己的光芒，而是针对那源于木制肖像的错误以及那虚幻形象的差异；同样地，\Quote{我是}，\Quote{我是\OriginalTerm{唯一天主}{Only God}}，\Quote{除我以外，没有别的}等语，亦即为使人们弃绝伪神，转而承认他是真天主。事实上，当天主说这话时，他是藉着自己的\OriginalTerm{圣言}{Word}所说的，除非现代犹太人再添上这一说：他不是藉着圣言说这话的；但无论这些魔鬼的追随者如何狂妄，他就是这样说了。因为\OriginalTerm{主}{Lord}的圣言临于先知，所听到的便是这话；天主无论说什么或做什么，无不在圣言内说，在圣言内做。所以，基督的仇敌啊，这些话不是针对他而说，而是针对那些与他无关、也不出于他的事物。按照上述比喻，如果太阳说了那些话，他必是纠正错误，这样说话，不是仿佛光芒在他之外，而是在光芒中彰显自己的光。因此，这类经文并非为否认\OriginalTerm{圣子}{Son}，也非针对圣子，而是为推翻虚妄。这样，天主在起初并没有对亚当说这样的话，虽然他的圣言与他同在，万物是藉着他而造成的；因为在偶像进入之前，并无此需要；但当人们起来反抗真理，为自己随意指定神时，才有需要说这样的话，为否认那些虚无之神。我还要补充说，这些话甚至是预见到这些反对基督者的愚妄而说的，好让他们知道，凡他们在\OriginalTerm{圣父}{Father}的\OriginalTerm{性体}{Essence}之外所设想的任何神，都不是\OriginalTerm{真天主}{True God}，也不是那\OriginalTerm{唯一者}{Only}和\OriginalTerm{第一者}{First}的肖像与圣子。\par

9. 那么，如果\OriginalTerm{圣父}{Father}被称为\OriginalTerm{唯一真天主}{only true God}，这话并非为否认那位说过\Quote{我是真理}\BibleRef{若 14:6}的\OriginalTerm{圣子}{Son}，反而是针对那些本性上不真的人，如同圣父和他的\OriginalTerm{圣言}{Word}是真的一样。因此，\OriginalTerm{主}{Lord}自己立刻加上：\Quote{并及你所派遣来的耶稣基督。}如果他是受造物，他就不会加上这话，把自己与造物主并列（因为真与不真之间有何相通？）；但事实上，他把自己与圣父相连，已表明他具有圣父的\OriginalTerm{性体}{Essence}；并使我们知道，他是由真父而生的真子嗣。同样，\Person{若望}也是如此，他按所学得的教导，在自己的书信中写道：\Quote{我们是在那真者内，即在他的子耶稣基督内；这一个\OriginalTerm{真天主}{True God}也是永生}\BibleRef{若一 5:20}。当先知论及创造说：\Quote{那独自展开诸天的}\EditorialNote{依 44:24}，当天主说：\Quote{我独自展开诸天}时，已向众人显明，在\Quote{唯独}中，也意指那\OriginalTerm{唯一者}{Only}的圣言，在祂内\Quote{万物受造}，离了祂\Quote{没有一样受造。}因此，如果诸天是藉着圣言受造，而祂却说\Quote{我独自}，并与那\Quote{唯独}一同理解那藉以造成诸天的圣子；那么，当经上说\Quote{一个天主}、\Quote{我唯独}和\Quote{我是最初的}时，在这一个、唯独和最初中，同样意指与之共存的圣言，如同光中的光芒。这一点除了圣言以外，不能用于任何其他的。因为其他万物是由虚无中藉着圣子而存在，其本性大不相同；但圣子自己是由圣父而生的自然真子嗣；这样，这些愚昧人认为可用来为他们的异端辩护的经文\Quote{我是最初的}，反倒揭露了他们的悖逆之心。因为天主说：\Quote{我是最初的，我是最末的}；那么，如果如同被列于祂以后的事物一样，说祂是它们中的第一，使得它们紧接在祂之后，那么你们肯定只能证明祂在时间上先于受造物而已；这暂不提，已是极大的不虔敬；但如果祂说\Quote{我是最初的}，是为了证明祂并非来自任何物，在祂以前也没有任何物，而祂是万物的根源和原因，并摧毁外邦人的神话，那么同样明显的是，当圣子被称为\OriginalTerm{首生者}{First-born}时，这并非为将祂与受造物同列，而是为证明万物是藉着圣子而构造和收养。因为正如圣父是第一位，同样，祂作为\OriginalTerm{第一者的肖像}{Image of the First}，也因为第一者在祂内，并且祂是来自圣父的子嗣，在祂内，整个受造界被创造并收养为子嗣。\par

% structure: p0013 source=h2 target=subsection reason=relative-heading-rank; base=h2

\subsection{第二十五章 经文解释；第九点，\BibleBook{若望}{福音}10:30；17:11等。\Person{亚略}派解释说，\OriginalTerm{圣子}{Son}与\OriginalTerm{圣父}{Father}在意志和判断上为一；但一切善人，甚至无生命之物也是如此；与圣子对比。他们之间的合一在于本性，因为工作的合一。天使不是祈祷的对象，因为他们不与天主一同工作，而圣子则不然；引证经文。看见天使，并非看见天主。\Person{亚略}派实际上持有两个天主，并倾向\OriginalTerm{外邦人}{Gentiles}的多神论。\Person{亚略}派解释说，圣父与子合一，如同我们与基督合一，但此说被\OriginalTerm{信仰准则}{Regula Fidei}推翻，并藉圣经中的比喻证明无效；比较的真实力量；所用词语的力量。\Quote{在我们内}的力量；\Quote{如同}的力量；由\Person{圣若望}证实。我们如何在\Quote{天主内}并成为他的\Quote{子女}的意义。}

10. 然而，他们在这里又提出自己私下的捏造，主张圣子和圣父并非如教会所宣讲的那样是\Quote{一}或\Quote{相似}，而是照他们自己的意愿。因为他们说，既然圣父所意愿的，圣子也意愿，且无论在思想上或在判断上都不相反，而是在各方面都与圣父一致，宣布相同的教义，发出与圣父的教导一致而合一的言语，因此祂与圣父乃是一体；并且其中有些人竟敢以笔端和言语宣扬这种说法。试问，还有比这更荒谬无理的事吗？如果因此圣子和圣父便是一体，且如果\OriginalTerm{圣言}{Word}就是这样相似圣父，那么天使、我们上面的其他存在——异能者、宰制者、上座者、统治权者，以及我们眼见的日、月、星辰——也应如圣子一样成为儿子；并且论及它们，也当说它们与圣父是一体，且每一个都是天主的\OriginalTerm{肖像}{Image}与圣言。因为凡天主所意愿的，它们也意愿；无论在审判上或教义上，它们都不矛盾，而是在一切事上顺服它们的创造者。因为它们若不是也意愿圣父所意愿的，就不会存留在自己的光荣中。例如那没有存留，而坠入歧途的，曾听见这话：\BibleQuote{啊，\Person{路西法}，清晨之子，你怎么从天坠落？}\BibleRef{依 14:12} 但若真是这样，为何只有祂是\OriginalTerm{独生子}{Only-begotten}、圣言与智慧？既然有这么多相似圣父的，为何只有祂才是肖像？因为就在人间，也会找到许多相似圣父的，例如众多的殉道者，在它们之前有宗徒和先知，更早又有圣祖。如今也有许多人遵守救主的命令，仁慈\BibleQuote{如同他们在天的父}，并遵从劝勉：\BibleQuote{所以你们应效法天主，如同蒙宠爱的儿女一样；又应在爱德中生活，就如\Person{基督}爱了我们。}\BibleRef{弗 5:1} 又有许多人效法了\Person{保禄}，如同保禄效法了基督\BibleRef{格前 11:1}。然而，这些人中没有一个是圣言、智慧、独生子或肖像；也没有人敢说：\BibleQuote{我与父是一体}，或\BibleQuote{我在父内，父在我内}；论到他们众人，反而说：\BibleQuote{上主，在众神之中，谁能像祢？在神的众子中，谁能与上主相比？} 论到祂，反而说：唯有祂是圣父真实而生自本性的肖像。因为我们虽是\BibleQuote{按照肖像}造成的，且被称为天主的肖像与光荣，但这并非由于我们自己，而是由于那居住在我们内的天主的真肖像与真光荣——就是祂的圣言，祂后来为了我们而取了血肉——我们才获得了这一称号的恩宠。\par

11. 他们的这种想法，既明显荒谬无理，又与其他想法一样，因此这相似和合一必须归于圣子的\OriginalTerm{本体}{Essence}；因为若非如此理解，便如前述，不能显示祂有任何超越受造物之处，也不能说祂与圣父相似，而只是与圣父的教义相似；祂与圣父的差异在于圣父是父，而教义和教导却是父的。因此，若就教义和教导而言，圣子与圣父相似，那么按照他们的说法，圣父将仅有父之名，而圣子并非真正的\OriginalTerm{肖像}{Image}，甚或显然毫无与圣父的特性或相似之处；因为一个与圣父如此截然不同的人，还能有什么相似或特性呢？\Person{保禄}的教导虽似救主，本质上却并非相似。既有此等想法，他们便说了谎；其实，圣子与圣父是一体的，正如前述，而且圣子正是如此相似于圣父自己并由祂所出，犹如我们看并理解子之于父，又如我们看\OriginalTerm{光辉}{Radiance}之于太阳。圣子既是如此，所以当圣子工作时，圣父便是工作者；圣子来到圣徒中间时，便是圣父在圣子内来临，一如祂所应许：\BibleQuote{我和父要来，并在他那里作我们的居所。}\BibleRef{若 14:23} 因此，也如我们刚才所言，当圣父赐予恩宠与平安时，圣子也同样赐予，正如保禄在每封书信中所指明的，写道：\BibleQuote{愿恩宠与平安由天主我们的父和主\Person{耶稣基督}赐予你们。} 因为同一恩宠源于圣父，在圣子内赐予，如同太阳与其光辉的光是一体的，且太阳的照明是借着光辉而成就的；同样，当他为\Place{得撒洛尼}人祈祷时说：\BibleQuote{愿天主我们的父自己，和我们的主耶稣基督，指引我们到你们那里去的道路。}\BibleRef{得前 3:11} 他保守了圣父与圣子的合一。因为他不说\Quote{愿他们指引}，好似从两个源头——这个和那个——赐下了双重的恩宠，而是说\Quote{愿祂指引}，以表明圣父借着圣子赐下恩宠；——面对这一点，这些不虔敬的人本应汗颜，却毫不羞愧。\par

12. 因为如果没有\OriginalTerm{合一}{unity}，如果\OriginalTerm{圣言}{Word}不是圣父\OriginalTerm{性体}{Essence}的\OriginalTerm{亲生后裔}{Offspring}，如同光明的\OriginalTerm{光耀}{radiance}，反是子的本性与父相分离，那么单由父赐予便已足够，因为\OriginalTerm{受造物}{originate things}中，没有谁可在\OriginalTerm{造物主}{Maker}的施予中与祂同为分享者；然而事实上，这种施予的方式正彰显了父与子的\OriginalTerm{合一}{oneness}。譬如，没有人会祈求从天主和天使，或任何其他受造物领受，也不会有人说\Quote{愿天主和天使赐予你}；而是从父和子祈求，因为祂们的合一及赐与的合一。因为凡所赐与的，皆是通过圣子而赐与；父除了借子行动以外，什么也不做；如此，领受者所得的\OriginalTerm{恩宠}{grace}才稳固。如果圣祖\Person{雅各伯}，祝福他的孙子\Person{厄弗辣因}和\Person{默纳舍}时说：\BibleQuote{一生牧养我直到今日的天主，救我脱离一切凶恶的天使，愿他祝福这两个孩子！}然而他并没有将受造的自然天使与造物主天主相连，也没有离弃那牧养他的天主，而向天使为其孙辈求祝福；而是当他说\BibleQuote{救我脱离一切凶恶的}时，便指明那并非受造的天使，而是他在祈祷中所联合于父的\OriginalTerm{圣言}{Word}，父正是透过祂，按自己的旨意施行拯救。因为知道祂也被称为父的\BibleQuote{伟大谋略的天使}\OriginalTerm{伟大谋略的天使}{Angel of great Counsel}，他就说，赐福者和救拔于邪恶者正是祂，而非其他。并非他为自身向天主求祝福，而为孙辈向天使求祝福，而是他亲自向祂恳求说：\BibleQuote{我若不蒙你祝福，就不放你走}（因为那正是天主，正如他亲口说：\BibleQuote{我面对面见了天主}），他也祈求这同一位祝福\Person{若瑟}的儿子们。天使的职责是奉命服役，他时常前去驱逐阿摩黎人，被派遣在途中保护百姓；但这些并非他自己的作为，而是派遣和命令他的天主的作为，拯救也是属於天主的，祂要拯救谁便拯救谁。因此，他所见的正是主天主自己，祂对他说：\BibleQuote{看，我与你同在；你不论往哪里去，我必护佑你}；他所见的正是天主，祂阻止了\Person{拉班}的诡计，吩咐他不可说恶言伤害\Person{雅各伯}；他所呼求的正是天主自己，说：\BibleQuote{求你救我脱离我哥哥\Person{厄撒乌}的手，因为我怕他}；因为他在与妻子们的谈话中也说：\BibleQuote{天主却不允许\Person{拉班}加害于我}。\par

13. 因此，\Person{达味}也为自己的救援呼求了天主自己，而非其他：\BibleQuote{我蒙难时呼求上主，他就俯听我；主啊，求你救我灵魂脱离说谎的口唇和诡诈的舌头！}他在上主救他脱离众仇敌和\Person{撒乌耳}之手的日子，也在第十七篇圣咏中向祂谢恩说：\BibleQuote{上主，我的力量，我爱慕你；上主是我的坚固磐石、我的堡垒和我的救主。}\Person{保禄}在经历许多迫害之后，也唯独向天主谢恩说：\BibleQuote{主已救我脱离这一切，我们信赖的那一位，祂还将继续拯救。}同样，祝福\Person{亚巴郎}和\Person{依撒格}的，不是别的，正是天主；依撒格为\Person{雅各伯}祈祷时说：\BibleQuote{愿天主祝福你，使你增长，使你繁殖，你要成为万民的多族；愿祂将我的父亲\Person{亚巴郎}的祝福赐与你。}但如果赐福和救拔唯独属于天主，而雅各伯的救主不是别的，正是上主自己，并且圣祖为孙辈所求的也正是这位救主，那么显然，他在祈祷中所联合于天主的，不是别的，就是天主的\OriginalTerm{圣言}{Word}，因此他称圣言为天使，因为唯独祂启示父。宗徒也是如此，当他说道：\BibleQuote{愿恩宠与平安，由我们的父天主和主耶稣基督赐与你们。}因为如此，祝福才稳固，这是由于圣子对父的\OriginalTerm{不可分离性}{indivisibility}，以及祂们所赐的恩宠是同一的。父施恩，是通过子才成；子被认为恩赐，却是父藉着子并在子内供应；因为宗徒致书给格林多人说：\BibleQuote{我时常为你们感谢我的天主，为天主在基督耶稣内赐与你们的恩宠，}\BibleRef{格前 1:4}。这一点在光与光耀的比喻中也可看见：光所照耀的，光耀便照射；光耀所照射的，其光照也来自光。因此，看见了子，便是看见了父，因为子是父的光耀；这样，父与子是一体。\par

然而，受造之物与受造物并非如此；因为当圣父工作时，并非任何天使或任何其他受造物在工作；因为其中没有一个堪称动力因，它们只是被造之物；而且它们与唯一天主分离、分开，本性相异，且是工程，既不能做天主所做的，也不能像我先前所说，当施予恩宠时，与祂一同施恩。人看见天使时，不会说他看见了圣父；因为天使，如经上记载，是\BibleQuote{奉职的神被派遣去服务}\BibleRef{希 1:14}，且是由祂借圣言赐予接受者的恩赐之传报者。天使显现时，自己承认是由主所派遣的；如同\Person{加俾额尔}向\Person{匝加利亚}，以及向天主之母\Person{玛利亚}所承认的。看见天使异象的人，知道自己看见了天使，而不是天主。因为\Person{匝加利亚}看见的是天使，\Person{依撒意亚}看见的却是主。\Person{三松}的父亲\Person{玛诺亚}看见的是天使，但\Person{梅瑟}看见的是天主。\Person{基德红}看见的是天使，但向\Person{亚巴郎}显现的却是天主。看见天主的，并没有看见天使；看见天使的，也不认为看见了天主；因为受造之物在本性上与造物主天主全然不同，或更确切地说，完全不同。但若有时，当天使被看见时，看见的人听到了天主的声音，如在荆棘丛中所发生的；因为\BibleQuote{主的天使在荆棘丛中的火焰里显现了，主在荆棘丛中呼唤\Person{梅瑟}说：\InnerQuote{我是你父亲的天主，\Person{亚巴郎}的天主，\Person{依撒格}的天主和\Person{雅各伯}的天主}}，然而，那天使并非\Person{亚巴郎}的天主，而是天主在天使内说话。所见的是天使；但天主在他内说话。因为正如祂在帐幕里借云柱向\Person{梅瑟}说话，天主也如此在天使内显现并说话。同样，祂也借一位天使向\Person{农的儿子}说话。但天主所说的，很明显是借圣言，而非借其他东西说的。而圣言，由于与圣父不分离，也不像或不同于圣父的本体\OriginalTerm{本体}{Essence}，祂所做的，就是圣父的工程，祂创造万物与圣父的是一；圣子所给予的，就是圣父的恩赐。看见圣子的人，知道在祂身上看见的并非天使，也非仅仅大于天使者，更非任何受造物，而是圣父自己。听见圣言的人，知道他所听见的是圣父；正如被光辉照耀的人，知道他是被太阳所照耀。\par

因为神圣圣经愿意我们如此理解这事，就给了我们如前所述的这些比喻；凭此，我们既能指证背信的犹太人，又能驳斥外邦人的指控，他们因天主圣三而主张并认为我们宣讲多神。因为，如这比喻所示，我们并不像\Person{马西昂}和\Person{摩尼}的追随者那样，引入三个根源或三个父；因为我们并未暗示三个太阳的形像，而是太阳与光辉。而光辉中的光来自太阳，且为一；因此我们承认只有一个根源；我们宣认创造万有的圣言，其天主性的形式，无异于唯一天主，因为祂由祂所生。那么，疯狂亚略派将理应被控以多神论，甚或无神论，因为他们妄称圣子为外在的、受造的，又妄称圣神出自虚无。因为他们要么会说圣言不是天主；要么虽然承认祂是天主（因为圣经如此记载），但并非属于圣父的本体\OriginalTerm{本体}{Essence}，于是因种类差异而引入多神（除非他们竟敢说，祂只是借分有，如同其他万物一样，被称为天主；倘若他们持此观点，他们的不虔仍是一样，因为他们视圣言为万物之一）。然而，我们决不可有此想法。因为天主性只有一个形式，这形式也在圣言内；只有一位天主，圣父，依其本性自存，超越万有之上，在圣子内显现，因其贯乎万物之中，并在圣神内，因其借圣言在万物中运行。因此我们宣认天主借圣三而为一，我们说，相信唯一的天主性在圣三内，远比异端者那多种类、多部分的天主性更为虔诚。\par

16. 因为如果不是这样，而\OriginalTerm{圣言}{Word}是受造物，是从虚无中造成的工程，那么祂要么不是真天主，因为祂自己也是受造物之一；要么他们若因圣经的缘故称祂为神，就不得不宣称有两位神，一位是创造者，另一位是受造物；并且侍奉两位主，一位是非受造者，另一位是受造者，是受造物；还必须有两种信德，一种信真天主，另一种信一个他们自己制造、塑造而称为神者。如此巨大的盲目，必然导致：当他们朝拜非受造者时，就背弃了受造者；而当他们转向受造物时，就背离了创造者。因为他们无法在一位内看见另一位，因为两者的\OriginalTerm{性体}{Essence}和运作都是外来的、截然不同的。抱持这种想法，他们当然会走向更多的神，因为这正是那些背离唯一真天主者的尝试。那么，当\OriginalTerm{亚略派}{Arians}持有这样的思考和观点时，他们岂不是将自己列于外邦人之中了吗？因为他们也像外邦人一样，朝拜受造物而非创造万有的天主；尽管他们回避外邦人的名称，以欺骗那些不明就里的人，但暗中却怀着与他们相似的念头。因为他们惯于提出那狡猾的说法：\Quote{我们不说有两位非受造者}，这显然是为了欺骗简单的人；因为就在他们宣称\Quote{我们不说有两位非受造者}时，已经暗示了有两位神，且具有不同的性体，一位是受造者，一位是非受造者。虽然希腊人朝拜一位非受造者和多位受造者，而这些人朝拜一位非受造者和一位受造者，但他们与外邦人并无区别；因为他们称为受造者的那位神，不过是众多中的一位，况且希腊人众多的神与这个神具有相同的性体，因为无论是他还是他们，都是受造物。他们是不幸的，更不幸的是，他们的伤害源于反对基督的思想；因为他们离弃了真理，比否认基督的犹太人更可耻的叛徒，他们与外邦人一同打滚，为天主所憎恶，因为朝拜受造物和众多神祇。因为只有一位天主，而非多位，祂的圣言也是唯一的，而非多位；因为圣言就是天主，唯独祂拥有\Emph{圣父}的形像。救主既是如此，便用这些话使犹太人惶惑不安：\BibleQuote{那派遣我来的父，亲自为我作证；你们从未听见过祂的声音，也从未见过祂的形像；并且你们并没有祂的圣言存留在你们内，因为你们不相信祂所派遣的那一位。}\BibleRef{若 5:37}祂恰当地将\Emph{圣言}与\Emph{形像}结合起来，为显示天主的圣言本身就是祂\Emph{圣父}的肖像、印记和形像；而犹太人因为没有接受那向他们说话的一位，也就没有接受圣言，那位就是天主的形像。这正是\Person{圣祖雅各伯}所看见的那一位，他从祂领受了祝福，并且他的名字由\Person{雅各伯}改为\Person{以色列}，正如圣经所见证的，说：\Quote{当他经过天主的形像时，太阳就升起照在他身上。}也正是这一位说过：\Quote{谁看见了我，就是看见了父}，又说\Quote{我在父内，父在我内}，并说\Quote{我与父原是一体}；因为这样，天主是唯一的，对圣父和圣子的信德也是唯一的；因为，虽然圣言是天主，但上主我们的天主是唯一的主；因为圣子是\Emph{圣父}本有的，按其性体的\OriginalTerm{本有}{proper}性和\OriginalTerm{特性}{peculiarity}，是不可分离的。\par

17. 然而，\OriginalTerm{亚略派}{Arians}甚至不因此感到惭愧，回答说：\Quote{并非如你所说，乃是如我们所主张的；}因为，你们既已推翻了我们先前的借口，我们又发明了一个新的，那就是：——圣子与圣父原是一体，而且圣父在圣子内，圣子在圣父内，就如我们也能在祂内成为一体一样。因为这在\WorkTitle{若望福音}中记载着，基督以下面的话为我们祈求：\BibleQuote{圣父啊！求祢因祢的名，保全那些祢所赐给我的人，使他们合而为一，正如我们一样。}\BibleRef{若 17:11}不久之后又说：\BibleQuote{我不但为他们祈求，而且也为那些因他们的话而信从我的人祈求，愿众人都合而为一！父啊！愿他们在我们内合而为一，就如祢在我内，我在祢内，为叫世界相信是祢派遣了我。我将祢赐给我的光荣赐给了他们，为叫他们合而为一，就如我们原为一体一样。我在他们内，祢在我内，使他们完全合而为一，为叫世界知道是祢派遣了我。}然后，这些诡诈的人自以为找到了逃避的路，便补充说：\Quote{如果正如我们在圣父内成为一体，祂与圣父也是一体，并且祂也在圣父内，那么，你怎能从祂所说的\InnerQuote{我与父原是一体}和\InnerQuote{我在父内，父在我内}这些话，就宣称祂是圣父性体本有的、相似的？因为其结论要么是，我们也是圣父性体本有的，要么是祂如同我们一样异于圣父的性体。}他们就这样愚妄地喋喋不休；但在他们的乖僻中，我看到的无非是从魔鬼来的无理性的胆大妄为和鲁莽，因为他们是按照魔鬼的模式说：\Quote{我们要升到天上，我们要与至高者相似。}因为人借着恩宠所获得的，他们竟要使之与施恩者的天主性相等。这样，他们听见人被称为儿子，就以为自己与本性为真子的那位相等了。如今，他们又听见救主说：\BibleQuote{愿他们合而为一，正如我们一样}\BibleRef{若 8:44}，他们就自欺欺人，竟狂妄地以为，自己也能如同圣子在圣父内，圣父在圣子内那样；却没有考虑他们的\Quote{父魔鬼}的堕落，那堕落正是因这样的妄想而发生的。\par

18. 那么，如果我们多次说过，天主的\OriginalTerm{圣言}{Word}与我们相同，除了时间以外与我们毫无分别，那就让祂像我们一样，并与我们一样在圣父前有同等的地位；不要称祂为\OriginalTerm{独生子}{Only-begotten}，也不要称祂为圣父的独一圣言或智慧；但让这同一个名称通用于所有像祂的我们。因为这是合理的，既然他们有同一本性，就应当共享名称，尽管他们因时间而彼此不同。因为亚当是人，保禄是人，现在出生的人也是人，而时间并不改变种族的本性。那么，如果圣言也仅在时间上与我们不同，我们就应当如同祂一样。但事实上，我们既不是圣言或智慧，祂也不是受造物或工程；不然，为什么我们都来自同一个，而祂却是独一圣言呢？然而，即使他们这样说可能合适，至少对我们而言，接纳他们的亵渎是不合适的。但虽然精细分析这些经文可能是多余的，考虑到它们如此清晰而虔敬的含意，和我们自己正统的信仰，然而为了在这里也揭露他们的不虔敬，来，让我们简略地，如同从教父们领受的那样，从经文揭示他们的异端。圣经有一种习俗，以自然界的事物作为人类的形像和说明；它这样做，是为了从这些物质对象解释人的道德冲动，从而显示他们的行为是恶的还是义。例如，在恶的方面，当它指责说：\BibleQuote{不要像那无知的骡马。}或者当它抱怨那些变成这样的人说：\BibleQuote{人处尊荣中而不理解，适与愚蠢的畜牲相似。}又说：\BibleQuote{他们像发情的马}\BibleRef{耶 5:8}。救主为揭露黑落德说：\BibleQuote{去告诉那个狐狸}\BibleRef{路 13:32}；另一方面，却吩咐祂的门徒：\BibleQuote{看，我派遣你们好像羊进入狼群中，所以你们要机警如同蛇，纯朴如同鸽子}\BibleRef{玛 10:16}。祂说这话，不是要我们在本性上变成驮畜，或变成蛇和鸽子；因为祂自己并没有把我们造成那样，因此本性也不允许；而是要我们避免那种畜牲的无理性冲动，并意识到那另一种动物的智慧，不被它所欺骗，且要穿戴鸽子的温良。\par

19. 再者，从神圣主题中为人取榜样，救主说：\BibleQuote{你们应当慈悲，就像你们的在天之父是慈悲的一样}\BibleRef{路 6:36}；又说：\BibleQuote{你们应当是成全的，如同你们的天父是成全的一样}\BibleRef{玛 5:48}。祂说这话，也不是要我们成为如同圣父一样；因为成为如同圣父一样，为我们是受造物，是从虚无中被造成的，是不可能的；而是正如祂命令我们：\BibleQuote{不要像马，}不是怕我们变成驮畜，而是要我们不效法它们的无理性；同样，祂说：\BibleQuote{你们应当慈悲，如同你们的天父，}不是要我们成为如同天主一样，而是要我们仰望祂的仁慈作为，我们所行的善事，不应为人的缘故，而是为祂的缘故，好使我们从祂那里，而不是从人那里获得赏报。因为，正如虽然按本性只有一位圣子，是真的、\OriginalTerm{独生子}{Only-begotten}，我们却也成为儿子，不是像祂那样在本性和真理上，而是依照召唤我们者的\OriginalTerm{恩宠}{grace}，并且尽管我们是从土而来的人，却被称为神，不是如同真天主或祂的\OriginalTerm{圣言}{Word}，而是按赐给我们这恩宠的天主所喜悦的；同样，我们也如同天主成为慈悲的，不是由于被等同于天主，也不是在本性和真理上成为施恩者（因为施恩不是我们的恩赐，而是属于天主），而是为了将我们从天主那里由恩宠所领受的，分施给他人，不加区分，而是慷慨地向众人扩展我们的爱心服务。因为只有这样，我们无论如何才能成为效法者，别无他法，当我们把从祂而来的事物供给他人时。而且，当我们对这些经文给予美好而正确的解释时，对\BibleBook{若望}{福音}中的那段经文同样作此解释。因为他不说，如同圣子在圣父内，我们也必须成为那样：——这怎么可能呢？当祂是天主的\OriginalTerm{圣言}{Word}和智慧，我们是从土中被造的，而祂本性本质上是圣言和真天主（因为若望说：\BibleQuote{我们知道天主子来了，赐给了我们理智，叫我们认识那真实者；我们确实在那真实者内，即在他的子耶稣基督内——他就是真实的天主和永生}\BibleRef{若一 5:20}），而我们却因祂藉着\OriginalTerm{义子名分}{adoption}和\OriginalTerm{恩宠}{grace}，分享祂的\OriginalTerm{圣神}{Spirit}，成为儿子（因为\BibleQuote{凡接受祂的，}他说，\BibleQuote{祂给他们权能，好成为天主的儿女，就是信祂名的人}\BibleRef{若 1:12}），因此祂也是真理（祂说：\BibleQuote{我是真理，}并在对父的致辞中说：\BibleQuote{求你以真理祝圣他们，你的话就是真理}）；而我们却是藉着效法成为有德的、成为儿子：因此，祂说\BibleQuote{使他们合而为一，如同我们一样}，并不是要我们成为祂那样；而是如同祂，作为圣言，在祂自己的父内，好使我们也，以祂为榜样，仰望祂，在和谐与精神的合一中彼此成为一体，不要像格林多人那样分争，而要同心合意，正如\BibleBook{宗徒}{大事录}中的那五千人\BibleRef{宗 4:4}，他们如同一人。\par

20. 因为我们是\Emph{儿子}，而非圣子；是\Emph{神}，而非祂自身；也不是如同圣父，而是\BibleQuote{如同圣父那样仁慈}。而且，如前所述，我们如此成为一体，如同圣父与圣子，我们也将如此，不是如同圣父在本性上在圣子内、圣子在本性上在圣父内那样，而是按照我们自己的本性，并按照我们能够藉此被塑造并学习我们应当如何成为一体的方式，正如我们也学会了仁慈一样。因为同类的事物在本性上彼此合一；这样，所有的血肉在种类上归为一类；但圣言与我们不同，却与圣父相似。因此，祂在本性和真理上与自己的圣父原是一体，而我们，由于彼此属于同一种类（因为所有人都是从一位而受造，所有人的本性也是一样），便在善的意向中彼此成为一体，以圣子与圣父的本性合一为我们的模板。因为祂从自己教导我们要良善心谦，说：\BibleQuote{你们跟我学吧！因为我是良善心谦的} \BibleRef{玛 11:29}，不是要我们变得与祂同等，那是不可能的，而是要我们仰望祂，常保良善心谦；同样在这里，祂希望我们彼此间的善的意向是真实、坚固且不可分解的，便从自己取了榜样，说：\BibleQuote{使他们合而为一，正如我们原为一体一样}，祂的合一性是不可分的；也就是说，他们从我们学习那不可分的本性，也能同样保持彼此之间的和谐一致。而如所说，这种对本性状态的模仿，为人尤为稳妥；因为那些本性状态恒常不变，而人的行为却是变化无常，人可以仰望那本性上不变的事物，避免恶事，并以最好的来改造自己。\par

21. 为此，\BibleQuote{使他们在我们内合而为一}这句话也有正确的意义。譬如，如果我们能够如同圣子在圣父内那样，那么话就该说\BibleQuote{使他们在你内合而为一}，正如圣子在圣父内一样；但实际上，祂并没有这样说；而是说\BibleQuote{在我们内}，借此指出了距离和差异；就是说，祂独自在独一的圣父内，作为唯一的圣言和智慧；而我们则在圣子内，并通过祂在圣父内。祂这样说是这个意思：\Quote{藉着我们的合一，他们也能彼此如此合一，如同我们在本性和真理上原是一体一样；因为除非他们从我们学习合一，否则不能成为一体。}而\BibleQuote{在我们内}具有这种含义，我们可以从保禄的话得知，他说：\BibleQuote{我为了你们的缘故，把这事贴在我自己和阿颇罗身上，好叫你们从我们学习\InnerQuote{不可超过所记载的}这话，免得你们自高自大} \BibleRef{格前 4:6}。那么，\BibleQuote{在我们内}这话，并不是如同圣子在圣父内那样\BibleQuote{在圣父内}；而是暗示一个榜样和肖像，等于说\BibleQuote{让他们向我们学习}。因为正如保禄对格林多人，圣子与圣父的合一也是一切人的模范和教训，他们借此仰望圣父与圣子那本性的合一，可以学习自己该怎样在心神上彼此成为一体。或者，若需别样解释这一说法，\BibleQuote{在我们内}也可意味着，在圣父和圣子的大能内，他们说同样的话，合而为一；因为没有天主，这是不可能的。这种说话方式也见于神圣经典，如\BibleQuote{我们靠天主奋勇作战}；\BibleQuote{我靠天主要跳过墙垣}；\BibleQuote{我们靠你的名要践踏我们的敌人}。因此，很明显，我们能够在圣父和圣子的名字内，成为一体，以坚牢爱德的联系。因为主仍坚持同样的思想说：\BibleQuote{我将在祢赐给我的光荣，赐给了他们，使他们合而为一，正如我们原为一体一样。}祂在这里说得也很合适，不是说\BibleQuote{使他们在我内，如同在祢内一样}，而是\BibleQuote{如同我们}；因为说\BibleQuote{如同}的人，不是指同一性，而是指所论事物的一个肖像和榜样。\par

22. 因此，圣言真实而真正地具有与圣父相同的本性；但对我们来说，则是被赐予去效仿，如前所述；因为祂立即又说：\BibleQuote{我在他们内，你在我内，使他们完全成为一体。}这里主最终为我们所求的，是更大更成全的事；因为圣言显然已来到我们内，因为祂穿上了我们的身体。\BibleQuote{父啊，你在我内}；\BibleQuote{因为我是你的圣言，既然你在我内，因为我是你的圣言，而我在他们内，是因为身体的缘故，并且因为你，人的救恩在我内得以成全，所以我祈求他们也成为一体，按照在我内的身体及其成全；使他们也能成为成全的，与这身体合一，并在它内成为一体；这样，如同所有人都由我承载，所有人都成为一个身体，一个心神，并长成为成全的人。}因为我们所有分享这同一者的人，成为一个身体，拥有同一的主在我们内。这段经文既有此意，基督之敌的谬说就更显被驳倒了。我再说：如果祂简单绝对地说\BibleQuote{使他们合而为一在你内}，或\BibleQuote{使他们和我合而为一在你内}，天主的敌人就有了一些借口，虽则是可耻的；但实际上祂没有简单地这样说，而是\BibleQuote{父啊，如同你在我内，我在你内，使他们也都合而为一。}再者，使用\BibleQuote{如同}一词，祂表明那些在远的意义上如同祂在父内一样的人；远不是在空间上，而是在本性上；因为在空间上，没有什么远离天主，但唯独在本性上，万物都远离祂。而且，像我以前说过的，凡使用小品词\BibleQuote{如同}的人，是暗示，不是指同一性，也不是指平等，而是指在某个方面所关涉之事的榜样。\par

23. 其实，我们也可以从救主自己那里学到，当祂说：\BibleQuote{就如约纳曾在大鱼腹中三天三夜，人子也要在地心三天三夜。} \BibleRef{玛 12:40} 因为约纳并非如同救主，约纳也没有下到\OriginalTerm{阴府}{hades}；大鱼也不是阴府；约纳被吞下时，并没有把先前被大鱼吞下的人带上来，只是当大鱼奉命时，他独自出来了。因此，\Quote{如同}这个词并不表示相同或相等，而是表示此物与他物；它表明在约纳身上有着某种平行，这是由于那三日的缘故。同样地，当我们使用主所说的\Quote{如同}时，我们既不会如圣子在圣父内那样，也不会如圣父在圣子内那样。因为我们是在心意和精神的合一上，如同圣父与圣子那样成为一体，而救主也将像约纳在地里那样；但正如救主不是约纳，祂下到阴府也不是像约纳那样被吞下，这仅仅是一种平行关系。同样，如果我们也在圣父内如同圣子那样成为一体，我们也不会像圣子，或与祂同等；因为祂和我们只是一种平行关系。为此，\Quote{如同}这个词应用在我们身上，因为那些在性体上彼此不同的事物，在某种关系中被看待时，就变得如同它们一样。因此，圣子自己单纯而无条件地在圣父内，因为这是祂由性体而具有的属性；但为我们这些并非本性如此的人来说，就需要一种形象和榜样，以致祂可以论到我们说：\BibleQuote{如同祢在我内，我在祢内。} \BibleQuote{当他们如此成全时，} 祂说，\BibleQuote{那时世界就知道是祢派遣了我，因为我若不来临，背负他们这个身体，他们中没有一个能得成全，反而全部仍处于腐朽之中。圣父啊，求祢在他们内工作，如同祢赐给我背负这身体一样，将祢的圣神赐给他们，好使他们在圣神内成为一体，并在我内得以成全。因为他们的成全显示祢的圣言曾寓居在他们中间；世界看见他们成全并充满天主，就会完全相信是祢派遣了我，我也曾寓居于此。因为他们的成全从何而来，岂不是由于我——祢的圣言，背负了他们的身体，降生成人，完成了圣父啊，祢交给我的工程吗？这工程得以完成，因为人从罪中得救赎，不再停留在死亡中；反而\OriginalTerm{被神化}{deified}，彼此之间，通过注视我，拥有爱德的联系。}\par

24. 我们为了对这节经文的表述提供一个粗浅的看法，已经引出了许多话语，但是圣若望将在他的书信中，比我们所能做的更简洁、更完美地阐明这些话语的意义。他既要驳斥这些不虔敬之人的解释，又要教导我们如何在主内，主如何在我们内；以及我们如何在祂内成为一体，圣子又如何在本性上与我们差异巨大，并且要阻止亚略派的人不再妄想他们会成为如同圣子一样，免得他们听见这样的话：\Quote{你只是人，并非天主，}和\Quote{你原是贫穷的，不要与富人为伍。} 接着圣若望这样写道：\BibleQuote{我们所以知道我们住在祂内，祂住在我们内，就是由于祂赐给了我们的圣神。} \BibleRef{若一 4:13} 因此，由于赐给我们的圣神的恩宠，我们得以在祂内，祂也在我们内；并且因为这是天主的圣神，所以借着祂居住在我们内，我们既拥有圣神，便理所当然地被视为在天主内，这样天主就在我们内。因此，我们并非如同圣子在圣父内那样，也在圣父内；因为圣子并非仅仅分享圣神才在圣父内；祂也不是领受圣神，而是自己将圣神供给所有的人；圣神也不将圣言与圣父结合，而是圣神由圣言领受。圣子在圣父内，如同祂自己的圣言和\OriginalTerm{光辉}{Radiance}；而我们，离了圣神，便是与天主疏远的陌生人，借着分享圣神，我们才被编织进\OriginalTerm{天主性}{Godhead}中；因此，我们在圣父内，并非出于我们自己，而是出于那住在我们内并常存于我们内的圣神，只要我们通过真确的宣认将这圣神保存在我们内，正如圣若望又说：\BibleQuote{谁若明认耶稣是天主子，天主就住在他内，他也住在天主内。} \BibleRef{若一 4:15} 那么，我们与圣子的相似和同等在哪里呢？难道亚略派的人不是从各方面被驳倒了吗？尤其是被圣若望驳倒，他表明圣子在圣父内是一种方式，而我们成为在祂内是另一种方式，并且我们绝不会如同祂一样，圣言也绝不会如同我们一样；除非他们竟敢，像往常一样，如今又说，圣子也是通过分享圣神和品行上的进步才得以自身也在圣父内。然而，仅仅接受这种想法就已经是极度不虔敬了。因为祂，如已说过的，将圣神赐予，而圣神所拥有的一切，都是从圣言来的。\par

25. 救主论到我们说：\BibleQuote{就如你，父啊，在我内，我在你内，使他们也在我们内合而为一}，这并不是表示我们要与祂同一；因为从\Person{约纳}的例子已显明了这点；而是如\Person{若望}所记载，是向父请求，愿圣神借着祂赐予信者，借着祂，我们得以在天主内，并在这方面与祂结合。因为\OriginalTerm{圣言}{Word}在父内，圣神又由圣言赐下，祂愿意我们领受圣神，好叫我们领受之后，因着那在父内的圣言的圣神，我们也因圣神而在圣言内，并借着祂在父内，成为一体。而祂说\BibleQuote{就如我们}，这也不过是请求那赐予门徒的圣神恩宠，能够永不衰退、绝不撤回。因为圣言按本性在父内所拥有的，如我所说，祂愿意借圣神无可挽回地赐予我们；宗徒深知此理，故说：\BibleQuote{谁能使我们与基督的爱隔绝？}又说：\BibleQuote{天主的恩赐和召选是决不会撤回的。}因此，是在天主内的圣神，而不是我们凭自己；而正如我们因着在我们内的圣言而成为儿子和神，同样，我们将在子和父内，并且要被视为在子和父内合而为一，因为那在父内的圣言内的圣神，也在我们内。几时人因某种邪恶而从圣神堕落，若他堕落之后悔改，恩宠仍为愿意者存留，绝不撤回；否则，那堕落者便不再在天主内（因为那在天主内的圣神和\OriginalTerm{护慰者}{Paraclete}已经离弃了他），罪人必在他所顺服的那一位内，就如发生在\Person{撒乌耳}身上的事一样；因为天主的神离弃了他，便有恶神来扰乱他。\BibleRef{撒上 16:14} 天主的仇敌听到这些，理当从此羞愧，不敢再伪装与天主同等。但他们既不明白（因为他说：\Quote{不虔敬的人，不明了知识}），也不能忍受虔敬的话，甚至觉得听来沉重。\par

% structure: p0030 source=h2 target=subsection reason=relative-heading-rank; base=h2

\subsection{第二十六章 从福音书论降生成人的经文导言。列举尚待解释的经文。亚略派与犹太人相比较。我们必须回归\OriginalTerm{信仰准则}{Regula Fidei}。我们的主并非进入人内，而是成为人，因此具有肉身的行动和情感。同一作为既是神性的又是人性的。如此，肉身得以净化，人类成为不朽的。参见\BibleRef{伯前 4:1}。}

26. 因为看啊，他们仿佛在反宗教的言论上不知疲倦，反而像\Person{法郎}那样心硬，当他们听见并看见福音中救主的人性表现时，竟完全像\Person{撒摩撒塔人}一样，忘记了子的父的天主性，并且狂妄大胆地说：\Quote{子怎会是本性地出于父、在本质上相似于父呢？他曾说：\BibleQuote{一切权柄都赐给了我}；又说：\BibleQuote{父不审判任何人，但把审判的全权交给了子}；又说：\BibleQuote{父爱子，并把一切交在他手中；那信从子的，便有永生}；又说：\BibleQuote{我父将一切交给了我；除了子，没有人认识父；除了子和子所愿意启示的人外，也没有人认识父}；又说：\BibleQuote{凡父交给我的，必到我这里来}。}对此，他们辩称：\Quote{如果他像你们说的那样是本性的子，他就无需领受，而是作为子本性地拥有。}\Quote{他又怎会是父的本性和真正的德能呢？他将近受难的时候却说：\BibleQuote{现在我心神烦乱，我该说什么才好？父啊！救我脱离这时辰吧！但我正是为此才面对这时辰的。父啊！光荣你的名吧！}于是有声音从天上说：\BibleQuote{我既已荣耀了它，还要再荣耀它。}\BibleRef{若 12:27} 另一次他也说了同样的话：\BibleQuote{父啊！如果可能，就让这杯离开我吧！}又说：\BibleQuote{耶稣说了这话，心神烦乱，就作证说：\InnerQuote{我实实在在告诉你们：你们中有一个要出卖我。}}}于是这些乖僻的人争辩说：\Quote{如果他是德能，他就不会害怕，反而会给别人提供力量。}他们接着说：\Quote{如果他按本性真是父固有的智慧，经上怎么写着：\BibleQuote{耶稣在智慧和身量上，并在天主和人前的恩宠上，渐渐增长}呢？}同样，当他来到\Place{凯撒勒雅斐理伯}一带时，他问门徒人们说他是谁；当他在\Place{伯达尼}时，他问\Person{拉匝禄}埋葬在哪里；他又曾对门徒说：\BibleQuote{你们有多少饼？}于是他们说：\Quote{他怎会是智慧呢？他既在智慧上增长，又曾不知道自己所问别人的事。}他们也力持此说：\Quote{他怎会是父固有的圣言、父若没有他就从未存在过的、如你们所认为父藉他创造万物的那位呢？他在十字架上曾喊说：\BibleQuote{我的天主，我的天主，你为什么舍弃了我？}在此之前又曾祈祷说：\BibleQuote{光荣你的名吧！}和\BibleQuote{父啊！请以在创世之前我与你同享的光荣来光荣我吧！}他常到旷野去祈祷，也嘱咐门徒祈祷，以免陷入诱惑；并曾说：\BibleQuote{心神固然切愿，但\Emph{肉身}却软弱。}又说：\BibleQuote{至于那日子和那时刻，没有人知道，连天上的天使和子都不知道。}}对此，那些可怜的人又说：\Quote{如果按你们的解释，子是永远与天主同在的，他就不会不知道那日子，而身为圣言理应知道；也不会因为是共存的而被舍弃；也不会要求得到光荣，因为他本来就在父内享有；也根本无需祈祷；因为身为圣言，他什么都不缺；但正因他是受造物，是受造物之一，他方才这样说话，且需要他所没有的；因为需要并缺乏自己所有，正是受造物的本性。}\par

27. 不虔敬者在其言论中所声称的正是如此；他们若如此辩论，就可能会愈发大胆地说：\Quote{圣言究竟为何要成为血肉？} 他们或许还会加上：\Quote{他既是天主，怎能成为人？} 或 \Quote{无形者怎能承受一个身体？} 或者他们可能还以更加犹太化的方式，随同\Person{盖法}说：\Quote{基督既是人，究竟为何要使自己成为天主？} 这正是犹太人在当时所见而暗自低语的，也是如今\OriginalTerm{亚略狂人}{Ario-maniacs}在阅读时不肯相信，反而堕落于亵渎之言中的。如果有人将这两种人的言论仔细对照，就必会发现他们达到同样的不信，其不虔敬的大胆程度相等，且都是与我们共同的争论。因为犹太人说：\Quote{他既是人，怎么能是天主？} 而亚略派说：\Quote{他若是出自天主的真天主，怎能成为人？} 犹太人当时反感并讥嘲说：\Quote{他若是天主子，就不会忍受十字架了。} 而亚略派站在他们一边，催促我们说：\Quote{你们怎敢说，那具有身体，以至忍受了这一切的，是属于父本体所固有的圣言？} 接着，当犹太人因主说过天主是他自己的父，并且使自己与天主平等，如同做父所做的工作，而想要杀害他时，亚略派也不仅学会了否认他与天主平等，且否认天主是圣言的自己本然之父，而且还寻求杀害那些持守此信理的人。再者，犹太人曾说：\Quote{这不是若瑟的儿子吗？他的父亲和母亲我们岂不认识？他怎么能说：\BibleQuote{在亚巴郎之前，我就在} 和 \BibleQuote{我从天降下} 呢？} 而亚略派也相应回应说：\Quote{那如同人一样睡眠、哭泣、询问的，怎么能是圣言或天主呢？} 这样，双方都因救主由于他所承担的血肉而摄取的那些人性特质，否认了圣言的永恒性和天主性。\par

28. 因此，这种错误既是犹太式的，且是背教者\Person{犹达斯}所怀的犹太式想法，就让他们公开承认自己是\Person{盖法}及\Person{黑落德}的弟子，而不要以基督宗教的名义掩盖犹太主义，并且，如我们先前所说，让他们干脆否认救主在肉身内的显现，因为这一学说与他们的异端相近。或者，如果他们因对\Person{君士坦提乌斯}的卑躬屈节，并为了他们所欺骗的那些人，而不敢公然犹太化、接受割损，那就不要讲犹太人所说的话；因为如果他们弃绝这个名称，就该公道地放弃这种教义。因为我们是基督徒，亚略派啊，我们是基督徒；我们的殊荣正是要好好认识有关救主的福音，既不像犹太人那样，一听说他的天主性与永恒性就拿石头砸他，也不像你们那样，因他为我们而说的人性卑微言论而跌倒。如果你们愿意成为基督徒，就脱去\Person{亚略}的疯狂，用宗教的言语清洗你们那被亵渎所污秽的耳朵；要知道，一旦不再作亚略派，你们也就不会再有现今犹太人的恶意了。那时，真理立刻就会从黑暗中照耀你们，你们将不再指责我们主张两个永恒者，而是会亲自承认主按本性是天主的真子，并非仅是永恒的，而是显明为与父的永恒共存者。因为有些被称为永恒的事物，是祂所创造的；因为\BibleBook{}{圣咏}第二十四篇写道：\BibleQuote{城门，请高举你们的门楣，古老的门户，请高抬门扉。} 明显这些事物都是藉着祂而造成的。但如果连这些永恒之事也是祂所创造的，我们当中谁还能争辩说，祂并不先于那些永恒之物，并因此证实祂是主，不是凭祂的永恒性，而是凭祂是天主子？因为祂是子，就与父不可分离，从来没有不存在之时，祂永远是存在的；并且祂是父的肖像和光辉，拥有父的永恒性。上文简要的论述已足以证明他们对其所引经文的错误理解；至于他们现今由福音书中所提出的，他们肯定给出了不正确的解释，如果我们现在思考我们基督徒所持守信仰的宗旨，并以之为准则，如宗徒所教导的，专心阅读受默感的圣经，就很容易看出。因为基督的仇敌，由于不认识这个宗旨，已偏离真理的道路，在绊脚石上跌倒了\BibleRef{罗 9:32}，所思所想超出了所当想。\par

29. 我们屡次说过，圣经的要旨和特色即在于此——它对救主有双重的叙述：祂始终是天主，是圣子，是圣父的\OriginalTerm{圣言}{Word}、\OriginalTerm{光辉}{Radiance}和\OriginalTerm{智慧}{Wisdom}；随后，为了我们，祂由童贞女、\OriginalTerm{天主之母}{Bearer of God}\Person{玛利亚}取了肉躯，而成为人。这要旨在整部受默感的圣经中处处可见，正如主亲自所说的：\BibleQuote{你们查考经典，以为其中有永生，正是这些为我作证}\BibleRef{若 5:39}。但为避免因汇集所有相关章节而使篇幅过长，此处仅提出几处作为例证，首先是\Person{若望}所说：\BibleQuote{在起初已有圣言，圣言与天主同在，圣言就是天主。这位在起初就与天主同在。万物是借着他而造成的；凡受造的，没有一样不是由他而造成的。}随后又说：\BibleQuote{圣言成了血肉，寄居在我们中间；我们见了他的光荣，正如父独生子的光荣。}\Person{保禄}也写道：\BibleQuote{他虽具有天主的形像，却并不以自己与天主同等为应当把持不舍的，反而倒空自己，取了奴仆的形像，成为人的样式；既具有人的形态，便贬抑自己，听命至死，且死在十字架上。}\BibleRef{斐 2:6}。若有人以这些经文为起点，并按其提供的解释通读整部圣经，就会领悟到：起初圣父如何对祂说：\BibleQuote{要有光！}\BibleQuote{要有穹苍！}\BibleQuote{让我们照我们的肖像造人！}但在时期一满，祂派遣祂来到世界，不是为审判世界，而是为叫世界借祂而得救；并且经上记载：\BibleQuote{看，一位贞女将怀孕生子，人将称他的名字为厄玛奴耳，意思是：天主与我们同在。}\BibleRef{玛 1:23}。\par

30. 于是，阅读圣经的人可以从这些古经章节中熟悉这些真理；另一方面，从福音书中他会知道主成了人；因为\Person{若望}说：\BibleQuote{圣言成了血肉，寄居在我们中间。}\BibleRef{若 1:14}。并且祂成了人，而非进入一个人内；这一点必须明白，免得这些不虔敬的人陷入这种想法，并欺哄人以为，正如在古时圣言惯常临到每一位圣人身上，如今祂也寓居在一个人内，同样地圣化他、显示自己，一如在其他圣人身上一样。因为若是如此，祂只是显现在一个人身上，那就毫无新奇之处，那些看见祂的人也不会惊讶地说：\Quote{祂从哪里来？你既是人，为什么把自己当作天主呢？}因为他们从\Quote{上主的话传于}这位或那位先知这类言辞中，早已熟悉这种观念。但是现在，既然天主圣言——万物是借祂而造成的——竟成了人子，自谦自卑，取了奴仆的形像，因此，为犹太人，基督的十字架是绊脚石，但为我们，基督却是\BibleQuote{天主的德能和天主的智慧}\BibleRef{格前 1:24}；因为正如\Person{若望}所说：\BibleQuote{圣言成了血肉}（因为圣经习惯以\Quote{血肉}称呼人，如借\Person{岳厄尔}先知说：\BibleQuote{我要将我的神倾注在一切有血肉的人身上}；又如\Person{达尼尔}对\Person{阿斯提阿革斯}说：\BibleQuote{我不敬拜人手所造的偶像，只敬拜生活的天主，祂创造了天地，对一切血肉具有主权}；因为\Person{岳厄尔}和\Person{达尼尔}都称人类为\Quote{血肉}）。\par

31. 在古时，祂惯常个别地临到圣人们身上，圣化那些正确接待祂的人；但是，那些圣人出生之时，并未曾说祂成了人；他们受苦之时，也并未曾说祂自己受了苦。但当祂为了废除罪恶，在末时代由玛利亚一次而永远地来到我们中间时（因为圣父乐意派遣自己的儿子，\BibleQuote{生于女人，生于法律之下}），这才说，祂取了肉躯而成了人，并在那肉躯内为我们受苦（正如\Person{伯多禄}所说：\BibleQuote{基督既然在肉身内为我们受了苦}\BibleRef{迦 4:4}，为显示并使众人相信，祂虽始终是天主，圣化那些祂所临到的人，并依照圣父的旨意安排一切，后来却为了我们成了人，并且，如宗徒所说，\OriginalTerm{天主性}{Godhead}\BibleQuote{有形地}\BibleRef{哥 2:9}寓居在肉身内；等于说：\Quote{祂既是天主，有祂自己的身体，并以这身体为工具，为了我们成了人。}）因此，肉身的特质——如饥饿、口渴、受苦、疲倦等肉躯所易感受的——都被说成是祂的，因为祂在肉身内；另一方面，圣言本有的工程，如复活死人、使盲者复明、治愈患血漏的妇人，祂都借祂自己的身体而行。圣言承担了肉身的疾苦，如同自己的，因为那肉身是祂的；肉身则服务于天主性的工程，因为天主性寓居在其内，因为那身体是天主的。先知说得好：\BibleQuote{祂背负的是}\BibleRef{依 53:4}；他没有说：\Quote{祂医治了我们的疾苦}，免得，如果祂像往常一样只从外部医治肉身，祂就仍让人臣服于死亡；但祂背负了我们的疾苦，亲自承担了我们的罪，为显示祂为了我们成了人，并且那在祂内承担这些的，正是祂自己的身体；而且，正如\Person{伯多禄}所说的，\BibleQuote{祂在自己身上，亲自承担了我们的罪过，上了木架}，祂自己丝毫没有受损，我们人类却从自己的私欲偏情中得蒙救赎，并充满了圣言的义德。\par

32. 因此，当肉身受苦时，圣言并不在它之外；所以这苦难被称为祂的。当祂以天主性方式做祂父的工作时，肉身并不在祂之外，而是主就在这身体内做了这些事。因此，当祂成为人后，祂说：\BibleQuote{如果我不做我父的工作，你们不必信我；但若是我做了，你们纵然不肯信我，至少要信这些工作，如此你们必定认出父在我内，我在父内。} 因此，当需要治好\Person{伯多禄}妻子的母亲，她患热症时，祂以人性伸出手，却以天主性阻止了疾病。在生来瞎眼之人的事例中，祂从肉身所吐的唾沫是属人的，但祂以天主性方式藉着泥开了他的眼睛。在\Person{拉匝禄}的事例中，祂作为人发出人声；但作为天主，以天主性使\Person{拉匝禄}由死者中复活。这些事之所以这样完成、这样彰显，是因为祂有一个身体，不是幻象，而是真实的；主在穿上人类肉身后，也完整地穿上肉身的固有感受，这是合宜的。这样，正如我们说这身体是祂自己的，我们也同样可以说这肉身的感受惟独属祂，尽管这些感受按祂的天主性并未触及祂。那么，如果这身体属于别人，这些感受也应归于那人；但如果这肉身是圣言的（因为\Quote{圣言成了血肉}），那么肉身的感受也必然归于肉身所属的那一位。而那位被归以这些感受——即被定罪、被鞭打、渴、十字架、死亡及其他肉身的软弱——的，祂也同样拥有胜利与恩宠。为此，这些感受一贯且合宜地不归给另一位，而归给主；这样就使得恩宠出自祂，我们亦能成为不朝拜任何别位，而是真正虔敬天主的人，因为我们所呼求的并非受造之物，亦非寻常之人，而是出自天主的本性真子，祂虽成为人，却并未因此稍减其为主、天主及救主。\par

33. 谁不赞叹此事？谁不认同这样的事真属天主呢？因为，如果圣言天主性的工程未曾藉着身体实现，人就不会被神化；反之，如果肉身的特性未被归于圣言，人就不能完全脱离它们；但正如我先前所说，虽然这些特性曾暂时止息，罪与腐朽却仍存留于他内，恰如在祂之前的人类那样；原因在此：——例如，很多人曾被圣化并清除所有罪恶；不仅如此，\Person{耶肋米亚}甚至从母胎就被祝圣，而\Person{若翰}还在母胎中，就因听到天主之母\Person{玛利亚}的声音而欢喜跳跃；然而，\BibleQuote{但死亡却作了王，从亚当直到梅瑟，连那些没有像亚当一样违法犯罪的人，也属他权下}\BibleRef{罗 5:14}；这样，人依旧如从前那样可朽可腐，易受其本性固有感受的影响。然而，如今圣言成为人，并将属于肉身的一切归为己有，这些事就不再触及身体，因为圣言进入了它，它们反被祂消灭；从此人们不再按自身的感受而处于罪恶与死亡中，而是凭借圣言的大能复活后，常存不朽不腐。因此，当肉身由天主之母\Person{玛利亚}所生，那赐予他人存在根源的祂自己，也被称为诞生了；好使祂能将我们的起源转移到祂自己内，我们便不再像单纯的尘土归于尘土，而是与来自天上的圣言结合，得以由祂携带升天。因此，同样并非没有原因，祂也将身体的其他感受都转移到自己身上；使得我们，不再仅是作为人，而是作为属圣言的人，能分享永生。因为我们不再按照在\Person{亚当}内的昔日起源而死亡；从此以后，我们的起源和肉身的一切软弱既已被转移到圣言，我们便从地中复活，罪恶的诅咒已因那在我们内、并为我们成了诅咒的那位而被除去。而且这是合理的；因为正如我们众人出自尘土，并在\Person{亚当}内死去，同样，我们由水和圣神从上重生，在基督内都得以活命；肉身不再属于地，从此反而成为圣言，因天主圣言为了我们\Quote{成了血肉}的缘故。\par

34. 为使人们能更确切地认识圣言本性中的\OriginalTerm{无痛性}{impassibility}，以及因肉身的缘故而归给祂的软弱，最好聆听\Person{圣伯多禄}；因为他将是关于救主的可靠见证。他在他的书信中如此写道：\BibleQuote{基督既在肉身内为我们受了苦难}\BibleRef{伯前 4:1}。因此，当圣经说祂饥渴、劳苦、不知、睡眠、哭泣、祈求、逃离、诞生、恳求免去这杯，以及一句话，经历所有属于肉身之事时，让我们在每种情况下都合宜地说：\Quote{基督那时饥渴，\InnerQuote{是为我们而在肉身内}；}又说：\Quote{祂不知、被打耳光、劳苦，\InnerQuote{是为我们而在肉身内}；}\Quote{祂被高举、诞生、成长，\InnerQuote{是在肉身内}；}\Quote{祂恐惧、躲藏，\InnerQuote{是在肉身内}；}并说：\BibleQuote{若是可能，就让这杯离开我吧！}\BibleRef{玛 26:39}，\Quote{以及被击打、被接纳，\InnerQuote{是为我们而在肉身内}；}总而言之，所有这一切\Quote{都是为我们而在肉身内。}为此，宗徒自己说了：\Quote{基督既在肉身内受苦}，并非在其天主性内，而是\InnerQuote{为我们而在肉身内}，好让人承认这些感受，并非圣言本身本性固有的，而是那肉身本身本性固有的。\par

因此，谁也不要因那属于人的事而跌倒，反而当知道，在性体上，\OriginalTerm{圣言}{Word}自身是\OriginalTerm{不能受苦的}{impassible}，但因为他穿上了那\OriginalTerm{血肉}{flesh}，这些事便归于他，因为它们是血肉所固有的，而那身体本身则是属于救主的。而他自己，虽然性体上不能受苦，却依然如故，不被这些感受所伤害，反而消除并毁灭它们；人，他们的欲情仿佛在\OriginalTerm{不能受苦者}{the Impassible}内改变且消灭了，从此自己也成为不能受苦的，且永远脱离它们，正如\Person{若望}所教导的，说：\BibleQuote{你们也知道，他显示自己，是为除去罪过，在他内并没有罪过。}\BibleRef{若一 3:5} 事实既是如此，就没有异端者会反驳说：\Quote{血肉本是可死的，为何复活？它若复活，为何不再饥渴、受苦、仍为可死的？因为它来自尘土，它本性的状况怎能脱离它？} 但如今，血肉能向这如此好辩的异端者回答说：\Quote{我来自尘土，本性是可死的，但后来我已成为圣言的血肉。} 而他\Emph{承担了}我的感受，尽管他自己并无这些；这样我也脱离了它们，不再受它们的奴役，因为那使我脱离它们的主。因为如果你反对我脱离本性的腐朽，就要小心，免得你反对天主圣言取了我的奴仆形像；因为正如主，穿上了身体，成了人，同样我们人因圣言而被\OriginalTerm{神化}{deified}，借着血肉被接纳到他那里，从此继承\Emph{永生}。\par

35. 我们认为必须先考察这些要点，以便当我们见他藉他自己的身体这一工具，做出或说出属神的事时，能知道他是天主，所以如此工作；并且，如果我们见他人性地说或受苦，也能知道他曾承受血肉，成为人，并因此他如此行动和说话。因为如果我们识别出各自的特点，并看出并明白这一切都是同一位所做的，我们的信德便是正确的，且永不会迷失。但如果有人看到\OriginalTerm{圣言}{Word}属神地所做之事，却否认身体，或者看到属于身体的事，却否认圣言临在于血肉内，或从人性的事中对圣言产生卑微的想法，这样的人，就像犹太酿酒人，将水与酒混合，必会以十字架为绊脚石，或像外邦人，以宣讲为愚妄。这正是天主的敌人\OriginalTerm{亚略派}{Arians}所遭遇的；因为他们看到救主人性的一面，就断定他是受造物。因此，他们若也观看圣言的属神工作，就当否认他的身体是有起源的，并从此将自己归于\OriginalTerm{摩尼教徒}{Manichees}之列。但对于他们，即使迟缓，也要学习：\BibleQuote{圣言成了血肉；} 而我们，持守信德的整体视野，就应承认他们所解释错误之处，本有正确的解释。\par

% structure: p0042 source=h2 target=subsection reason=relative-heading-rank; base=h2

\subsection{第27章 经文解释：第十点，\BibleBook{玛窦}{福音}11:27；\BibleBook{若望}{福音}3:35，等等。这些经文旨在排除\OriginalTerm{撒伯里乌}{Sabellian}派对圣子的观念；它们与天主教会关于圣子的教义相吻合；它们由\BibleBook{若望}{福音}5:26中的\Quote{照样}所解释。（预先提及下一章。）再者，它们也被用来指我们主的人性；是为了我们的缘故，使我们借着他得以领受而不失落，如同在他内领受一样。这也与圣经的其他部分一致，表明他在领受之前便已有权能等等。他既是天主又是人，他的行为往往同时是属神的又是人性的。}

35（\Emph{续}）。因为，\BibleQuote{父爱子，已将万有交在他手中；} 和 \BibleQuote{我由我父承受了一切；} 以及 \BibleQuote{我由我自己不能做什么；我所听得的，我就这样审判，} 这类经文并不表明圣子曾有一时没有这些特恩——（因为在\OriginalTerm{本质}{essence}上就是\OriginalTerm{父的唯一圣言与智慧}{Only Word and Wisdom of the Father}的那一位，岂不永远拥有父所有的一切？他也曾说：\BibleQuote{凡父所有的一切，都是我的；} 和 \BibleQuote{我的，就是父的。} 因为如果父的一切都是子的，而父永远拥有它们，那么显然子所拥有的，既是父的，就永远在子内），——所以他说这话，并非因为他曾没有这些，而是因为，虽然子永远拥有他所拥有的，但他却是从父而拥有的。\par

36. 因为，免得人因看到子拥有父所有的一切，并由于祂所有之物的完全相似与同一，而陷入\Person{撒伯里乌}的不虔，以为祂就是父，所以祂说了\BibleQuote{赐给我了}和\BibleQuote{我领受了}以及\BibleQuote{交付给我了}，\BibleRef{若 10:18}，只是为了表明祂不是父，而是父的圣言、永生的子，祂因为与父相似，永恒地从父那里拥有祂所有的一切，并因为祂是子，永恒地从父拥有。再者，从这些经文本身我们可以得知，\BibleQuote{赐给了}和\BibleQuote{交付了}之类的话并不减损子的\OriginalTerm{天主性}{Godhead}，反而更显明祂真是子。因为，如果万有都交付给祂，首先，祂与祂所领受的万有不同；其次，祂是万有的承继者，唯独是子，按父的\OriginalTerm{本质}{Essence}专有的。因为如果祂是万有中的一个，那么祂就不是\BibleQuote{万有的承继者}，\BibleRef{希 1:2}，而是每个都按父所愿和所赐的领受了。但如今，祂既然领受了万有，就与它们都不同，且唯独为父所专有。再者，\BibleQuote{赐给了}和\BibleQuote{交付了}并不表示祂曾一度没有，我们可以从类似的经文，并以同样的方式对所有这类经文得出结论；因为救主自己说：\BibleQuote{就如父在自己内有生命，同样祂也赐给子在自己内有生命}，\BibleRef{若 5:26}。从\BibleQuote{赐给了}这句话，祂表明自己不是父；但说\BibleQuote{同样}，则显示子与父在本性上的相似和专有关系。如果父曾一度没有，那么显然子也曾一度没有；因为如父\BibleQuote{同样}，子也拥有。但如果这样说是不虔的，并且相反，说父历来拥有是虔敬的，那么当子说，如父拥有，\BibleQuote{同样}子也拥有时，他们却说子没有\BibleQuote{同样}而是别样，岂不是不当的吗？因此，圣言是信实的，祂所说所领受的一切，祂一直都拥有，却是从父而来；父诚然不是从任何其他得来，但子是从父而来。如同光芒的例子：如果光芒本身说：\BibleQuote{光把一切地方都赐给我照亮，我并非从自己照亮，而是照光所愿}，然而，这样说并不表示它曾一度没有，而是意指：\BibleQuote{我是专属于光的，光的一切都是我的}；同样，甚至更应如此理解子的情况。因为父将一切赐给了子，在子内仍拥有一切；而子拥有了，父仍拥有它们；因为子的\OriginalTerm{天主性}{Godhead}就是父的\OriginalTerm{天主性}{Godhead}，这样，父就在子内对万有施行祂的\OriginalTerm{天主眷顾}{Providence}。\par

37. 虽然这些话的含义如此，但那些以人性方式论及救主的话，也容许有虔敬的意义。因为我们预先考察它们就是为了这个目的：如果我们听到祂问\Person{拉匝禄}安放在何处，或者当祂来到\Place{凯撒勒亚}一带时问\BibleQuote{人们说我是谁？}，或\BibleQuote{你们有多少饼？}以及\BibleQuote{你愿意我给你做什么？}时，我们可以根据前面所说的，知道这些经文的正确意义，不致像基督的仇敌亚略派那样跌倒。首先，我们必须向这些不虔的人提出这个问题：他们为什么认为祂无知？因为一个发问的人，未必是由于无知才发问；而是有可能一个知道的人，仍然询问祂所知道的事。正如\Person{若望}知道，当基督问\BibleQuote{你们有多少饼？}时，祂并非不知道，因为若望说：\BibleQuote{祂说这话是为试探他，因为祂自己知道要做什么}，\BibleRef{若 6:6}。但如果祂知道自己在做什么，那么祂就不是出于无知，而是以知识在问。从这个例子我们可以理解类似的；就是，当主发问时，祂并非因无知才问拉匝禄在哪里，或人们说祂是谁；而是知道自己所问的事，也清楚自己将要做的事。如此，他们的诡辩就轻易地被驳斥了；但如果他们仍因祂的发问而坚持，那么必须告诉他们，在\OriginalTerm{天主性}{Godhead}中确实没有无知，但无知属于\OriginalTerm{血肉}{flesh}，正如已经说过的。而且情况确实如此，请注意：那位询问拉匝禄安放在何处的主，当祂不在现场而是相距甚远时，自己却说\BibleQuote{拉匝禄死了}，\BibleRef{若 11:14}，并且知道他在哪里死的；以及那位被他们认为无知者，正是那预知门徒们的心思，知道每个人心里有什么，\BibleQuote{在人内有什么}，并且更重要的，惟独认识父，说\BibleQuote{我在父内，父在我内}的那一位。\par

38. 因此，这一点对所有人都是清楚的：肉身确实是无知的，但\OriginalTerm{圣言}{Word}本身，就其作为圣言而言，甚至在万物未成之前就已知道一切。因为当祂成为人时，并没有停止是天主；也不是因为祂是天主，就回避属人的事；绝无此意！相反，祂身为天主，却取了自己的肉身，且在肉身内\OriginalTerm{神化}{deifies}了肉身。因为正如祂在肉身内发问，同样也在肉身内复活死人；祂向众人显示：那位赐生命给死人、召回灵魂的主，更认识万有的隐秘。祂知道\Person{拉匝禄}躺在哪里，却仍发问；因为那为我们的缘故忍受一切的天主至圣之言做了这事：祂这样背负我们的无知，是为赐给我们认识祂自己唯一而真实的父，以及认识祂自己——祂是为了万民的\OriginalTerm{救恩}{salvation}，因我们而被派遣的那一位；没有比这更大的\OriginalTerm{恩宠}{grace}了。因此，当救主使用他们据以自辩的那些话——\BibleQuote{权柄赐给了我}、\BibleQuote{光荣你的儿子}——并且\Person{伯多禄}说\BibleQuote{权柄交给了他}时，我们都在同一意义上理解所有这些经文：祂说这一切，是\OriginalTerm{以人性方式}{humanly}，因着\OriginalTerm{身体}{body}的缘故。因为虽然祂原无所需，但祂仍被说成以人性方式领受了祂所领受的，另一方面，既然主已领受，而所赐的恩惠已寄存在祂内，恩宠便得以稳固。因为人若仅仅是领受，就可能再失去（正如\Person{亚当}的情形所示：他领受了，却失去了），但为使恩宠不可收回，且能在人内稳固保存，因此祂自己将恩赐归于己有；祂说，作为人，祂领受了权能，而这权能，祂作为天主，向来拥有，并且祂说\BibleQuote{光荣我}（而祂本是光荣他人的那一位），为显示祂拥有一个需要这些事的肉身。因此，当肉身领受时，由于那领受者是在祂内，且祂因取了肉身而成为人，所以祂自己被说成是领受了。\par

39. 那么，如果（正如已多次说过的）圣言未曾成为人，你们就会按自己的意思，把领受、需要光荣和无知归于圣言；但如果祂已成为人（祂确实已成为人），而领受、需要和无知是属人的事，我们为什么将施与者视为领受者，怀疑那分施于人的主有所欠缺，并将圣言与父分离，视其为不完全和匮乏者，同时却剥夺了人性的恩宠呢？因为如果圣言本身，就其作为圣言而言，为自己领受了光荣并受到了光荣，如果按祂的\OriginalTerm{天主性}{Godhead}，祂是那被\OriginalTerm{圣化}{hallowed}并复活了的，那么人还有什么希望呢？因为他们仍保持原状，赤身、可怜、死了，与赐给圣子的事物无份。圣言又为什么来到我们中间，成为肉身呢？如果是为了领受祂所说的那些事物，那么在此之前祂就没有这些，并且必然反而应感谢身体：因为在祂进入身体之后，才从父那里领受了这些，而这是祂降生成人之前所没有的。照这种说法，祂似乎因身体而自身被提升了，而不是身体因祂而被提升。但这种想法是犹太式的。然而，如果是为了救赎人类，圣言才来到我们中间；并且为了圣化并神化人类，圣言成为肉身（祂正是为此而成为肉身），谁看不出，随之而来的结论就是：祂所说当祂成为肉身时领受的一切，并非为了自己，而是为了肉身提及的呢？因为在祂内说话的那个位格，藉着祂从父所领受的恩赐，本就属于肉身。但让我们看看祂所求的是什么，以及祂所说领受的事物究竟是哪些，这样他们或许也能有所感受。祂曾求光荣，然而祂也曾说：\BibleQuote{一切都交给了我}\BibleRef{路 10:22}。复活后，祂说自己领受了一切权能；但在这之前，祂早已说过\BibleQuote{一切都交给了我}，祂本就是万有的主，因为\BibleQuote{万物是藉着祂而造成的}；并且\BibleQuote{只有一个主，万物却是出于祂}\BibleRef{格前 8:6}。当祂求光荣时，祂本来就是那光荣的主，正如\Person{保禄}所说：\BibleQuote{如果他们认识了，决不至将光荣的主钉在十字架上}\BibleRef{格前 2:8}；因为祂在说\BibleQuote{在世界未有以前，我在祢前所有的光荣}\BibleRef{若 17:5}时，祂所祈求的，正是祂本来就拥有的光荣。\par

40. 同样，祂所说复活后领受的权能，祂在领受之前、在复活之前就已拥有。因为祂曾亲自斥责\Person{撒殚}，说：\BibleQuote{撒殚，退到我后面去！}\BibleRef{路 4:8}；祂也曾赐给门徒权能制伏撒殚，当他们回来时祂说：\BibleQuote{我看见撒殚如同闪电从天跌下}\BibleRef{路 10:18}。再者，祂所说领受的，在领受之前就已拥有，这从祂驱逐魔鬼，解开撒殚所束缚的人——如同祂在\Person{亚巴郎}的女儿身上所做的——；从祂赦免罪过，对瘫痪者和为祂洗脚的妇人说：\BibleQuote{你的罪赦了}；以及从祂复活死人，并修复瞎子的原始本性，恩赐他看见，都可以看出。祂做这一切时，并没有等待领受，而是\Emph{拥有权能}。从这一切可以清楚看出，祂身为圣言所拥有的，当祂成为人并复活后，祂说祂以人性方式领受了；这是为了使世人从此以后，在地上能有制伏魔鬼的权能，因为已分享了\OriginalTerm{天主性}{divine nature}；并在天上，脱离败坏，永远为王。因此我们必须彻底认清：祂所说领受的一切，并非祂原先没有而领受的；因为圣言身为天主，原本就拥有一切；在这些章节中，祂被说成以人性方式领受，乃是为了：当肉身在祂内领受之后，从此恩赐便能为我们稳固地存留。因为\Person{伯多禄}所说的\BibleQuote{祂由天主领受了尊敬和光荣，天使也屈伏在祂权下}就有这个意思。正如祂以人性方式询问，却以天主的方式复活了\Person{拉匝禄}，同样\BibleQuote{祂领受}是论祂的人性方式而言，而天使的屈伏则显明圣言的天主性。\par

41. 所以，\OriginalTerm{被天主所憎恶者}{abhorred of God}啊，住口吧！不要贬低\OriginalTerm{圣言}{Word}，也不要减损祂的\OriginalTerm{天主性体}{Godhead}——这性体原属于圣父——仿佛祂有所需要或无知似的；免得你们像那曾用石头砸祂的犹太人一样，把自己的论点反过来攻击基督。因为这一切并不属于作为圣言的圣言，而是属于人；正如当祂吐唾沫、伸手、呼唤拉匝禄时，我们并不说那些奇事是人性的，尽管是通过肉身实现的，而是天主的作为；同样，另一方面，虽然福音中将人的事归于\OriginalTerm{救主}{Saviour}，但我们应考虑到所言之事的性质，知道这些事与天主不相称，而不要将它们归给圣言的天主性体，却要归给祂的\OriginalTerm{人性}{manhood}。因为虽然\BibleQuote{圣言成了血肉}，但情感却是属于肉身的；虽然肉身在圣言内为天主所据有，但恩宠和权能仍属于圣言。祂于是通过肉身行圣父的工；而同样真实地，肉身的各种情感也在祂身上显示出来；例如，祂询问，然后使拉匝禄复活；祂责备母亲，说：\BibleQuote{我的时辰还没有到}，随即又变水为酒。因为祂是在肉身内的\OriginalTerm{真天主}{Very God}，又是在圣言内的真肉身。因此，祂通过自己的作为，既启示了自己是\OriginalTerm{天主子}{Son of God}，也启示了祂自己的圣父；同时，通过肉身的各种情感，祂表明自己具有真实的肉身，且是祂自己的肉身。\par

% structure: p0050 source=h2 target=subsection reason=relative-heading-rank; base=h2

\subsection{二十八、经文解释；第十一，\BibleRef{谷 13:32}和\BibleRef{路 2:52}。\OriginalTerm{亚略派}{Arians}对前一段经文的解释违反\OriginalTerm{信理准则}{Regula Fidei}，也与上下文相悖。我们的主因着祂的\OriginalTerm{人性}{human nature}，说祂不知道那日子。若\OriginalTerm{圣神}{Holy Spirit}知道那日子，那么\OriginalTerm{圣子}{Son}也知道；若圣子认识\OriginalTerm{圣父}{Father}，那么祂也知道那日子；若祂拥有圣父的一切，那么也就知道那日子；若祂在圣父内，祂就在圣父内知道那日子；若祂创造并维系万物，祂就知道它们何时终止。从\BibleRef{玛 24:42}的论证可知，祂作为人而不知道。正如祂询问拉匝禄的坟墓等，虽明知却仍询问，同样祂知道；又如圣保禄说：\Quote{或在肉身内，我不知道}等，虽明知却如此说，同样祂知道。祂说不知道，是为了我们的益处，叫我们不要好奇\BibleRef{宗 1:7}。正如\OriginalTerm{全能者}{Almighty}询问亚当和加音，虽明知却仍询问，同样圣子作为天主也知道。再者，祂作为人在智慧上也日渐增长，否则祂会使\OriginalTerm{天使}{Angels}在自己之前就臻于完美。祂增长，是因为\OriginalTerm{天主性体}{Godhead}在祂内随时间越发彰显出来。}

42. 这些事既是如此，来吧，让我们现在细查\BibleQuote{至于那日子和那时刻，没有人知道，连天主的天使也不知道，子也不知道}这段经文；因为那些人对这话极其无知，又为之惊愕，便以为从中找到了支持他们\OriginalTerm{异端}{heresy}的重要论据。但在我看来，当那些\OriginalTerm{异端者}{heretics}提出这话并据以自恃时，我看到的是那些再次与天主为敌的巨人。因为那创造了天地、万有借祂而成的主，竟要在他们面前为日子和时刻而争讼；那通晓万有的\OriginalTerm{圣言}{Word}，竟被他们指控对某一日子无知；那认识\OriginalTerm{圣父}{Father}的子，竟被说成对某日某时无知；如今还有什么话比这更违背常理，有什么疯狂能与此相比呢？万有原是借着圣言造成的，时间和季节、黑夜和白昼以及整个受造界；难道万物的创造者反而对祂的工程无知吗？而且，这段经文的上下文本身就显示\OriginalTerm{天主子}{Son of God}知道那时刻和那日子，尽管亚略派因自己的无知而一头栽倒。因为在说完\BibleQuote{子也不知道}之后，祂便向\OriginalTerm{门徒}{disciples}讲述那日子之前的事，说：\BibleQuote{这些事将要发生，然后是终结}。那讲述那日子之前之事的主，必定也知道那日子，那日子将在所预言的事之后显现。如果祂不知道那时刻，祂就不能指明它之前的事件，因为不知道它何时会发生。正如有人要为一群不认识某房子或某城市的人指示那地方，就先讲述那房子或城市前面的情景，一一描述之后说：\Quote{随后立刻就看见那城或那房子了}，他当然知道那房子或城市在哪里（因为他若不知道，就不会描述它前面的情景，免得因无知把听者引到错路上，或在讲述时不知不觉越过目标）；同样，主既然说出那日子和那时辰之前的事，就确切知道那时辰和那日子何时来到，并非无知。\par

43. 至于祂虽知道，当时却不明确告诉门徒的原因，既然祂保持沉默，谁也不要好奇追问；因为\BibleQuote{谁曾知道主的心意？谁曾作过祂的谋士？}\BibleRef{罗 11:34}然而，祂虽知道，却说\BibleQuote{不，连圣子也不知道}，我认为信友中无人不知，祂这样说乃是如同其他类似的话，因\OriginalTerm{肉身}{flesh}而以人的身份说的。因为这不是\OriginalTerm{圣言}{Word}的缺失，而是那本以无知为特性的\OriginalTerm{人性}{human nature}的缺失。若诚实地考察救主说这话的时机，何时以及向谁说的，就能清楚看出这一点。这话不是在祂创造天之时，也不是在祂与\OriginalTerm{圣父}{Father}同在、\OriginalTerm{圣言}{Word}\Quote{措置万物}之时，更不是在祂成为人之前说的，而是在\BibleQuote{圣言成了血肉}\BibleRef{若 1:14}之后。因此，将祂成为人之后以人的方式所说的一切，都归于祂的人性，是合理的。因为圣言理当知晓所造之物，不会不知道这些事物的起始与终结（因为这些工程是祂的），祂知道自己造了多少物，以及它们维持的界限。既然知道各个事物的起始与终结，祂必定知道万物的普遍与共同的终局。当然，当祂在福音中论及自己人性的身份时说：\BibleQuote{父啊！时辰来到了，光荣你的圣子}，显然，作为圣言，祂也知道万物终局的时辰，虽然作为人祂对此无知，因为无知属于人，尤其是对这些事的无知。再者，这也合乎救主对人的爱；因为祂既成了人，就不因那无知的肉身而羞于说\BibleQuote{我不知道}，为要显示祂作为天主知晓，但按肉身却是无知的。因此祂没有说\BibleQuote{不，连天主子也不知道}，免得\OriginalTerm{天主性}{Godhead}显得无知，而只说\BibleQuote{不，连圣子也不知道}，好使那无知属于由人而生的圣子。\par

44. 因此，祂提到了天使，却没有进一步说\Quote{不，连圣神也不知道}；祂保持沉默，具有双重暗示：首先，如果\OriginalTerm{圣神}{Spirit}知道，那么圣言作为圣神由之所受者，更该知道；其次，祂对圣神保持沉默，清楚表明祂是就自己的\OriginalTerm{人的职务}{human ministry}而说\BibleQuote{不，连圣子也不知道}的。证据就是：当祂以人的方式说了\BibleQuote{不，连圣子也不知道}之后，却显示了自己按天主性知道一切。因为那位祂宣称不知道那日子的圣子，祂却宣称他认识父；因为祂说：\BibleQuote{除了圣子外，没有人认识父}\BibleRef{玛 11:27}。除了\OriginalTerm{亚略派}{Arians}之外，所有人都会一致承认，那位认识父的，必然更加认识整个受造界；而在整个受造界中，也认识其终局。如果日子和时辰已经由圣父所定，那么显然它们是通过圣子而定的，祂自己也知道那通过祂所定的事，因为没有一样事物不是通过圣子而存在和确定的。因此，祂作为\OriginalTerm{宇宙的塑造者}{Framer of the universe}，知道圣父愿意创造的万物具有何种性质、何等规模，并带有何种界限；而在这种\Quote{多少}和\Quote{多远}之中，也包含了它的期限。再者，如果凡属于圣父的，也都属于圣子（这是祂自己所说的，\BibleRef{若 16:15}），而知道那日子是圣父的属性，那么显然圣子也知道，因为祂从圣父得到这一专属于己的。再者，如果圣子在圣父内，圣父在圣子内，而圣父知道那日子和时辰，那么显然，圣子既在圣父内，并知道圣父的事，祂自己也就知道那日子和时辰。如果圣子也是圣父的\OriginalTerm{真像}{Very Image}，而圣父知道那日子和时辰，那么显然，圣子也具有与圣父同样的知道它们的能力。如果那创造万有、宇宙赖以存在的祂，自己知道所造之物，以及个别与全部的终局何时来临，这并不稀奇；稀奇的是，这种狂妄，虽然与\OriginalTerm{亚略狂热派}{Ario-maniacs}的疯狂相称，竟迫使我们进行如此冗长的辩护。因为他们将天主之子、永恒的圣言列入受造之物，就离妄称圣父次于受造物不远了；因为如果那认识父的不知道日子和时辰，我担心他们对受造物的认识，或者更确切地说对最低级受造物的认识，会如他们疯狂所称，比对父的认识更大了。\par

45. 但对于他们，当他们如此亵渎\Person{\OriginalTerm{圣神}{the Spirit}}时，正如\Person{主}所说的，他们必永不得赦免这种不虔；但让我们这些热爱\Person{基督}、内心怀有\Person{基督}的人知道，\Person{\OriginalTerm{圣言}{the Word}}，并非无知，就作为\Person{圣言}而言，他说了\BibleQuote{我不知}，因为他确实知道，而是为显示他的\OriginalTerm{人性}{manhood}，因为无知是人所固有的，而他穿上了无知的血肉之躯，因着这躯体，他按肉身说：\BibleQuote{我不知}。为此，在说了\BibleQuote{连圣子也不知道}，并提到\Person{诺厄}时代人们的无知之后，他立即加上说：\BibleQuote{所以，你们要醒寤，因为你们不知道：你们的主人在那一天要来}；又说：\BibleQuote{在你们不料想的时辰，人子就来了}\BibleRef{玛 24:42}。因为我，为了你们而成为你们一样，也说过\BibleQuote{连圣子也不知道}。因为，如果他\OriginalTerm{以天主性}{divinely}无知，他就必定说：\BibleQuote{所以，你们要醒寤，因为我不知道}，以及\BibleQuote{在我想不到的时辰}；但事实上他并没有这样说；而是说\BibleQuote{你们不知道}和\BibleQuote{在你们想不到的时辰}，这表明无知属于人；正是为了人的缘故，他取了与他们一样的肉身，成为人，说\BibleQuote{连圣子也不知道}，因为他在肉身中不知，尽管作为\Person{圣言}他知道。此外，\Person{诺厄}的例证揭露了\Person{基督}仇敌的无耻；因为在那里他也没说：\BibleQuote{我不知道}，而是说：\BibleQuote{他们仍然没有觉察，直到洪水来了，把他们全部卷了去}\BibleRef{玛 24:39}。因为人们不知道，但那位带来洪水的（就是\Person{救主}自己）却知道那日子和时辰，他开启了天上的闸门，裂开了深渊，并对\Person{诺厄}说：\BibleQuote{你和你全家进入方舟}\BibleRef{创 7:1}。因为如果他无知，他就不会预先告诉\Person{诺厄}：\BibleQuote{再过七天，我要使洪水泛滥在地上}。但如果他在描述那日子时，用了\Person{诺厄}时代的类比，而且他确实知道洪水那日，因此他也知道自己来临的那日。\par

46. 再者，在讲述了童女的比喻之后，他更清楚地显示，谁是对那日子和时辰无知的人，他说：\BibleQuote{所以，你们要醒寤，因为你们不知道那日子，也不知道那时辰}\BibleRef{玛 25:13}。他刚才不久说：\BibleQuote{没有人知道，连圣子也不知道}，现在他却不说\BibleQuote{我不知道}，而是说\BibleQuote{你们不知道}。同样，当他的门徒询问末世时，他恰当地按照肉身，因身体之故说：\BibleQuote{连圣子也不知道}；为显示作为人，他不知道；因为无知是人所固有的。然而，如果他是\Person{\OriginalTerm{圣言}{the Word}}，如果他就是那将要来者，那审判者，那新郎，他就知道他何时和哪个时辰来临，以及何时要说：\BibleQuote{你这睡眠的，醒起来吧！从死者中起来吧！基督必要光照你！}\BibleRef{弗 5:14}。因为如同他成为人时，会饥饿、口渴并与人们一同受苦，同样，作为人，他与人们一同不知道；然而，\OriginalTerm{按天主性}{divinely}，作为在\Person{圣父}内的\Person{圣言}与智慧，他知道，且没有他不知道的事。同样，关于\Person{拉匝禄}，他以人的方式询问，而他正要去复活他，并且知道从哪里召回\Person{拉匝禄}的灵魂；知道灵魂在哪里，是比知道身体躺在哪里更大的事；但他以人的方式询问，是为了以天主的方式复活他。同样，当来到\Place{凯撒勒雅}境内时，他询问门徒，虽然他在\Person{伯多禄}回答之前就已经知道。因为如果\Person{圣父}向\Person{伯多禄}启示了对\Person{主}的问题的答案，很明显那启示是通过\Person{圣子}而来的，因为他说：\BibleQuote{除了父外，没有人认识子；除了子和子所愿意启示的人外，也没有人认识父}\BibleRef{路 10:22}。但如果对\Person{圣父}和\Person{圣子}的认识是通过\Person{圣子}启示的，那么毫无疑问，那位询问的\Person{主}，先从\Person{圣父}那里启示给\Person{伯多禄}，然后以人的方式询问；为要显示，他按肉身询问，却按\OriginalTerm{天主性}{divinely}知道\Person{伯多禄}将要说什么。因此，\Person{圣子}知道，因为知道一切，且知道自己的\Person{父}，没有比这知识更大或更完美的了。\par

47. 这足以驳斥他们；但为了进一步显明他们敌视真理，是基督的仇敌，我仍愿向他们提出一个问题。宗徒在\BibleBook{格林多}{后书}中写道：\BibleQuote{我认得一个在基督内的人，十四年前，他是否在肉身内，我不知道，或是否在肉身外，我也不知道；惟有天主知道。}现在你们怎么说？宗徒在神视中所经历的事，他是知道的，尽管他说\BibleQuote{我不知道}，还是他并不知道？如果他并不知道，那么要小心，免得你们因惯于谬误，而陷入\OriginalTerm{弗吕家人}{Phrygians}的过犯中，因为他们说先知们和圣言的其他仆人既不知道他们所做的是什么，也不知道他们所传报的是什么。但如果他说\Quote{我不知道}时是知道的——因为他内有基督向他启示一切——那么天主仇敌的心岂不是真的乖僻而\Emph{自定其罪}吗？因为当宗徒说\Quote{我不知道}时，他们说他其实知道；但当主说\BibleQuote{我不知道}时，他们却说祂不知道。因为如果保禄因基督在他内，便知道自己所说\Quote{我不知道}的事，那么基督自己虽也说\BibleQuote{我不知道}，岂不更知道了吗？所以宗徒因主启示给他，知道他所经历的事；正因如此他才说：\BibleQuote{我认得一个在基督内的人}。而既认得那人，也就知道那人怎样被提去。同样，厄里叟看见了厄里亚，也知道他怎样被提去；但他虽然知道，当先知的弟子们以为圣神将厄里亚投在某座山上时，他因从起初就知道自己所见的，便试图说服他们；但当他们强求时，他便静默不语，任凭他们去寻找。难道他因静默便不知吗？他确实知道，却如同不知一样，任凭他们去，好叫他们因亲身证实，对厄里亚的被提不再疑惑。因此，保禄自己是那被提去的人，更知道自己是怎样被提的；因为厄里亚知道；若有人问他，他便会说出经过。然而保禄却说\Quote{我不知道}，这有两个原因，至少我是这样看的：其一，正如他自己所说的，免得因所得的启示极高超，别人对他的估计超过所见的事实；其二，因为我们的救主说了\BibleQuote{我不知道}，所以他宜于也说\Quote{我不知道}，免得仆人看似高过主，门徒超过师傅。\par

48. 因此，将那知识赐给保禄的，自己更远远知道；因为祂既论到那日子以先的事，正如我先前所说，祂也知道那日子和那时辰，但祂虽然知道，却说：\BibleQuote{不，连圣子也不知道。}为什么那时作为主，明知的事，祂却说\BibleQuote{我不知道}呢？我们可以通过探索猜测，至少我认为，祂这样做是为我们的益处；愿祂恩准我们当前所提出的解释能有真义！救主在我们两方面都保障了利益：祂已经指明了末世以先要发生的事，好叫我们——如祂亲口所言——当这些事发生时，不致惊慌失措，反可由这些事期待其后的终局。至于那日子和时辰，按祂的天主性，祂不愿说\BibleQuote{我知道}，而是按人性说\BibleQuote{我不知道}，这是为了那无知的人性的缘故，正如我先前所说：免得他们再追问下去，那时祂若缄口不言，就会使门徒痛苦；若发言，则可能对门徒和我们所有人造成损伤。因为祂所做的一切，全然是为了我们，既然圣言也是为我们而\BibleQuote{成为血肉}。所以为我们祂说\BibleQuote{不，连圣子也不知道}；祂这样说并非不真实（因为祂是作为人，按人性说\BibleQuote{我不知道}），同时也没有让门徒强逼祂发言，因为祂借着说\BibleQuote{我不知道}阻止了他们的追问。因此在\BibleBook{宗徒}{大事录}上记载着，当祂超越诸天使、以人的身份上升，并带着祂所背负的肉身升上高天时，门徒看见，又追问说：\Quote{何时是终局，祢何时临在？}祂更清楚地对他们说：\BibleQuote{父以自己的权柄所定的时间和时节，不是给你们知道的}\BibleRef{宗 1:7}。那时祂没有像先前按人性说话那样说\BibleQuote{不，连圣子也不知道}，而是说\BibleQuote{不是给你们知道的}。因为如今肉身已复活，脱下可朽，并被\OriginalTerm{被神化}{deified}；当祂正要进入诸天之际，已不再适合按人性作答；从今以后，祂要以属神的方式教训：\BibleQuote{父以自己权柄所定的时间或时节，不是你们应当知道的；但你们将领取德能。}而这圣父的德能不就是圣子吗？因为基督就是\BibleQuote{天主的德能和天主的智慧}。\par

49. 所以子知道，作为他是圣言；因为他说\Quote{我知道，但这不是为你们知道的}时，暗示了这一点；因为我是为了你们的缘故，才按肉身说\Quote{不，连子也不知道}，这是为了你和众人的益处。因为对于你们来说，听到有关天使和子的这些事，是有益的，是由于以后将有迷惑人者；即使魔鬼伪装成天使，试图谈论末日，你们也不要信，因为他们是无知的；并且，如果\OriginalTerm{假基督}{Antichrist}也伪装自己，说\Quote{我是基督}，并试图谈论那日子和末日，以欺骗听众，你们因有我说的\Quote{连子也不知道}这句话，也就可以不信他。再者，不知道末日何时来临，或末日之日何时到来，对人是有益的，免得他们知道了，就会忽略中间的时间，只等待临近末日的日子；因为他们会辩称只有那时才必须留意自己。因此，他也对每个人死亡的时刻保持沉默，免得人因知道而自高，就立即在大部分时间里忽略自己。因此，万物的结局和每个人的终末，圣言都向我们隐藏了（因为在万物的结局中有每个人的结局，在每个人的结局中也包含了万物的结局），这样，既然它不确定且一直悬在面前，我们就可以如同被召唤一样，日复一日地前进，向着前面的目标奔跑，忘记后面的。因为，谁知道末日的日子，不会在间隔中懈怠呢？但若不知道，岂不日复一日地准备着？正是为此，救主补充说：\BibleQuote{所以，你们要醒寤，因为你们不知道主来的时辰；在你们不料想的时辰，人子就来了。}\BibleRef{玛 24:42} 因此，他说这话是为这出于无知而来的益处；因为他说这话，是愿意我们时常准备着；他说：\Quote{因为你们不知道；但我，主，知道我来的时候，尽管\OriginalTerm{亚略派}{Arians}不等候我，我是父的圣言。}\par

50. 因此，主知道什么对我们胜于我们自己的益处，便这样保全了门徒；而门徒们，因着这样的教导，纠正了\Place{得撒洛尼}的信友，因为他们在这方面很可能陷入错误。然而，既然基督的仇敌连这些考虑都不让步，我愿意——虽然知道他们的心比\Person{法郎}更硬——就这一点再问他们。在乐园里，天主问：\BibleQuote{亚当，你在哪里？}他也查问\Person{加音}说：\BibleQuote{你的兄弟\Person{亚伯尔}在哪里？}对此你们怎么说呢？因为如果你们认为他是无知的，所以才问，你们就已经属于\OriginalTerm{摩尼教徒}{Manichees}一党了，因为这是他们狂妄的想法；但如果你们因害怕公开的名号而勉强说，他是在知道的情况下询问，那么，在教义中有什么过分或奇怪的呢，你们竟跌倒，竟然发现子——当时天主就在他里面询问，如今这同一位子穿了血肉，便以人的身份询问门徒？除非，很不幸，你们已经变成了摩尼教徒，你们愿意谴责当时对\Person{亚当}的询问以及一切，好给你们邪恶的心思以充分的空间。因为你们在各方面都暴露了，却依然对\Person{路加}的话窃窃私语，那些话说得正确，但被你们理解错了。而这是怎么回事，我们必须说明，这样，他们败坏的见解也就可以揭露出来。\par

51. 现在\Person{路加}说：\BibleQuote{耶稣在智慧和身量上，并在天主和人前的恩宠上，渐渐地增长。}\BibleRef{路 2:52} 这就是那段经文，因为他们在这里跌倒，我们不得不像法利塞人和撒杜塞人那样，问他们路加所说的是指哪一位。情况是这样的。耶稣基督是人，像其他人一样，还是他是背负血肉的天主？如果他是一个普通人像其他人一样，那么，就让他作为人而增长吧；然而这是\OriginalTerm{萨莫萨特派}{Samosatene}的见解，实际上你们也持这种见解，虽然在人们面前你们在名义上否认它。但如果他是背负血肉的天主，正如他确实是，并且\BibleQuote{圣言成了血肉}，身为天主下降到地上，那么，他本与天主同等，有什么增长呢？或者，子如何增长，因为他永远在父里？因为如果他永远在父里却增长了，那么，请问，还有什么超出父之外可以使他增长的呢？接下来，这里适宜重复一下关于他接受和被光荣的论点。如果他在成为人时增长了，那么显然，在成为人之前，他是不圆满的；而且，与其说他为血肉带来了圆满，不如说血肉成了他获得圆满的原因。再者，如果作为圣言，他增长了，那么，他还有什么可以成为的，胜过圣言和智慧、子和天主的能力呢？因为圣言就是这一切，如果一个人能以某种方式分享它，就像分享一道光线一样，这样的人就在众人当中成为完全的人，并与天使同等。因为天使、总领天使、宰制者、所有异能者和上座者，由于分享圣言，常仰视他父的圣容。那么，那位向别人提供圆满的，自己怎么会比他们更晚才增长呢？因为天使甚至在他人性出生时服事了他，而路加的那段记载是在天使的服事之后。那么，这怎么会进入人的思想呢？或者，智慧怎么会在智慧上增长呢？或者，那位赐给别人恩宠的（正如\Person{保禄}在每封书信中所说的，知道藉着他赐下恩宠，\BibleQuote{愿我们的主耶稣基督的恩宠与你们众人同在}，却自己在恩宠上增长了呢？因为要么他们得说宗徒的话不真，妄称子不是智慧，要么如果他是智慧，如\Person{撒罗满}所说，且\Person{保禄}写道，\BibleQuote{基督是天主的德能和天主的智慧}，那么智慧还能有什么进一步的增长呢？\par

52. 因为人既是受造物，便能在某种方式上向前迈进、修德行善。比如，\Person{哈诺客}是这样被提升的，\Person{梅瑟}日益增长、达到圆满；\Person{依撒格}\Quote{因着进步而成为伟大}；宗徒也说，他\BibleQuote{向前的目标每日迈进}\BibleRef{斐 3:13}。因为每个人都有前进的余地，注目于前一步。但是天主子，祂是唯一独生子，有什么余地能令祂向前迈进呢？因为万有都因注目于祂而前进；祂，作为唯一独生子，就在唯一圣父内，并不从圣父再向前迈进，而是永远住在祂内。所以，进步是属于人的；天主子，因为祂不能进步，祂在圣父内是完美的，却为了我们贬抑自己，好使我们在祂的贬抑中，反而能够增长。我们的增长无非是弃绝感官的事物，而来到\OriginalTerm{圣言}{Word}本身；因为祂的贬抑无非是祂摄取我们的肉躯。所以，并非作为圣言的圣言在进步；祂由完美的圣父而来，完美无缺，无所缺乏，反而使他人得以进步；但在这里，按人性而论祂也被说成进步，因为进步属于人。因此，圣史以谨慎准确的口吻，提到了身量上的进步；可是身为圣言和天主，祂并不以身量来衡量，身量属于肉体。所以，进步是属于肉身的；因为肉身进步时，\OriginalTerm{天主性}{Godhead}对那些看见的人也在显示上进步了。而且，随着天主性愈益彰显，祂作为人的恩宠在众人面前也愈益增长。因为当祂还是婴孩时，就被抱进圣殿；当祂成为少年时，就留在那里，向司祭们询问法律。随着祂的身体逐渐成长，圣言在祂内显现，从此先由\Person{伯多禄}宣认，继而是所有的人：\BibleQuote{这真是天主子}；尽管犹太人，无论是古代的，还是现今的这些，都故意紧闭双眼，不愿看清智慧上的进步，并不是智慧本身的进步，而是人性在智慧内的进步。因为\BibleQuote{耶稣在智慧和恩宠上增长了}；并且，若我们可以说一些既解释又真实的话，那就是：祂在自己内进步了；因为\BibleQuote{智慧为自己建造了居所}，并在自身内使那居所得以进步。\par

53. （何况，所说的这个进步，除了如我之前所说，就是智慧赋予人的\OriginalTerm{神化}{deifying}和恩宠，罪恶在他们内被涂抹，他们内在的腐朽也照着与圣言肉身的相似和亲缘关系被除掉，还有什么呢？）因为这样，随着身体在身量上增长，天主性的彰显也在其中发展起来，并向众人显示：那身体是天主的\OriginalTerm{圣殿}{Temple}，天主在身体内。若有人强辩说，\BibleQuote{圣言成了血肉}就被称为\Person{耶稣}，并将\Quote{进步}一词归给祂，那么必须告诉他们：这并不减损\OriginalTerm{圣父的光明}{the Father's Light}，即圣子，反而仍表明圣言成了人，并承受了真实的血肉。并且，正如我们说祂在肉躯内受苦，在肉躯内饥饿，在肉躯内疲乏，同样，也合理地说，祂在肉躯内进步；因为这样的进步，并非发生在圣言脱离肉躯的情况下，因为进步的血肉就在祂内，并且被称为祂的，如前所述，好使人的进步得以持久不衰，因为有圣言与之同在。所以，进步既不属于圣言，血肉也不是智慧，但血肉成了智慧的形体。因此，如我们已经说过的，不是智慧，就其为智慧而言，在自身方面进步；而是人性在智慧上进步，逐渐超越人性，被神化，并对众人成为并显得如同智慧施行天主性之运作与光辉的器官。所以，圣史并没有说\Quote{圣言进步了}，而是说\Person{耶稣}——主降生成人时被称呼的名字；因此，进步是属于人性的，一如我们上面所解释的。\par

% structure: p0063 source=h2 target=subsection reason=relative-heading-rank; base=h2

\subsection{第二十九章 经文释义；第十二点：\BibleRef{玛 26:39}；\BibleRef{若 12:27}等。\OriginalTerm{亚略派}{Arian}的推论违背\OriginalTerm{信仰准则}{Regula Fidei}，如前所述。祂哀哭等等，是因人而哀哭。其他经文证明祂是天主。天主不能畏惧。祂畏惧，是因祂的肉躯畏惧。}

54. 所以，正如当肉躯成长时，祂被称为成长了，因为身体本是祂自己的；同样，在祂死的时候所记载的，祂烦乱、祂哭了，也必须按同一意义去理解。他们四处奔走，好像要借此重新推举他们的异端，说：\Quote{看啊，\InnerQuote{祂哭了}，又说，\InnerQuote{如今我的灵魂烦乱}，祂还求那杯离开祂；那么，祂若如此说，怎么可能是天主，是圣父的圣言呢？} 是的，圣经也记载祂哭了，天主的仇敌啊，也记载祂说：\BibleQuote{我的灵魂烦乱}，在十字架上祂说：\BibleQuote{厄里、厄里，肋玛撒巴黑塔尼？}意思是\BibleQuote{我的天主，我的天主，你为什么舍弃了我？}，祂还求那杯离开祂。是的，圣经确实这样记载；但我还要再问你们（因为对你们的每一项反对，都必然要做出同样的答复）：如果说这话的只是人，就让他哭泣并畏惧死亡，因为他是人；但如果祂是圣言在肉躯内（我们不必忌讳重复），祂既是天主，又有何可畏惧？祂就是生命，又从死亡中拯救他人，何以畏惧死亡？祂曾说：\BibleQuote{不要害怕那杀害肉身的}（\BibleRef{路 12:4}），自己何以害怕？祂曾对亚巴郎说：\BibleQuote{不要害怕，因为我是你的盾牌}，鼓励梅瑟对抗法郎，并对农的儿子说：\BibleQuote{你应勇敢坚决}，自己何以在黑落德和比拉多面前恐惧呢？再者，那位帮助别人战胜恐惧者（圣经说：\BibleQuote{主偕同我，我不畏惧世人能对我做什么}），自己却畏惧总督、畏惧世人么？那亲自前来战胜死亡的主，自己却恐惧死亡吗？说祂畏惧死亡或畏惧阴府，难道不荒谬不虔敬吗？阴府的门卫见到祂尚且战栗。但如果如你们所主张的，圣言在恐惧中，那么祂为何早已预言犹太人的阴谋，却不逃走，反而在被寻找时说：\BibleQuote{我就是}呢？因为祂本可以避免死亡，正如祂说的：\BibleQuote{我有权舍掉我的生命，也有权再取回它}，以及\BibleQuote{没有人能夺去我的生命}。\par

55. 但是，这些受苦之情并不属于圣言的本性，就祂之为圣言而言；但圣言就在那如此受苦的肉躯内，基督的仇敌和忘恩负义的犹太人哪！因为祂在尚未取得肉躯之前并没有说这一切；而是在\BibleQuote{圣言成了肉躯}，成为了人之后，圣经才记载祂说了这些，即以人性的方式说的。圣经所记载的这一位，正是那使拉匝禄从死者中复活、变水为酒、赐胎生瞎子看见，并说\BibleQuote{我与父原是一体}的那一位。那么，他们若以祂的人性表现作为轻视天主子、甚至把祂看做完全出于地而非出于天的根据，为何不因祂的天主的工程而认识那在圣父内的圣言，从而从此抛弃他们任意妄为的不虔敬呢？因为他们本该看到，那行这些工程的正与那显示自己的身体能受苦——允许身体哭泣、饥饿并表现身体其他特性——的是同一位。因为祂通过这类事表明，祂虽然是不能受苦的天主，却取了能受苦的肉躯；但借着所行的工程，祂显示自己正是天主的圣言，后来成了人，说：你们即使不信我，只凭我披戴人性的身体，也该信我的工程，好叫你们知道并相信\BibleQuote{我在父内，父在我内}。在我看来，基督的仇敌已显出赤裸裸的厚颜无耻和亵渎；因为他们一听到\BibleQuote{我与父原是一体}，就凶狠地扭曲意义，割裂父与子的合一；但当读到祂的眼泪、汗水或苦难时，他们却不顾及祂的身体，反因这些缘故，将创造万物的那位列于受造物中。那么，他们与犹太人还有什么区别呢？因为犹太人不虔敬地将天主的工程归给贝尔则步，同样，这些人将行那些工程的主列于受造物中，也必遭受与犹太人相同的不可饶恕的惩罚。\par

56. 然而，当听到\BibleQuote{我与父原为一体}时，他们理应看出，在基督身上有天主性的一体以及圣父性体的特质；当听到\BibleQuote{他哭了}等类似经句时，则应说明这些事是专属于身体的；尤其因为两方面都有可理解的依据，即一处是论及天主而记，另一处则是针对基督的人性。因为在无形体者内，本不会有身体的属性，除非祂取了一个可朽坏且有死的身体；因为至圣玛利亚是有死的，而祂的身体正是从她而来。因此，当祂带着身体受苦、哭泣、劳碌时，这些专属肉躯的事，连同身体便被归于祂。倘若祂哭泣并忧闷，那哭泣忧闷的并不是就其本身而论的圣言，而是专属肉躯的事；同样，倘若祂祈求这杯离开祂，也不是天主性在惊惧，这情绪也是专属人性的事。至于\BibleQuote{你为什么舍弃了我？}这句话是祂说的，按前文的解释（虽然祂未受任何痛苦，因为圣言是无苦的），仍由圣史们记述了下来；因为主降生成人，这些事都是按一个人的方式去行去说的，为的是祂能亲自减轻肉躯的这些痛苦，并使之摆脱这些痛苦。因此，主也不可能被圣父舍弃，因为祂常住在圣父内，无论是在说话之前，还是发出这呼喊之时。同样，说主在惊惧也是不合常理的，因为地狱的门卫一见祂就战栗，并开启了地狱，坟墓裂开，许多圣者的身体复活起来，显现给他们的同城人。因此，一切异端者都当闭口，绝不敢将惊惧归于主：死亡像蛇一样逃避祂，魔鬼在祂面前战栗，海洋惊骇；诸天为祂裂开，一切威能都动摇。因为看，当祂说出\BibleQuote{你为什么舍弃了我？}时，圣父便显示祂常在自己内，即使那时亦然；因为大地一认出说话的是它的主，立即震动，圣所的帐幔裂开，太阳隐蔽，岩石崩裂，坟墓——正如我已说过的——自行裂开，其中的死人都复活了；而且奇妙的是，当时在场、先前否认过祂的人，一看到这些征兆，便承认说：\BibleQuote{这人真是天主子！}\par

57. 至于祂所说的\BibleQuote{如果可能，就让这杯离开我吧！}，请注意，虽然祂这样说，祂却责备伯多禄说：\BibleQuote{你所体会的，不是天主的事，而是人的事。}因为祂所祈求的，正是祂所愿意的，祂正是为此而来；但意愿是出于祂（因为祂正是为此而来），而惊惧则属于肉躯。因此，祂也是以人的身份说出这话的，然而两者都由同一位说出，为的是显示祂是天主，在自己内愿意，但当祂降生成人时，便有一个会惊惧的肉躯。正是为了这肉躯，祂将自己的意愿与人性的软弱结合，为的是摧毁这情感，进而使人面对死亡不再畏惧。看，真是一件奇特的事！基督的仇敌将惊惧之语归咎于祂，而祂却藉着所谓惊惧的言辞，使人变为无畏无惧。于是，真福宗徒们因祂的这些话语，便对死亡生出如此大的轻蔑，以致于他们毫不理会那审问他们的人，却回答说：\BibleQuote{听从天主胜过听人，是应该的。}\BibleRef{宗 5:29}而其他圣殉道者如此勇敢，竟以为自己是走向生命，而非经历死亡。那么，若一方面赞赏圣言之仆人的勇气，另一方面却说圣言本身——他们正是藉着祂而轻看死亡——曾在惊惧中，这岂不是荒唐吗？但从圣殉道者极其坚忍的志向和勇气可见，天主性并不在惊惧中，而是救主消除了我们的惊惧。因为正如祂以死亡废除了死亡，以人性的方式废除了人性的一切灾祸，同样，祂也藉着所谓惊惧消除了我们的惊惧，使人从此不再畏惧死亡。祂的言语与行动并行。因为\BibleQuote{让这杯离开}和\BibleQuote{你为什么舍弃了我？}这些话是出于人性，而使太阳失光和死人复活的行动，则是出于天主性。再者，祂以人性说：\BibleQuote{我的灵魂现在烦乱；}又以天主性说：\BibleQuote{我有权舍弃我的生命，也有权再取回它。}因为烦乱属于肉躯，而有权舍弃生命并随意再取回，这并非人的特性，而是圣言的权能。因为人死不是凭自己的权能，而是出于本性的必然，且非自愿；但主本身虽是不死的，却有可死的肉躯，作为天主，祂有权随时离开身体并再次取回。关于这一点，达味也在\WorkTitle{圣咏}中说：\BibleQuote{你绝不会将我的灵魂遗弃在阴府，也绝不让你的圣者见到腐朽。}因为按理说，肉躯既是可朽坏的，便不该再按其本性保存有死，却要因穿上它的圣言而常存不朽。因为就如祂来了，带着我们的身体，适应了我们的处境，同样，我们接受祂，便分享那来自祂的不朽。\par

58. 因此，这些\OriginalTerm{亚略狂人}{Ario-maniacs}关于\OriginalTerm{圣言}{Word}的绊脚石借口是徒劳的，而且他们的理解是狭隘的，因为经上记载：\BibleQuote{祂心神烦乱，}\BibleQuote{祂哭了。} 因为他们似乎连人的情感都没有，如果他们对人的本性和属性如此无知；这反而更令人惊讶，圣言竟住在如此受苦的肉躯内，既不阻止那些密谋害祂的人，也不报复那些置祂于死地的人，尽管祂有能力，祂曾阻止一些人死亡，并使另一些人从死者中复活。祂让自己的身体受苦，因为正如我以前说过的，祂正是为此而来，使祂在肉躯内受苦，从此肉躯成为无苦难和不朽的，并且如我们多次说过的，侮辱和其他苦难定要加于祂身上，并且因着祂而完全废去，从此以后不再临到后人；从此人们能永远存留不朽，作为圣言的宫殿。如果基督的仇敌坚持这些思想，并认识到\OriginalTerm{教会的目标}{ecclesiastical scope}乃是信仰的锚，他们就不会使信仰触礁，也不会如此无耻地抗拒那些想要挽回他们免于跌倒的人，并将那些劝诫他们虔敬的人视为仇敌。\par

% structure: p0069 source=h2 target=subsection reason=relative-heading-rank; base=h2

\subsection{第三十章 续前反对（如第七至十章）。圣子是否由圣父的意志所生？这实际上与\Quote{是否曾有一时他不存在？}相同。并为\OriginalTerm{亚略派}{Arians}引入后一问题。\OriginalTerm{信仰规范}{Regula Fidei}立即以相反经文予以否定。亚略派随从\OriginalTerm{华伦提努派}{Valentinians}，主张存在先行的意志；而这种意志仅由天主施于受造物。圣经中的例证。\Person{阿斯忒里乌}的不一致。如果圣子因意志而存在，则在他之前必另有圣言。如果天主因意志而是善的或存在，那么圣子也因他的意志而存在。如果他愿意有理性或智慧，那么他的圣言和智慧就是随他的意志的。圣子是活的意志，并拥有一切表示同性同体的称号。圣父对圣子所有的意志，圣子同样对圣父有。圣父愿意圣子，圣子愿意圣父。}

58. \Emph{续} 然而，看来异端者确实是邪恶的，在各方面他的心灵都败坏而不虔敬。因为你看，尽管他们在所有论点上都被驳倒，并被证明完全失去了理智，他们仍不感到羞耻；正如\OriginalTerm{外教寓言}{Gentile fable}中的\OriginalTerm{海德拉}{hydra}，当它原有的蛇被杀死后，生出新蛇，以产生新蛇来对抗杀旧蛇者，照样，他们对天主心怀敌意和憎恨，如同海德拉，在他们提出的反对中失去生命，又为自己制造出其他犹太式的和愚妄的问题，以及新的权宜之计，仿佛真理是他们的仇敌，以此越发表明他们在一切事上是基督的敌人。\par

59. 在如此多的明证反对他们，甚至连他们的父魔鬼也自觉羞愧而退后之后，他们又从自己邪恶的心中嘀咕出其他的权宜之计，时而低声细语，时而如蚊蚋嗡嗡作响；他们说：\Quote{就这样吧。你们这样解释这些经文，在推理和证据上得胜；但你们仍必须承认圣子是由圣父的意志和喜好而获得存在。}因为他们就这样欺骗许多人，提出天主的意志和喜好。现在如果有任何信仰正确的人在单纯中这样说，就没有理由怀疑这表达，正确的意向胜过这有些简单的用词。但既然这说法出自异端者，而异端者的话是可疑的，并且正如经上所载：\BibleQuote{恶人是诡诈的，}\BibleQuote{恶人的言语是诡诈，}即使他们只是示意，因为他们的心是败坏的，来吧，让我们也考察这个说法，免得他们虽然四面受敌，却仍像海德拉一样，制造出一个新词，通过这种巧妙的语言和似是而非的推脱，以另一种方式再次散播他们的不虔敬。因为说\Quote{圣子由天主的意志而来}的人，与说\Quote{曾有一时他不存在}、\Quote{圣子从无中而来}、\Quote{他是受造物}的人，意思相同。但既然他们现在以这些说法为耻，这些狡猾的人便试图以另一种方式表达他们的意思，提出\Quote{意志}一词，如同\OriginalTerm{墨鱼}{cuttlefish}喷出墨汁，以此来蒙蔽单纯的人，并谨记他们特有的异端。因为他们从哪里提出\Quote{由意志和喜好}呢？或从哪段圣经？让他们说吧，他们如此猜疑自己的言语，如此捏造不虔敬。因为从天上启示自己圣言的圣父宣布说：\BibleQuote{这是我的爱子；}并且藉着达味说：\BibleQuote{我的心涌出美善的圣言；}并且命若翰说：\BibleQuote{在起初已有圣言；}达味在圣咏中说：\BibleQuote{在你那里有生命的泉源，藉你的光明我们才能看见光明；}宗徒写道：\BibleQuote{他是天主光荣的辉耀，}又说：\BibleQuote{他虽具有天主的形体，}以及\BibleQuote{他是不可见的天主的肖像。}\par

60. 众人到处都告诉我们\OriginalTerm{圣言}{Word}的\OriginalTerm{存有}{being}，却没有人说祂是\Quote{\OriginalTerm{出于意志}{by will}}而存有，也根本没有说祂是受造的；但我请问他们，他们在哪里找到了先于天主圣言的\OriginalTerm{意志}{will}或\OriginalTerm{愉悦}{pleasure}，且是\Quote{先存的}呢？除非他们离弃了圣经，模仿\Person{瓦伦廷}的乖谬。因为瓦伦廷派的\Person{托勒密}说，\OriginalTerm{无始者}{Unoriginate}有一对属性，\OriginalTerm{思想}{Thought}和\OriginalTerm{意志}{Will}，祂先思想，然后定意；而祂所思想的，若非加上意志的能力，就无法推出。\Person{亚略派}从此学了一课，希望意志和愉悦先于圣言。那么，就让他们去与瓦伦廷的教义争胜吧；而我们，当我们阅读\OriginalTerm{神圣训言}{divine discourses}时，发现\Quote{祂是}适用于圣子，但我们只听见祂是在圣父内、是圣父的\OriginalTerm{肖像}{Image}；唯独在\OriginalTerm{受造物}{things originate}身上，由于这些物按本性原不存在，后来才得以出现，我们才认识到有先存的意志和愉悦，正如达味在第一百一十三篇圣咏中所说：\BibleQuote{我们的天主在天上，祂愿意做什么，就做什么}\BibleRef{咏 113:11}，在第一百一十篇中说：\BibleQuote{上主的工程何其伟大，凡喜爱祂的必须究察}\BibleRef{咏 110:2}；又在第一百三十四篇中说：\BibleQuote{上主愿意做什么，就在天上、在地上、在海中、在一切的深处成就}\BibleRef{咏 134:6}。如果祂是工程和受造物，与其它物一样，那么，像其它物一样，祂也可说是\Quote{出于意志}而得以存在，圣经也表明这些物就是这样产生的。异端的辩护者\Person{阿斯忒里乌斯}同意这一点，他这样写道：\Quote{如果万有的\OriginalTerm{造物者}{Framer}凭愉悦而造是不相宜的，那么，为确保祂的尊威不受损害，就该把愉悦同样地从一切受造物那里除去。或者，如果天主愿意是合宜的，那么就容这更好的方式适用于首生\OriginalTerm{子嗣}{Offspring}。因为同一位天主，既凭祂的愉悦而造，却不凭祂的意志而造，这是不可能的。}尽管这位\OriginalTerm{诡辩家}{Sophist}在言辞中引入了大量的不虔敬，即\OriginalTerm{子嗣}{Offspring}和受造物是相同的，而圣子不过是所有\OriginalTerm{子嗣}{Offspring}中的一个，他却得出结论说，工程是出于意志和愉悦才是合宜的。\par

61. 因此，如果祂异于万有，正如上面所证明的，并且万有更是借着祂而得以存在，那么就不应将\Quote{出于意志}归给祂，否则祂就会像那些借着祂而存在的受造物一样，以同样的方式存在。因为\Person{保禄}，他先前并不存在，后来按天主的\OriginalTerm{意志}{will}成了宗徒；而我们自己的召叫，本身曾不存在，如今后来才发生，也是为意志所先导，并且正如保禄自己再次说的，是\Quote{按祂意志的\OriginalTerm{良好意愿}{good pleasure}}而成就的\BibleRef{弗 1:5}。至于\Person{梅瑟}所叙述的：\BibleQuote{要有光！}、\BibleQuote{大地出现！}、\BibleQuote{让我们造人！}，我认为，按照前文，这表明了行动者的意志。因为那些曾不存在、后来由外在原因而出现的事物，\OriginalTerm{造物者}{Framer}商议去造它们；但祂自己的\OriginalTerm{圣言}{Word}，由祂本性所生，关于祂，祂并没有预先商议；因为在圣言内，圣父造作并形成祂所商议的其他一切；正如\Person{雅各伯}宗徒教导说：\BibleQuote{祂自愿用\OriginalTerm{真理的圣言}{Word of truth}生成了我们}\BibleRef{雅 1:18}。因此，天主对万有的意志，无论是被重新生成还是最初被带入存在，都在祂的圣言内，在圣言内祂既做又重生了祂认为适宜的事；正如宗徒在致得撒洛尼人书中再次指明\BibleRef{得前 5:18}：\BibleQuote{因为这是天主在基督耶稣内对你们所有的意志}。但是，如果圣父在谁内造，意志也在谁内，并且圣父的愉悦是在基督内，那么祂怎能像其他受造物一样，凭意志和愉悦而存在呢？因为如果祂也如你们所坚持的，是凭意志而存在，那么随后的就是：关于祂的意志存在于另一个圣言内，祂再借着那圣言而得以存在；因为已经证明，天主的意志并不在祂所带入存在的那些事物内，而是在祂所借以并在其内使一切受造物得以存在的（那一位）内。再者，既然说\Quote{出于意志}和说\Quote{祂曾经不存在}是完全一样的，就让他们下定决心说\Quote{祂曾经不存在}吧，好让他们羞惭地认识到，后者指的是时间，从而明白说\Quote{出于意志}就是把时间置于圣子之前；因为商议先于那些曾不存在的事物，就如所有受造物那样。但是，如果圣言是受造物的造物者，并且祂与圣父\OriginalTerm{共存}{coexists}，那么商议怎能先于\OriginalTerm{永恒者}{Everlasting}，好像祂不存在一样呢？因为如果商议在先，万有又怎能借着祂呢？相反，祂也如同其他（受造物）一样，是凭意志受生为圣子，就如我们也是借着真理的圣言成为儿子一样；并且，如前述，剩下的就是寻求另一个圣言，借着祂，祂也得以产生，并与一切依照天主愉悦而有的万物一同受生。\par

62. 那么，如果另有天主的圣言，圣子便是藉着话语而生的；但如果没有——事实正是如此——万物倒是藉着祂而有，就是圣父所愿意的一切，这岂不暴露了这些人多头的诡诈吗？他们羞于说\Quote{工程}、\Quote{受造物}，以及\Quote{天主的圣言在祂受生之前并不存在}，却以另一种方式，他们主张祂是受造物，提出\OriginalTerm{意志}{will}的说法，并说：\Quote{除非祂是藉意志而有，天主便有了一个圣子，出于必然，且违反祂的善意。}那么，是谁将必然性强加给祂呢？啊，最邪恶的人哪，你们把一切都拽向你们异端的目的！因为他们看到的是违反意志的事物；但那超越它、比它更大的，却逃过了他们的理解。因为凡是出乎定意的，是违反意志的；而按本性的，则超越了计议，且先于计议。人藉着\OriginalTerm{计议}{counsel}建造房屋，但按\OriginalTerm{本性}{nature}生育儿子；建造的东西，是按意志开始存在的，且外在于建造者；但儿子是父亲\OriginalTerm{性体}{essence}的本有亲生，并不外在于他；因此，父亲也不为儿子而计议，免得他像是替自己计议。那么，正如圣子超越受造物，按本性的事物同样超越意志。他们在论及祂时，本不该用意志去衡量那按本性的事物；然而，他们忘记了他们是在谈论天主的圣子，竟敢将人间的矛盾——即\Quote{必然}与\Quote{出乎定意}——用在天主身上，以便能由此否认有真实的圣子。让他们自己告诉我们——天主是善的、仁慈的，这是藉着意志属于祂，还是并非如此呢？如果是藉着意志，我们就必须认为祂开始成为善的，并且祂可能不是善的；因为计议与选择意味着两种倾向的可能，且是理性本性的属性。但若说天主是善的、仁慈的，是基于意志，那太不成体统，那么他们自己的话必须回敬他们——\Quote{所以，是出于必然，而非按祂的善意，祂才是善的}；并且，\Quote{是谁将这必然性强加给祂呢？}但如果谈论天主意谓着\Quote{必然}是不成体统的，因而祂是按本性而为善，那么，就更加确凿而真实的是，祂是按本性，而非按意志，成为圣子的父。\par

63. 此外，让他们回答我们这一点：——（为了针对他们的无耻，我想提出一个更进一步的问题，这问题虽然大胆，却怀着虔诚的心意；主啊，求祢垂怜！）——圣父自己，祂是先有了计议，然后乐意，抑或先于计议而存在呢？因为既然他们对圣言如此放肆，他们也必得着同样的答复，好叫他们知道，他们的这种狂妄甚至上达圣父自己。那么，若他们自己就意志来计议，并说天主也是出于意志，那么，祂在计议之前是什么呢？或者，如你们所想的，在计议之后，祂获得了什么呢？但如果这样的问题是不成体统的、自我毁坏的，甚至提出来都令人震惊（因为只要听到天主之名，就足以让我们知道并理解祂是\OriginalTerm{自有者}{He that Is}），那么，对天主的圣言怀有类似的想法，并编造出意志与善意之说，岂不也违反理性吗？因为同样地，只要听到圣言之名，就足以知道并理解：那并非藉意志而是天主的那一位，并非藉意志，而是按本性，拥有祂自己的圣言。并且，只要存想一下，天主自己在计议、思考、选择并怀有善意，为的是使祂不致没有圣言和没有智慧，而是拥有二者——这岂不已超出了一切可以想象的疯狂吗？因为，那为属于祂性体的事物而计议的，显得是在为自己而盘算。那么，既然这种想法中含有莫大的亵渎，虔诚的说法将是：万有是由\Quote{恩惠和意志}而有的，但圣子却不是意志的产物，也不是像受造物那样后来才有的，而是按本性，作为天主性体的亲生后代。因为，作为圣父自己的圣言，祂不容许我们将意志视为先于祂自己，因为祂自己就是圣父生活的计议与能力，以及对于圣父所视为善的事物的制造者。而这就是祂在\WorkTitle{箴言}中论及自己时所说的：\BibleQuote{计议属我，保障属我；智慧属我，能力属我}\BibleRef{箴 8:14}。因为，尽管祂自己就是\Quote{智慧}——祂在其中预备了诸天，祂自己也是\Quote{能力与力量}（因为基督是\BibleQuote{天主的能力与天主的智慧}\BibleRef{格前 1:24}），祂在这里却变换了说法，说\BibleQuote{智慧属我}和\BibleQuote{能力属我}；同样，当祂说\BibleQuote{计议属我}时，祂自己必是圣父生活的计议；正如我们也从先知那里得知，祂成为\BibleQuote{伟大计议的使者}\BibleRef{依 9:6}，并被称作圣父的善意；因此，我们必须用有关天主的人间比喻来驳斥他们。\par

64. 所以，如果受造之物都是\Quote{凭旨意和恩惠}而存在，整个受造界\Quote{照天主的乐意}而被造，\Person{保禄}蒙召为宗徒是\Quote{凭天主的旨意}，我们的召叫也是\Quote{凭祂的乐意和旨意}而来，且万物都是借圣言而有的，那么圣言就外于那些凭旨意而有的受造物了；相反，祂本身就是圣父的\OriginalTerm{永活的谋略}{Living Counsel}，万物都凭这谋略而得以存在；\Person{达味}也在第七十二篇圣咏中为此感谢，\BibleQuote{你握着我的右手，以你的规劝引导我。}\BibleRef{咏 73:24}那么，圣言既是圣父的谋略和乐意，祂自己怎能像其他受造物一样，\Quote{凭乐意和旨意}而存在呢？除非——如我之前所说——他们狂乱地再推说，祂是凭自己或凭别的什么而存在。那么祂是凭谁而存在的呢？让他们另造一个圣言吧；让他们再命名另一位基督吧，这无异于\Person{瓦伦廷}的学说；因为这绝非出自圣经。就算他们再另造一个来，那一个也必定是凭某个他者而存在；这样一来，当我们层层推算、追问这连串的源头时，便显明这无神派的多头异端，最终导致多神论和无边的疯狂；在这种异端中，他们既想使圣子成为受造物、从虚无而来，就另说一套，用\Quote{旨意}和\Quote{乐意}这些字眼来暗示同样的意思，而这些字眼本应归于起初之物和受造之物。那么，将受造之物的特性归给万物的创造者，岂不是不虔？硬说在圣言之先，圣父内已有旨意，岂不是亵渎？因为如果圣父内先有旨意，圣子所说的\BibleQuote{我在父内}\BibleRef{若 14:20}这句话就不真实了；或者即使祂在父内，也只得屈居次位，并且祂不应当说\BibleQuote{我在父内}，因为照你们的看法，旨意是在祂之先，万物都是在这旨意内被造的，祂自己也是在这旨意内存在的。纵使祂荣耀超绝，也仍不过是凭旨意而有的受造物之一。而且，如前所说，如果真是这样，祂怎能是主，而它们是仆役呢？但祂是万有的主，因为祂与圣父的主权原为一体；而受造界则都处于奴役之下，因为它们外于圣父的一体性，且它们原不存在，是被造出来的。\par

65. 而且，如果他们说圣子是凭旨意而有的，那他们也应该说祂是凭理智而有的；因为我认为理智和旨意本是同一回事。因为人谋划什么，他对此也就有理智；他理智什么，他也就谋划什么。的确，救主自己也显明这两者是相称而同源的，因为祂说：\BibleQuote{谋略和稳妥属我；理智和力量属我。}\BibleRef{箴 8:14}因为力量和稳妥既是一回事（因为它们指的是同一个属性），那么我们也可以说，理智和谋略也是一回事，而这正是主。可是这些不虔的人不愿意圣子成为圣言和永活的谋略；他们倒编造说，在天主那里有一个仿佛是习性的东西，像凡人一样，时而出现，时而消失，那就是理智、谋略、智慧；他们挖空心思，提出\Person{瓦伦廷}的\Quote{思想}和\Quote{意志}，好叫他们能把圣子与圣父分开，称祂为受造物，而不是圣父本有的圣言。那么，对他们就该说那曾对术士\Person{西满}说过的话：\Quote{愿\Person{瓦伦廷}的不虔同你们一起丧亡！}\BibleRef{宗 8:20}；让人人宁愿信服\Person{撒罗满}吧，他说圣言就是智慧和明智。因为他写道：\BibleQuote{上主以智慧奠定了大地，以明智巩固了诸天。}\BibleRef{箴 3:19}并且，像这里说凭明智，在\BibleBook{}{圣咏}中却又说\BibleQuote{因上主的一句话，诸天造成}\BibleRef{咏 33:6}。而又如诸天由圣言创造，同样\BibleQuote{凡祂所喜爱的，样样都成功了}\BibleRef{咏 115:3}。又像宗徒致书\Place{得撒洛尼}人时所说的：\BibleQuote{天主的旨意就在基督耶稣内}\BibleRef{得前 5:18}。所以天主圣子，祂就是\Quote{圣言}和\Quote{智慧}；祂就是\Quote{明智}和永活的\Quote{谋略}；在祂内有\Quote{圣父的乐意}。祂是圣父的\Quote{真理}、\Quote{光明}和\Quote{能力}。但如果天主的旨意就是智慧和明智，而圣子就是智慧，那么说圣子\Quote{凭旨意}而来的人，实际上就是说，智慧是在智慧中产生的，圣子是在一个子内被造，圣言是借着圣言受造；这既与天主的本性不相称，也与圣经相悖逆。因为宗徒宣告圣子乃是父的光辉和印记，不是出于父的旨意，而是出于父的性体本身，说：\BibleQuote{他是天主光荣的辉耀，是天主本体的真像}\BibleRef{希 1:3}。可是，如果如我们之前所说，圣父的性体和位格并不来自旨意，那么很明显，属于圣父位格的本有之物也不是来自旨意；因为那位蒙福的位格是怎样的，以怎样的方式存在，那由祂而出的本有之子也就必然怎样。所以圣父自己并没有说：\Quote{这是我凭我的旨意所生的子}，也没有说：\Quote{这是出于我的恩惠而有的子}，祂只单单说：\Quote{我的子}，并进一步说：\Quote{我所喜悦的}；祂的意思是：这是本质意义上的子；\Quote{凡我所喜悦的一切，都寄托在祂内}。\par

66. 既然圣子出自本性而非出自意志，祂是否就没有圣父的悦乐，也不符合圣父的意志呢？不，绝非如此；圣子正是与圣父的悦乐同在，并且正如祂自己所说：\BibleQuote{父爱子，并将一切显示给祂。}\BibleRef{若 5:20} 因为正如祂并非\Quote{出于意志}才开始是善，然而没有意志和悦乐也并非善（因为祂之所是，便是祂的悦乐），同样，圣子的存在，虽非\Quote{出于意志}，却也不缺乏祂的悦乐，也不违背祂的旨意。因为正如祂自己的\OriginalTerm{存在体}{Subsistence}是出于祂的悦乐，同样，圣子既属于祂的\OriginalTerm{本质}{Essence}，便不缺祂的悦乐。因此，愿圣子成为圣父悦乐和爱的对象；这样，让每一个人虔诚地理解天主的悦乐和心甘情愿。因为正是借着圣子成为圣父悦乐之对象的那善意，圣父也成了圣子爱、悦乐和尊崇的对象；而这善意是一，从圣父而在圣子内，因此在这里我们也可瞻仰圣子在圣父内、圣父在圣子内。因此，谁也不要像\Person{瓦伦廷}那样引入一个在先的意志；谁也不要借着\Quote{谋略}的借口，在独一的圣父与独一的圣言之间进行干预；因为在它们之间安置意志和考虑乃是疯狂。因为说\Quote{祂出于意志而存在}是一回事，而说圣父对祂那本性属己的圣子怀有爱和善意是另一回事。因为说\Quote{祂出于意志而存在}，首先意味着祂曾一度不存在；其次，如前所述，这意味着一种双向的倾向，以致有人可以设想圣父甚至可能不意愿圣子。但对圣子说\Quote{祂可能不曾存在}，是一种不虔的臆断，甚至触及圣父的本质，仿佛那属祂自己的可能不曾存在。因为这等同于说\Quote{圣父可能不曾是善的}。正如圣父总是本性为善，祂也总是本性为父；而说\Quote{圣父的善意即是圣子}，以及\Quote{圣言的善意即是圣父}，并非暗示一个在先的意志，而是本性的纯真、所有权的归属与本质的相似。因为正如就光辉和光而言，人们可以说，在光中光辉并不先于意志，但它却是光的自然产物，是出于生它的光的悦乐，并非借着意志和考虑，而是借着本性和真理；同样，就圣父和圣子而言，人们也可以正确地说，圣父对圣子怀有爱和善意，圣子对圣父怀有爱和善意。\par

67. 因此，不要称圣子为善意的工程；也不要把\Person{瓦伦廷}的学说引入教会；却要承认祂是活生生的旨意，是真理与本性中的\OriginalTerm{后裔}{Offspring}，如同出自光的光辉。因为圣父这样说过：\BibleQuote{我心涌溢美善的圣言。}\BibleRef{咏 45:1} 而圣子也相应地说：\BibleQuote{我在父内，父在我内。}\BibleRef{若 14:10} 但如果圣言在心内，意志何在？如果圣子在父内，善意何在？如果祂就是意志本身，意志内又怎会有谋略？这是不适宜的；免得圣言在言语中产生，圣子在儿子中产生，智慧在智慧中产生，正如一再所说。因为圣子是圣父的全部；在圣言之前，圣父内绝无一物；但在圣言内有意志，并且借着祂，意愿的对象得以成就，正如圣经所显示的。我但愿那些不虔的人，既已陷入如此缺乏理智的境地，竟要考虑意志，现在不再去询问他们常去询问的那些生儿育女的女人：\Quote{你还没有怀孕前，就有了儿子吗？}，而是去问父亲：\Quote{你们成为父亲，是由于谋略，还是由于你们意志的自然法则？}或\Quote{你们的儿女，是否相似你们的本性和本质？}，好叫他们即使从父亲那里也能学到羞耻，因为他们正是从父亲那里借用了关于生育的命题，也希望从父亲那里得到相关的知识。因为他们会回答：\Quote{我们所生的，不是相似我们的善意，而是相似我们自己；我们成为父母，也不是由于预先的谋略，生育是我们的本性所固有的；因为我们也是我们父亲的肖像。}那么，或者让他们谴责自己，停止向女人询问天主的圣子，或者让他们向女人学习，圣子并非由意志所生，而是在本性和真理中所生。对他们而言，来自人间实例的反驳是适宜且恰当的，因为这些心思悖逆的人以人间的方式争论天主性的事。那么，\Person{基督}的仇敌为什么仍然疯狂呢？因为这个，连同他们其他的借口，都被显明并证实为纯粹的幻想和寓言；因此，他们应该，即使迟了，也要凝视他们跌落其中的愚妄悬崖，从深处重新站起，逃避魔鬼的罗网，正如我们所劝诫的。因为真理是爱人的，并不断地呼喊：\Quote{如果因我取了肉身的衣服你们不信我，也当信这些事，好叫你们知道：\BibleQuote{我在父内，父在我内}\BibleRef{若 14:10}，\BibleQuote{我与父原是一体}\BibleRef{若 10:30}，\BibleQuote{看见我的，就是看见了父。}\BibleRef{若 14:9}} 但是主照祂的常情是爱人的，并乐意\BibleQuote{扶助跌倒的人}\BibleRef{咏 145:14}，正如\Person{达味}的赞美诗所说；然而那些不虔的人，不愿听从主的声音，也不忍见到祂被众人承认为天主及天主子，可怜的人，如同甲虫一般，同他们的父魔鬼四处奔走，为不虔寻找借口。那么，他们还能从何处找到什么借口呢？除非他们借取犹太人和\Person{盖法}的亵渎，并采纳外邦人的无神论？因为圣经对他们封闭了，从圣经的每一部分，他们都被驳斥为无知的人及基督的仇敌。\par

\SourceRecord{Translated by John Henry Newman and Archibald Robertson.；Nicene and Post-Nicene Fathers, Second Series；Vol. 4.；Edited by Philip Schaff and Henry Wace.；Buffalo, NY: Christian Literature Publishing Co.,；1892.；<http://www.newadvent.org/fathers/28163.htm>.}

% structure: p0001 source=h1 target=section reason=page-title

\section{\WorkTitle{驳亚略派第四篇讲论}}

% structure: p0002 source=h2 target=subsection reason=relative-heading-rank; base=h2

\subsection{1-5. 自圣经证圣言的实体性。若那唯一元始是实体的，则其圣言是实体的。除非圣言与圣子为另一元始，或为一工程，或为一属性（因而天主成为复合体），或同时即是圣父，或意味着天主内存在另一本性，否则祂便出自圣父的本质且与圣父有别。按\BibleBook{申命}{纪} 4:4阐释\BibleBook{若望}{福音} 10:30之例证。}

1. 圣言是由天主而来的天主；因为\BibleQuote{圣言就是天主}\BibleRef{若 1:1}，又说\BibleQuote{圣祖是他们祖先，并且基督按血统说也是从他们来的，他是在万有之上，世世代代应受赞颂的天主！阿们}\BibleRef{罗 9:5}。既然基督是由天主而来的天主，且是天主的圣言、智慧、圣子与德能，因此圣经中只宣示一位天主。因为圣言是那独一天主的圣子，祂与那使祂存在者相联；所以圣父与圣子是两位，但天主性的\OriginalTerm{一元}{Monad}却是不可分割、不可分离的。这样，我们也就保存了天主性只有一个\OriginalTerm{元始}{Beginning}，而非两个元始，由此严格说来便是\OriginalTerm{唯一神权}{Monarchy}。而这元始的圣言，按本性即是圣子，并非如同另一个自立自存的元始，也不是在那元始之外生出，以免因这种差异而产生\OriginalTerm{二神权}{Dyarchy}与\OriginalTerm{多神权}{Polyarchy}；相反，祂是那唯一元始的亲生子，自有智慧，自有圣言，从祂而存在。因为按若望所说，\BibleQuote{在元始中就有圣言，圣言与天主同在}，因为元始就是天主；既然祂出自元始，所以\BibleQuote{圣言就是天主}。既然只有一个元始，因而只有一位天主，同样，那真实、真正、确实的\OriginalTerm{本质}{Essence}与\OriginalTerm{自立体}{Subsistence}也只有一位，即说\BibleQuote{我是自有者}\BibleRef{出 3:14}的那位，免得有两个元始；从这唯一者，按本性与真理而生的圣子，就是祂的独有圣言、独有智慧、独有德能，且与祂不可分离。正如没有另一个本质，以免有两个元始，同样，出自那唯一本质的圣言也不会消解，也不是有意义的声息，而是一位本质的圣言、本质的智慧，亦即真正的圣子。因为如果祂不是本质的，天主就会\BibleQuote{向虚空说话}\BibleRef{格前 14:9}，并具有身体，与人毫无区别；但祂既不是人，故祂的圣言也不随从人的软弱。因为正如元始是一个本质，同样，祂的圣言是一个本质的、自立的，祂的智慧亦然。因为正如祂是由天主而来的天主，由智慧者而来的智慧，由理性者而来的圣言，由圣父而来的圣子，同样，祂是由自立体而来的自立者，由本质而来的本质与实体，由存在而来的存在。\par

2. 因为如果祂不是本质的智慧、实体的圣言，且作为圣子实际存在，而仅是圣父内的智慧、圣言与子，那么圣父本身的性体便由智慧与圣言组合而成。若是这样，就会产生前述的荒谬；祂将是自身的父，圣子由自己生也为自己所生；或者圣言、智慧、圣子只是个名号，而那拥有或更好说是这些名号者并不自立存在。如果祂并不自立存在，这些名号便是虚空而无实质的，除非我们说天主就是\OriginalTerm{至智慧者}{Very Wisdom}、\OriginalTerm{至圣言者}{Very Word}。但若是这样，祂既是自身的父也是子；当祂智慧时是父，当祂是智慧时是子；但这些事并非作为某种属性存在于天主内；应摒弃这种不敬的思想；因为它将引致这样的结论：天主由本质与属性组合而成。因为一切属性都寓于本质，显然，那不可分割的神性\OriginalTerm{一元}{Monad}，若被分为本质与\OriginalTerm{依附体}{accident}，就必然是复合的。因此我们必须质问那些执拗的人：圣子被宣示为天主的智慧与圣言；那么祂是如何如此呢？若是作为一种属性，其荒谬已如前述；但若天主就是那至智慧者，那就落入\Person{撒伯里乌}的谬论；因此，祂是真正意义上的出自圣父本身的子，犹如光的比喻。因为正如光生于火，同样，圣言生于天主，智慧生于智慧者，圣子生于圣父。这样，一元保持着不可分割与完整，其子是圣言，并非\Emph{无本质}，亦非\Emph{不自立}，而是真真实实的本质的。若不是这样，所论的一切都仅成为意念与言语之说。但我们必须避免这种荒谬，那么真圣言便是本质的。因为正如确有真实的圣父，同样也确有真实的智慧。就这方面而言，他们是两位；不是因为像\Person{撒伯里乌}所说，父与子是同一，而是因为父是父，子是子，他们是一体，因为子按本性是出自圣父本质的圣子，作为父的自身圣言而存在。这就是主所说的：\BibleQuote{我与父原是一体}\BibleRef{若 10:30}；因为圣言既非与父分离，父也从未、亦非现在无圣言；为此祂说：\BibleQuote{我在父内，父在我内}。\par

3. 再者，基督是天主的圣言。那么，祂是凭自己存有，并在存有之后与圣父结合，还是天主造了祂或称之为自己的圣言？如果是前者，我的意思是如果祂是凭自己存有且是天主，那么就有两个\OriginalTerm{开端}{Beginnings}；而且很明显，祂不是圣父自己的，因为祂不是出于圣父，而是出于自己。但如果相反，祂是从外部被造的，那么祂就是受造物。因此只能说是从天主自己而来；但如果是这样，那么从另一位而来的是一件事物，而那来源又是另一件事物；按照此说，就是两个。但如果它们不是两个，而是同一者有不同的称呼，那么原因与结果就会相同，生者与受生者也会相同，这在\Person{撒伯流}的例子中已显为荒谬。但如果祂出于祂，却又不是另一位，那么祂将同时是生者和非生者；生，因为祂从自己产生；非生，因为祂不是自己以外的任何事物。但如果是这样，同一个就在概念上被称为圣父与圣子。然而，如果说这不合宜，那么圣父与圣子必须是两个；但他们是一，因为圣子不是从外面来，而是由天主所生。但是，如果任何人不敢说\OriginalTerm{子}{Offspring}，而只说圣言与天主同在，那么他要当心，因为不敢说圣经所言，反会陷入荒谬，使天主成为具有双重本性的存在。因为不承认圣言来自\OriginalTerm{单体}{Monad}，却简单地以为祂是与圣父结合的，他就引入了双重本质，而两者都不是对方的父。对于\OriginalTerm{大能}{Power}也是如此。如果我们参照圣父来思考，这一点就更清楚了；因为圣父只有一位，不是两位，而圣子出于那一位。因此正如没有两位圣父，而是一位，照样也没有两个\OriginalTerm{开端}{Beginnings}，而是一个，而圣子是由那一位\OriginalTerm{本质地}{essential}所生。\par

4. 相反，我们必须问\OriginalTerm{亚略派}{Arians}：（因为\OriginalTerm{撒伯流派}{Sabellianisers}必须从圣子的概念去反驳，而\OriginalTerm{亚略派}{Arians}必须从圣父的概念去反驳：）那么，让我们说——天主是智慧的，而非无圣言的：还是反过来，祂是无智慧的，也是无圣言的？如果是后者，立刻就有荒谬；如果是前者，我们必须问，祂如何是智慧的，却又不是无圣言的？祂是从外部拥有圣言和智慧，还是从自己？如果从外部，那么必定有谁先将它们赐给祂，而在领受之前，祂是无智慧、无圣言的。但如果从自己，那么显然圣言并非从虚无而来，也不是曾经不存在；因为圣言永远存在；因为圣言是祂的\OriginalTerm{肖像}{Image}，而祂永远存在。但如果他们说，天主确实是智慧的，不是无圣言的，但祂在自己内有自己的智慧和自己的圣言，而这并非基督，而是用来造了基督的那个，我们必须回答说，如果基督是在那个圣言内被造成的，那么显然万物也都是如此；那一定是\Person{若望}所说的那位：\BibleQuote{万物是藉着祂而造成的}，以及圣咏作者所说的：\BibleQuote{你以智慧创造了万物}。而基督的话\BibleQuote{我在圣父内}就会被发现是不真实的，因为有另一位在圣父内。而\BibleQuote{圣言成了血肉}\BibleRef{若 1:14}按照他们也就不是真实的。因为，如果那位\BibleQuote{万物是藉着祂而造成的}，祂自己成了血肉，但基督作为圣言不在圣父内，那么基督就没有成为血肉，而或许基督只是被称为圣言。但若是这样，首先，在名称之外还有另一位，其次，万物就不是藉着祂造成的，而是在那另一位内造成的，基督也在那一位内被造。但如果他们说，智慧是圣父内的一种属性，或说圣父就是\OriginalTerm{真智慧}{Very Wisdom}，那么前面提到的荒谬就会接踵而至。因为祂就成了\OriginalTerm{复合的}{compound}，并且自证为自己的圣子和圣父。此外，我们必须基于以下理由驳斥并使他们无言：在天主内的圣言不可能是受造物，也不可能从虚无而来；但如果圣言在天主内，那么祂必定是基督，祂说：\BibleQuote{我在圣父内，圣父也在我内}\BibleRef{若 14:20}，因此祂也就是\OriginalTerm{独生子}{Only-begotten}，因为没有其他从祂所生者。这就是独一的圣子，是圣言、\OriginalTerm{智慧}{Wisdom}、\OriginalTerm{大能}{Power}；因为天主并非由这些组合而成，而是生发它们。因为正如祂藉着圣言形成受造物，照样，按祂自己\OriginalTerm{本质}{Essence}的本性，祂有圣言作为\OriginalTerm{子}{Offspring}，并藉着祂形成、创造、安排万物。因为藉着圣言和\OriginalTerm{智慧}{Wisdom}，万物得以存在，并按照祂的定命一同存立。对于\Emph{圣子}这个称呼也是如此；如果天主没有圣子，祂就没有工；因为圣子是祂的\OriginalTerm{子}{Offspring}，藉着祂而工作；但若非如此，同样的问题和同样的荒谬将随他们的胆大而来。\par

5. \WorkTitle{申命记}记载：\BibleQuote{你们凡是依附上主你们天主的，今日都还活着}\BibleRef{申 4:4}。由此我们可以看出区别，并知道天主子不是受造物。因为圣子说：\BibleQuote{我与父原是一体}，也说：\BibleQuote{我在圣父内，圣父也在我内}；但受造物，当它们进步时，是依附于上主的。因此，圣言是在圣父内，因为祂是圣父自己的；然而受造物，由于是外在的，便依附于祂，因为它们在本性上是陌生的，且是藉着自由抉择而依附。因为按本性而生的儿子，与生他者是一体；但那从外面来、被立为儿子的，则依附于家族。因此圣子立刻接着说：\BibleQuote{哪有这样伟大的民族，天主竟这样亲近他们？}别处又说：\BibleQuote{我是近处的天主}；因为对受造物，祂亲近，因为它们对祂是陌生的；但对圣子，祂不仅是亲近，而是在祂内。而且圣子不是依附于圣父，而是与祂同在；因此同一部\WorkTitle{申命记}中，\Person{梅瑟}又说：\BibleQuote{你们应听从祂的话，且依附祂}\BibleRef{申 13:4}；然而凡是依附的，是从外面依附。\par

% structure: p0008 source=h2 target=subsection reason=relative-heading-rank; base=h2

\subsection{6, 7. 当圣言与圣子饥饿、哭泣、疲惫时，祂作为我们的中保，承担我们所有，为将祂所有分赐于我们。}

6. 然而，为答复亚略派软弱而属人的观念，即他们认为主有所匮乏，因为祂说\BibleQuote{已赐给我了}，和\BibleQuote{我接受了}，以及\Person{保禄}说\BibleQuote{因此天主举扬了祂}，和\BibleQuote{使祂坐在自己的右边}等话，我们必须指出：我们的主，既是圣言和圣子，穿上了肉躯，成了人子，好成为天主与人之间的中保，将天主的事交付我们，并将我们的事献于天主。当祂被说成饥饿、哭泣和疲倦，并呼喊\BibleQuote{厄里，厄里}时，这些原是我们人性的感受，祂从我们接受，并献于圣父，为我们转求，好使它们在祂内得以消除。当说\BibleQuote{天上地下的一切权柄都赐给了我}和\BibleQuote{我接受了}以及\BibleQuote{因此天主举扬了祂}时，这些恩赐是经由祂而由天主赐给我们的。因为圣言从未匮乏，也未有过起头；而人也没有能力为自己获得这些事，而是经由圣言赐给了我们；所以，这些恩赐犹如赐给了祂，实际是分施给我们。这正是祂\OriginalTerm{降生成人}{Incarnation}的理由：好使恩赐犹如赐给了祂，便能传递给我们。因为这样的恩赐，单纯的人是不配接受的；同样，单纯的圣言也不需要它们；于是圣言与我们结合，然后分施给我们能力，并举扬我们。因为圣言居于人内，就举扬了人本身；当圣言在人内时，人本身就领受了。那么，既然圣言在肉身内，人本身被举扬并领受了能力，这些事便归于圣言，因为它们是为了祂的缘故而赐下的；因为正是由于在人内的圣言，这些恩赐才得以赐予。正如\BibleQuote{圣言成了血肉}\BibleRef{若 1:14}，同样，人本身也领受了经由圣言而来的恩赐。因为人所领受的一切，都说成是圣言领受了；为要显示，人本身就其本性而言，本不堪当领受，却因圣言成了血肉而领受了。因此，若有什么说是赐给主的，或类似的，我们必须理解为，恩赐并非赐给需要它的祂，而是经由圣言赐给人本身。因为凡为他人转求的，是以本人的身份领受恩赐，并非自己有需要，而是为了他所转求的人。\par

7. 因为正如祂承受我们的软弱，自己并非软弱；祂饥饿，却非为自己饥饿，而是献上我们所有的，好使它们被消除；同样，从天主来的恩赐，取代我们的软弱，祂自己也确实领受了，为使人与祂结合后，能分享这些恩赐。因此主说：\BibleQuote{凡祢交给我的，我都交给了他们}，又说：\BibleQuote{我为他们祈求}\BibleRef{若 17:7}。祂为我们祈求，承担了我们所有，而祂也将所领受的分施。既然圣言与人本身结合，圣父垂顾祂，惠赐人得以被举扬，拥有一切权能等；因此，凡我们经由祂所领受的一切，都归于圣言自身，并好像赐给了祂。因为正如祂为了我们降生成人，我们也为了祂而被举扬。那么，如果祂为了我们自谦自卑，同样为了我们而被说成受举扬，这并非荒谬。所以\BibleQuote{祂赐给了祂}，意思就是\BibleQuote{为了祂赐给了我们}；\BibleQuote{祂举扬了祂}\BibleRef{斐 2:9}，意思就是\BibleQuote{在祂内举扬了我们}。而圣言自身，当我们被举扬、领受、获得扶助时，犹如祂自己被举扬、领受和获得扶助一样，便向圣父称谢，将我们所有的归于自己，并说：\BibleQuote{凡祢交给我的，我都交给了他们}\BibleRef{若 17:7}。\par

% structure: p0011 source=h2 target=subsection reason=relative-heading-rank; base=h2

\subsection{八、亚略派将圣子的起头定得比\Person{马赛路斯}更早，等等。}

8. \Person{欧瑟比乌}及其同党，即亚略派狂人，赋予圣子一个存在的起头，却佯称不愿祂有一个王权的起头。但这实属可笑；因为凡将存在的起头归给圣子的，也极其显然地将为王的起头归给了祂；他们如此盲目，竟承认自己所否认的。再者，那些声称圣子仅是一个名号，说天主子，即圣父的圣言，是无\OriginalTerm{本质}{unessential}且非\OriginalTerm{自立}{non-subsistent}的，却假装对说\BibleQuote{曾有一时祂不存在}的人发怒。这同样可笑；因为那些完全不赋予祂存在的人，反倒对那些至少承认祂存在于时间内的人发怒。如此一来，这些人在指责别人时，也承认了自己所否认的。再者，\Person{欧瑟比乌}及其同党，承认圣子，却否认祂就本性而言是圣言，并要将圣子称为名目上的圣言；而另一派人承认祂是圣言，却否认祂是圣子，并要将圣言称为名目上的圣子；同样都是毫无根基。\par

% structure: p0013 source=h2 target=subsection reason=relative-heading-rank; base=h2

\subsection{9-10. 除非圣父与圣子只是名目上的两个，或是作为部分而各自不完全，或是两位神，否则他们是\OriginalTerm{同性同体}{coessential}，在天主性内为一，而圣子出于圣父。}

9. \BibleQuote{我与父原是一体}\BibleRef{若 10:30}。你说这两者是一，或者说那一位有两个名号，又或者说那位一被分裂为二。如果那位一被分裂为二，那么被分裂的必然是一物体，而每一部分都不是完美的，因为每一部分只是部分，不是整体。但如果那位一有两个名号，这便是\Person{撒贝里}的策略，他说圣子与圣父相同，且在有一位时便取消了另一位：有圣子时无圣父，有圣父时无圣子。但如果这两者为一，那么必然有二，但就天主性而言，按照圣子与圣父的\OriginalTerm{同性同体}{coessentiality}，以及圣言出于圣父自身而言，他们是二，因为有一位圣父，和圣子，即圣言；而又是一，因为只有一个天主。若非如此，祂就会说：\BibleQuote{我是父，}或\BibleQuote{我与父是}；但事实上，在\BibleQuote{我}字中祂指圣子，在\BibleQuote{与父}中则指生祂的那位；在\BibleQuote{一体}中则指同一个天主性和祂的同性同体。因为并非如外邦人所主张，那同一位是智慧与智慧，或圣父与圣言相同；因为祂作自己的圣父是不合理的，然而属神的教训认识圣父和圣子，智慧和智慧，天主和圣言；同时始终保守祂在各方面不可分割、不可分离、不可分解。\par

10. 但若有人，一听说圣父与圣子是两位，便错误地指责我们宣讲两位天主（因为有些人如此臆想，并立刻讥讽说：\Quote{你们持有两位天主。}），我们对这样的人必须回答：若承认圣父和圣子就是持有两位天主，那么紧接着若只承认一位天主，我们就必须否认圣子并陷入撒伯流主义。因为如果说两位便是陷入异教主义，那么说一位，我们必然陷入撒伯流主义。但这并非如此；绝无此理！但正如我们说圣父和圣子是两位时，仍承认一位天主，同样当我们说只有一位天主时，让我们思想圣父和圣子是两位，而他们在天主性内为一，且在圣父的圣言与祂不可分离、不可分割、不可分开方面为一。且让火及其光芒作为比喻，它们在实体与表象上是二，但因光芒不可分割地源于火，故为一。\par

% structure: p0016 source=h2 target=subsection reason=relative-heading-rank; base=h2

\subsection{11，12. 玛尔凯路斯及其门徒，与亚略派一样，声称圣言并非受造，而是被发出，为要创造我们，仿佛天主的静默是一种无所作为的状态，而当天主藉圣言说话时，祂才行动；或说圣言有出发与回归；这种道理暗示圣父与圣子中有变化与不完美。}

11. 他们陷入了与亚略派同样的愚昧；因为亚略派也声称祂是受造为了我们，好让祂创造我们，仿佛天主等待我们的受造，祂才能发出——如一派所言——或受造——如另一派所言。于是，亚略派对我们比对圣子更为慷慨；因为他们说，我们不是为了祂，而是祂为了我们而存在；就是说，如果祂因此受造并存在，为让天主藉着祂创造我们。而这些人，同样不虔诚甚至更甚，给予天主的少于给予我们的。因为我们常常，即使在静默时，也活跃地思想，以至将我们思想的成果形成概念；但他们想要天主在静默时一无所为，当祂说话时才施展力量；即是说，如果静默时祂不能创造，说话时才开始了创造。我们理当问他们，当圣言在圣父内时，祂是否已是完美的，足以能够创造？如果一方面，祂在圣父内时是不完美的，但藉着受生而成为完美的，那么我们就是祂完美的原因，也就是说，如果祂是为了我们而受生；因为我们，祂才领受了创造的能力。但若祂在圣父内已是完美的，足以创造，那么祂的受生便是多余的；因为即使尚在圣父内时，祂也能构造世界；所以，要么祂未曾受生，要么祂受生，并非为了我们，而是因为祂永远出自圣父。因为祂的受生见证的不是我们被造，而是祂出自天主；因为在我们的受造以先，祂就已存在。\par

12. 针对圣父，也能对他们证明同样的假设；因为如果静默时，祂不能创造，那么祂必然通过受生，即说话，获得了能力。祂从哪里获得这能力？又为何？如果在祂内拥有圣言时，祂便能创造，那么祂受生便是多余的，因为即使在静默中也能创造。其次，如果圣言在祂受生以前已在圣父内，那么受生之后，祂便在圣父之外、在圣父外部了。但若如此，祂如今怎么说：\BibleQuote{我在圣父内，圣父在我内}\BibleRef{若 14:10}？但若祂如今在圣父内，那么祂就永远在圣父内，正如现在一样，那么说\Quote{祂为了我们而受生，并在我们被造后回归，好恢复祂原来的状态}就是多余的了。因为祂并非现今不是祂之前所是，也不是现在不是祂将来所是；而是祂过去如何，如今也如何，且是在同一状态和同一方面；否则祂就会显得不完美和可变。因为如果祂之前是什么，后来还要成为什么，仿佛现在不是如此，那么显然现在祂不是祂过去所是、将来所是。我的意思是，如果祂以前在圣父内，后来还要再次在圣父内，那么现在圣言不在圣父内。但主斥责这样的人，祂说：\BibleQuote{我在圣父内，圣父在我内}；因为祂现在如何，便永远如何。但若祂现在如何，便永远如何，那么结论就不是祂在某时受生，另一时不受生，也不是天主时而有静默，然后说话，而是永远有一位圣父，和一位圣子，就是祂的圣言，不只是在名义上是圣言，也不只是概念中的子，而是存在着的、与圣父同实体的，不是为我们而受生，因为我们是为祂而受造。因为，如果祂为我们而受生，在祂的受生中我们被造，在祂的受生中受造物成立，然后祂回归以恢复原来的状态，那么首先，那受生者将再次成为非受生。因为如果祂的出来是受生，那么祂的回归将是那受生的终结，因为当祂再次在圣父内时，圣父将再次沉默。但若圣父沉默，那么将有静默时的一切，静止而无创造，因为受造物将停止存在。因为正如圣言出来时，受造物形成并存在，同样圣言退隐时，受造物将不复存在。那么它成为存在又有什么用，如果它要停止？或者天主为什么说话，然后又沉默？为什么祂发出一位却又召回？为什么祂受生了一位，却又让祂的受生停止？再次，祂将来成为什么也不确定。因为要么祂将永远沉默，要么祂将再次受生，并设计一个不同的受造界（因为祂不会造同样的，否则已造成的将存留，而是另一个）；到了时候祂也会终结那另一个，并设计又一个，如此无穷无尽。\par

% structure: p0019 source=h2 target=subsection reason=relative-heading-rank; base=h2

\subsection{13，14. 这样的道理排斥了在天主性中一切真正位格的区分。从\BibleBook{格林多}{后书}\BibleRef{格后 6:11}等处阐示圣经的道理。}

13. 这或许是他从斯多噶学派借来的，他们主张他们的神随着创造而收缩又扩张，然后无止境地休息。因为被扩张的东西先是收缩的；被扩展的东西起初是收缩的；它原本是那样，只是经历了一种感受。如果那\OriginalTerm{一元}{Monad}被扩张而成了\OriginalTerm{三元}{Triad}，而一元就是圣父，三元则是圣父、圣子及圣神，那么首先，一元在被扩张时，经历了一种感受，变成了它原先不是的东西；因为它被扩张了，而原先并未扩张。其次，如果一元本身被扩张成三元，即圣父、圣子和圣神，那么圣父、圣子和圣神就证明是同一的，正如\Person{撒伯里乌}所主张的，除非他所说的那\OriginalTerm{一元}{Monad}是圣父之外的什么，那么他就不该谈论扩张，因为一元要成为三，这样就有一个一元，然后才有圣父、圣子和圣神。因为如果一元被扩张，并扩展自己，那么它自己必定就是那被扩展的。而三元在被扩张时就不再是一元，当其是一元时则还不是三元。这样，那曾是圣父的还不是圣子和圣神；但当成为这些时，就不再只是圣父了。而一个如此说谎的人，必然将形体归于天主，并把祂描绘成有情欲的；因为扩张是什么，不就是那被扩张之物的感受吗？或者说，那被扩张的是什么，不就是先前并非如此，而是确实狭窄的吗？因为它还是同一个，只是在时间上与自己不同。\par

14. 这位神圣的宗徒知道这一点，当他写信给格林多人时，他说：\BibleQuote{你们在我们内并不狭窄，而是你们自己狭窄，格林多人啊，你们宽大吧！}\BibleRef{格后 6:12}；因为他劝勉同一群人从狭窄变为宽大。正如假设格林多人曾是狭窄的，转而宽大后，他们并没有变成别人，仍是格林多人，同样，如果圣父被扩张成三元，那么这三元就再次仅仅是圣父。他又说同样的话：\BibleQuote{我们的心是宽大的}；而\Person{诺厄}说：\BibleQuote{愿天主使\Person{耶斐特}扩张}，因为同样的心和同样的\Person{耶斐特}处于扩张之中。那么，如果一元扩张，它就会为别的而扩张；但如果它为自己扩张，那么它就会成为那被扩张的；而那被扩张的除了是圣子和圣神之外，还能是什么呢？当这样说的时候，最好问问他，这扩张的作用是什么？或者，实际上，它到底为何发生？因为凡是不能保持一样，而在时间进程中扩张的东西，必然有一个扩张的原因。如果是为了让圣言和圣神与祂同在，那么说\Quote{先是一元，然后扩张}就不合目的了；因为圣言和圣神并非后来才有，而是永远就有，否则天主就成了没有圣言的，正如亚略派所主张的。因此，如果圣言和圣神永远就有，那么它就永远扩张了，而不是起初是一元；但如果它是后来扩张的，那么后来才有圣言。但如果它是为了降生成人而扩张，然后成了三元，那么在降生成人之前就还没有三元。甚至看起来圣父成了血肉，也就是说，如果圣父就是那一元，并且在人内扩张了；如此，或许就只会有一元、血肉和第三位圣神；也就是说，如果祂自己扩张的话；那么三元就只是名义上的了。说它为了创造而扩张也是荒谬的；因为它保持为一元也能创造万物；因为一元不需要扩张，在扩张之前也不缺乏能力；在天主的事上这样想或说，实在是荒谬而不虔敬的。还会产生另一个荒谬。因为如果它是为了创造而扩张，而它是一元时创造并不存在，但在完成之后扩张停止，它将再次成为一元，那么创造也将归于虚无。因为正如它为了创造而扩张，同样，扩张停止时，创造也将停止。\par

% structure: p0022 source=h2 target=subsection reason=relative-heading-rank; base=h2

\subsection{15-24. 既然圣言来自天主，祂必是圣子。既然圣子从永恒就有，祂必是圣言；否则要么祂高于圣言，要么圣言就是圣父。\WorkTitle{新约}中提出圣子与圣父合一的经文；因此圣子是圣言。驳斥三种假说——1. 那人就是圣子；2. 圣言与人结合后才是圣子；3. 圣言在降生成人时才成为圣子。\WorkTitle{旧约}中谈论圣子的经文。如果它们仅是预言性的，那么那些关于圣言的经文也可能如此。}

15. 说一元扩张为三元，必然产生这些荒谬。但是因为说这些话的人胆敢分开圣言与圣子，说圣言是一回事，圣子是另一回事，并且先有圣言后有圣子，那么就让我们也来审查这个学说。他们的臆测形式多样；因为有些人说救主所取的人就是圣子；另一些人则说人和圣言在结合时才成为圣子。还有一些人说圣言自己在成为人时才成为圣子；因为他们说，祂从圣言变成了圣子，之前不是圣子，而只是圣言。这两种说法都是斯多噶的学说，不论说天主扩张还是否认圣子，但尤其荒谬的是，称呼圣言却否认祂是圣子。因为如果圣言不是来自天主，那麼他们否认祂是圣子还算合理；但如果祂来自天主，他们怎么看不出凡是出自某物的，就是那生出他的那一位的儿子呢？其次，如果天主是圣言的父，为什么圣言不是自己父的圣子？因为一个人之所以是父并被如此称呼，是因为他有儿子；而一个人之所以是某人的儿子并被如此称呼，是因为他有父。那么，如果天主不是基督的父，圣言也就不是圣子；但如果天主是父，那么圣言理所当然也是圣子。但如果先是天主，后来才有父，这就是亚略派的想法。再者，说天主改变是荒谬的，因为改变属于形体；但如果他们辩称，在创造的事上，祂后来成了造物主，那么他们当知道，改变存在于后来生成的事物中，而不在天主内。\par

16. 如果圣子也是受造物，那么天主完全可能像对其他受造物那样，开始成为祂的圣父；但如果圣子不是受造物，那么圣父永远存在，圣子也永远存在。但如果圣子永远存在，祂就必然是\OriginalTerm{圣言}{Word}；因为如果圣言不是圣子——有人竟敢如此主张——那么他要么认为圣言就是圣父，要么认为圣子高于圣言。因为圣子 \BibleQuote{在父怀里}\BibleRef{若 1:18}，因此必然的结果是：要么圣言不先于圣子（因为在父内的那位之前没有任何存在），要么如果圣言不同于圣子，那么圣言就必然是圣父，而圣子就在圣父内。但如果圣言不是圣父而是圣言，那么圣言必然在圣父之外，因为正是圣子\BibleQuote{在父怀里}。因为并非圣言和圣子都在父怀里，而是只有一位在其中，那就是圣子，\OriginalTerm{独生子}{Only-begotten}。从另一个理由也可推出，如果圣言是一回事，圣子是另一回事，那么圣子就高于圣言；因为\BibleQuote{除了圣子，没有人认识父}\BibleRef{玛 11:27}，而不是圣言。那么，要么圣言不知道，要么如果祂知道，那么\Quote{没有人认识}这句话就不真实了。同样地，\BibleQuote{谁看见了我，就是看见了父}，以及\BibleQuote{我与父原是一体}，因为这是由圣子说出的，而非圣言，正如他们所主张的那样，而福音书清楚表明了这一点；因为根据若望记载，当主说\BibleQuote{我与父原是一体}时，犹太人便拿起石头来要砸死祂。\BibleQuote{耶稣回答他们说：\Quote{我赖我父给你们显示了许多善事，为了哪一件，你们要砸死我呢？}犹太人回答说：\Quote{为了善事，我们不是砸死你，而是为了亵渎的话，因为你是人，却把自己当作天主。}耶稣回答他们说：\Quote{在你们的法律上不是记载着：\InnerQuote{我说过：你们是神}么？如果那些承受天主话的人，天主尚且称他们为神——而圣经是不能废弃的——那么，父所祝圣并派遣到世界上来的，因为说过：我是天主子，你们就说：你说亵渎的话么？假使我不做我父的工作，你们就不必信我；但若是我做了，你们纵然不肯信我，至少要信这些工作，如此你们必定认出父在我内，我在父内。}}然而，就字面意思而言，祂既没有说\Quote{我是天主}，也没有说\Quote{我是天主子}，而是说\BibleQuote{我与父原是一体}。\par

17. 犹太人听到\Quote{一体}时，便像\Person{萨贝里乌斯}那样认为祂声称自己是圣父，但我们的救主用以下论据指出了他们的错误：\Quote{虽然我若说了\InnerQuote{天主}，你们也应该记得经上所载的：\BibleQuote{我说过：你们是神}；}然后为了澄清\BibleQuote{我与父原是一体}，祂用这些话解释了圣子与圣父的\OriginalTerm{合一性}{oneness}：\BibleQuote{因为我说了：我是天主子}。因为即使祂没有用言语说出，祂仍将\BibleQuote{原是一体}的寓意指向了圣子。因为除了出自圣父的那位，没有任何存在与圣父是一体。那出自圣父的，除圣子外，还能是谁呢？因此祂补充说：\BibleQuote{好使你们知道我在父内，父也在我内}。因为，在解释\Quote{一体}时，祂说合一与不可分离性不在于此位格即是与之为一的彼位格，而在于祂在父内，父也在祂内。因为如此祂既推翻了\Person{萨贝里乌斯}——在说\BibleQuote{我是}时，祂并没有说\Quote{圣父}，而是说\BibleQuote{我是天主子}——又推翻了\Person{亚略}，通过说\BibleQuote{原是一体}。那么，如果圣子和圣言不相同，那么与圣父为一体的就不是圣言，而是圣子；也不是看见圣言的\BibleQuote{看见了父}，而是\BibleQuote{看见}圣子的。由此可推出，要么圣子大于圣言，要么圣言并不超越圣子。因为有什么能比\Quote{一体}、\BibleQuote{我在父内，父在我内}以及\BibleQuote{谁看见了我，就是看见了父}更伟大或更完美呢？因为这些话语都属于圣子。因此，同一部若望福音记载祂说：\BibleQuote{谁看见了我，就是看见了派遣我来的；}\BibleQuote{谁接纳我，就是接纳派遣我来的；}\BibleQuote{我身为光明，来到了世界上，使凡信我的，不留在黑暗中。无论谁，若听我的话而不遵行，我不审判他，因为我来不是为审判世界，乃是为拯救世界。他所听见的话，要在末日审判他，因为我往父那里去。}祂说，这宣讲要审判那不遵守诫命的人；\BibleQuote{因为如果我没有来，没有教训他们，他们就没有罪；但现在他们对自己的罪，是无可推诿的，}\BibleRef{若 15:22}祂说，听我的话，那些遵行的，必将获得救恩。\par

18. 也许他们会如此厚颜无耻地说，这句话不属于圣子而属于圣言；但从上文可清楚看出，说话者是圣子。因为那在这里说\BibleQuote{我不是来审判世界，而是要拯救世界}\BibleRef{若 12:47}的那一位，正是同一位\Person{若望}先前所言，\BibleQuote{天主竟这样爱了世界，甚至赐下了自己的独生子，使凡信他的人不至丧亡，反而获得永生。因为天主没有派遣子到世界上来审判世界，而是为叫世界藉着他而获救。那信从他的，不受审判；那不信的，已受了审判，因为他没有信从天主独生子的名字。审判就在于此：光明来到了世界，世人却爱黑暗甚于光明，因为他们的行为是邪恶的。}的那一位，正是天主的\OriginalTerm{独生子}{Only-begotten}。那么，若那位说\BibleQuote{我不是来审判世界，而是要拯救世界}的，与那位说\BibleQuote{看见我的，也就是看见了派遣我来的}是同一位；且那来拯救世界而非审判它的是天主的独生子，那么显然，说\BibleQuote{看见我的，也就是看见了派遣我来的}的，就是这同一位圣子。因为那曾说\BibleQuote{信我的}，以及\BibleQuote{若有人听我的话，我不审判他}的，就是圣子本人，圣经论及他说：\BibleQuote{那信从他的，不受审判；那不信的，已受了审判，因为他没有信从天主独生子的名字。}又说：\BibleQuote{审判就在于此}，即那不信从子的，\BibleQuote{光明来到了世界}，他们却不相信他，也就是子；因为他必须是\BibleQuote{那普照每人的真光，正在进入这世界}。只要他尚在地上，就降生成人而言，他就是世上的光，正如他自己说：\BibleQuote{你们几时有光，就应当信从光，好成为光明之子。}因为他说：\BibleQuote{我身为光明，来到了世界上}。\par

19. 这一点既已显明，则圣言就是圣子。但如果圣子就是那来到世界的光，那么毫无疑问，世界是藉圣子造成的。因为在这部福音的开头，这位圣史论及\Person{若翰洗者}时说：\BibleQuote{他不是那光，只是为给那光作证。}因为我们之前已说过，基督本人就是那普照每人的真光，正在进入这世界。因为若\BibleQuote{他已在世界上，世界原是藉他造成的}，那么他必是天主的圣言，关于他，这位圣史也见证万物是藉他而造成的。因为他们要么被迫讲论两个世界，一个是藉圣子造成的，另一个是藉圣言造成的；要么，如果世界只有一个，受造界也只是一个，那么圣子和圣言在一切受造之前就是同一位，因为世界是藉他而有的。所以，如果说万物如何藉圣言而有，同样也藉圣子而有，那么说\BibleQuote{在起初已有圣言}，或说\BibleQuote{在起初已有圣子}，并不矛盾，甚至完全相同。但如果因为若望没有说\BibleQuote{在起初已有圣子}，他们就坚持圣言的属性不适用于圣子，那么反过来立即得出，圣子的属性也不适用于圣言。但已证明，\BibleQuote{我与父原是一体}这一句属于圣子；他也是那位\BibleQuote{在父怀里的}；以及\BibleQuote{看见我的，也就是看见了派遣我来的}；而\BibleQuote{世界原是藉他造成的}是圣言与圣子所共有的；由此可见，圣子在创世之前就已存在；因为创造者必然先于受造之物。并且按照他们，对\Person{斐理伯}所说的话必定不属于圣言，而是属于圣子。因为圣经记载：\BibleQuote{耶稣说：\InnerQuote{斐理伯！这么长久的时候，我和你们在一起，而你还不认识我吗？谁看见了我，就是看见了父；你怎么说：把父显示给我们呢？你不信我在父内，父在我内吗？我对你们所说的话，不是凭我自己讲的；而是住在我内的父，作他自己的事业。你们要相信我：我在父内，父也在我内；若不然，你们至少该因那些事业而相信。我实实在在告诉你们：凡信我的，我所做的事业，他也要做，并且还要作比这些更大的事业，因为我往父那里去。你们因我的名无论求什么，我必要践行，为叫父在子身上获得光荣。}}所以，如果父在子身上获得光荣，那么圣子必定是那位说\BibleQuote{我在父内，父在我内}的，也是那位说\BibleQuote{谁看见了我，就是看见了父}的。因为他，就是如此说话的这一位，藉着加上\BibleQuote{为叫父在子身上获得光荣}这句，显示了他自己就是圣子。\par

倘若他们主张，\OriginalTerm{圣言}{Word}所穿戴的人，而非圣言本身，是天主的\OriginalTerm{独生子}{Only-begotten}，那么，按此推论，那人必须是在圣父内，圣父也在他内的那一位；那人必须是那与圣父为一的，是在圣父怀中的，是真光。他们就不得不承认，世界是藉着这人造成的，这人降来不是为审判世界，而是为拯救世界；并且他就是那在\Person{亚巴郎}出现以前就已存在的。因为圣经记载，\Person{耶稣}向他们说：\BibleQuote{我实实在在告诉你们：在\Person{亚巴郎}出现以前，我就是} \BibleRef{若 8:58}。说一个出自\Person{亚巴郎}后裔、相隔四十二代的人，竟在\Person{亚巴郎}出现之前就已存在，这岂不荒谬吗？若圣言所承担的肉躯本身就是圣子，却说那由\Person{玛利亚}而来的肉躯是造成世界的根源，岂不荒谬吗？他们又将如何解释\BibleQuote{祂已在世界上}这句经文呢？因为圣史为表明圣子的先存性，且在降生成人之前，便接着说：\BibleQuote{祂已在世界上}。再者，若非圣言，而是那人为圣子，祂自己既是世界中的一员，又如何能拯救世界呢？如果这还不能使他们羞愧，那么当那人在圣父内时，圣言又在哪里？当那人与圣父为一时，圣言相对圣父又处于何等位置？但若那人是独生子，圣言的位置又将何在？要么必须说圣言次于独生子，要么若圣言高于独生子，圣言就必须是圣父自己了。因为正如圣父是唯一的，从祂而来的独生子也是唯一的；若圣言不是圣子，祂又有什么高于那人的地方呢？因为尽管圣经说世界是藉着圣子和圣言而造成的，创造世界是圣言与圣子共有的工，然而圣经却接着指出，看见圣父的，不在圣言，而在于圣子；拯救世界的，也不归于圣言，而归于独生子。因为，如圣经所说，\Person{耶稣}说过：\BibleQuote{我和你们在一起这么长久，你还不认识我吗，\Person{斐理伯}？谁看见了我，就是看见了父。} 圣经也不说圣言认识圣父，而说圣子；不说圣言看见圣父，而说那在圣父怀中的独生子。\par

再者，若如他们所主张，圣子是一位，圣言是另一位，那么圣言为我们的救恩所作的贡献，能比圣子多些什么呢？因为诫命是要我们信从圣子，而非信从圣言。因为\Person{若望}说：\BibleQuote{信从子的，有永生；不信从子的，必见不到生命} \BibleRef{若 3:36}。而蕴涵着全部信仰实质的圣洗，也不是以圣言之名，而是以圣父、圣子及圣神之名施行的。那么，若按他们所持，圣言是一位，圣子是另一位，圣言并非圣子，圣洗便与圣言毫无关联。那么，在施洗时圣言不与圣父同在，他们又如何能持守圣言与圣父同在呢？或许他们会说，在圣父之名中已包含了圣言？那为何不包含圣神呢？难道圣神在圣父之外吗？这样一来，（若圣言不是圣子）那人固然是依圣父得名，但圣神却是依那人得名吗？于是，那\OriginalTerm{独一}{Monad}不但没有扩展为\OriginalTerm{三一}{Triad}，反而按照他们的说法，扩展成了\OriginalTerm{四一}{Tetrad}，即圣父、圣言、圣子和圣神。他们因这一论点而自觉羞愧，于是转向另一说法，声称不是主所承受的那人单独成为圣子，而是圣言与人两者合一才是圣子；因为他们说，二者结合才被称为圣子。那么，二者中谁是原因，谁是被导致圣子者？或者更明白地说：圣言是因肉躯而成为圣子呢？还是肉躯因圣言而被称作圣子？还是两者都不是原因，只是二者的结合？若圣言是因肉躯而成为圣子，则肉躯必然就是圣子，那么先前从主张那人是圣子的说法所推导出的一切荒谬结论，便都随之而来。但若肉躯是因圣言而被称作圣子，那么，在未有肉躯之前，圣言本身既已是如此，则必定已经是圣子了。因为，如果他自己不是圣子，又如何能使他人成为圣子呢？尤其当有一位父存在时？若他是为自己造圣子，那他自身就是父；若是为父造圣子，那他必是圣子，而且更应是那位使其他人都得以成为子民的圣子。\par

22. 因为若祂不是\OriginalTerm{圣子}{Son}，而我们却是儿子，天主便是我们的父，而不是祂的父。那么祂怎能反而自称为父，说\Quote{我的父}，又说\Quote{我从父而来}呢？因为倘若祂是众人共同的父，祂就不只是祂的父，而祂也不唯独从\OriginalTerm{圣父}{Father}而来。但他们却说，祂有时也被称为我们的父，因为祂自己分享了我们的人性。为此，\OriginalTerm{圣言}{Word}成了血肉，好使祂——因为圣言是圣子——借着居住在我们内的圣子，也能被称为我们的父；因为圣经说：\BibleQuote{他派遣了自己儿子的圣神，到我们心里喊说\InnerQuote{\OriginalTerm{阿爸}{Abba}，父啊}}\BibleRef{迦 4:6}。因此，在我们内的圣子呼求自己的父，使得祂也被称为我们的父。显然，谁的心里没有圣子，天主也就不能称为他们的父。但如果那人是因着圣言而被称为圣子，那么必然可推知，既然古人在降生成人之前就被称为儿子，圣言在祂暂居我们中间之前就已经是圣子了；因为经上说：\BibleQuote{我生了儿子们}；在\Person{诺厄}的时代，\BibleQuote{天主的儿子们看见}；以及在那首\BibleBook{}{雅歌}中，\BibleQuote{祂不是你的父吗？}因此，也有那位\OriginalTerm{真圣子}{True Son}，正是为了祂的缘故，他们才成了儿子。但如果，如他们再次声称的，两者都不是圣子，而是取决于两者的结合，那么结果就是两者都不是圣子；我是说，既不是圣言，也不是那人，而是他们结合所凭借的某个原因；因此，那使成为圣子的原因便先于结合而存在。因此，按这种方式，圣子在血肉之前就已存在。当这一点被追问时，他们就要寻求另外的借口说：那人不是圣子，两者结合也不是，而是圣言起初确实是纯然如字的圣言，但当祂成了人，那时才被称为圣子；因为在祂显现之前，祂不是圣子，仅仅是圣言；正如\BibleQuote{圣言成了血肉}，先前不是血肉，同样，圣言成为圣子，先前不是圣子。这就是他们的虚言；但很容易被驳倒。\par

23. 因为如果仅仅在成了人之后，祂才成为\OriginalTerm{圣子}{Son}，那么成为人本身就是原因。而如果那人是祂成为圣子的原因，或两者结合是原因，那么同样荒谬的结论就会出现。再者，如果祂先是\OriginalTerm{圣言}{Word}，后是圣子，那么祂就不是本有，而是后来才认识\OriginalTerm{圣父}{Father}；因为祂不是作为圣言认识圣父，而是作为圣子。因为\BibleQuote{除了子，没有人认识父}。这也会导致：祂后来才\BibleQuote{在父怀里}\BibleRef{玛 11:27}，后来才与圣父成为一体；后来才有\BibleQuote{看见我的，就是看见了父}\BibleRef{若 14:9}。因为所有这些话都是论及圣子的。因此，他们被迫要说，圣言只是一个名字而已。因为那在我们内与父同在的就不是祂，看见圣言的也并非就是看见了父，父也根本不为任何人所知——因为父是借着子才被认识的（因为经上记着：\BibleQuote{子所愿意启示的人}），而如果圣言还不是子，就没有人认识父。那么，祂怎么会被\Person{梅瑟}看见，怎么会被先祖看见呢？因为祂自己在\BibleBook{列王}{纪}中说：\BibleQuote{我不是向你父家清楚显现了吗？}但是，如果天主曾显现，就必定有一位子来启示，正如祂自己说的：\BibleQuote{子所愿意启示的人}。因此，说圣言是一回事，圣子是另一回事，不仅不虔敬，而且愚蠢，至于他们这种想法从何而来，倒要好好问问他们。他们回答说，因为在\OriginalTerm{旧约}{Old Testament}中没有提到子，只提到圣言；因此他们确信，子晚于圣言，因为只有在\OriginalTerm{新约}{New Testament}中，而非旧约中，才谈到子。这就是他们不虔敬的说法；因为首先，把两个约分开，使得彼此不相关联，正是摩尼教徒和犹太人的伎俩，前者反对旧约，后者反对新约。再者，按他们的说法，如果旧约所含的内容年代更早，新约的年代更晚，而时间取决于写作的时间，那么\BibleQuote{我与父原为一体}、\OriginalTerm{独生子}{Only-begotten}、以及\BibleQuote{看见我的，就是看见了父}等话，就都是较晚的了，因为这些见证并非引自旧约，而是引自新约。\par

然而事实并非如此；因为在旧约中关于圣子所说的话也确实很多，如在\BibleBook{圣咏}{第二篇}中：\BibleQuote{你是我的儿子，我今日生了你}；在\BibleBook{圣咏}{第九篇}的标题中：\BibleQuote{交与乐官。关于儿子隐藏之事的训诲诗。达味诗歌。}；在\BibleBook{圣咏}{第四十四篇}中：\BibleQuote{交与乐官，属于科辣黑子孙的训诲歌，关于爱子的歌。}；在\BibleBook{依撒意亚}{}书中：\BibleQuote{我要歌唱我的爱子，唱一首关于他葡萄园的情歌。我的爱子有一座葡萄园。}\BibleRef{依 5:1}；这\Quote{爱子}是谁，岂不就是\OriginalTerm{独生子}{Only-begotten}吗？同样在\BibleBook{圣咏}{第一百零九篇}中说：\BibleQuote{在晓明之前，我由胎中生了你}，关于这一点我稍后再讲；在\BibleBook{箴言}{}中：\BibleQuote{远在诸山未曾奠定以前，我已受生}；在\BibleBook{达尼尔}{}书中：\BibleQuote{并且那第四位的相貌，好像神子}；以及许多其他经文。因此，如果古老性源自旧约，那么圣子必然也是古老的，他在旧约中多处被清楚描述。他们却说：\Quote{是的，是这样，但必须按先知性意义理解。}因此，也必须说\OriginalTerm{圣言}{Word}是先知性地被论及；因为这不能一种方式理解这样，另一种方式理解那样。因为如果\BibleQuote{你是我的儿子}指的是将来，那么\BibleQuote{因上主的话，诸天造成}也指的是将来；因为经上并非说\InnerQuote{被造成}，也不是\InnerQuote{祂造了}。但\InnerQuote{造成}指的是将来，这在另一处写道：\BibleQuote{上主为王}，紧接着说\BibleQuote{祂奠定了大地，使它永不动摇}。而如果\BibleBook{圣咏}{第四十四篇}中的\BibleQuote{关于我的爱子}指的是将来，那么紧随其后的\BibleQuote{我的心灵涌溢优美的圣言}也指的是将来。如果\BibleQuote{从胎中}指的是人，那么\BibleQuote{从心中}也指的是人。因为如果胎是属人的，那么心也是肉体的。但如果从心而发者是永恒的，那么\InnerQuote{从胎而生}者也是永恒的。如果\Quote{独生子}是\Quote{在父怀里}，那么\Quote{爱子}也是\Quote{在父怀里}。因为\InnerQuote{独生子}和\InnerQuote{爱子}是相同的，正如经上说：\BibleQuote{这是我的爱子}。因为祂说\InnerQuote{爱子}，并非为表示对祂的爱，好像显得祂恨别人似的，而是借此表明祂是独生子，为显示唯独祂是由祂而来的。因此，圣言为向亚巴郎传达\InnerQuote{独生子}的观念，说：\BibleQuote{带上你的儿子，你心爱的独生子}\BibleRef{创 22:2}；但任何人都清楚，依撒格是撒辣唯一的儿子。所以，圣言是子，不是后来才成为子或被称为子的，而是永为子。因为如果不是子，祂也就不是圣言；如果不是圣言，祂也就不是子。因为由父而来的就是子；而由父而来的，除了那从心中发出、由胎而生的圣言，还能是什么？因为父不是圣言，圣言也不是父；而一位是父，另一位是子；一位生，另一位受生。\par

% structure: p0033 source=h2 target=subsection reason=relative-heading-rank; base=h2

\subsection{\OriginalTerm{马赛里安派}{Marcellian}引用\BibleRef{格前 12:4}的例证，予以驳斥。}

亚略于是妄言圣子是从虚无中来的，并且曾有一时不存在；同样，撒伯里乌也妄言父即是子，子又是父，在\OriginalTerm{位格}{subsistence}上是一位，在名称上是两位；他还妄用圣神的恩宠作为例证。因为他声称：\Quote{正如神恩虽有区别，却是同一圣神}，同样，父也是一位，但扩展为子和圣神。然而这完全是荒诞的；因为如果天主如同圣神那样，那么父就会成为圣言和圣神，对这人成为父，对那人成为子，对另一人成为圣神，以适应每人的需要；在名称上虽是子与圣神，但实质上只是父；祂在成为子时便有了开始，然后又停止被称为父，并且在名义上降生成人，但事实上甚至没有来到我们中间；祂说\BibleQuote{我与父原是一体}也是不真实的，因为在现实中祂自己就是父，还有其他从撒伯里乌的例子而来的荒谬后果。而且当需要被满足之后，子与圣神的名号必然会停止；所发生的一切将完全只是伪装，因为它仅仅是名义上的表现，而非真实的。而且当他们主张子的名号停止时，那么圣洗圣事的恩宠也将停止；因为它是因着圣子而赐下的。岂止于此，接下来岂不就是受造物的毁灭？因为如果圣言出来了是为使我们受造，而当祂出来时我们便存在了，那么当祂如他们所说退回到父内时，我们显然也就不再存在了。因为祂将回到祂原先的状态；我们也将不再存在，如同那时我们不存一样；因为当祂不再出来时，受造物也就不复存在。因此，这实在是荒谬的。\par

% structure: p0035 source=h2 target=subsection reason=relative-heading-rank; base=h2

\subsection{由新约论证圣子是永与圣父同在的圣言。续论旧约经文；尤其\BibleRef{咏 110:3}。此外，旧约中的圣言可以是新约中的圣子，正如旧约中的圣神是新约中的\OriginalTerm{护慰者}{Paraclete}。\BibleRef{宗 10:36}的异议；通过类似经文回应，如\BibleRef{格前 1:5}、\BibleRef{肋 9:7}等。圣言取肉身的必要性，即为了圣化肉身而不毁灭肉身。}

26. 然而，圣子并没有存在之始，在他降生成人之前，他常与圣父同在，\Person{若望}在他的第一封书信中清楚表明，写道：\BibleQuote{论到那从起初就有的生命之言，就是我们听见过、亲眼看见过、瞻仰过、并且用手摸过的；这生命已显现出来，我们看见了，也为他作证，且把这原与圣父同在，且已显现给我们的永生，传报给你们}\BibleRef{若一 1:1}。他在这里说\Quote{这生命}不是\Quote{成为}，而是\Quote{与圣父同在}，而在书信的末尾，他说圣子就是生命，写道：\BibleQuote{我们是在那真实者内，即在他的子耶稣基督内，他即是真实的天主和永生。} 但若圣子是生命，生命又与圣父同在，且圣子与圣父同在，同一位圣史又说：\BibleQuote{圣言与天主同在}\BibleRef{若 1:1}，那么圣子必是那常与圣父同在的圣言。而既然\Quote{圣子}即是\Quote{圣言}，那么\Quote{天主}必是\Quote{圣父}。再者，根据\Person{若望}，圣子不仅是\Quote{天主}，而是\Quote{真天主}；因为同一位圣史说：\BibleQuote{圣言就是天主}；而圣子说：\BibleQuote{我就是生命。} 因此，圣子就是那与圣父同在的圣言和生命。此外，同一\Person{若望}书中说：\BibleQuote{那在圣父怀里的独生子}，表明圣子永恒常在。因为\Person{若望}称为圣子的那位，\Person{达味}在圣咏中称他为\OriginalTerm{天主的手}{God's Hand}，说：\BibleQuote{你为何不将你的右手从怀中伸出？} 那么，若手在怀中，圣子在怀中，圣子便是手，手便是圣子，圣父藉着他创造万物；因为经上记载：\BibleQuote{你的手创造了这一切，} 又说：\BibleQuote{他用自己的手引领他的百姓}；所以是藉着圣子。而若说\Quote{这就是至高者右手的改变}，又说：\BibleQuote{到了末了，论及那将要改变的事，为我的爱子之歌}；那么爱子就是那改变了的手；关于他，天父的声音也说：\BibleQuote{这是我的爱子。} 因此，\Quote{\OriginalTerm{我的手}{My Hand}}等同于\Quote{\OriginalTerm{我的儿子}{My Son}}。\par

27. 但有一些受教不良的人，他们抗拒圣子的道理，轻视那说\BibleQuote{在晓星之前，我从母胎中生了你}的话，以为这是指他与\Person{玛利亚}的关系，声称他是\Person{玛利亚}\Quote{在晓星之前}所生，因为\Quote{母胎}一词不可能指他与天主的关系，我们在此不得不说几句话。如果因为\Quote{母胎}是属人的，所以它与天主格格不入，那么显然\Quote{心}也有属人的意思，因为有心的也有母胎。既然两者都是属人的，我们必须两者都否定，或设法两者都解释。正如言语出于心，后裔也出于母胎；正如论及天主的心时，我们并不以人意来理解，同样，若圣经说\Quote{从母胎}，我们也不可按肉体的意义来理解。因为圣经通常以人的方式来说并意指那超乎人性之事。例如论及创造时说：\BibleQuote{你的手造成了我，塑造了我}，又说：\BibleQuote{你的手创造了这一切}，又说：\BibleQuote{他一命，它们便受造。} 因此，圣经的用语对一切都合适；它将\Quote{本有}和\Quote{真实}归于圣子，将\Quote{存在之始}归于受造物。因为天主造且创万物；但祂从自己生成了圣子，即\OriginalTerm{圣言}{Word}或\OriginalTerm{智慧}{Wisdom}。而\Quote{母胎}与\Quote{心}清楚表明了本有和真实；因为我们也是从母胎而有此；但我们用双手制造作品。\par

28. 那么，他们说，\Quote{在晓星之前}是什么意思？我要回答，如果\Quote{在晓星之前}表明他由\Person{玛利亚}的诞生是奇妙的，那么许多其他人在星升起之前也已出生。那么，在他身上所说的，有什么如此奇妙，竟使它被记载为某种特选的特恩，既然这是许多人共有的？其次，\Quote{生}与\Quote{产生}不同；因为\Quote{生}包含首原的基础，而\Quote{产生}无非是使已有者出现。如果这词语属于身体，就当注意，他并非在夜间牧羊人得到喜讯时才获得存在的开始，而是在天使对童贞女说话的时候。而那并非夜间，因为经上并未说那是夜间；相反，当他从母胎中出来时才是夜间。圣经作出了这种区分，一方面说他在晓星之前受生，另一方面又说他从母胎中出来，如第二十一篇圣咏中说的：\BibleQuote{是你使我从母胎中出来。} 此外，他没有说\Quote{在晓星升起之前}，而只简单地说\Quote{在晓星之前}。那么，如果这句话必须指身体，那么身体要么必须在\Person{亚当}之前，因为星辰在\Person{亚当}之前，要么我们就不得不追究字面含义。而这\Person{若望}使我们能够做到，他在\BibleBook{默示录}{}中说：\BibleQuote{我是阿尔法和奥默加，肇始和终末，元始和终末。那些使自己衣裳阔大的，是有福的！他们可以对生命树有权柄，并从门进入城。在城外有狗、邪术的、行奸的、杀人的、拜偶像的，以及一切喜爱并实行欺诈的人。我耶稣派遣了我的使者，给你们在众教会内证实这些事。我是\Person{达味}的根苗和后裔，又是那颗光亮的晓星。\OriginalTerm{圣神}{Spirit}和新娘说：来罢！听见的也说：来罢！渴的，也来罢！凡愿意的，都可以白白取用生命的水。} 那么，如果\Quote{\Person{达味}的后裔}就是\Quote{那颗光亮的晓星}，显然，救主的血肉被称为\Quote{晓星}，而自\Person{天主}而来的后裔是在此之前；因此，那圣咏的意思就是：\BibleQuote{在你的血肉显现之前，我从自己生了你}；因为\Quote{在晓星之前}等同于\Quote{在\OriginalTerm{圣言}{Word}降生成人之前}。\par

29. 同样在旧约中，关于子也有明确的陈述；同时，辩论这一点是多余的；因为如果旧约没有提及的都是后来才有的，那么那些好辩的人，请指出旧约中哪里提到了圣神，即护慰者？因为旧约虽然提到了圣神，却没有一处提到护慰者。难道圣神是一位，护慰者是另一位，而护慰者是后来才有的，因为旧约没有提及吗？但断乎不可说圣神是后来的，或将圣神与护慰者分开；因为圣神是唯一且同一的，无论过去还是现在，祂都圣化并安慰那些领受祂的人；正如唯一且同一的圣言和子，那时也引导配得的人成为义子。因为旧约下的众子，正是通过子而成为子。因为如果甚至在玛利亚之前，还没有出自天主的子，那么当众子先于祂时，祂怎能先于万有呢？如果祂在许多人之后，又怎能是\Quote{首生者}呢？然而，护慰者也不是第二位的，因为祂先于万有；子也不是后来的，因为\BibleQuote{在起初已有圣言}\BibleRef{若 1:1}。正如圣神与护慰者是同一的，子与圣言也是同一的；救主论到圣神说：\BibleQuote{但那护慰者，就是父因我的名所要派遣来的圣神}\BibleRef{若 14:26}，是指着同一位，并不作区分；同样，若望也类似地描述说：\BibleQuote{于是，圣言成了血肉，寄居在我们中间；我们见了祂的光荣，正如父独生者的光荣}\BibleRef{若 1:14}。因为在此他也没有区分，而是见证同一性。正如护慰者并不与圣神有别，而是同一位，同样，圣言并不与子有别，而是同一的，因为圣言就是独生子；因为圣经所说的并不是血肉本身的光荣，而是圣言的光荣。因此，谁若敢在圣言与子之间作区分，就让他先在圣神与护慰者之间作区分；但如果圣神不可区分，那么圣言也同样不可区分，因为祂也是子、智慧与大能。此外，就连精通辞令的希腊人也知道，\Quote{所爱的}一词与\Quote{独生的}等同。荷马在\WorkTitle{奥德赛}第二卷中，就这样描述\Person{尤利西斯}的独生子\Person{忒勒玛科斯}：\par

\Quote{亲爱的青年，你为何奔波于广袤大地，\newline{}你可是独生子，是所爱的儿子？}\par

\Quote{你所哀悼的那位，神圣的\Person{尤利西斯}，已倒下，\newline{}远离故土，在异邦人居住之地。}\par

因此，一个人若是父亲的独生子，就被称为\Quote{所爱的}。\par

30. 撒摩撒塔派的某些追随者，将圣言与子区分开来，宣称子是基督，而圣言是另一位；他们以\Person{伯多禄}在\WorkTitle{宗徒大事录}中的话为根据，这些话他讲得好，他们却解释得糟。这话是：\BibleQuote{祂将圣言传给以色列子民，藉着耶稣基督宣讲和平；这一位是万有的主}\BibleRef{宗 10:36}。因为他们说，既然圣言是藉着基督说话，如同先知的情形，\BibleQuote{上主这样说}，那么先知是一位，上主是另一位。但对此，我们可以同样用致\WorkTitle{格林多人前书}的话来反驳：\BibleQuote{等待我们的主耶稣基督的出现，他也必要坚固你们到底，使你们在我们的主耶稣基督的日子上，无瑕可指}\BibleRef{格前 1:7}。因为正如一位基督并不坚固另一位基督的日子，而是祂自己，在自己的日子里坚固那些等候祂的人，同样，父派遣成了血肉的圣言，为使祂成了人，能藉着自己宣讲。因此，他紧接着说：\BibleQuote{这一位是万有的主；}但万有的主就是圣言。\par

31. \BibleQuote{梅瑟对\Person{亚郎}说：你到祭台前，奉献你的赎罪祭和全燔祭，为你自己与人民赎罪；再奉献人民的祭品，为他们赎罪，正如上主吩咐\Person{梅瑟}的}\BibleRef{肋 9:7}。看，这里虽然梅瑟是一位，但梅瑟自己却说\Quote{正如上主吩咐梅瑟}，仿佛是在说另一位梅瑟。同样地，如果真福\Person{伯多禄}也论及天主圣言，说是藉着\Person{耶稣}基督传给以色列子民，我们不必理解圣言是一位，基督是另一位，而应理解他们因着在祂出于天主性的慈爱俯就并成了人的过程中所发生的结合，是一位且是同一位。即使可以从两个角度来理解祂，这仍然没有分割圣言，正如受默感的\Person{若望}所说：\BibleQuote{于是，圣言成了血肉，寄居在我们中间}\BibleRef{若 1:14}。所以，真福\Person{伯多禄}所说的话既美善又正确，而撒摩撒塔派的追随者却理解得既恶劣又错误，不站在真理中。因为在圣经中，基督可以从两种角度来理解，例如圣经说基督是\BibleQuote{天主的德能和天主的智慧}\BibleRef{格前 1:24}。因此，如果\Person{伯多禄}说圣言是藉着耶稣基督传给以色列子民，我们应该理解他的意思是，成了血肉的圣言显现给了以色列子民，这样就能与\BibleQuote{于是，圣言成了血肉}\BibleRef{若 1:14}相符。但如果他们另有理解，虽然承认圣言是神圣的，正如祂原本那样，却将祂所取的人性与祂分离——而我们相信祂正因此人性与祂合而为一——并说祂是藉着耶稣基督而被派遣，那么他们就在不知不觉中自相矛盾了。因为那些在此处将天主圣言与天主的降生成人分开的人，似乎对祂成了血肉的教义持有低劣的理解，且怀有外邦人的思想，竟认为天主的降生成人是圣言的一种变化。但事情并非如此；断乎不可！\par

32. 因为正如\Person{若望}在此所宣讲的不可理解的结合：\BibleQuote{有死的被生命吞灭}\BibleRef{格后 5:4}，不，是被身为生命本身的那位所吞灭（正如主对\Person{玛尔大}所说的：\BibleQuote{我就是生命}\BibleRef{若 11:25}），同样，当真福\Person{伯多禄}说圣言是通过耶稣基督被派遣的，他也暗示了这天主性的结合。因为正如当一个人听到\BibleQuote{圣言成了血肉}时，他不会认为圣言不再存在——这是荒谬的，正如先前所说的——同样，当他听闻与肉身结合的圣言时，也让他理解那唯一而单纯的天主性奥迹。然而，比一切推理都更清楚且无可争辩的，是总领天使对\OriginalTerm{天主之母}{Bearer of God}本人所说的话，它显示了天主圣言与人的合而为一。因为他说：\BibleQuote{圣神要临到你，至高者的能力要庇荫你，因此，那要由你诞生的圣者，将被称为天主之子}\BibleRef{路 1:35}。那么，\OriginalTerm{萨莫撒塔人}{Samosatene}的追随者无理地将那被清楚宣告与出自\Person{玛利亚}的人合而为一的圣言分开。因此，圣言不是通过那个人被派遣；而是祂在他内派遣了，说：\BibleQuote{你们要去，教训万民！}\BibleRef{玛 28:19}。\par

33. 这是圣经惯用的方式，用朴实简单的言辞来表达。例如，在\BibleBook{}{户籍纪}中我们会发现，\Person{梅瑟}对\Place{米德杨}人，\Person{梅瑟}的岳父\Person{勒乌耳}说。因为并不是有一个说话的\Person{梅瑟}，另有一个岳父是\Person{勒乌耳}的\Person{梅瑟}，而是同一个\Person{梅瑟}。同样，如果天主圣言被称为\OriginalTerm{智慧}{Wisdom}、\OriginalTerm{能力}{Power}、\OriginalTerm{右手}{Right-Hand}、\OriginalTerm{臂膀}{Arm}等等，并且如果祂因着对人的爱而与我们合而为一，穿上我们的初果并与它融合，那么，其他的名号自然也成了圣言的属性。因为\Person{若望}说过：\BibleQuote{在起初已有圣言，圣言与天主同在，圣言就是天主，万物是藉着祂而造成的，没有祂，什么也没有造成}，这清楚地表明，连人也是天主圣言的塑造。如果祂在接纳了那衰弱的人到自己内后，藉着那确实的更新，使他重新达至无穷的持久，并因此与他合而为一，为使他获得更神圣的命运，我们怎能说圣言是通过那出自\Person{玛利亚}的人而被派遣，并竟把祂——\OriginalTerm{宗徒之主}{Lord of Apostles}——等同于其他宗徒，即那些由祂派遣的先知呢？基督又怎能被称为一个纯粹的人呢？相反，由于与圣言合而为一，祂理应被称为基督和天主之子，远在旧约，先知就已大声且清楚地将圣父的位格归于祂，曾说：\Quote{我要派遣我的儿子基督}，并在\Place{约旦河}畔说：\BibleQuote{这是我的钟爱之子}。因为当祂履行了祂的许诺时，便恰如其分地显示出，祂就是祂所说要派遣的那一位。\par

34. 那么，让我们从两面来看基督：在\Person{玛利亚}内与出自\Person{玛利亚}的那位合而为一的天主圣言。因为圣言在她胎中为自己建造了居所，一如起初祂从泥土造了\Person{亚当}；或者更神圣地说，关于此事，\Person{撒罗满}也公开说过——他知道圣言也被称为\OriginalTerm{智慧}{Wisdom}——\BibleQuote{智慧为自己建造了房屋}\BibleRef{箴 9:1}；宗徒解释这话时说：\BibleQuote{这房屋就是我们}\BibleRef{希 3:6}，在别处又称我们为圣殿，因为天主居于圣殿是合宜的；这圣殿的肖像，就是用石头建造的，是祂曾命古代人民藉\Person{撒罗满}所建；因此，当真理显现时，肖像就终止了。因为当那些残忍的人想要证明肖像是真理，并毁灭我们深信为祂与我们结合的真实居所时，祂并没有威胁他们；但祂知道他们的罪行是对付自己，就对他们说：\BibleQuote{你们拆毁这座圣殿，三天之内我要把它重建起来}\BibleRef{若 2:19}，我们的救主，由此确实表明，人们所忙碌的事物本身带着毁灭。因为若非上主建造房屋，护守城堡，建造者的劳碌便是枉然，守卫者的警醒也是徒劳。因此，犹太人的工程被拆毁了，因为它们只是影子；但教会则稳固建立；它\Quote{建在磐石上}，\Quote{阴间的门决不能胜过它}。那些说\Quote{祢既是人，怎么把自己当作天主呢？}的人，他们的门徒就是\OriginalTerm{萨莫撒塔人}{Samosatene}；因此，他向他的追随者传授其异端，是理所当然的。但是，\Quote{如果我们听过}基督，并从祂受了教，\Quote{就应脱去那照欺骗的私欲而败坏的旧人}，穿上\Quote{那照着天主，在正义和真理的圣善中受造的新人}。那么，让我们以虔敬的方式从两面来看基督。\par

但若圣经屡次以基督之名来称呼其身体，如\Person{圣伯多禄}对\Person{科尔乃略}所说，教导他有关\BibleQuote{\Place{纳匝肋}人耶稣，天主以圣神傅了祂}；又对犹太人说：\BibleQuote{\Place{纳匝肋}人耶稣，天主为你们所举扬的人}；\Person{圣保禄}对\Place{雅典}人也说：\BibleQuote{祂已立定了一人，使众人确信，使祂从死者中复活了}\BibleRef{宗 17:31}（我们可见立定与派遣常与傅油同义；由此，凡愿意的人可以学到，圣作者们的话并无矛盾，他们不过是用不同名称称呼\OriginalTerm{天主圣言}{God the Word}与由\Person{玛利亚}所生之人之间的合一，有时称傅油，有时称派遣，有时称立定）。由此可知，\Person{圣伯多禄}所言属实，并纯粹宣讲独生子的\OriginalTerm{天主性}{Godhead}，毫不将天主圣言的\OriginalTerm{自立性}{subsistence}与由\Person{玛利亚}所生之人分开（\Emph{万万不可！}因为他既多次听闻\BibleQuote{我与父原是一体}及\BibleQuote{谁看见我，就是看见了父}，又怎能如此？）在那人内，复活后亦然，门扇紧闭时，我们知道祂来到众宗徒中间，藉着祂的话驱散一切难以相信的事：\BibleQuote{触摸我看看，因为鬼魂没有肉和骨，如同你们看我却是有的}\BibleRef{路 24:39}。祂没有说\Quote{这}或\Quote{我所取的这人}，而是说\Quote{我}。因此，\OriginalTerm{萨摩萨塔派}{Samosatene}将毫无凭据，既有如此多论据证实天主圣言的合一，更有\OriginalTerm{天主圣言}{God the Word}亲自给众人带来讯息，藉进食而使他们确信，并许可他们当时触摸祂。因为给人食物者，与给祂食物者，手必接触。因为圣经说\BibleQuote{他们给了祂一块烤鱼和蜜房，祂当着他们的面吃了，然后拿起剩下的，给了他们}。看，虽未允许\Person{多默}，却以另一种方式使他们完全确信，就是被他们触摸；但你若想见伤痕，则当求教于\Person{多默}。\BibleQuote{伸出你的手探入我肋旁，伸出你的指头看我的手}\BibleRef{若 20:27}；此乃天主圣言，论及自己的肋旁与手，论及自身兼真人真天主，首先若我们认为，藉紧闭门户而入，是使圣人透过身体而感知圣言；继而藉祂的身体立于他们中间，给予完全的确信。凡此种种，为坚固信者、纠正不信者，如此说明颇为合宜。\par

因此，\Person{萨摩萨塔的保禄}亦当受纠正，因他听到那说\BibleQuote{我的身体}的天主之声，祂不是说\Quote{基督在我这圣言之外}，而是说\Quote{祂与我同在，我也在祂内}。因为我\OriginalTerm{圣言}{Word}是\OriginalTerm{圣化之油}{chrism}，而从我领受圣化之油的那位是人；若没有我，祂就不能被称为基督，但祂与我同在，我也在祂内。所以，提及圣言的派遣，便显明了那在生于\Person{玛利亚}的耶稣身上所发生的结合，祂的名字意为救世主，这并非由于别的，而是由于那人与\OriginalTerm{天主圣言}{God the Word}成为一体。这一段与\BibleQuote{派遣我的父}及\BibleQuote{我不是由我自己来的，而是父派遣了我}意义相同。因为他把与那人结合称为派遣，藉此人，\OriginalTerm{不可见的天主性}{Invisible nature}可通过可见者为人所知。因为天主并不像我们隐藏在地方那样变换位置，当祂以我们卑微的方式，在肉身中显示自己的存在时；那充满天地者怎能如此？但因着在肉身内的临在，义人论及祂的派遣。所以，\OriginalTerm{天主圣言}{God the Word}自己就是由\Person{玛利亚}所生的基督，\OriginalTerm{天主}{God}而\OriginalTerm{人}{Man}；不是另一位基督，而是同一者；万世之前由\OriginalTerm{父}{Father}所生，末期由\OriginalTerm{童贞女}{Virgin}所生；先前连天上的圣善之神体也看不见，如今因祂与可见之人合一而可见；我说，不是在不可见的\OriginalTerm{天主性}{Godhead}上被看见，而是藉着人性与全部人而运作的天主性，这人性已因与祂相结合而被更新。愿\OriginalTerm{朝拜}{adoration}与\OriginalTerm{钦崇}{worship}归于祂，往昔、现今及将来，直至万世万代。阿们。\par

\SourceRecord{Translated by John Henry Newman and Archibald Robertson.；Nicene and Post-Nicene Fathers, Second Series；Vol. 4.；Edited by Philip Schaff and Henry Wace.；Buffalo, NY: Christian Literature Publishing Co.,；1892.；<http://www.newadvent.org/fathers/28164.htm>.}
